%\RequirePackage[notext,nomath]{stix}
\documentclass[12pt]{amsart}
\usepackage{amsmath,amsthm,amssymb}
\newcommand\hmmax{0}
\newcommand\bmmax{0}
\usepackage{dsfont}
\usepackage{mathrsfs} %use \mathscr
\usepackage[all,cmtip]{xy}
\usepackage{patchcmd}
\input{diagxy}
\usepackage{hyperref}
\usepackage{bm}
\usepackage{mathtools}
\usepackage{enumerate}
\usepackage{stmaryrd}
%\usepackage{color}
\usepackage[usenames,dvipsnames,svgnames,table]{xcolor}

% BEGIN REFEREE INSERTION PREAMBLE
% Purple referee remarks; all original color commands are unchanged.
\definecolor{RefereePurple}{RGB}{103,34,147}
\newcommand{\RefereeRemark}[4]{%
  \par\medskip\begingroup\normalfont\small\color{RefereePurple}%
  \emergencystretch=1.5em\noindent\textbf{Referee R#1 [#2] -- #3.}\par\nobreak\noindent #4\par
  \endgroup\medskip
}
% Author-approved resolutions are green; the original referee text is retained in gray.
\definecolor{ResolutionGreen}{RGB}{0,100,55}
\newcommand{\ResolvedRefereeRemark}[5]{%
  \par\medskip\begingroup\normalfont\small\color{gray}%
  \emergencystretch=1.5em\noindent\textbf{Resolved referee R#1 [#2] -- #3.}\par\nobreak\noindent #4\par
  \medskip\color{ResolutionGreen}\noindent\textbf{Author resolution.} #5\par
  \endgroup\medskip
}
\newcommand{\PartlyResolvedRefereeRemark}[5]{%
  \par\medskip\begingroup\normalfont\small\color{gray}%
  \emergencystretch=1.5em\noindent\textbf{Partly resolved referee R#1 [#2] -- #3.}\par\nobreak\noindent #4\par
  \medskip\color{ResolutionGreen}\noindent\textbf{Author resolution.} #5\par
  \endgroup\medskip
}
% END REFEREE INSERTION PREAMBLE
\usepackage[normalem]{ulem}
\usepackage[top=2.3cm, bottom=2.3cm, left=2cm, right=2cm]{geometry}
%\usepackage{showkeys}
\allowdisplaybreaks

%MYlogic

\newcommand{\arthur}[1]{{\color{red}#1}}
\newcommand{\yimu}[1]{{\color{blue}#1}}
\newcommand{\memph}[1]{{\color{magenta}\emph{#1}}}

%table of contents------------------------------------------------------------------
\makeatletter
\patchcommand\@starttoc{\begin{quote}}{\end{quote}}
\def\@tocline#1#2#3#4#5#6#7{\relax
  \ifnum #1>\c@tocdepth % then omit
  \else
    \par \addpenalty\@secpenalty\addvspace{#2}%
    \begingroup \hyphenpenalty\@M
    \@ifempty{#4}{%
      \@tempdima\csname r@tocindent\number#1\endcsname\relax
    }{%
      \@tempdima#4\relax
    }%
    \parindent\z@ \leftskip#3\relax \advance\leftskip\@tempdima\relax
    \rightskip\@pnumwidth plus4em \parfillskip-\@pnumwidth
    #5\leavevmode\hskip-\@tempdima
      \ifcase #1
       \or\or \hskip 1em \or \hskip 2em \else \hskip 3em \fi%
      #6\nobreak\relax
    \dotfill\hbox to\@pnumwidth{\@tocpagenum{#7}}\par
    \nobreak
    \endgroup
  \fi}
\makeatother

\newcommand{\quotnum}[1]{ \vspace{.1in} \refstepcounter{equation}%
\noindent (\thesection.\arabic{equation}) \hspace{.6cm} \pbox{0.85\textwidth}{#1} \vspace{.1in}}

% THEOREM Environments ---------------------------------------------------
 \theoremstyle{plain}
 \newtheorem{thm}{Theorem}[section]
 \newtheorem{cor}[thm]{Corollary}
 \newtheorem{lem}[thm]{Lemma}
 \newtheorem{fact}[thm]{Fact}
 \newtheorem{prop}[thm]{Proposition}
 \newtheorem{cla}[thm]{Claim}
 \newtheorem{conj}[thm]{Conjecture}
\theoremstyle{definition}
 \newtheorem{defn}[thm]{Definition}
 \newtheorem{test}[thm]{Test}
 \newtheorem{hyp}[thm]{Hypothesis}
 \newtheorem{crit}[thm]{Criterion}
 \newtheorem{axio}[thm]{Axiom}
\theoremstyle{remark}
 \newtheorem{rem}[thm]{Remark}
 \newtheorem{ques}[thm]{Question}
 \newtheorem{quess}[thm]{Questions}
 \newtheorem{ter}[thm]{Terminology}
 \newtheorem{nota}[thm]{Notation}
 \newtheorem{conv}[thm]{Convention}
 \newtheorem{exam}[thm]{Example}

 \numberwithin{equation}{section}
 \renewcommand{\theequation}{\thesection.\arabic{equation}}

\newtheorem*{incla}{Claim}
\newtheorem*{sta}{\textbf{\textup{Claim}} $(\star)$}

\theoremstyle{plain}
\newtheorem*{theo}{Main Theorem}
\newtheorem*{thm*}{Theorem}

% MATH objects ---------------------------------------------------
\newcommand{\VF}{\textup{VF}}
\newcommand{\ACVF}{\textup{ACVF}}
\newcommand{\RV}{\textup{RV}}
\newcommand{\DC}{\textup{DC}}
\newcommand{\MM}{\mdl M}
\newcommand{\LL}{\mathcal{L}}
\newcommand{\OO}{\mdl O}
\newcommand{\II}{\mathcal{I}}
\newcommand{\CC}{\mathcal{C}}
\newcommand{\UU}{\mathfrak{U}}
% MATH operators---------------------------------------------------
\DeclareMathOperator{\rk}{rk}
\DeclareMathOperator{\otp}{otp}
 \DeclareMathOperator{\ann}{Ann}
 \DeclareMathOperator{\ran}{ran}
 \DeclareMathOperator{\dom}{dom}
 \DeclareMathOperator{\cod}{cod}
 \DeclareMathOperator{\id}{id}
 \DeclareMathOperator{\ad}{AD}
 \DeclareMathOperator{\ed}{ED}
 \DeclareMathOperator{\cd}{CD}
 \DeclareMathOperator{\Th}{Th}
 \DeclareMathOperator{\tp}{tp}
 \DeclareMathOperator{\Hom}{Hom}
 \DeclareMathOperator{\im}{im}
 \DeclareMathOperator{\lh}{lh}
 \DeclareMathOperator{\lcm}{lcm}
 \DeclareMathOperator{\spec}{Spec}
 \DeclareMathOperator{\ck}{coker}
 \DeclareMathOperator{\aut}{Aut}
 \DeclareMathOperator{\ext}{Ext}
 \DeclareMathOperator{\maxlim}{max\,lim}
 \DeclareMathOperator{\td}{tr\,deg}
 \DeclareMathOperator{\tc}{tr\,cl}
 \DeclareMathOperator{\ord}{Ord}
 \DeclareMathOperator{\cha}{char}
 \DeclareMathOperator{\ac}{\overline{ac}}
 \DeclareMathOperator{\supp}{supp}
 \DeclareMathOperator{\quot}{Quot}
 \DeclareMathOperator{\acl}{acl}
 \DeclareMathOperator{\dcl}{dcl}
 \DeclareMathOperator{\pr}{pr}
 \DeclareMathOperator{\alg}{ac}
 \DeclareMathOperator{\avg}{avg}
 \DeclareMathOperator{\tr}{tr}
 \DeclareMathOperator{\mgl}{GL}
\DeclareMathOperator{\msl}{SL}
\DeclareMathOperator{\jcb}{Jcb}
\DeclareMathOperator{\K}{\mathds{k}}
\DeclareMathOperator{\dvn}{div}
\DeclareMathOperator{\Tr}{Tr}
\DeclareMathOperator{\Jac}{Jac}
\DeclareMathOperator{\GL}{GL}
\DeclareMathOperator{\res}{res}  % map into the residue field

%---integrals-----------------------------------------------------------
\def\Xint#1{\mathchoice
{\XXint\displaystyle\textstyle{#1}}%
{\XXint\textstyle\scriptstyle{#1}}%
{\XXint\scriptstyle\scriptscriptstyle{#1}}%
{\XXint\scriptscriptstyle\scriptscriptstyle{#1}}%
\!\int}
\def\XXint#1#2#3{{\setbox0=\hbox{$#1{#2#3}{\int}$}
\vcenter{\hbox{$#2#3$}}\kern-.5\wd0}}
\def\ddashint{\Xint=}
\def\dashint{\Xint-}

% shorthand for fonts and compounds-------------------------------------

\newcommand{\Z}{\mathds{Z}}
\newcommand{\A}{\mathds{A}}
\newcommand{\C}{\mathds{C}}
\newcommand{\D}{\mathds{D}}
\newcommand{\F}{\mathds{F}}
\newcommand{\G}{\mathds{G}}
\newcommand{\I}{\mathds{I}}
\newcommand{\KK}{\mathds{K}}
\newcommand{\Q}{\mathds{Q}}
\newcommand{\N}{\mathds{N}}
\newcommand{\M}{\mathds{M}}
\newcommand{\R}{\mathds{R}}
\newcommand{\dS}{\mathds{S}}
\newcommand{\dO}{\mathds{O}}
\newcommand{\PP}{\mathds{P}}
\newcommand{\p}{$p$\nobreakdash}
\newcommand{\omin}{$o$\nobreakdash}
\newcommand{\bmin}{$b$\nobreakdash}
\newcommand{\cmin}{$C$\nobreakdash}
\newcommand{\vmin}{$V$\nobreakdash}
\newcommand{\T}{$T$\nobreakdash}
%\newcommand{\M}{$M$\nobreakdash}
\newcommand{\NN}{$N$\nobreakdash}
\newcommand{\aldash}{$\alpha$\nobreakdash}
\newcommand{\bedash}{$\beta$\nobreakdash}
\newcommand{\omdash}{$\omega$\nobreakdash}
\newcommand{\gadash}{$\gamma$\nobreakdash}
\newcommand{\lexp}[2]{{\vphantom{#2}}^{#1}{#2}}
\newcommand{\dand}{\quad \text{and} \quad}
\newcommand{\tand}{~\text{and}~}

%----fonts---------------------------------------------------------------

\newcommand{\gP}{\mathfrak{P}}
\newcommand{\gR}{\mathfrak{R}}
\newcommand{\gA}{\mathfrak{A}}
\newcommand{\gB}{\mathfrak{B}}
\newcommand{\gC}{\mathfrak{C}}
\newcommand{\gD}{\mathfrak{D}}
\newcommand{\gE}{\mathfrak{E}}
\newcommand{\gG}{\mathfrak{G}}
\newcommand{\gH}{\mathfrak{H}}
\newcommand{\gI}{\mathfrak{I}}
\newcommand{\gJ}{\mathfrak{J}}
\newcommand{\gM}{\mathfrak{M}}
\newcommand{\gN}{\mathfrak{N}}
\newcommand{\gO}{\mathfrak{O}}
\newcommand{\gZ}{\mathfrak{Z}}
\newcommand{\ga}{\mathfrak{a}}
\newcommand{\gb}{\mathfrak{b}}
\newcommand{\gc}{\mathfrak{c}}
\newcommand{\gd}{\mathfrak{d}}
\newcommand{\gee}{\mathfrak{e}}
\newcommand{\gf}{\mathfrak{f}}
\newcommand{\frakg}{\mathfrak{g}}
\newcommand{\gh}{\mathfrak{h}}
\newcommand{\gi}{\mathfrak{i}}
\newcommand{\gj}{\mathfrak{j}}
\newcommand{\gk}{\mathfrak{k}}
\newcommand{\gl}{\mathfrak{l}}
\newcommand{\gm}{\mathfrak{m}}
\newcommand{\gn}{\mathfrak{n}}
\newcommand{\go}{\mathfrak{o}}
\newcommand{\gp}{\mathfrak{p}}
\newcommand{\gq}{\mathfrak{q}}
\newcommand{\gr}{\mathfrak{r}}
\newcommand{\gs}{\mathfrak{s}}
\newcommand{\gt}{\mathfrak{t}}
\newcommand{\gu}{\mathfrak{u}}
\newcommand{\gv}{\mathfrak{v}}
\newcommand{\gw}{\mathfrak{w}}
\newcommand{\gx}{\mathfrak{x}}
\newcommand{\gy}{\mathfrak{y}}
\newcommand{\gz}{\mathfrak{z}}

\newcommand{\0}{\emptyset}

%-------------------------------

\DeclareMathAlphabet{\mathpzc}{OT1}{pzc}{m}{it}

\newcommand{\disa}{\mathpzc a}
\newcommand{\disb}{\mathpzc b}
\newcommand{\disc}{\mathpzc c}
\newcommand{\disd}{\mathpzc d}
\newcommand{\dise}{\mathpzc e}
\newcommand{\dish}{\mathpzc h}
\newcommand{\diso}{\mathpzc o}
\newcommand{\disr}{\mathpzc r}
\newcommand{\diss}{\mathpzc s}
\newcommand{\pM}{\mathpzc M}
\newcommand{\pF}{\mathpzc F}
\newcommand{\pL}{\mathpzc L}

% double bars etc.-----------------------
\newcommand{\dbra}[1]{% \llrrbracket{..}
  \,[\mkern-6.8mu[\, #1 \,]\mkern-6.8mu]}
\newcommand{\dpar}[1]{% \llrrparen{..}
   (\mkern-4mu (#1 )\mkern-4mu )}

% double bars etc.-----------------------
 \DeclarePairedDelimiter\abs{\lvert}{\rvert}
%\newcommand{\abs}[1]{\left\vert#1\right\vert}
 \newcommand{\ceil}[1]{\left\lceil#1\right\rceil}
% \newcommand{\set}[1]{\left\{#1\right\}}
\newcommand{\seq}[1]{\left<#1\right>}
 \DeclarePairedDelimiter\norm{\lVert}{\rVert}
 \newcommand{\wt}[1]{\widetilde{#1}}
 \newcommand{\wh}[1]{\widehat{#1}}
 \newcommand{\usub}[2]{#1_{\textup{#2}}}
 \newcommand{\dhat}[1]{\Hat{\Hat{#1}}}
 \newcommand{\dlbr}{( \! (}
 \newcommand{\drbr}{) \! )}
 \newcommand{\btd}{\bigtriangledown}
 \newcommand{\btu}{\bigtriangleup}
 \newcommand{\lan}[3]{\mathcal{L}_{#1 \textup{#2} #3}}

% special fonts for theory names, functions, axioms, etc.
\newcommand{\mdl}[1]{\mathcal{#1}}  % model; e.g. M
\newcommand{\bo}[1]{\mathbf{#1}}    % boldface;
\newcommand{\bb}[1]{\mathbb{#1}}
\newcommand{\comp}{\thicksim}



% logical connectives
\newcommand{\limplies}{\rightarrow}
\newcommand{\lfollows}{\rightarrow}
\newcommand{\liff}{\leftrightarrow}
\newcommand{\ex}[1]{\exists #1 \;} % exists x ...
\newcommand{\fa}[1]{\forall #1 \;} % forall x ...
\newcommand{\lImplies}{\Rightarrow}
\newcommand{\lFollows}{\Leftarrow}
\newcommand{\lIff}{\Longleftrightarrow}

\newcommand{\inot}{\mathord{\sim}}
\newcommand{\lam}[1]{\lambda #1 \;} % lambda x ...

% useful symbols
\newcommand{\proves}{\vdash}
\newcommand{\forces}{\Vdash}
\newcommand{\nforces}{\nVdash}

\newcommand{\wtbo}[1]{\widetilde{\mathbf{#1}}}
\newcommand{\rest}{\upharpoonright}
\newcommand{\lbar}{\vec}
\newcommand{\ulin}{\underline}
\newcommand{\fun}{\longrightarrow}
\newcommand{\efun}{\longmapsto}
\newcommand{\sub}{\subseteq}
\newcommand{\esub}{\preceq}
\newcommand{\mi}{\smallsetminus}
\newcommand{\xfun}{\xrightarrow}
\newcommand{\dbar}[1]{\bar{\bar{#1}}}
\newcommand{\ddlim}[1]{\underset{#1}{\varprojlim} \,\,}
\newcommand{\dclim}[1]{\underset{#1}{\varinjlim} \,\,}
\newcommand{\dlim}{\varprojlim}  %inverse limit or projective limit or limit
\newcommand{\clim}{\varinjlim}  %direct limit or colimit
\newcommand{\dia}{\diamond}
\newcommand{\colim}[1]{\underset{#1}{\text{colim}}\,}  %direct limit or colimit

\newcommand{\dash}{\mathalpha{\mbox{-}}} % e.g. for Sigma^1_1-AC
\newcommand{\goedel}[1]{\ulcorner #1 \urcorner}

\newcommand{\la}{\langle}
\newcommand{\ra}{\rangle}
\newcommand{\half}{\frac{1}{2}}
%\newcommand{\st}{ \; | \; } % x such that ...
%\newcommand{\res}{\upharpoonright}
\newcommand{\concat}{\mathord{\hat{\;}}}

%set-------------------------------------------------
% just to make sure it exists
\providecommand\given{}
% can be useful to refer to this outside \set
\newcommand\SetSymbol[1][]{%
\nonscript \: #1 \vert
\allowbreak
\nonscript\:
\mathopen{}}
\DeclarePairedDelimiterX\set[1]\{\}{%
\renewcommand\given{\SetSymbol[\delimsize]}
#1
}
%-------------------------------------------------------------------
%+++ The following code is (essentially) from Sam Buss +++
% It computes the height of the ulcorner and urcorner correctly.

\newbox\gnBoxA
\newdimen\gnCornerHgt
\setbox\gnBoxA=\hbox{$\ulcorner$}
\global\gnCornerHgt=\ht\gnBoxA
\newdimen\gnArgHgt

\def\code #1{%
        \setbox\gnBoxA=\hbox{$#1$}%
        \gnArgHgt=\ht\gnBoxA%
        \ifnum \gnArgHgt<\gnCornerHgt
                \gnArgHgt=0pt%
        \else
                \advance \gnArgHgt by -\gnCornerHgt%
        \fi
        \raise\gnArgHgt\hbox{$\ulcorner$} \box\gnBoxA %
                \raise\gnArgHgt\hbox{$\urcorner$}}
%+++ End{Sam Buss} +++

%-------------------------------------------------------------------
\newcommand{\mVF}{{\mu}{\VF}}
%\DeclareMathOperator{\vVF}{_1 \negthickspace \VF}
\newcommand{\mgVF}{{\mu} {\VF}}
\newcommand{\mRV}{{\mu}{\RV}}
\newcommand{\mgRV}{{\mu}{\RV}}
\newcommand{\mG}{{\mu}{\Gamma}}
\newcommand{\RES}{\textup{RES}}
\newcommand{\mgRES}{{\mu}{\RES}}
\newcommand{\mRES}{{\mu}{\RES}}
\newcommand{\isp}{\textup{I}_{\textup{sp}}}
\newcommand{\dsp}{\textup{D}_{\textup{sp}}}
\newcommand{\misp}{\mu \isp}
\newcommand{\ispt}{\isp^{\times}}
\newcommand{\mgisp}{\mu \isp}
\newcommand{\mispdb}{\mu \dsp}
\newcommand{\mgispdb}{\mu_{\Gamma} \dsp}
\newcommand{\mgD}{{\mu}{\bb D}}
\newcommand{\mgE}{{\mu}{\bb E}}
\newcommand{\mgL}{{\mu}{\bb L}}
\newcommand{\mL}{{\mu}{\bb L}}
\newcommand{\vgRES}{{\vol_{\Gamma}}{\RES}}
\newcommand{\maRV}[1]{{\mu_{#1}}{\RV}}
\newcommand{\maG}[1]{{\mu_{#1}}{\Gamma}}
\newcommand{\maRES}[1]{{\vol_{#1}}{\RES}}

%-------------------------------------------------------------------

\DeclareMathOperator{\rv}{rv}
\DeclareMathOperator{\sn}{sn}
\DeclareMathOperator{\stan}{st}
\DeclareMathOperator{\csn}{csn}
\DeclareMathOperator{\rcsn}{\overline {csn}}
\DeclareMathOperator{\vv}{val}
\DeclareMathOperator{\orb}{orb}
\DeclareMathOperator{\vs}{vs}
\DeclareMathOperator{\mor}{Mor}
\DeclareMathOperator{\gsk}{\mathbf{K}_+}
\DeclareMathOperator{\ggk}{\mathbf{K}}
\DeclareMathOperator{\sgsk}{!\mathbf{K}_+}
\DeclareMathOperator{\sggk}{!\mathbf{K}}
\DeclareMathOperator{\ob}{Ob}
\DeclareMathOperator{\fn}{FN}
\DeclareMathOperator{\fib}{fib}
\DeclareMathOperator{\fin}{fin}
\newcommand{\vol}{\textup{vol}}
\DeclareMathOperator{\mR}{\mathsf{R}}
\DeclareMathOperator{\mC}{\mathsf{C}}
\DeclareMathOperator{\hen}{HEN}
\DeclareMathOperator{\rad}{rad}
\DeclareMathOperator{\vcr}{vcr}
\DeclareMathOperator{\vrv}{vrv}
\DeclareMathOperator{\RVH}{RVH}
\DeclareMathOperator{\MON}{MON}
\DeclareMathOperator{\can}{\mathbf{c}}
\DeclareMathOperator{\wgt}{\bm w}
\DeclareMathOperator{\pvf}{pr_{\VF}}
\DeclareMathOperator{\prv}{pr_{\RV}}
\DeclareMathOperator{\der}{d}
\DeclareMathOperator{\mds}{mod}
\DeclareMathOperator{\gal}{Gal}
\DeclareMathOperator{\ito}{int}
\DeclareMathOperator{\cl}{cl}
\DeclareMathOperator{\sym}{Sym}
\newcommand{\RCVF}{\textup{RCVF}}
\newcommand{\RCF}{\textup{RCF}}
\DeclareMathOperator{\db}{d \! b}
\DeclareMathOperator{\bdd}{b \! d}
\newcommand{\var}{\textup{Var}}
\DeclareMathOperator{\TVF}{TVF}
\DeclareMathOperator{\TRV}{TRV}
\DeclareMathOperator{\TRES}{TRES}
\DeclareMathOperator{\TCVF}{TCVF}
\DeclareMathOperator{\RVV}{RV_\infty}
%local-------------------------------------------------------------------
\newcommand{\ol}[1]{\overline{#1}}
\DeclareMathOperator{\lant}{\lan{T}{}{}}
\DeclareMathOperator{\dist}{dist}
\newcommand{\LT}{$\lan{T}{}{}$\nobreakdash}
\DeclareMathOperator{\len}{length}
\DeclareMathOperator{\gra}{Gr}
\newcommand{\pard}[1]{\tfrac{\partial}{\partial x_{#1}}}
\newcommand{\pardy}[1]{\tfrac{\partial}{\partial y_{#1}}}
\newcommand{\ddx}{\tfrac{d}{d x}}
\DeclareMathOperator{\bfk}{\bf k}
\DeclareMathOperator{\BSA}{BSA_{\R}}
\DeclareMathOperator{\puR}{\tilde \R}
\DeclareMathOperator{\puC}{\tilde \C}
\DeclareMathOperator{\abv}{\abs{\vv}}
\newcommand{\ticvf}[1]{\text{T}_{#1}\text{CVF}}
\DeclareMathOperator{\ssin}{\Yright}
\newcommand{\mzG}{{\mu_0}{\Gamma}}
\newcommand{\mzRES}{{\mu_0}{\RES}}
\newcommand{\mzRV}{{\mu_0}{\RV}}
\newcommand{\loc}{\textup{loc}}
\DeclareMathOperator{\tor}{Tor}
\DeclareMathOperator{\dsh}{\sharp\sharp}
\DeclareMathOperator{\Vol}{Vol}
\newcommand{\Var}{\textup{Var}}
\DeclareMathOperator{\Def}{Def}
\DeclareMathOperator{\sgn}{sgn}
\DeclareMathOperator{\tbk}{tbk}
\newcommand{\gdv}{\ggk^{\hat \delta}\var_{\R}}
\newcommand{\gmv}{\ggk^{\hat \mu}\var_{\C}}
\newcommand{\kuq}{{\K^{\times}} \cup \Q}
\newcommand{\ccirc}{\circ \circ}
%local-------------------------------------------------------------------




\DeclareMathOperator{\Cl}{Cl}


\begin{document}

\title[Bounded integral and motivic Milnor fiber]{Bounded integral and motivic Milnor fiber}

\author[Arthur Forey]{Arthur Forey}
\address[A. Forey]{D-MATH, ETH Z\"urich, R\"amistrasse 101, CH-8092 Z\"urich, Switzerland}
\email{arthur.forey@math.ethz.ch}

\author[Yimu Yin]{Yimu Yin}
\address[Y. Yin]{Pasadena, California}
\email{yimu.yin@hotmail.com}

\thanks{Many thanks to Ehud Hrushovski and Fran\c cois Loeser for encouragements and enlightening discussions during the preparation of this manuscript. We also thanks Johannes Nicaise for suggesting to work on the Fubini theorem. Thanks also to Raf Cluckers and Martin Hils for inspiring discussions and comments. A.F.\ is supported by a DFG-SNF lead agency program grant
  (grant 200020L\_175755). }

%\keywords{Motivic integration}
%\subjclass[2000]{}



\begin{abstract}
We construct a new Hrushovski-Kazhdan style motivic integral, the so-call bounded integral, that interpolates the integral with volume forms and that without volume forms. This is done within the framework of model theory of algebraically closed valued fields of equicharacteristic zero. As the main application, we rectify and refine some results of Hrushovski and Loeser on how to construct motivic Milnor fiber without resolution of singularities.
\end{abstract}



\maketitle

% BEGIN REFEREE INSERTION OVERVIEW
{\normalfont\small\color{RefereePurple}
\noindent\textbf{Technical referee report: audited second version.}
Every one of the first report's 161 remarks has been rechecked against the manuscript, its conventions, and the relevant cited results. The text has also been reread for omissions. Of the original findings, 103 are confirmed in substance, 52 are revised, and 6 are withdrawn. Twelve new remarks bring the active total to 167. Revisions include corrections to the report, narrower classifications, and fuller proofs of proposed repairs; they do not all indicate that the original objection was false.
\par\smallskip
The original numbers are retained. Revised findings are marked \textbf{revised}; new findings are \textbf{R162--R173} and marked \textbf{new}. Withdrawn remarks R017, R052, R096, R097, R113, and R129 are omitted from the text. R017 misread the stated scope of Section~5; R052, R096, R113, and R129 overlooked the explicit singleton/fiber conventions; R097 added no independent issue beyond the already stated parameter convention and the descent objection R088. Unwarranted parts of R002, R005, R118, and R155 are also corrected in place. A separate audit log records the disposition and reason for all 161 original remarks.
\par\smallskip
Purple remarks distinguish \textbf{Error} (a false assertion, invalid formula, or counterexample), \textbf{Gap} (an argument not established as written), \textbf{Clarification} (including optional exposition), and \textbf{Typo}. Each gives its original source line. The authors' text and Yimu's blue notes are preserved.
\par\smallskip
\noindent\textbf{Recommendation: major revision on technical grounds.}
The principal counterexamples concerning Jacobians, continuity, residue-field descent, and summability survive the audit. Newly identified points include an RV-hull leaving the stated finite-fiber category and the missing compatibility of truncation cutoffs in the coefficient comparison. Several earlier proof-gap remarks now include explicit repairs. The scope and dependency checks attached to each proposed repair remain part of the requested revision.
\par}
% END REFEREE INSERTION OVERVIEW

% BEGIN REFEREE INSERTION R001
\RefereeRemark{001}{Typo}{Acknowledgements and abstract}{%
\textit{Original source line 507.}\quad In the acknowledgements, replace ``We also thanks'' by ``We also thank'' and preferably ``for encouragements'' by ``for their encouragement''. In the abstract, ``so-call'' should be ``so-called'', and ``construct motivic Milnor fiber'' should read ``construct the motivic Milnor fiber''. These are corrections to source lines 493 and 502.
}
% END REFEREE INSERTION R001


% BEGIN REFEREE INSERTION CONTENTS PAGE BREAK
\clearpage
% END REFEREE INSERTION CONTENTS PAGE BREAK
\tableofcontents


%\section*{Notes}
%Many things in the notes are cut-pasted from other sources, so there are  some notational idiosyncrasies, all recorded in \S~\ref{sec:notation}, though.
%
%\S~\ref{sec:HK:main} is of course largely redundant (we want to establish an analog of \cite{hru:loe:lef} for the real algebraic case, the excessive details are intended for that paper), but just in case there are differences as to what the categories are.
%
%The main issue is the failure of (8.1.4) and (8.1.7) in \cite{hru:loe:lef}. So the homomorphism $h_m$ needs to be constructed in a different way. This is in \S~\ref{sec:spec:pui}.
%
%To make \S~7 work, however, we actually need a variation of \cite[Theorem~10.11]{hrushovski:kazhdan:integration:vf}. This is \S~\ref{sec:retract}.
%
%I also found issues concerning (8.5.4) and (8.5.5) in \cite{hru:loe:lef}. Firstly $\Upsilon$ does not vanish on $\isp$ and hence the expression in \cite[Corollary~8.5.3]{hru:loe:lef} is not well-defined. We can force it by modding out $[\A]$, but this is detrimental to our purpose. Secondly the integrand in \cite[Corollary~8.5.3]{hru:loe:lef} does not a priori produce a class in the 2bdd $\RV$-category, hence \cite[Proposition~8.5.1]{hru:loe:lef} cannot be applied directly. The modified definition of $\Upsilon$ now appears towards the end of \S~\ref{sec:retract}.  To solve the second issue, a new integral for the $\VF$-category of invariant objects and covariant morphisms is constructed. This is \S~\ref{sec:inv:cov}.
%
%The work in \S~\ref{sec:inv:cov} is a variation on \cite[Proposition~3.2.2]{hru:loe:lef}. The theorem there can be established, but the length of the construction will  not be desirable. Maybe there are other routes that are more direct for our purpose here? See the last few remarks in \S~\ref{sec:milnor}.

\pagebreak

\section{Introduction}
Recent years have seen significant development in applying the Hrushovski-Kazhdan integration theory to the  study of a variety of topics in geometry and topology, especially those around Denef-Loeser's motivic Milnor fiber. Such an approach  was first articulated in  \cite{hru:loe:lef} for the purpose of  finding a more conceptual (read ``resolution-free'') construction of the complex  motivic Milnor fiber, among other things; unfortunately, it contains a substantial technical oversight. The main goal of this paper is to present a new Hrushovski-Kazhdan style integral,  the so-call \memph{bounded integral}, that interpolates (or refines) the integral with volume forms and that without volume forms, and thereby fully rectify the results in  \cite{hru:loe:lef}.
The overhaul taken up here constitutes a comprehensive expansion of the theoretical framework of the Hrushovski-Kazhdan style integration that should prove amenable to new applications in the future.

% BEGIN REFEREE INSERTION R002
\RefereeRemark{002}{Typo; revised}{Wording}{%
\textit{Original source line 529.}\quad Replace ``so-call'' by ``so-called''. The first report's additional request to replace ``rectify'' by ``rectifies'' is withdrawn: ``to present ... and thereby rectify ...'' is a valid coordination of infinitives. Recasting the long sentence may improve readability, but no correction of that verb is required.
}
% END REFEREE INSERTION R002


We now explain our construction in more detail.
%The Grothendieck semigroup $\gsk \mdl C$ of a category $\mdl C$ endowed with a binary operation ``$\mi$'' and a binary relation ``$\subseteq$'' (subobject) is the free semigroup generated by the isomorphism classes of $\mdl C$, subject to the usual scissor relation
%$[A \mi B] + [B] = [A]$, when $B\subseteq A$,
%where $[A]$, $[B]$ denote the isomorphism classes of the objects $A$, $B$, and ``$\mi$'' and ``$\subseteq$'' are compatible with isomorphisms between objects, usually they are just set subtraction and set inclusion. Sometimes $\mdl C$ is also equipped with a binary operation (for example, cartesian product) that induces multiplication in $\gsk \mdl C$, in which case $\gsk \mdl C$ becomes a (commutative) semiring. The formal groupification of  $\gsk \mdl C$, which is then a ring, is denoted by $\ggk \mdl C$. If a group $G$ acts on the objects of $\mdl C$ and the morphisms of $\mdl C$ are $G$-equivariant, that is, they commute with $G$-actions, then the corresponding $G$-equivariant Grothendieck ring is denoted by $\ggk^G \mdl C$. If $G = \lim_n G_n$ is profinite then we shall always impose the condition that a $G$-action factor through some $G_n$-action.
In this introduction, for concreteness, we work in  the field $\puC = \bigcup_{m \in \Z^+} \C \dpar{ t^{1 / m} }$ of complex Puiseux series. This field is the algebraic closure of the field $\C \dpar t$ of complex Laurent series. A typical element  takes the form $x = \sum_{n \in \Z} a_n t^{n/m}$ for some $m \in \Z^+$ such that its support $\supp(x) = \set{n/m \in \Q \given a_n \neq 0 }$ is well-ordered, in other words, there is a $q \in \Q$ such that $a_n = 0$ for all $n/m < q$. We think of $\bfk \coloneqq \C$ as a subfield of $\puC$ via the embedding $a \efun at^0$. The map $\vv: \tilde \C^\times \fun \Q$ given by $x \efun \min\supp(x)$ is indeed a valuation, and its valuation ring $\OO$ consists of those series $x$ with $\min\supp(x) \geq 0$ and its maximal ideal $\MM$ of those series $x$ with $\min\supp(x) > 0$. Its residue field $\K$ admits a section onto $\bfk$ and hence is isomorphic to $\C$. For a series $x = \sum_{n \in \Z} a_n t^{n/m} \in \puC$ with $\vv(x) = p/m$, let $\rv(x) = a_p t^{p/m}$, which is called the \memph{leading term} of $x$.

% BEGIN REFEREE INSERTION R003
\RefereeRemark{003}{Clarification}{Zero series convention}{%
\textit{Original source line 536.}\quad Set $\vv(0)=\infty$ explicitly, or add $0$ separately to the displayed descriptions of $\OO$ and $\MM$: $\min\supp(0)$ is undefined unless this convention is imposed. The leading-term formula here applies to nonzero series.
}
% END REFEREE INSERTION R003

It is well-known that $(\puC, \OO)$ is an algebraically closed valued field of equicharacteristic zero. Thus, we may  think of it as an $\lan{}{RV}{}$-model of the first-order theory $\ACVF(0,0)$ of such fields, with the parameter space (the ``ground field'') $\bb S = \C \dpar t$.  Here  the two sorts $\VF$, $\RV$ of the language $\lan{}{RV}{}$ are  interpreted, respectively, as $\puC$, $\tilde \C^\times / 1 + \MM$ (or, equivalently, the set of leading terms) and the cross-sort function $\rv : \VF^\times \fun \RV$ as the quotient map (or the leading term map described above). The obvious epimorphism from $\RV$ onto the value group $\Gamma = \Q$, also referred to as the $\Gamma$-sort, with the kernel $\K^\times$, is denoted by $\vrv$. See \cite[\S~2]{Yin:special:trans} for the precise definitions. The Galois group $\gal(\puC / \C \dpar t)$, which  is also the group of automorphisms over $\C \dpar t$ of $\puC$ as a model of $\ACVF(0,0)$, may be identified with the procyclic group $\hat \mu$ of roots of unity (the limit of the inverse system of groups $\mu_n$ of $n$th roots of unity).

% BEGIN REFEREE INSERTION R004
\RefereeRemark{004}{Typo}{Quotient parentheses}{%
\textit{Original source line 538.}\quad Write $\tilde\C^{\times}/(1+\MM)$: without parentheses, the displayed quotient is ambiguous.
}
% END REFEREE INSERTION R004



The category $\VF_*$ essentially consists of the definable subsets of $\VF^n$, $n \geq 0$, as objects and the definable bijections between them as morphisms. The category $\RV[k]$ essentially consists of the definable subsets of $\RV^k$ as objects and the definable bijections between them as morphisms. The category  $\RV[*]$  is the coproduct of $\RV[k]$, $k \geq 0$, and hence is equipped with a gradation by ambient dimensions. The actual definitions will be recalled in \S~\ref{sec:HK:main}.

% BEGIN REFEREE INSERTION R162
\RefereeRemark{162}{Clarification; new}{The introductory description excludes finite covers}{%
\textit{Original source line 542.}\quad The phrase ``essentially consists of the definable subsets of $\RV^k$'' is too restrictive, even up to definable bijection. For example,
\[
U=\{(x,y)\in(\K^\times)^2:y^2=x^3-x\}
\]
is an affine open of a genus-one curve, and its projection to $x$ makes it an object of $\RES[1]\subseteq\RV[1]$. It is not definably bijective to a subset of $\RV$: orthogonality forces such an image to have finitely many value-group fibers; after naming lifts of their values, those fibers are constructible subsets of $\K^\times$. In characteristic zero, a definable bijection would then give a birational parametrization of a dense open of the genus-one curve by a subset of the affine line, which is impossible. Use the description in Definition~\ref{defn:c:RV:cat}: an object is a pair $(U,f)$, with $U$ in an arbitrary finite RV power and $f:U\to\RV^k$ finite-to-one. The grading is the chosen target rank $k$, not the ambient dimension of $U$.
}
% END REFEREE INSERTION R162

To such categories one associates Grothendieck rings $\ggk\VF_*$ and $\ggk\RV[*]$, roughly constructed as the free group generated by isomorphism classes of objects, subject to the usual scissor relation $[A \mi B] + [B] = [A]$ for any objects $A \sub B$, and the product being given by fiber product. See  Definition~\ref{def:Kgroup} for details.

% BEGIN REFEREE INSERTION R005
\RefereeRemark{005}{Error; revised}{Scissor relation}{%
\textit{Original source line 544.}\quad The inclusion is reversed: the relation $[A\mi B]+[B]=[A]$ requires $B\subseteq A$. With the displayed $A\subseteq B$, taking $A=\varnothing$ would force every $[B]$ to be zero. The multiplication intended here is Cartesian product. The first report's parenthetical description using a terminal object is withdrawn: the categories here admit only bijections, so that categorical description was inappropriate.
}
% END REFEREE INSERTION R005

Pulling back from $\RV$ to $\VF$ along the map $\rv: \VF^\times \fun \RV$ yields a ring homomorphism between the   Grothendieck rings
\begin{equation}\label{lifting:hom}
\bb L : \ggk\RV[*] \fun \ggk\VF_*
\end{equation}
One of the  main results of \cite{hrushovski:kazhdan:integration:vf} states that $\bb L$ is surjective with the kernel $(\bm P-1)$, where $\bm P$ stands for the element $[\{1\}] - [\rv(\MM \mi 0)]$  in $\ggk \RV[1]$ (so the principal ideal $(\bm P - 1)$ is not homogenous), and hence, inverting, we get a canonical isomorphism

% BEGIN REFEREE INSERTION R006
\RefereeRemark{006}{Typo}{Spelling}{%
\textit{Original source line 550.}\quad Replace ``homogenous'' by ``homogeneous'' here and at its subsequent occurrences in this introduction.
}
% END REFEREE INSERTION R006
\[
\int : \ggk\VF_* \fun \ggk\RV[*]/(\bm P-1).
\]

The structure of $\ggk \RV[*]$ can be significantly elucidated. To wit, it is isomorphic to a tensor product of two other Grothendieck rings $\ggk \RES[*]$ and $\ggk \Gamma[*]$, where  $\RES[*]$ is  the category of twisted constructible sets in the residue field $\K$  and $\Gamma[*]$ is the category of definable sets in the value group $\Gamma$ (as an \omin-minimal group), both are graded by ambient dimensions. The objects of $\RES[*]$ are twisted because the short exact sequence
\[
1 \fun \K^\times \fun \RV \fun \Gamma \fun 0
\]
does not admit a natural splitting, and $\ggk \Gamma[*]$ is not the Grothendieck ring of \omin-minimal groups because not all definable bijections are admitted as morphisms. Anyway, we have two retractions from $\ggk \RV[*]$ onto a quotient $\sggk \RES$ of $\ggk \RES$ (the gradation is forgotten) that vanish on $(\bm P-1)$, reflecting the fact that there are two Euler characteristics in the $\Gamma$-sort; these are denoted by $\bb E_{b}$, $\bb E_{g}$.

% BEGIN REFEREE INSERTION R007
\RefereeRemark{007}{Error}{Different targets for the two retractions}{%
\textit{Original source line 559.}\quad The two targets should be distinguished already here. The bounded Euler-characteristic construction $\bb E_b$ lands in $\sggk\RES$, but the geometric one $\bb E_g$ requires affine-line localization: its dimension-free image lies in $\sggk\RES[[\A]^{-1}]$. For example, $\bb E_g([1])=[\A]^{-1}$ after forgetting the grading, as computed in Remark~\ref{eb:no:eg}. Thus both maps cannot be described here as retractions onto the same unlocalized ring.
}
% END REFEREE INSERTION R007
%We shall only need $\bb E_{b}$; see the left column of the diagram (\ref{diag-diamond-interpol}).


The objects of $\VF_*$, $\RV[*]$ may be equipped with $\Gamma$-volume forms, and the resulting categories are denoted by $\mgVF[*]$, $\mgRV[*]$. There is again a  canonical isomorphism
\[
\int^\mu : \ggk\mgVF[*] \fun \ggk\mgRV[*]/(\bm P),
\]
where the principal ideal $(\bm P)$ is now homogenous. The ring $\ggk\mgRV[*]$ can  be expressed as a tensor product of  $\ggk \mgRES[*]$ and $\ggk \mG[*]$, which induces two retractions $\mgE_ {b}$, $\mgE_{g}$ from $\ggk\mgRV[*]$ onto $\sggk \RES[*]$ that vanish on $(\bm P)$; see the right column of the diagram (\ref{diag-diamond-interpol}).

% BEGIN REFEREE INSERTION R008
\RefereeRemark{008}{Error}{Measured retraction target}{%
\textit{Original source line 567.}\quad The target stated here omits both the volume forms and the quotient. Proposition~\ref{prop:eu:retr:k} gives $\sggk\mgRES[*]/([\A])$. Also, the current diagram~\eqref{diag-diamond-interpol} does not display either retraction: its former bottom row has been commented out. Correct the target and remove or update the reference to its ``right column''. 
}
% END REFEREE INSERTION R008

We think of $\int$ as the universal additive invariant (or generalized Euler characteristic) of definable sets and $\int^\mu$ as a motivic integral. They are not directly related, though: in  $\ggk\VF_*$, two objects are identified if there is a definable bijection between them, whereas in $\ggk\mgVF[*]$, two objects are identified if there is a measure-preserving bijection between two dense subobjects.


The main result of this paper is  that the two isomorphisms $\int$, $\int^\mu$ can be interpolated by a third one, the bounded integral (Theorem \ref{main:prop}):
\[
\int^\diamond : \ggk \mgVF^\diamond[*] \fun \ggk  \mgRV^{\db}[*] / (\bm P_\Gamma).
\]
The category $\mgVF^\diamond[*]$ consists of those \memph{proper invariant} objects of $\mgVF[*]$, which roughly means a bounded definable set that is a union of open boxes of valuative radius $\gamma$, $\gamma \in \Gamma$. The category $\mgRV^{\db}[*]$ consists of those \memph{doubly bounded} objects of $\mgRV[*]$, in the sense that its image in the value group is doubly bounded. See \S~\ref{sec:inv:cov} for the precise definitions. Note that $\mgVF^{\diamond}[*]$  is also graded since, as in classical measure theory, gradation by ambient dimensions is a necessity in the presence of volume forms (a curve has no volume if considered as a subset of a surface). The ideal $(\bm P_\Gamma)$ is homogenous but no longer principal; it is generated by the elements $\bm P_\gamma \in \ggk \mgRV^{\db}[1]$, one for each $\gamma \in \Z^+$, defined as follows. Let $\MM_\gamma$ be the open  disc around $0$ with valuative radius $\gamma$ and $t_\gamma \in \vrv^{-1}(\gamma)$ a definable element. Then $\bm P_\gamma$ is the element
\[
[\rv(\MM \mi \MM_{\gamma})] + [\{t_\gamma\}] - [\{1\}]
\]
with the constant volume form $0$, and it does not depend on the choice of $t_\gamma$. We can express $\ggk  \mgRV^{\db}[*]$ as a tensor product of $\ggk \mRES[*]$ and $\ggk \mG^{\db}[*]$. Since the objects of $\mG^{\db}[*]$ are doubly bounded, the two Euler characteristics coincide and consequently there is only one retraction $\bb E^{\diamond}$ from $\ggk  \mgRV^{\db}[*]$ onto $\sggk \RES$ that vanishes on $(\bm P_\Gamma)$.

% BEGIN REFEREE INSERTION R009
\RefereeRemark{009}{Clarification}{Uniqueness claim}{%
\textit{Original source line 580.}\quad The agreement of $\chi_b$ and $\chi_g$ on doubly bounded sets shows that these two constructions yield the same retraction; it does not establish uniqueness among all possible retractions. Replace ``there is only one retraction'' by ``the two constructions give the same retraction''. 
}
% END REFEREE INSERTION R009


The situation is summarized in the  commutative diagram (\ref{diag-diamond-interpol}). More precisely speaking, the bounded integral $\int^\diamond$ is a common refinement of $\int$ and $\int^\mu$: it is finer than $\int$ because it carries volume forms, and it is finer than $\int^\mu$ because it only admits those measure-preserving morphisms that are defined everywhere (as opposed to on dense subobjects).
\begin{figure}[htb]
\begin{equation}\label{diag-diamond-interpol}
\bfig
  \hSquares(0,0)/<-`->`->`->`->`<-`->/<400>[{\ggk \VF_*}`{\ggk \mgVF^\diamond[*]}`{\ggk \mgVF[*]}`{\ggk \RV[*] / (\bm P - 1)}`{\ggk \mgRV^{\db}[*] / (\bm P_\Gamma)}`{\ggk \mgRV[*] / (\bm P)}; ``\int`\int^{\diamond}`\int^\mu``]
\efig
\end{equation}
%\begin{equation}\label{diag-diamond-interpol}
%\bfig
%  \hSquares(0,0)/<-`->`->`->`->`<-`->/<400>[{\ggk \VF_*}`{\ggk \mgVF^\diamond[*]}`{\ggk \mgVF[*]}`{\ggk \RV[*] / (\bm P - 1)}`{\ggk \mgRV^{\db}[*] / (\bm P_\Gamma)}`{\ggk \mgRV[*] / (\bm P)}; ``\int`\int^{\diamond}`\int^\mu``]
%  \hSquares(0,-400)/<-`->`->`->`->`<-`./<400>[{\ggk \RV[*] / (\bm P - 1)}`{\ggk \mgRV^{\db}[*] / (\bm P_\Gamma)}`{\ggk \mgRV[*] / (\bm P)}`{\sggk \RES}`{\sggk \RES[*]}`{\sggk \mgRES[*] / ([\A])}; ``\bb E_b`\mgE^{\db}`\mgE_{b}`\phi`]
%\efig
%\end{equation}
\end{figure}
%with $\bb E^{\diamond} = \phi \circ \mgE^{\db}$,  where $\phi$ is the forgetful homomorphism, and the dotted line means ``no meaningful relation'' insofar as the results in this paper are concerned.



% \bfig
% \hSquares(0,0)/->`->`->`->`->`->`->/<400>[{\ggk \mgVF^\diamond[*]}`{\ggk  \mgRV^{\db}[*] / (\bm P_\Gamma)}`{\sggk \RES}`{\ggk \VF_*}`{\ggk \RV[*] / (\bm P - 1)}`{\sggk \RES}; \int^{\diamond}`\bb E^{\diamond}```\id`\int`\bb E_b]

%
%
%\bfig
%\node a(0,800)[\ggk {\mgVF^\diamond[*]}]
%\node b(1000,800)[\ggk {\mgRV^{\db}[*]} / (\bm P_\Gamma)]
%\node c(0,0)[\ggk \VF_*]
%\node d(1000,0)[\ggk {\RV[*]} / (\bm P - 1)]
%%\node e(2000,0)[{\ggk^{\hat \mu} \var_{\C}[[\A]^{-1}]}]
%\node f(700,400)[\ggk {\mgVF[*]}]
%\node g(1700,400)[\ggk {\mgRV[*]} / (\bm P)]
%%\node h(2700,400)[{\ggk^{\hat \mu} \var_{\C}[[\A]^{-1}][Y]_\dag}]
%\arrow|a|/->/[a`b;{\int^{\diamond}}]
%\arrow|l|/->/[a`c;]
%\arrow|l|/->/[b`d;]
%\arrow|l|/->/[b`g;]
%\arrow|l|/->/[a`f;]
%\arrow|l|/->/[c`d;{\int}]
%\arrow|l|/->/[f`g;{\int^\mu}]
%%\arrow|l|/->/[g`h;Z]
%%\arrow|l|/->/[d`e;\Theta \circ \bb E_b]
%%\arrow|b|/->/[h`e;{-\lim}]
% \efig


The general strategy for the construction of the bounded integral $\int^\diamond$ follows roughly that of the original Hrushovski-Kazhdan integral. One first defines a lifting homomorphism from $\ggk  \mgRV^{\db}[*]$ to $\ggk \mgVF^\diamond[*]$,  shows that it is surjective, then studies its kernel.
On the technical side, in order to stay within the realm of proper invariant sets at each step of the construction, one is led to the notion of a \memph{proper covariant} bijection, first formulated in \cite{hru:loe:lef}. It goes as follows (see Definition~\ref{prop:pseu} for more details). For $\alpha \in \Gamma$, let $\pi_\alpha : \VF\fun \VF/\MM_\alpha$ be the projection map;  if $\alpha\in \Gamma^n$ then $\pi_{\alpha}$ denotes the product of the maps $\pi_{\alpha_i}$. A definable bijection $f$ from $A\sub \VF^n$ onto  $B \sub \VF^m$ is said to be proper covariant if there are $\alpha \in \Gamma^n$, $\beta \in \Gamma^n$, and a function $f_\downarrow : \pi_{\alpha}(A) \fun \pi_{\beta}(B)$ such that $A$ is $\alpha$-invariant, $B$ is $\beta$-invariant, and  $\pi_\beta \circ f = f_\downarrow \circ \pi_\alpha$. In order to show the existence of such proper covariant bijections of a very special kind, we establish in \S~\ref{sec:cont:fib:prop} some general properties of definable continuous functions, reminiscent of those of definable continuous functions in \omin-minimal theories (see, for instance,  \cite[Chapter~6]{dries:1998}).

% BEGIN REFEREE INSERTION R010
\RefereeRemark{010}{Error}{Proper covariance in the introduction}{%
\textit{Original source line 629.}\quad Change $\beta\in\Gamma^n$ to $\beta\in\Gamma^m$. More substantially, include boundedness of $A$ and $B$ to match Definition~\ref{prop:pseu}; the displayed quotient condition alone is satisfied by the identity on the whole valued field, which is not proper covariant under that definition.
}
% END REFEREE INSERTION R010

Next, we turn to the construction of the motivic Milnor fiber attached to a nonconstant polynomial function $f : (\C^d, 0) \fun (\C, 0)$,  where $0$ is assumed to be a singular point, that is, $\nabla f (0) = 0$.

By an (algebraic) variety over a field $\bfk$, we mean a reduced separated $\bfk$-scheme of finite type. We denote by $\Var_{\bfk}$ the category of  varieties over $\bfk$.
Let $\ggk \var_{\C}$ be the Grothendieck ring of complex algebraic varieties and $\ggk^{\hat \mu} \var_{\C}$ the corresponding Grothendieck ring with good $\hat \mu$-actions. Here, in addition to the factorization requirement for the $\hat \mu$-actions, the qualifier ``good'' also means that a $\C$-vector space endowed with a linear $\hat \mu$-action is identified with the same  $\C$-vector space but endowed with the trivial $\hat \mu$-action.

% BEGIN REFEREE INSERTION R011
\RefereeRemark{011}{Error}{Good actions and the affine-action relation}{%
\textit{Original source line 634.}\quad Two different notions are conflated. A good finite-group action means that every orbit is contained in an affine open subset; for $\hat\mu$ one additionally requires factorization through some $\mu_n$. Identifying different linear actions on affine space is a separate defining relation in the modified equivariant Grothendieck ring. State the goodness condition and the affine-action relation separately, including compatibility with products with an arbitrary equivariant variety. See \href{https://arxiv.org/pdf/math/0006050}{Denef--Loeser, Geometry on arc spaces, Section 2.4}.
}
% END REFEREE INSERTION R011

Let $\mathscr L$ be the space of formal arcs on $\C^d$ at $0$. So each element in $\mathscr L$ is of the form $\gamma(t) = (\gamma_1(t), \ldots, \gamma_d(t))$, where $\gamma_i(t)$ is a complex formal power series and $\gamma_i(0) = 0$. For each $m \in \Z^+$, let $\mathscr L_m$ be the space of such arcs modulo $t^{m+1}$ (also referred to as ``truncated arcs''). Consider the following locally closed subset of $\mathscr L_m$:
\begin{equation*}
  \mdl X_{f,m}  = \set{\gamma(t) \in \mathscr L_m \given f(\gamma(t)) = t^m \mod t^{m+1} }.
\end{equation*}
It may be viewed in a natural way as the set of closed points of an algebraic variety over $\C$ and, as such, carries a natural $\mu_{m}$-action. The \memph{motivic zeta function} attached to $f$ is then the generating series whose coefficients are in effect the ``$\hat \mu$-equivariant motivic volumes'' of the sets of truncated arcs above:
\[
 Z_f(T) = \sum_{m \in \Z^+} [\mdl X_{f,m}] [\A]^{-nd}T^m \in \gmv[[\A]^{-1}] \dbra{T}.
\]

% BEGIN REFEREE INSERTION R012
\RefereeRemark{012}{Error}{Zeta coefficient exponent}{%
\textit{Original source line 643.}\quad Replace $[\A]^{-nd}$ by $[\A]^{-md}$: $m$ is the summation and truncation index; $n$ is not defined here. The normalization in \eqref{z:coeff:change} is the correct one.
}
% END REFEREE INSERTION R012
It is  shown in \cite{denefloeser:arc, den:loe:2002} that $Z_f(T)$ is rational and belongs to $\gmv[[\A]^{-1}][T]_\dag$, where the subscript ``$\dag$'' indicates that the ring is localized by  a certain multiplicative family of elements. The \memph{motivic Milnor fiber}
\[
\mathscr S_{f} = - \lim_{T \limplies \infty} Z_{f}(T)
\]
is then extracted from this rational expression via a formal process of sending the variable $T$ to infinity (this process is also summarized in \cite[\S~8.4]{hru:loe:lef}). Of course, to justify calling $\mathscr S_{f}$ a ``Milnor so-and-so'' one needs to show, at the very least, that invariants of the \memph{topological Milnor fiber} $F_{f}$ attached to $f$ can be recovered from it.  This is indeed the case for, say, the Euler characteristic and the Hodge characteristic.

Originally,  the proofs that $Z_f(T)$ is rational and  the Euler (or Hodge) characteristics coincide both rely on resolution of singularities.  More recently, in \cite{hru:loe:lef}, these results are established by way of a more conceptual construction. To begin with, $Z_{f}(T)$ is expressed as
\begin{equation*}\label{intro:zeta:mot}
\sum_{m \in \Z^+} H_m([\mdl X_f]) T^m,
\end{equation*}
where the coefficients $H_m([\mdl X_f])$ are certain integrals of definable sets that take values in the ring $\gmv [[\A]^{-1}]$,  and the so-called \memph{nonarchimedean Milnor fiber} of $f$ is the definable set
\[
\mdl X_f = \set{x \in \MM^d \given \rv(f(x)) = \rv(t) }.
\]
Formulated in this way, the rationality of $Z_{f}(T)$ essentially follows from  certain computation rules of (convergent) geometric series. That the Euler characteristics of $\mathscr S_{f}$ and $F_{f}$ coincide follows from the fact that we can express both the Euler characteristic of each coefficient of $Z_{f}(T)$ and that of $F_{f}$ in terms of traces of the monodromy action on the cohomological groups of $F_{f}$, where the first expression relies on  the resolution-free proofs of the A'Campo-Denef-Loeser formula (this is the main point of \cite{hru:loe:lef}) and quasi-unipotence of local monodromy (see \cite[Remark~8.5.5]{hru:loe:lef}). However, technical difficulties, arising from the fact that none of the morphisms (8.1.4), (8.1.6), and (8.1.7) in \cite{hru:loe:lef} is surjective, contrary to what is claimed in \cite[Proposition~10.10(2)]{hrushovski:kazhdan:integration:vf}, make their construction of the integral $H_m$ incomplete.

% BEGIN REFEREE INSERTION R013
\RefereeRemark{013}{Clarification}{Identify the failed surjectivity explicitly}{%
\textit{Original source line 658.}\quad Since this failure motivates the paper, give an explicit object outside the image of the cited tensor morphism, or a precise cross-reference to such an example with its obstruction explained. Merely listing the incorrect surjectivity claims does not let the reader verify why the replacement is necessary. This request in the human report remains relevant.
}
% END REFEREE INSERTION R013

To remedy this, as an application of our main result, we complete their construction by recasting it in the formalism of the bounded integral $\int^\diamond$. More specifically, recall from \cite[\S~4.3]{hru:loe:lef} the canonical isomorphism
\[
\Theta : \sggk \RES \fun \gmv,
\]
which is essentially the operation of ``twisting back'' via a reduced cross-section (Definition~\ref{twistback}). Let  $\gmv[[\A]^{-1}]_{\loc}$ be the localization of $\gmv[[\A]^{-1}]$ by the multiplicative family generated by the elements $1 - [\A]^{-i}$, $i \in \Z^+$. Let $\ggk^{\hat \mu} \var_{\C}[[\A]^{-1}][T]_{\dag}$ be the localization of $\gmv[[\A]^{-1}][T]$ by the multiplicative family generated by the elements $1 - [\A]^{a}T^b$, where $a \in \Z$ and $b \in \Z^+$; it may be regarded as a subring of  $\ggk^{\hat \mu} \var_{\C}[[\A]^{-1}] \dbra T$. Then we have a  commutative diagram (\ref{diag-diamond-zeta-function}).
\begin{figure}[htb]
\begin{equation}\label{diag-diamond-zeta-function}
\bfig
  \iiixiii(0,400)/<-`->`->`->`->`<-`->`->`->`->``->/<1200,400>[{\ggk \VF_*}`{\ggk^\natural \mgVF^\diamond[*]}`{\ggk^\natural \mgVF[*]}`{\ggk \RV[*] / (\bm P - 1)}`{\ggk^\natural \mgRV^{\db}[*] / (\bm P_\Gamma)}`{\ggk^\natural \mgRV^{\bdd}[*] / (\bm P)}`\gmv`\gmv{[[\A]^{-1}][T]_{\dag}}`\gmv{[[\A]^{-1}]_{\loc}[T]_{\dag}}; ``\int`\int^{\diamond}`\int^\mu```\Theta \circ \bb E_b`Z`Z``]
  \square(0,0)|aamb|/`->`->`=/<1200,400>[\gmv`\gmv{[[\A]^{-1}][T]_{\dag}}`\gmv{[[\A]^{-1}]}`\gmv{[[\A]^{-1}]};
  ``-\lim`]
  \square(1200,0)/``->`->/<1200,400>[\gmv{[[\A]^{-1}][T]_{\dag}}`\gmv{[[\A]^{-1}]_{\loc}[T]_{\dag}}`
  \gmv{[[\A]^{-1}]}`\gmv{[[\A]^{-1}]_{\loc}};``-\lim`]
%  \hSquares(0,400)/<-`->`->`->`->`<-`->/<400>[{\ggk \VF_*}`{\ggk^\natural \mgVF^\diamond[*]}`{\ggk^\natural \mgVF[*]}`{\ggk \RV[*] / (\bm P - 1)}`{\ggk^\natural \mgRV^{\db}[*] / (\bm P_\Gamma)}`{\ggk^\natural \mgRV^{\bdd}[*] / (\bm P)}; ``\int`\int^{\diamond}`\int^\mu``]
%  \hSquares(0,0)/<-`->`->`->`->``->/<400>[{\ggk \RV[*] / (\bm P - 1)}`{\ggk^\natural \mgRV^{\db}[*] / (\bm P_\Gamma)}`{\ggk^\natural \mgRV^{\bdd}[*] / (\bm P)}`\gmv{[[\A]^{-1}]}`\gmv{[[\A]^{-1}][T]_{\dag}}`\gmv{[[\A]^{-1}]_{\loc}[T]_{\dag}}; ``\Theta \circ \bb E_b``Z``]
%  \morphism(0,0)/<-/<2240,0>[\gmv{[[\A]^{-1}]}_{\loc}`\gmv{[[\A]^{-1}]}_{\loc};-\lim]
 \efig
\end{equation}
\end{figure}

% BEGIN REFEREE INSERTION R014
\RefereeRemark{014}{Error}{Domain of the limit arrow}{%
\textit{Original source line 678.}\quad The arrows labeled $-\lim$ are not defined on all of $R[T]_{\dag}$, where $R$ denotes the relevant coefficient ring: the polynomial $T$ belongs to this localization and has no finite value at infinity. Use the subring of rational functions regular at infinity and with zero constant coefficient in their expansion at $T=0$, equipped with Hadamard multiplication; equivalently specify the actual image of $Z$ and prove it has these properties. On this domain $-\lim$ is compatible with Hadamard product. See \href{https://arxiv.org/pdf/1111.1954}{Hrushovski--Loeser, Section 8.4 and Lemma 8.4.1}. The domain issue for $Z$ itself is addressed immediately below.
}
% END REFEREE INSERTION R014
The superscript ``$\natural$'' means that the integral is restricted to the subring of a sort of ``integrable'' sets. The map $Z$ assigns, to each object $[\bm U] \in \ggk^\natural \mgRV^{\db}[*]$,  a power series \yimu{this is not quite right: the power series can indeed be assigned to every $[\bm U] \in \ggk^\natural \mgRV^{\db}[*]$, this is how $\ggk^\natural$ is defined, but to show that it belongs to  $\gmv[[\A]^{-1}][T]_{\dag}$ we need to further restrict to objects of the form $[\bm U] = [(U, f, l)]$ with $l \in \Z$}
\[
Z([\bm U])(T) = \sum_{m \in \Z^+} \bm H_m([\bm U]) T^m \in \ggk^{\hat \mu} \var_{\C}[[\A]^{-1}] \dbra T,
\]
that is, the motivic zeta function of $\bm U$, where $\bm H_m$ is the rectified version of $H_m$. Note that $Z$ is a ring homomorphism with respect to the Hadamard product in the target. It turns out that $Z([\bm U])(T)$ is rational and belongs to $\gmv[[\A]^{-1}][T]_{\dag}$. The formal process $T \limplies \infty$ mentioned above then lands us in $\gmv[[\A]^{-1}]$; this is the map $-\lim$. If we take $[\bm U] / (\bm P_\Gamma) = \int^{\diamond}([\mdl X_f])$ then the end result is the motivic Milnor fiber $\mathscr S_{f}$.

% BEGIN REFEREE INSERTION R015
\RefereeRemark{015}{Gap; revised}{Incorporate the blue restriction into the statement}{%
\textit{Original source line 683.}\quad Yimu's blue note is mathematically substantive and is still contradicted by the following unrestricted rationality claim and diagram. The definition of $\ggk^{\natural}$ in Remark~\ref{dia:to:db:dag} ensures integral exponents for each $m$; it does not itself establish rationality as $m$ varies. The proof of Theorem~\ref{direct:mil} establishes rationality for integer constant volume forms and finite sums of those classes. Restrict the general claim and diagram to the corresponding precisely specified image subring, checking the arrows on that subring, or prove the broader assertion. Also revise the unrestricted bounded-category claim near the end of Section~\ref{sec:spec:pui}. For the theorem about $\mdl X_f$ itself, the needed constant-zero representative is supplied by Lemma~\ref{L:measure:surjective} and Remark~\ref{dia:to:db:dag}, subject to the preparation repairs below; it is not an additional missing hypothesis of that theorem.
}
% END REFEREE INSERTION R015

Since $\mdl X_f$ also lives in $\VF_*$, we can recover   $\mathscr S_{f}$ directly as $(\Theta \circ \bb E_b \circ \int)([\mdl X_f])$ in $\gmv$, without even inverting $[\A]$. This is indeed a salient point. However, for comparison with the Denef-Loeser construction, we have to go down to  $\gmv[[\A]^{-1}]$, since that is where the coefficients of the motivic zeta function live.

% BEGIN REFEREE INSERTION R016
\RefereeRemark{016}{Clarification}{Class before and after localization}{%
\textit{Original source line 685.}\quad The formula defines a canonical class in the unlocalized ring whose image is the Denef--Loeser motivic Milnor fiber. Equality after localization does not uniquely determine an unlocalized class, since localization at $[\A]$ need not be injective (as acknowledged in Remark~\ref{rem-comp-milnor-fib}). Say ``a canonical unlocalized lift of $\mathscr S_f$'' here, or explicitly adopt a separate notation for that lift.
}
% END REFEREE INSERTION R016

The composite  $\Theta \circ \bb E_b \circ \int$ is referred to as the \memph{motivic volume operator}, usually denoted by $\Vol$, and is already used implicitly in \cite{hrushovski:kazhdan:integration:vf}. Since its reemergence in \cite{hru:loe:lef}, it has been featured in a growing list of applications, see \cite{thuong-int-conj, nicaise_ps_tropical, Nic:Pay:trop:fub, forey_virtual_2017, NS-deg-stab-rat}.

The nonarchimedean Milnor fiber $\mdl X_f$ is closely related to the analytic Milnor fiber $\mathcal{F}^{\mathrm{an}}_{f}$ introduced in \cite{nicaise:sebag:motivic:serre}. The $\puC$-rational points of $\mathcal{F}^{\mathrm{an}}_{f}$ may be understood as an object of  $\VF_*$ (but not $\mgVF^\diamond[*]$) and hence, as such, admits a motivic volume $\Vol(\mathcal{F}^{\mathrm{an}}_{f})$. This is indeed the motivic Milnor fiber $\mathscr S_{f}$. This fact is already proven  in \cite{nicaise_ps_tropical}, but without taking the $\hat \mu$-actions into account, and also in \cite{Nic:Pay:trop:fub, forey_virtual_2017}. The arguments there all rely on resolution of singularities.

The bounded integral $\int^\diamond$ is also used in \cite{FY:Mot-Int-Mil-Fib} to obtain similar results on real motivic Milnor fiber.

In \S~\ref{csndag}, as a computational tool, we introduce the notion of a twistoid. There we prove two Fubini theorems that are not related to the rest of the paper. The first is  a version of the  tropical motivic Fubini theorem of \cite{Nic:Pay:trop:fub}. In a  simplified form (but of central interest), the second states the following.

\begin{thm*}
Let $f :A \fun B$, $f : A' \fun B$ be morphisms of varieties over $\C$. If $\Vol_b([f^{-1}(b)])=\Vol_b([f'^{-1}(b)])$ for all $b\in B$  then $\Vol([A])=\Vol([A'])$. Consequently, if $\Vol_b([f^{-1}(b)])= 0$ for all $b\in B$ then $\Vol([A]) = 0$.
\end{thm*}

% BEGIN REFEREE INSERTION R018
\RefereeRemark{018}{Typo}{Second morphism}{%
\textit{Original source line 697.}\quad The second morphism must be $f':A'\longrightarrow B$, not a second $f$.
}
% END REFEREE INSERTION R018

Here the operator $\Vol_b$ is similar to $\Vol$, but over the parameter space $\C\dpar t (b)$. It is essential to consider $b$ in a sufficiently saturated model $\bb U$ of $\ACVF(0,0)$ instead of $\puC$, allowing the hypothesis to be tested at generic points.  A more general version, with $A$, $B$ possibly arbitrary definable sets, appears as Theorem \ref{fubini}, but one needs to be more careful about the definition of $\Vol_b$.

% BEGIN REFEREE INSERTION R019
\RefereeRemark{019}{Gap}{Meaning of the relative volume}{%
\textit{Original source line 699.}\quad The saturated-point requirement is appropriate, but the promised construction of $\Vol_b$ in Remark~\ref{rem:volb} needs correction for residue-field descent; see the detailed remark there. Until that is supplied, state this introductory theorem using the equality in the parameter-relative $\sggk\RES$ ring appearing in Theorem~\ref{fubini}, rather than using an unestablished isomorphism $\Theta_b$.
}
% END REFEREE INSERTION R019

Some property of this form is needed in \cite{loeser2019motivic}.
We expect to find more applications of these results in motivic Donaldson-Thomas theory, along the line of \cite{Nic:Pay:trop:fub}.



The paper is organized as follows. In \S~\ref{sec-tech} we recall and extend certain facts on definable sets in $\ACVF(0,0)$ that are needed later on. In particular,  some new results about continuous definable functions in this environment are obtained, which could be of independent interest, such as fusion of continuous fibers (Proposition~\ref{prop:continuity:fiberwise}) and curve selection (Lemma~\ref{curve:select}), which closely resemble their counterparts in the \omin-minimal setting. In the short \S~\ref{sec:HK:main} is introduced notations and results from the original Hrushovski-Kazhdan paper \cite{hrushovski:kazhdan:integration:vf}. In \S~\ref{sec:retract} we show how to replace classes of polytopes with numbers  and thereby construct a retraction map from the Grothendieck ring in $\RV$ to that of varieties. Invariance and covariance are defined and investigated  in \S~\ref{sec-pscb}, where the surjectivity of the lifting homomorphism (\ref{lifting:hom}) is also established. Then in \S~\ref{section:int} we study its kernel and finish the construction of the bounded integral $\int^\diamond$. Finally, \S~\ref{sec:spec:pui} is devoted to the construction of motivic zeta function and motivic Milnor fiber, rectifying the result in \cite{hru:loe:lef}.

% BEGIN REFEREE INSERTION R020
\RefereeRemark{020}{Typo}{Grammar}{%
\textit{Original source line 706.}\quad Replace ``In the short Section ... is introduced notations and results'' by ``The short Section ... introduces notation and results''. Also write ``the construction of the motivic zeta function and the motivic Milnor fiber''. 
}
% END REFEREE INSERTION R020



\section{Technicalities}
\label{sec-tech}

\subsection{Setting and notation}\label{sec:notation}

\begin{nota}[Coordinate projections]\label{indexing}
For each $n \in \N$, let $[n]$ denote the set $\{1, \ldots, n\}$. For any $E \sub [n]$, we write $\pr_E(A)$, or even $A_E$, when there is no danger of confusion, for the projection of $A$ into the coordinates contained in $E$. It is often more convenient to use simple standard descriptions as subscripts. For example, if $E$ is a singleton $\{i\}$ then we shall always write $E$ as $i$ and $\tilde E \coloneqq [n] \mi E$ as $\tilde i$; similarly, if $E = [i]$, $\set{k \given  i \leq k \leq j}$, $\set{k \given  i < k < j}$, $\{\text{all the coordinates in the sort $S$}\}$, etc., then we may write $\pr_{\leq i}$, $\pr_{[i, j]}$, $A_{(i, j)}$, $A_{S}$, etc.; in particular, we shall frequently write $A_{\VF}$ and $A_{\RV}$ for the projections of $A$ into the $\VF$-sort and $\RV$-sort coordinates.


Unless otherwise specified, by writing $a \in A$ we shall mean that $a$ is a finite tuple of elements (or ``points'') of $A$, whose length, denoted by $\lh(a)$, is not always indicated. If $a = (a_1, \ldots, a_n)$ then, for $1 \leq i < j \leq n$, following the notational scheme above, $a_i$, $a_{\tilde i}$, $a_{\leq i}$, $a_{[i, j]}$, $a_{\VF}$, etc., are shorthand for the corresponding subtuples of $a$.

We shall write $\{t\} \times A$, $\{t\} \cup A$, $A \mi \{t\}$, etc., simply as $t \times A$, $t \cup A$, $A \mi t$, etc., when it is clearly understood that $t$ is an element and hence must be interpreted as a singleton in these expressions.

For $a \in A_{\tilde E}$, the fiber $\{b : ( b, a) \in A \} \sub A_E$ over $a$ is often denoted by $A_a$. Note that, in the discussion below, the distinction between the two sets $A_a$ and $A_a \times \{a\}$ is usually immaterial and hence they may and  shall be tacitly identified. In particular, given a function $f : A \fun B$ and $b \in B$, the pullback $f^{-1}(b)$ is sometimes written as $A_b$ as well. This is a special case since functions are identified with their graphs. This notational scheme is especially useful when the function $f$ has been clearly understood in the context and hence there is no need to spell it out all the time.
\end{nota}

% BEGIN REFEREE INSERTION R021
\ResolvedRefereeRemark{021}{Clarification}{Elements versus tuples; scalar hypotheses}{%
\textit{Original source line 724.}\quad The sentence beginning `by writing $a\in A$' should say that an element of $A$ may itself be a tuple, rather than that it is a tuple of elements of $A$. More importantly, please explicitly override this convention whenever a result requires a single coordinate; the exchange principle and Lemma~\ref{finite:fib:def} below are false for arbitrary tuples.
}{The explicitly stated tuple convention is retained. The author has clarified that literal subsets of $\VF$ or $\RV$ specify one coordinate; the tuple convention does not enlarge that scope. Scalar exchange is explicit in Lemma~2.14, and Lemma~\ref{finite:fib:def} is already unary as written. The suggested rephrasing of the convention is editorial; no further main-text amendment is adopted.}
% END REFEREE INSERTION R021

From here on we work in a sufficiently saturated model $\bb U$ of $\ACVF(0,0)$, together with a parameter  set $\mathbb{S}$, which is supposed to be a substructure of $\bb U$. So by a definable set  we shall mean an $\bb S$-definable set, unless indicated otherwise. For the rest of this paper, we assume that $\Gamma(\mathbb S)$ is nontrivial.

\begin{nota}
There is a special element $\infty = \rv(0)$ in the $\RV$-sort. For simplicity, we shall write $\RV$ to mean the $\RV$-sort without the element $\infty$, and $\RV_\infty$ otherwise --- that is, $\RV = \rv(\VF^\times)$ and $\RV_\infty = \rv(\VF)$ --- although the difference rarely matters (when it does we will of course provide further clarification).
\end{nota}

\begin{ter}[Sets and subsets]\label{nota:sub}
By a definable set in $\VF$ we mean a definable subset in the sort $\VF$, by which we just mean a subset of $\VF^n$ for some $n$, unless indicated otherwise; similarly for other sorts or structures  in place of $\VF$ that have been clearly understood in the context, such as $\RV_\infty$, $\K$, $\MM$, or any substructure $\bb M$ of $\bb U$. In particular,  a definable set without further qualification means a definable set in $\bb U$, that is, a  definable subset of $\VF^n \times \RV^m_\infty$ for some $n, m \in \N$.
\end{ter}

\begin{ter}[Relations as functions]
We may think of a relation $C$ between two sets $A$ and $B$, that is, a subset $C \sub A \times B$, as a function from $A$ into the powerset $\mdl P(B)$. For instance, if $A \sub \VF^n \times \RV^m$ is a definable set then it may be thought of as a \memph{definable} function from $A_{\VF}$ into  $\mdl P(\RV^m)$. Thus a definable function $f : A \fun \mdl P(B)$ is ontologically the same as the definable set $\bigcup_{a \in A} a \times f(a)$. We say that $f$ is \memph{finitary} if $f(a)$ is finite for every $a \in A$.
\end{ter}

The default topology in $\bb U$ is the valuation topology.

\begin{nota}\label{defn:disc}
A set $\gp \sub \VF^n \times \RV_\infty^m$ of the form $\prod_{i \leq n} \gb_i \times t$ is called an (\memph{open, closed, $\RV$-}) \memph{polydisc} if each $\gb_i$ is an (open, closed, $\RV$-) disc. The \memph{radius} $\rad(\gp)$ of the polydisc $\gp$ is the tuple $(\rad(\gb_1), \ldots, \rad(\gb_n))$. The open polydisc centered at $(a,t) \in \VF^n \times \RV_\infty^m$ with radius $\gamma \in \Gamma^n$ is denoted by $\go((a,t), \gamma)$; the closed polydisc with radius $\gamma \in \Gamma_\infty^n$ is denoted by $\gc((a,t), \gamma)$. Radius $\infty$ is used only for closed discs, with $\gc(a,\infty)=\{a\}$; open discs have finite valuative radius.

The subdiscs $\go(0, \gamma)$, $\gc(0, \gamma)$ of $\VF$ are more suggestively denoted by $\MM_\gamma$, $\OO_\gamma$, respectively, and $\MM_0$, $\OO_0$ are even more suggestively denoted by $\MM$, $\OO$. Write $\RV^{\ccirc}_\infty = \rv(\MM)$ and $\RV^{\ccirc} = \RV^{\ccirc}_\infty \mi \infty$. More generally, for $\gamma \in \Gamma$, denote by $\RV^{\ccirc}_\gamma$ the set $\rv(\MM_{\gamma} \mi 0)$.

An \memph{$\RV$-disc} is a fiber $\rv^{-1}(t)$ for $t\in\RV_\infty$, including $\rv^{-1}(\infty)=\{0\}$. In a mixed ambient space, an $\RV$-polydisc keeps the existing $\RV_\infty$-coordinates fixed. The \memph{$\RV$-hull} of a set $A$, denoted by $\RVH(A)$, is the union of all the $\RV$-polydiscs whose intersections with $A$ are nonempty. If $A$ equals $\RVH(A)$ then $A$ is called an \memph{$\RV$-pullback}.
\end{nota}

% BEGIN REFEREE INSERTION R022
\ResolvedRefereeRemark{022}{Clarification}{RV-discs and the zero fiber}{%
\textit{Original source line 748.}\quad Define an $\RV$-disc explicitly as $\rv^{-1}(t)$ for $t\in\RV_\infty$, including $\rv^{-1}(\infty)=\{0\}$. This convention is needed for the equivalence between being an $\RV$-pullback and having empty $\RV$-boundary; otherwise sets meeting coordinate hyperplanes are not covered by their stated $\RV$-hulls. Also specify that taking the hull keeps the existing $\RV$-coordinates fixed.
}{This is the intended convention. The definition now explicitly includes the singleton zero fiber and preserves the existing $\RV$-coordinates under the hull.}
% END REFEREE INSERTION R022


\begin{ter}\label{rvfiber}
A \memph{$\VF$-fiber} of a set $A$ is a set of the form $A_t$, where $t \in A_{\RV}$; in particular,  a $\VF$-fiber of a function $f : A \fun B$ is a set of the form $f_t$ for some  $t \in f_{\RV}$ (here $f$ also stands for its own graph) which is indeed (the graph of) a function.  We say that $A$ is open if every one of its $\VF$-fibers is, $f$ is continuous if every one of its $\VF$-fibers is, and so on. Composition is understood fiberwise, retaining the intermediate $\RV$-coordinates: for every fixed input, intermediate, and output $\RV$-tuple, the corresponding $\VF$-fiber maps compose continuously on their natural domain. Thus, if $f$, $g$ are continuous functions with $\ran(f)\sub\dom(g)$, the composite with these intermediate labels retained is continuous in this fiberwise sense.
\end{ter}

% BEGIN REFEREE INSERTION R023
\ResolvedRefereeRemark{023}{Error}{Fiberwise continuity is not closed under composition}{%
\textit{Original source line 753.}\quad The last assertion is false with the stated definition. Put $r=\rv(1)$ and $s=\rv(-1)$, let $f:\VF\to\{r,s\}$ send $0$ to $r$ and every nonzero point to $s$, and let $g(r)=0$, $g(s)=1$ in $\VF$. Every $\VF$-fiber of both graphs is continuous in the prescribed sense, whereas $g\circ f$ is discontinuous at $0$. A sufficient repair is to require, in addition, that the output $\RV$-coordinates be locally constant on each input $\VF$-fiber; equivalently use the valuation topology on $\VF$ and the discrete topology on $\RV$. If the weaker convention is retained, delete the composition assertion and check every later use of composition separately.
}{The composition remark concerns composition on fixed $\RV$-fibers, with the intermediate $\RV$-labels retained. This intention is now explicit. The original fiberwise continuity convention is retained; no stronger topology is imposed. The example discards the intermediate label, so it does not contradict the clarified assertion.}
% END REFEREE INSERTION R023


\begin{nota}\label{gamwhat}
Semantically we shall treat the value group $\Gamma$ as a definable sort (the $\Gamma$-sort). However, syntactically any reference to $\Gamma$ may be eliminated in the usual way and we can still work with $\lan{}{RV}{}$-formulas for the same purpose. This implies that $\gamma\in \Gamma$ is definable means that $\vv^{-1}(\gamma)$ is definable.

We shall write $\gamma^\sharp$, $\gamma \in \Gamma$, when we want to emphasize that it is the set $\vrv^{-1}(\gamma) \sub \RV$  that is being considered. More generally, if $I$ is a set in $\Gamma$ then we write $I^\sharp = \bigcup\set{\gamma^\sharp \given \gamma \in I}$. Similarly, if $U$ is a set in $\RV$ then $U^\sharp$ stands for $\bigcup\set{\rv^{-1}(t) \given t \in U}$. In particular, $\gamma^{\sharp \sharp}=\set{ x \in \VF \given \vv(x)=\gamma}$.
\end{nota}


\begin{nota}\label{jcbG:op}
Let $U \sub \RV^n \times \Gamma^m$, $V \sub \RV^{n'} \times \Gamma^{m'}$, and $C \sub U \times V$. The \memph{$\Gamma$-Jacobian} of $C$ at $((u, \alpha), (v, \beta)) \in C$, written as $\jcb_{\Gamma} ((u, \alpha), (v, \beta))$, is the element
\[
\jcb_{\Gamma} ((u, \alpha), (v, \beta)) =  - \Sigma (\vrv(u), \alpha) + \Sigma (\vrv(v), \beta),
\]
where $\Sigma (\gamma_1, \ldots, \gamma_n) = \gamma_1 + \ldots + \gamma_n$. If $C$ is the graph of a function then we just write $\jcb_{\Gamma}C(u, \alpha)$ instead of $\jcb_\Gamma((u, \alpha), C(u, \alpha))$.
\end{nota}


%\begin{conv}\label{how:gam}
%
%\end{conv}

\begin{conv}\label{topterm}
We may and shall assume that in any $\lan{}{RV}{}$-formula, every $\VF$-sort term (polynomial) occurs in the scope of an instance of the function symbol $\rv$. For example, if $f(x)$, $g(x)$ are polynomials then the formula $f(x) = g(x)$ is equivalent to $\rv(f(x) - g(x)) = \infty$. The polynomial $f(x)$ in $\rv(f(x))$ is referred to as a \memph{top term} (of the $\lan{}{RV}{}$-formula in question).
\end{conv}


\begin{nota}[The definable sort $\DC$ of discs]\label{disc:exp}
At times it will be more convenient to work in the traditional expansion $\bb U^{\textup{eq}}$ of $\bb U$ by all definable sorts. However, for our purpose, a much simpler expansion $\bb U^{\bullet}$ suffices. This expansion has only one additional sort $\DC$ that contains, as elements, all the open and closed discs. Since each point in $\VF$ may be regarded as a closed disc of valuative radius $\infty$, for convenience, we may and occasionally do think of $\VF$ as a subset of $\DC$. Heuristically, we may think of a disc that is properly contained in an $\RV$-disc as a ``thickened'' point of certain stature in $\VF$. There are two uniform cross-sort maps $\VF\times\Gamma\fun\DC$ in $\bb U^{\bullet}$, sending $(a,\gamma)$ to the open disc $\go(a,\gamma)$ and to the closed disc $\gc(a,\gamma)$, respectively. The closed-disc map extends to $\VF\times\Gamma_\infty$ by $(a,\infty)\mapsto\{a\}$. The radius is an argument of these maps, rather than an index for separately named function symbols.

The expansion $\bb U^{\bullet}$ can help reduce the technical complexity of our discussion. However, as is the case with the definable $\Gamma$-sort, it is conceptually inessential since, for the purpose of this paper, all allusions to discs as (imaginary) elements may be eliminated in favor of objects already definable in $\bb U$.

Whether parameters in $\DC$ are used or not shall be indicated explicitly, if it is necessary, by means of  introducing parameters of the form $\code \ga$ or the phrase ``in $\bb U^{\bullet}$,'' in which case the substructure $\mathbb{S}$ may contain names for discs that may or may not be definable from $\VF(\mathbb{S}) \cup \RV(\mathbb{S})$.

A disc centered at $0$ may be identified with its valuative radius together with its open/closed designation ($\Gamma$ is already treated as a definable sort). At a fixed finite radius the open and closed discs are distinct; the radius $\infty$ designates only the closed singleton disc.

For a disc $\ga \sub \VF$, the corresponding imaginary element in $\DC$ is denoted by $\code{\ga}$ when notational distinction makes the discussion more streamlined; $\code{\ga}$ may be heuristically thought of as the ``name'' of $\ga$. Conversely, a set $D \sub \DC$ is often identified with the set $\set{\ga \given \code \ga \in D}$, in which case $\bigcup D$ denotes a subset of $\VF$.
\end{nota}

% BEGIN REFEREE INSERTION R024
\ResolvedRefereeRemark{024}{Clarification}{A uniform language for the disc sort}{%
\textit{Original source line 791.}\quad For a literal language expansion, use two uniform maps $(a,\gamma)\mapsto\go(a,\gamma)$ and $(a,\gamma)\mapsto\gc(a,\gamma)$, with domain $\VF\times\Gamma$ (and the specified degenerate cases). A separate named function for every element $\gamma$ of the monster would name those radii and change the parameter structure. Also retain an open/closed tag when identifying discs centered at $0$ with radii: the two discs with a fixed finite radius are distinct.
}{The original wording was informal shorthand for uniform maps with the radius as an argument. The maps and the open/closed designation are now explicit. Open discs have finite valuative radius; only closed discs admit the singleton radius-$\infty$ case.}
% END REFEREE INSERTION R024



\subsection{Of points, discs, and atoms}
The substructure $\bb S$ is \memph{$\VF$-generated} if it is generated solely by elements in $\VF$, similarly for being $(\VF, \RV)$-generated, and so on. If $\bb S$ is $\VF$-generated then the map $\rv$ is surjective in $\bb S$, and the standing assumption makes $\Gamma(\mathbb S)$ nontrivial. Then the model-theoretic algebraic closure $\acl \mathbb{S}$ of $\mathbb{S}$ is a model of $\ACVF(0,0)$; see, for instance, \cite[Corollary~4.18]{Yin:QE:ACVF:min}. In that case, we also know what the definable closure $\dcl \mathbb{S}$ of $\mathbb{S}$ is:

\begin{lem}\label{cut:to:hensel:substru}
Suppose that $\mathbb{S}$ is $\VF$-generated and $\Gamma(\mathbb{S})$ is nontrivial. Then $\dcl \mathbb{S} = \mathbb{S}$ if and only if the underlying valued field $(\VF(\mathbb{S}), \OO(\mathbb{S}))$ of $\mathbb{S}$ is henselian.
\end{lem}
\begin{proof}
We have $\VF(\acl \mathbb{S}) = \VF(\mathbb{S})^{\alg}$, where the latter is the field-theoretic algebraic closure of $\VF(\mathbb{S})$. Since the valued field automorphisms of $(\VF(\mathbb{S})^{\alg}, \OO(\acl \mathbb{S}))$ over $(\VF(\mathbb{S}), \OO(\mathbb{S}))$ are in one-to-one correspondence with the $\lan{}{RV}{}$-automorphisms of $\acl \mathbb{S}$ over $\mathbb{S}$ and $\VF(\mathbb{S})$ is the fixed field with respect to these automorphisms if and only if $(\VF(\mathbb{S}), \OO(\mathbb{S}))$ is henselian, we see that $\VF (\dcl \mathbb{S}) = \VF(\mathbb{S})$ if and only if $(\VF(\mathbb{S}), \OO(\mathbb{S}))$ is henselian. By quantifier elimination and routine syntactical inspection, $\RV(\dcl \mathbb{S}) = \RV(\mathbb{S}) $.  The lemma follows.
\end{proof}

We shall impose at a later stage the assumption that $\mathbb{S}$ is $\VF$-generated and (without loss of generality) definably closed,  that is, $\dcl \mathbb{S} = \mathbb{S}$.

\begin{lem}[{\cite[Lemma~4.10]{Yin:QE:ACVF:min}}]\label{finite:VF:project:RV}
For any finite definable set $A$ in $\VF$ there is a definable bijection $f : A \fun U$ for some set $U$ in $\RV$.
\end{lem}

The phenomenon exemplified by this simple lemma is dubbed ``internalizing finite sets.'' Its failure in positive residue characteristics belies the deep obstacles one faces in developing a Hrushovski-Kazhdan integration theory therein; see the discussion around \cite[Lemma~3.9]{hrushovski:kazhdan:integration:vf}.

\begin{lem}[{\cite[Lemma~3.3]{Yin:special:trans}}]\label{dcl:to:ac}
For $a, b \in \VF$ and $t \in \RV$, if $b$ is algebraic over $(a, t)$ then it is algebraic over $a$.
\end{lem}

\begin{lem}[{\cite[Lemma~4.3]{Yin:QE:ACVF:min}}]\label{exchange}
The exchange principle holds for single elements in both sorts, with all algebraic closures taken over the same fixed parameter base:
\begin{itemize}
 \item For any $a$, $b \in \VF$, if $a \in \acl(b) \mi \acl(\0)$ then $b \in \acl(a)$.
 \item For any $t$, $s \in \RV$, if $t \in \acl(s) \mi \acl(\0)$ then $s \in \acl(t)$.
\end{itemize}
\end{lem}

% BEGIN REFEREE INSERTION R025
\ResolvedRefereeRemark{025}{Error; revised}{Exchange requires singleton elements}{%
\textit{Original source line 823.}\quad Here $a,b$ (and $t,s$) must be single elements of the indicated sort, with exchange allowed over an additional common parameter set. Under the tuple convention above, take independent $x,y$, $a=x$, and $b=(x,y)$: then $a\in\acl(b)\setminus\acl(\varnothing)$ but $b\notin\acl(a)$. State the scalar version explicitly. Tuple applications should use the corresponding pregeometry argument, with all closures taken over the same fixed base.
}{The exchange assertions concern single VF or RV elements over the same fixed parameter base. This intended scope is now explicit; tuple arguments use scalar exchange coordinate by coordinate.}
% END REFEREE INSERTION R025

\begin{lem}\label{finite:fib:def}
Let $f : A \fun \mdl P(B)$ be a definable finitary function, where $A$, $B$ are subsets of $\VF$. Then there is a finite definable subset $B' \sub B$ such that, for every $b \in B$, the fiber $A_b = \set{a \in A \given b \in f(a)}$ is infinite if and only if  $b \in B'$. The same holds if $A$, $B$ are subsets of $\RV$.
\end{lem}
\begin{proof}
Let $b \in B$ such that $A_b$ is infinite. Then, by compactness, there is an element $a \in A_b \mi \acl(b)$. Since $f$ is finitary, which means $b \in \acl(a)$, we must have, by Lemma~\ref{exchange},  $b \notin \acl(a) \mi \acl(\0)$ and hence $b \in \acl(\0)$. It follows that there is an algebraic subset $B' \sub B$ such that if $A_b$ is infinite then $b \in B'$, for otherwise, by compactness again, there would be a $b \in B$ that lies outside every algebraic subset of $B$ and yet the size of $A_b$ is not bounded by any natural number. Clearly we can adjust $B'$ so that it contains exactly those $b \in  B$ with $A_b$ infinite. The second case is similar.
\end{proof}

% BEGIN REFEREE INSERTION R026
\ResolvedRefereeRemark{026}{Error; revised}{The finite-exception lemma is unary}{%
\textit{Original source line 830.}\quad Require $A\subseteq\VF$ with one source coordinate, and state the corresponding unary-source restriction in the $\RV$ version. Under Terminology~\ref{nota:sub}, the assertion is false for $A=\VF^2$, $B=\VF$, $f(x,y)=\{x\}$: every inverse fiber is infinite. The target may still be a tuple: if an infinite source fiber exists, scalar exchange applied to each target coordinate shows that each is algebraic over the base. The compactness argument then gives the required finite exceptional set. Clarify that all closures in this proof are over $\mathbb S$.
}{The author confirms that the original wording $A,B\subseteq\VF$ (and likewise $\RV$) already specifies single coordinates; the tuple convention does not apply to this wording. The statement is unchanged, and all closures in the proof are over the fixed base $\mathbb S$.}
% END REFEREE INSERTION R026


\begin{defn}
We set $\dim_{\VF}(\varnothing)=-\infty$. For a nonempty definable set $A$, the \memph{$\VF$-dimension},  denoted by $\dim_{\VF}(A)$, is the largest natural number $k$ such that, possibly after re-indexing of the $\VF$-coordinates, $\pr_{\leq k}(A_t)$ has nonempty interior for some $t \in A_{\RV}$.
\end{defn}

It is a fact that if $A \sub \VF^n$ is definable then $\dim_{\VF}(A)$ equals the Zariski dimension of the Zariski closure of $A$, see for example \cite[\S~3.8]{hrushovski:kazhdan:integration:vf} for more  on dimensions.

\begin{lem}[{\cite[Lemma~4.6]{Yin:special:trans}}]\label{full:dim:open:poly}
Let $A$ be a definable set with $A_{\VF} \sub \VF^n$. Then $\dim_{\VF} (A) = n$ if and only if there is a $t \in A_{\RV}$ such that $A_t$ contains an open polydisc.
\end{lem}

\begin{lem}\label{ball:to:ac}
Let $b \in \VF$ and $\ga \sub \VF^n$ be a polydisc with $\dim_{\VF}(\ga) = n$. If $b$ is algebraic over $\code \ga$ then it is indeed algebraic.
\end{lem}
\begin{proof}
Let $\alpha = \rad(\ga)$. Then, for any $a \in \ga$, $b$ is algebraic over $(a, \alpha)$ and hence, by Lemma~\ref{dcl:to:ac}, is algebraic over $a$. Since there is an $a \in \ga$ such that $a_i \notin \acl(a_{\tilde i}, b)$ for all $i \in [n]$, by Lemma~\ref{exchange},  $b$ must be algebraic.
\end{proof}


\begin{lem}[{\cite[Lemma~4.15]{Yin:QE:ACVF:min}}]\label{effectiveness}
If $\mathbb{S}$ is $(\VF, \RV,\Gamma)$-generated then every algebraic (respectively, definable) closed disc  contains an algebraic (respectively, definable) point.
\end{lem}

\begin{lem}[{\cite[Lemma~4.17]{Yin:QE:ACVF:min}}]\label{algebraic:balls:definable:centers}
If $\mathbb{S}$ is $(\VF, \Gamma)$-generated then every algebraic (respectively, definable) disc  contains an algebraic (respectively, definable) point.
\end{lem}

A pillar of the structure of definable sets in $\bb U$ is \emph{\cmin-minimality}, meaning that every definable subset of $\VF$ is a finite boolean combination of valuative discs, allowing additional parameters to name the individual discs.

% BEGIN REFEREE INSERTION R027
\ResolvedRefereeRemark{027}{Clarification}{Parameters of the balls in C-minimality}{%
\textit{Original source line 859.}\quad The balls in a finite Boolean decomposition need not be individually $\mathbb S$-definable. For example, a definable pair of conjugate algebraic points need not split into two definable singleton balls. Say `balls with additional parameters (whose finite canonical family is algebraic over the defining parameters)' and retain the Galois-invariant family where individual components cannot be chosen.
}{Unqualified definability remains relative to the fixed ground parameter set. The individual discs in a finite Boolean decomposition may require additional parameters to be named; the sentence now makes that distinction explicit.}
% END REFEREE INSERTION R027

\begin{cor}\label{open:disc:def:point}
Suppose that $\mathbb{S}$ is $(\VF, \RV,\Gamma)$-generated.
Let $\ga \sub \VF$ be a definable disc and $A$ a definable subset of $\VF$. If $\ga \cap A$ is a nonempty proper subset of $\ga$ then $\ga$ contains a definable point.
\end{cor}
\begin{proof}
It is not hard to see that, by \cmin-minimality, if $\ga \cap A$ is a nonempty proper subset of $\ga$ then $\ga$ contains a definable closed disc and hence the claim is immediate by Lemma~\ref{effectiveness}.
\end{proof}

\begin{defn}
Let $D$ be a set of parameters. We say that a (not necessarily definable) nonempty set $A$ \memph{generates a (complete) $D$-type} if, for every $D$-definable set $B$, either $A \sub B$ or $A \cap B = \0$. In that case, $A$ is \memph{$D$-type-definable} if no set properly contains $A$ and also generates a $D$-type. If $A$ is $D$-definable and generates a $D$-type, or equivalently, if $A$ is both $D$-definable and $D$-type-definable then we say that $A$ is \memph{$D$-atomic} or \memph{atomic over $D$}.
\end{defn}

We simply say ``atomic'' when $D =\0$.

In the literature, a type could be a partial type and hence a type-definable set may have nontrivial intersection with a definable set. In this paper, since partial types do not play a role, we shall not carry the superfluous qualifier ``complete'' in our terminology.

\begin{lem}[{\cite[Lemma~3.3]{Yin:int:acvf}}]\label{atomic:base:change}
Suppose that $\Gamma(\mathbb S)$ is nontrivial. Let $A$ be an atomic set. Then $A$ is $\gamma$-atomic for all $\gamma \in \Gamma$.
\end{lem}

% BEGIN REFEREE INSERTION R028
\ResolvedRefereeRemark{028}{Error}{Atomic base change needs a nontrivial base value group}{%
\textit{Original source line 879.}\quad The unrestricted statement remains false over a trivially valued base, as the human referee observed. Over a trivially valued algebraically closed field, $A=\MM\setminus\{0\}$ is atomic; a parameter $\gamma>0$ cuts it into the nonempty proper subset $\{x:0<\vv(x)<\gamma\}$. Impose a nontrivial base value group here and propagate that assumption through uses of this lemma. In the cited proof, the essential step is that a nonempty definable subset of $\Gamma$ contains a definable point; this fails for $(0,\infty)$ over the trivial base.
}{The author adopts nontrivial $\Gamma(\mathbb S)$ as a standing assumption throughout the paper, matching the blanket hypothesis in [18]. It persists under parameter extension and supplies the required hypothesis here and in subsequent uses, including Lemma~\ref{atom:self} and Corollary~\ref{part:rv:cons}. The contrary permission at the opening of Section~\ref{sec-pscb} is removed.}
% END REFEREE INSERTION R028

\begin{lem}\label{at:ball:t}
Let $\ga \sub \VF^n$ be an atomic open polydisc. Then $\ga$ is $t$-atomic for all $t \in \RV$.
\end{lem}
\begin{proof}
{\color{ResolutionGreen}If disc parameters have previously been named, the unary assertion follows directly from \cite[Lemma~3.9]{Yin:int:acvf}, in its stated $\bb U^{\bullet}$ context: for any formula $\varphi(x,t)$ over this base, the definable map $x\mapsto\{s\in\RV:\varphi(x,s)\}$ is constant on the atomic open disc. Thus naming $t$ does not split its type. For a real/$\Gamma$ ground base the following argument also applies.} For the case $n=1$, if $\ga$ were not $t$-atomic then, by Corollary~\ref{open:disc:def:point} and Lemma~\ref{dcl:to:ac}, it would contain a point algebraic over the fixed ground base. This point belongs to a finite ground-definable set. Atomicity would then force the whole infinite open disc $\ga$ to lie in that finite set, a contradiction. For the case $n>1$, by induction on $n$, $\ga_1$ is $t$-atomic and for every $a \in \ga_1$, $\ga_{>1}$ is $(a, t)$-atomic. So $\ga$ cannot have a nonempty proper $t$-definable subset.
\end{proof}

% BEGIN REFEREE INSERTION R029
\ResolvedRefereeRemark{029}{Gap}{Atomicity proof: algebraic versus definable}{%
\textit{Original source line 886.}\quad Lemma~\ref{dcl:to:ac} only gives an algebraic point over the original base, not necessarily a definable point. This suffices: an infinite atomic disc cannot meet a finite base-definable set. Replace `definable point' at this step by that argument. The final sentence also needs `nonempty proper $t$-definable subset'. Finally, indicate why Corollary~\ref{open:disc:def:point} remains applicable whenever previously named disc parameters occur in the base; its stated hypothesis excludes such extra parameters.
}{The author approves both proof corrections: the point is algebraic over the fixed ground base, so atomicity would place the infinite disc inside a finite ground-definable set; and the final sentence now says ``nonempty proper.'' The statement is unchanged. For previously named disc parameters, the unary case now cites the accepted result \cite[Lemma~3.9]{Yin:int:acvf} in its $\bb U^{\bullet}$ context: set-valued $\RV$ constancy prevents a new $\RV$ name from splitting the atomic disc. The original real/$\Gamma$ proof and product induction are retained. This resolves the additional disc-parameter issue without extending Corollary~\ref{open:disc:def:point} beyond its stated hypotheses.}
% END REFEREE INSERTION R029

\begin{lem}\label{atom:self}
Suppose that $\Gamma(\mathbb S)$ is nontrivial. Let $A \sub \VF^n$ be a set that generates a type and $\gb \sub A$ an open (or closed) polydisc. Then $\gb$ is $\code \gb$-atomic.
\end{lem}
\begin{proof}
We do induction on $n$. The base case $n=1$ is just \cite[Lemma~3.4]{Yin:int:acvf}. For the case $n>1$, write $\gb = \prod_i \gb_i$, $\gb' = \prod_{i>1} \gb_i$, and $\gamma=\rad(\gb_1) $. For any $a \in \gb_1$, $A_a$ generates a type. So, by Lemma~\ref{atomic:base:change} and the inductive hypothesis, $\gb'$ is $(a, \gamma, \code{\gb'})$-atomic and hence $\code \gb$-atomic. Similarly $\gb_1$ is $\code \gb$-atomic. It follows that $\gb$ must be $\code \gb$-atomic.
\end{proof}


%there is a definable set $\{\gb_1, \ldots, \gb_k\}$ of open discs of the same radius such that $f(\ga) = \bigcup_i \gb_i$ and for every $\gb_i$,    $\gb_i \cap f(a)$ is a singleton.

\begin{lem}\label{at:rv:split}
Let $\ga \sub \VF^n$ be an atomic open polydisc and $f : \ga \fun \mdl P(\VF)$ a definable finitary function such that $\rv \rest f(a)$ is injective for every $a \in \ga$.  Then  $\rv \circ f$ is constant.
\end{lem}
\begin{proof}
This is clear since, by Lemma~\ref{at:ball:t}, for every $t \in \rv(\bigcup f(\ga))$, $\ga$ is $t$-atomic and hence $f$ induces a $t$-definable function $\ga \fun t^\sharp$. The lemma follows.
\end{proof}

A \memph{thin annulus} is a set of the form $\gc(a,\gamma)\setminus\go(a,\gamma)=\{x\in\VF:\vv(x-a)=\gamma\}$, where $a\in\VF$ and $\gamma\in\Gamma$. Equivalently, it is a closed disc with one maximal open subdisc removed.

\begin{lem}[{\cite[Lemma~3.8]{Yin:int:acvf}}]\label{hae:component:11}
Let $A \sub \VF$ be an atomic open disc or an atomic closed disc or an atomic thin annulus. Let $f : A \fun \VF$ be a definable function. Then $f(A)$ is also atomic of one of these three forms.
\end{lem}

% BEGIN REFEREE INSERTION R163
\ResolvedRefereeRemark{163}{Clarification}{Define thin annuli before using them}{%
\textit{Original source line 907.}\quad The term ``thin annulus'' has not been defined in this manuscript. Specify a closed disc with one maximal open subdisc removed; after choosing the center of the removed subdisc, it is a shell $\{x:\vv(x-a)=\gamma\}$. This makes the alternatives here and the later component argument unambiguous.
}{The author approves the finite-radius shell definition $\gc(a,\gamma)\setminus\go(a,\gamma)=\{x\in\VF:\vv(x-a)=\gamma\}$, with $\gamma\in\Gamma$. It is now stated before the first use in Lemma~\ref{hae:component:11}.}
% END REFEREE INSERTION R163

\begin{lem}[{\cite[Lemma~3.13]{Yin:int:acvf}}]\label{atomic:open:ball:open}
Let $A \sub \VF$ be an atomic open disc and $f : A \fun \VF$ a definable function. Then either $f$ is constant or $f(A)$ is also an atomic open disc.
\end{lem}

\begin{lem}\label{atomic:open:multidisc}
Let $\ga$, $f$ be as in Lemma~\ref{at:rv:split} and suppose that $f$ is not constant. Then there are $t_1, \ldots, t_k \in \RV_\infty$ and for each $i$, a $t_i$-definable function $f_i : \ga \fun t^\sharp_i$ such that $f = \bigcup_i f_i$ and each $f_i(\ga)$ is either a point or a $t_i$-atomic open disc.
\end{lem}
\begin{proof}
By Lemma~\ref{at:rv:split} and injectivity of $\rv$ on each finite fiber, $f$ decomposes into finitely many $t$-definable single-valued branches $g:\ga\fun t^\sharp$. A branch with $t=\infty$ is the constant zero branch. Any constant branch already has the required form, so fix a nonconstant branch, for which $t\in\RV$. Lemma~\ref{at:ball:t} makes $\ga$ atomic over $t$.

Every coordinate line of $\ga$ is atomic over its fixed coordinates and $t$. By Lemma~\ref{atomic:open:ball:open}, its image under $g$ is a point or an open disc. Any two points of $\ga$ can be joined by changing one coordinate at a time. Starting at a point with image $p$ and ending at a point with different image, the first nonconstant line image is an open neighborhood of $p$ contained in $g(\ga)$. Thus $g(\ga)$ is open. It is a proper definable subset of $\VF$, since it lies in $t^\sharp$.

By \cmin-minimality, every point of a proper definable open subset of $\VF$ lies in a maximal open disc contained in that subset. Choose such a disc $D\sub g(\ga)$ containing one value of $g$. An open line image that meets $D$ is contained in $D$: intersecting open discs are nested, and maximality excludes a larger one inside $g(\ga)$. A point image also stays in $D$ when its initial value is in $D$. The coordinate chains therefore keep every value of $g$ in $D$, so $g(\ga)=D$. The definable image of an atomic set is atomic over the same base, hence this open disc is $t$-atomic, as required.

\end{proof}

% BEGIN REFEREE INSERTION R030
\ResolvedRefereeRemark{030}{Gap; revised}{A unary image lemma is used on a polydisc}{%
\textit{Original source line 918.}\quad For $n>1$, Lemma~\ref{hae:component:11} does not directly apply to $f:\ga\to\VF$, since it has a unary domain. Supply the missing argument that $f(\ga)$ has just one primitive component. A possible route is the canonical finite component decomposition of an atomic unary set: every coordinate slice of $\ga$ is atomic over its fixed coordinates, and its image is a point or an open disc by induction. Show that such an image must lie in one primitive component, then connect any two source points by changing coordinates one at a time. The same coordinate-chain argument excludes a closed disc or thin annulus as the whole image, because each slice image lies in a single maximal open subdisc. This fills the role of the inapplicable unary citation, once the component claim is spelled out.
}{The author approves the coordinate-chain proof. Unary coordinate-line images are points or open discs; maximal open discs in a proper definable open image then force the entire nonconstant image to be one open disc. This replaces the inapplicable unary citation on the polydisc. The tag ranges here and in Corollary~\ref{part:rv:cons} are also enlarged to $\RV_\infty$, retaining the constant zero branch.}
% END REFEREE INSERTION R030

\begin{cor}\label{part:rv:cons}
Let $A \sub \VF^n$ be a definable set and $f : A \fun \mdl P(\VF)$ a definable finitary function as in Lemma~\ref{at:rv:split}. Then there are a definable finite partition $(A_i)_i$ of $A$ and, for each $i$ and all open polydiscs $\ga \sub A_i$, finitely many $(\code\ga,t_{\ga j})$-definable functions $f_{\ga j} : \ga \fun t^\sharp_{\ga j}$, where $t_{\ga j} \in \RV_\infty$, such that  $f \rest \ga = \bigcup_j f_{\ga j}$  and each $f_{\ga j}(\ga)$ is either a point or an open disc.
\end{cor}
\begin{proof}
By compactness, there is a uniform finite bound $N$ on the cardinalities of the fibers $f(a)$. For an open-polydisc code $d$, write $\ga_d$ for its polydisc and let $P(d)$ say that $\rv\circ(f\rest\ga_d)$ is constant and that the image of each of its single-valued $\rv$-tagged branches is a point or an open disc. This is a first-order condition: use the finite disjunction over the possible numbers $0\leq k\leq N$ of distinct tags, and express a point or an open disc by its usual uniform codes. Injectivity of $\rv$ on each fiber gives uniqueness of each tagged branch.

For $a\in A$, let $D_a\sub A$ be the complete type-definable set containing $a$. By Lemma~\ref{atom:self}, every open polydisc $\ga_d\sub D_a$ is $\code\ga_d$-atomic. Lemma~\ref{atomic:open:multidisc}, over the ground base together with this disc parameter, therefore gives $P(d)$. If no finite part of the type of $a$ had this property for every polydisc contained in it, compactness would produce an open-polydisc code $d$ with $\ga_d\sub D_a$ and $\neg P(d)$, a contradiction. Thus there is a ground-definable set $A_a\sub A$ containing $a$ such that $P(d)$ holds whenever $\ga_d\sub A_a$. A second application of compactness gives finitely many such sets covering $A$; take the partition they generate.

For a polydisc $\ga$ in a partition piece, choose its finitely many distinct tags $t_{\ga j}\in\RV_\infty$. The full graph of its $j$th branch is
\[
\{(x,y):x\in\ga,\ y\in f(x),\ \rv(y)=t_{\ga j}\}.
\]
It is $(\code\ga,t_{\ga j})$-definable, the branches reconstruct $f\rest\ga$, and $P(d)$ gives the asserted images.
\end{proof}

% BEGIN REFEREE INSERTION R031
\ResolvedRefereeRemark{031}{Error}{The local branch functions need the disc parameter}{%
\textit{Original source line 925.}\quad The functions $f_{\ga j}$ must be allowed parameters $(\code\ga,t_{\ga j})$, not just $t_{\ga j}$. Their domains are $\ga$, so literal $t_{\ga j}$-definability would imply that every such disc is $t_{\ga j}$-definable. Already for $f(a)=\{1\}$ on $A=\VF$, this is impossible for arbitrary subdiscs, since $t_{\ga j}=\rv(1)$ is fixed. The proof itself introduces $\code\ga$, which explains the missing parameter. Also replace the appeal to the unary Lemma~\ref{atomic:open:ball:open} by the corrected multidimensional Lemma~\ref{atomic:open:multidisc}, and explicitly formulate the uniform first-order property used in the compactness step.
}{The author approves explicitly naming the disc together with its RV tag in each full branch graph. The proof now invokes the corrected multidimensional Lemma~\ref{atomic:open:multidisc} and states a uniform first-order property, using a finite bound on fiber size and the usual point and open-disc codes, to justify both compactness steps. The zero-tag and empty-fiber cases are retained.}
% END REFEREE INSERTION R031

\begin{defn}\label{defn:dtdp}
Let $f : A \fun B$ be a bijection between two sets $A$ and $B$, each with exactly one $\VF$-coordinate. We say that $f$ is \memph{concentric} if, for every $\VF$-fiber $f_{t}$ of $f$ and all open polydiscs $\ga \sub \dom(f_{t})$, $f_{t}(\ga)$ is also an open polydisc; if both $f$ and $f^{-1}$ are concentric then $f$ has the \memph{disc-to-disc property}.
\end{defn}

\begin{lem}[{\cite[Proposition~3.19]{Yin:int:acvf}}]\label{open:pro}
Let $f : A \fun B$ be a definable bijection between two sets $A$ and $B$, each with exactly one $\VF$-coordinate. Then there is a definable finite partition $(A_i)_i$ of $A$ such that every restriction $f \rest A_i$ has the disc-to-disc property.
\end{lem}



\subsection{Continuity and fiberwise properties}\label{sec:cont:fib:prop}

\begin{defn}
We set $\dim_{\RV}(\varnothing)=-\infty$. For a nonempty definable subset $U \sub \RV^m_\infty$, the \memph{$\RV$-dimension}, denoted by $\dim_{\RV}(U)$, is the smallest number $k$ such that there is a definable finite-to-one function $f: U \fun \RV^k$ ($\RV^0$ is taken to be the singleton $\{\infty\}$).
\end{defn}

% BEGIN REFEREE INSERTION R164
\ResolvedRefereeRemark{164}{Clarification; new}{Dimension of the empty set}{%
\textit{Original source line 941.}\quad Specify $\dim_{\VF}(\varnothing)=\dim_{\RV}(\varnothing)=-\infty$ (or another negative convention) separately. The VF definition otherwise has no largest natural number for the empty set, whereas the literal RV definition gives it dimension $0$. With that latter convention the boundary lemma fails when $n=0$, since its boundary is empty but $0<0$ is false. Alternatively restrict that lemma to $n\geq1$, while still specifying the empty VF dimension.
}{The author approves $\dim_{\VF}(\varnothing)=\dim_{\RV}(\varnothing)=-\infty$. Both definitions now state this convention separately and retain their original clauses for nonempty sets, including zero-coordinate cases.}
% END REFEREE INSERTION R164

\begin{defn}
Let $A$ be a definable subset of $\VF^n$. The \memph{$\RV$-boundary} of $A$, denoted by $\partial_{\RV}A$, is the definable subset of $\rv(A)$ such that $t \in \partial_{\RV} A$ if and only if $t^\sharp \cap A$ is a proper nonempty subset of $t^\sharp$. The definable set $\rv(A) \mi \partial_{\RV}A$, denoted by $\ito_{\RV}(A)$, is called the \memph{$\RV$-interior} of $A$.
\end{defn}

% BEGIN REFEREE INSERTION R032
\ResolvedRefereeRemark{032}{Typo}{Definability of the RV-boundary}{%
\textit{Original source line 945.}\quad Insert `definable' before `subset of $\VF^n$'; without that assumption the boundary and interior need not be definable.
}{The author approves explicitly assuming that $A$ is definable, so that its RV-boundary and RV-interior have the stated definability.}
% END REFEREE INSERTION R032

Obviously, $A \sub \VF^n$ is an $\RV$-pullback if and only if $\partial_{\RV} A$ is empty.


\begin{lem}\label{RV:bou:dim}
Let $A$ be a definable subset of $\VF^n$. Then $\dim_{\RV}(\partial_{\RV} A) < n$.
\end{lem}
\begin{proof}
For $n=0$, the boundary is empty and has dimension $-\infty$. For $n=1$, take a finite Boolean disc decomposition of $A$. A disc containing $0$ is an $\RV$-pullback, whereas a disc avoiding $0$ lies in one $\RV$-disc. A partially met $\RV$-fiber therefore contains a member of this finite disc family that avoids $0$, so there are only finitely many such fibers.

For the inductive step, fix $i\in[n]$ and put $X_i=\pr_{\tilde i}(A)$ and $T_a=\partial_{\RV}A_a$ for $a\in X_i$. Each $T_a$ is finite. For each $t\in T_a$, Corollary~\ref{open:disc:def:point}, over the parameters $(a,t)$, gives a definable point in $t^\sharp$. By compactness there are finitely many definable partial choice functions covering all such pairs $(a,t)$; order this list and use the first available choice to obtain a definable function $c(a,t)\in t^\sharp$. Set $f_i(a)=\{c(a,t):t\in T_a\}$ if $T_a\ne\varnothing$, and $f_i(a)=\{0\}$ otherwise. Thus every fiber is finite and contains exactly one representative of each boundary tag when that boundary is nonempty; in particular $\rv\rest f_i(a)$ is injective.

First split $X_i$ by the pattern of its zero coordinates, and apply Corollary~\ref{part:rv:cons} on the remaining coordinates of each piece. This gives a definable finite partition $(A_{ij})_j$ of $X_i$ such that, whenever $t^\sharp\sub A_{ij}$, the finite set $\rv(f_i(a))$ is independent of $a\in t^\sharp$. The zero-pattern split allows $t$ to have $\infty$ coordinates: after deleting them, $t^\sharp$ is an open polydisc in the remaining coordinates. Write $f_{ij}=f_i\rest A_{ij}$ and let $C_{ij}(t)$ be this common finite set for $t\in\ito_{\RV}(A_{ij})$.

Regard each $C_{ij}$ as its graph in $\RV_\infty^n$, retaining the input tags and placing its output tag in coordinate $i$. Put
\[
 C=\bigcup_{i,j}C_{ij},\qquad
 B=\bigcup_{i,j}\{s\in\rv(A):\pr_{\tilde i}(s)\in\partial_{\RV}A_{ij}\}.
\]
Each graph $C_{ij}$ is finite over its $n-1$ input coordinates, so $\dim_{\RV}(C)<n$. By induction, $\dim_{\RV}(\partial_{\RV}A_{ij})<n-1$; with only one remaining coordinate, this gives $\dim_{\RV}(B)<n$. Finally, for a partially met $\RV$-polydisc, start in $A$ and change coordinates one at a time to reach a point outside $A$. The first exit identifies a boundary tag of a unary section. Its other tags either lie in some $\partial_{\RV}A_{ij}$, placing the full tuple in $B$, or in $\ito_{\RV}(A_{ij})$, placing it in $C$. Hence $\partial_{\RV}A\sub B\cup C$, as required.
\end{proof}

% BEGIN REFEREE INSERTION R033
\ResolvedRefereeRemark{033}{Gap}{RV-boundary proof needs uniform branches}{%
\textit{Original source line 962.}\quad When constructing $f_{ij}$, ensure that each fiber contains exactly one representative of each boundary $\RV$-disc; injectivity of $\rv$ on these finite fibers is a hypothesis of Corollary~\ref{part:rv:cons}. The sentence about constancy should mean that the finite set $\rv(f_{ij}(a))$ is independent of $a\in t^\sharp$, not that $\rv$ is constant on the union of all those finite sets. The argument also depends on the preceding corollary after its parameter and multidimensional gaps are repaired. A shorter alternative is the human referee's generic-lift argument: over a model containing the parameters, a tuple of full $\RV$-dimension has a lift realizing one complete type, so its lift cannot meet $A$ partially; state the dimension/type fact precisely if using this route.
}{The author approves choosing exactly one VF representative of each boundary RV tag and stating constancy of the finite RV-image set along each source RV fiber. The proof now splits zero-coordinate faces before applying Corollary~\ref{part:rv:cons}, retains all input tags in the graphs, and gives the zero-dimensional base case and the coordinate-chain coverage argument.}
% END REFEREE INSERTION R033

Let $f : \VF^n \fun \VF^m$ be a definable function.

\begin{lem}[{\cite[Lemma~8.7]{Yin:special:trans}}]\label{fun:almost:cont}
There is a definable closed set $A \sub \VF^n$ with $\dim_{\VF}(A) < n$ such that $f \rest (\VF^n \mi A)$ is continuous.
\end{lem}

Per the usual terminology, this lemma says that $f$ is continuous \memph{almost everywhere}.

\begin{defn}\label{defn:diff}
In this definition, for a nonzero tuple $x\in\VF^k$ we set $\vv(x)=\min_i\vv(x_i)$, and we set $\vv(0)=\infty$. This minimum convention is separate from the coordinatewise valuation image used in Terminology~\ref{bounded}.


For any $a \in \VF^n$, we say that $f$ is \memph{differentiable at $a$} if there is a linear map $\lambda : \VF^n \fun \VF^m$ (of $\VF$-vector spaces) such that, for any $\epsilon \in \Gamma$, if $b \in \VF^n\setminus\{0\}$ and $\vv(b)$ is sufficiently large then
\[
\vv(f(a + b) - f(a) - \lambda(b)) - \vv(b) > \epsilon.
\]
It is straightforward to check that if such a linear function $\lambda$ (represented by a matrix with entries in $\VF$) exists then it is unique and hence is called the \memph{derivative of $f$ at $a$}.

Write $f = (f_1, \ldots, f_m)$. For $a = (a_i, a_{\tilde i}) \in \VF^n$, if the derivative of the function $f_j \rest (\VF \times a_{\tilde i})$ at $a_i$ exists then we call it the \memph{$ij$th partial derivative of $f$ at $a$}. Whenever $f$ is differentiable at $a$, the $(j,i)$-entry of its derivative matrix is this $ij$th partial derivative.
\end{defn}

% BEGIN REFEREE INSERTION R034
\ResolvedRefereeRemark{034}{Clarification; revised}{Valuation of a tuple and the zero increment}{%
\textit{Original source line 980.}\quad Define $\vv(b)=\min_i\vv(b_i)$ for a nonzero tuple, with $\vv(0)=\infty$, and require $b\ne0$ in the difference quotient. Otherwise the display includes the undefined expression $\infty-\infty$ at $b=0$. Later, where valuations of tuples are taken coordinatewise (Terminology~\ref{bounded}), distinguish that use from this minimum convention. The input-first indexing of partial derivatives is legitimate, but please keep the corresponding matrix convention explicit.
}{The author approves the minimum coordinate valuation for differentiability, with $\vv(0)=\infty$, and excludes the zero increment from the difference quotient. This avoids $\infty-\infty$. The input-first $ij$th partial-derivative convention is retained and its derivative-matrix position $(j,i)$ is explicit. The separate boundedness convention remains coordinatewise.}
% END REFEREE INSERTION R034

We differentiate functions between arbitrary definable sets by ``forgetting'' the $\RV$-coordinates. More precisely, let
\[
f : \VF^n \times \RV^m \fun \VF^{n'} \times \RV^{m'}
\]
be a definable function. For each fixed input $\RV$-tuple $t$, generic local constancy of definable $\RV$-valued maps \cite[Lemma~8.9]{Yin:special:trans} implies that, outside a definable set of $\VF$-dimension less than $n$, the output $\RV$-tuple is constant on a neighborhood of each input $a$. If this tuple is $s$, then $\dom(f_{(t,s)})$ contains a neighborhood of $a$. At such a point $(a,t)$, we define the partial derivatives of $f$ to be those of this local $\VF$-map $f_{(t,s)}$. If $n = n'$ and all the partial derivatives exist, the \memph{Jacobian of $f$ at $(a,t)$} is the determinant of the matrix whose $(j,i)$-entry is the $ij$th partial derivative, and is denoted by $\jcb_{\VF} f(a,t)$.

\begin{lem}\label{diff:almost:every}
Assume $n=n'$. The Jacobian of $f$ exists and is a continuous function almost everywhere.
\end{lem}
\begin{proof}
First remove the definable locus where the output $\RV$-tuple is not locally constant in the input $\VF$-coordinates. On each resulting open fiber with both input and output $\RV$-tuples fixed, extend the $\VF$-map by zero outside its domain and apply \cite[Corollary~9.9]{Yin:special:trans} to its partial derivatives. Their nonexistence and discontinuity loci are definable and have $\VF$-dimension less than $n$ in every such fiber. Retaining or projecting away $\RV$-coordinates does not increase $\VF$-dimension, so the combined exceptional locus still has $\VF$-dimension less than $n$. Off this locus all partial derivatives, and hence their determinant, are continuous. In arity zero the Jacobian is the empty determinant $1$.
\end{proof}

% BEGIN REFEREE INSERTION R035
\ResolvedRefereeRemark{035}{Gap}{Equal VF dimensions and generic RV constancy}{%
\textit{Original source line 993.}\quad State $n=n'$ in Lemma~\ref{diff:almost:every}; otherwise the determinant called the Jacobian has not been defined. Moreover, the preceding existence of one full-dimensional fiber for each input $t$ does not show that such fibers cover almost every input point. Use the standard generic local constancy of a definable map from $\VF^n$ to $\RV^{m'}$ first, and then apply the cited differentiability theorem on its open fibers. This provides the missing almost-everywhere coverage.
}{The author approves the square-dimension hypothesis $n=n'$ and the generic local-constancy argument for the output $\RV$-labels. The proof now applies the cited partial-differentiability theorem on open fixed-label charts and explains the smaller-dimensional exceptional locus. The arity-zero Jacobian is the empty determinant $1$.}
% END REFEREE INSERTION R035

\begin{ter}\label{bounded}
We say that a set $I \sub \Gamma_{\infty}^n$ is \memph{$\gamma$-bounded}, where $\gamma \in \Gamma$, if it is contained in the box $[\gamma, \infty]^n$, and is \memph{doubly $\gamma$-bounded} if it is contained in the box $[-\gamma, \gamma]^n$. More generally, let $A$ be a subset of $\VF^n \times \RV_\infty^m \times \Gamma_\infty^l$ and
\[
\operatorname{valim}(A) = \set{(\vv(a), \vrv(t), \gamma) \given (a, t, \gamma) \in A} \sub \Gamma_{\infty}^{n+m+l},
\]
where $\vv(a)$ and $\vrv(t)$ are taken coordinatewise. We say that $A$ is \memph{$\gamma$-bounded} if $\operatorname{valim}(A)$ is, and likewise for \memph{doubly $\gamma$-bounded}. A set is \memph{bounded} if it is $\gamma$-bounded for some $\gamma\in\Gamma$, and \memph{doubly bounded} if it is doubly $\gamma$-bounded for some $\gamma\in\Gamma$. The notation $A_{\Gamma}$ continues to mean projection to existing $\Gamma$-coordinates.
\end{ter}

% BEGIN REFEREE INSERTION R036
\ResolvedRefereeRemark{036}{Clarification}{A valuation image is not a coordinate projection}{%
\textit{Original source line 1001.}\quad The notation $A_\Gamma$ here conflicts with Notation~\ref{indexing}, where it denotes projection to existing $\Gamma$-coordinates. Use a separate notation such as $\operatorname{valim}(A)$ for this coordinatewise valuation image. Also explicitly say that `bounded' means $\gamma$-bounded for some $\gamma$, and similarly for `doubly bounded'.
}{The author approves $\operatorname{valim}(A)$ for the coordinatewise valuation image, reserving $A_{\Gamma}$ for projection to existing $\Gamma$-coordinates. The existential finite bound in ``bounded'' and ``doubly bounded'' is now explicit. This coordinatewise convention is distinct from the minimum tuple valuation used in differentiation.}
% END REFEREE INSERTION R036

\begin{defn}[Contractions]\label{defn:corr:cont}
A function $f : A \fun B$ is \memph{$\rv$-contractible} if there is a (necessarily unique) function $f_{\downarrow} : \rv(A) \fun \rv(B)$, called the \memph{$\rv$-contraction} of $f$, such that
\[
(\rv \rest B) \circ f = f_{\downarrow} \circ (\rv \rest A).
\]
Similarly, it is \memph{$\res$-contractible} (respectively, \memph{$\vv$-contractible}) if the same holds in terms of $\res$ (respectively, $\vv$ or $\vrv$, depending on the coordinates) instead of $\rv$.
\end{defn}


\begin{lem}\label{db:to:db}
Let $U \sub \RV^k$ be a doubly bounded set and $f: U \fun \RV$ a definable function. Then $f(U)$ is also doubly bounded. The same holds if the codomain of $f$ is $\Gamma$.
\end{lem}
\begin{proof}
For an $\RV$-valued $f$, put $h=\vrv\circ f$; for a $\Gamma$-valued $f$, put $h=f$. We retain the induction on $k$, with the case $k=0$ immediate. For $k>1$ and $t\in\pr_1(U)$, the inductive hypothesis makes $h(U_t)$ doubly bounded. Its finite bound
\[
\gamma_t=\sup\set{|h(t,x)|\given x\in U_t}
\]
exists by definable completeness of $\Gamma$, and the least-upper-bound formula makes $t\efun\gamma_t$ definable. The projection $\pr_1(U)$ is doubly bounded, so the unary $\Gamma$-valued case bounds all $\gamma_t$ and hence bounds $h(U)$. Thus both claims reduce to the case $k=1$ and a definable function $g:U\fun\Gamma$. Then, by \cmin-minimality, we may assume that $g$  $\vrv$-contracts to a function $g_{\downarrow} : \vrv(U) \fun \Gamma$. Since definable functions in $\Gamma$ are piecewise $\Q$-linear, the range of $g_{\downarrow}$ must be doubly bounded.
\end{proof}

% BEGIN REFEREE INSERTION R037
\ResolvedRefereeRemark{037}{Gap}{Uniformity in the bounded RV-image lemma}{%
\textit{Original source line 1017.}\quad The induction on $k$ must control the bounds uniformly in the remaining coordinates; boundedness on each fiber alone does not give a bound on their union. A precise route is to consider the relation of possible values $\vrv(f(u))$ over $\vrv(u)\in\Gamma^k$. Residue-field/value-group orthogonality makes each such fiber finite, uniformly; stable embeddedness and elimination in divisible ordered abelian groups then give finitely many piecewise rational-affine branches. These have bounded images on a bounded box. Include this justification, or an equivalent explicit uniform induction.
}{The author retains the original induction. Its fiber bounds are chosen definably as finite suprema of the absolute output values; the unary $\Gamma$-valued case then bounds this definable bound function uniformly. The proof now spells out that step, without replacing the induction by a relation of all input and output values.}
% END REFEREE INSERTION R037


% the following lemma has been moved to the section on proper covariants bijecetions, since it is not used before and is a corollary of the sujectivity of the lifting morphism
%Recall that, by \cmin-minimality, every definable set $A \sub \VF$ may be written as a unique boolean combination of discs $\ga_i$. Any atomic subset of these discs is called an \memph{atomic band} of $A$. If $\gB$ is an atomic band of $A$ then every disc in it has the same attributes: radius, open or closed, a positive or a negative boolean component, etc. So they may be considered as attributes of $\gB$. The totality of these attributes is called the \memph{signature} of $\gB$.
%
%\begin{lem}\label{invar:db}
%Let $A \sub \VF^n \times \RV^m$ be a definable set such that $A_{\VF}$ is bounded and $A_{\RV}$ is doubly bounded. Let $f : A \fun \Gamma$ be a definable function  that is constant on $A \cap \ga$ for all open polydiscs $\ga$ of radius $\gamma \in \Gamma$. Then $f(A)$ is doubly bounded.
%\end{lem}
%\arthur{I suspect one can get this for free as a corollary of lemma \ref{L:measure:surjective} and \ref{db:to:db}, if we can postpone the proof. It seems to be the case, since the lemma is only used in the last section.}
%\begin{proof}
%By Lemma~\ref{db:to:db} and \omin-minimality in the $\Gamma$-sort, it is enough to consider the case $A \sub \VF^n$. It is easy to see that, by induction on $n$, this may be further reduced to the case $n=1$. By the assumption, we may assume that, for every open disc $\ga$ of radius $\gamma$, either $\ga \cap A = \0$ or $\ga \sub A$, and moreover $A \cap \MM_\gamma = \0$, that is, $A$ is doubly bounded.
%
%Suppose for contradiction that $f(A)$ contains a definable interval $I$ of the form $(\beta, \infty)$. By \cmin-minimality, for  each $\alpha \in I$, if $\delta^{\dsh} \cap f^{-1}(\alpha)$ is a nonempty proper subset of $\delta^{\dsh}$ for some  $\delta \in \Gamma$ then some atomic band $\gB_\alpha$ of $f^{-1}(\alpha)$ is contained in $\delta^{\dsh}$. Therefore, shrinking $I$ if necessary, we can construct a definable function $g$ on $I$ such that
%\begin{itemize}
%  \item either $g(\alpha)  = \gB_\alpha$ of the same signature for all $\alpha \in I$ or
%  \item $g(\alpha) \in \Gamma$ and $g(\alpha)^{\dsh} \sub f^{-1}(\alpha)$ for all $\alpha \in I$.
%\end{itemize}
%In the latter case, we may assume that $g$ is $\Q$-linear, but then its image cannot be doubly bounded since $g$ is injective, contradicting the choice of $I$. The former case requires further argument.
%
%Consider the function $\bar g$ on $I$ given by $\alpha \efun \rad(\gB_\alpha)$. Since $\rad(\gB_\alpha) \leq \gamma$ for all $\alpha \in I$ and $A$ is bounded,  $\bar g(I)$ must be doubly bounded. By \omin-minimality in the $\Gamma$-sort, shrinking $I$ if necessary, we may assume that $\bar g$ is constant, say $\bar g(I) = \delta$, and either every $\gB_\alpha$ is closed or every $\gB_\alpha$ is open. In the former case, by Lemmas~\ref{effectiveness} and \ref{dcl:to:ac}, $g(I)$ must be finite; but this is impossible since either every $\gB_\alpha$ is a positive boolean component or every $\gB_\alpha$ is a negative boolean component. In the latter case, for the same reason, we may assume that $g(I)$ is infinite. Then a moment of reflection shows that there would be a closed disc $\gb$ of radius $\delta$ and a $\code \gd$-definable function, induced by the ``inverse'' of $g$, from the set of maximal open subdiscs of $\gd$ onto an infinite subset of $\Gamma$, which is easily seen to contradict \cmin-minimality.
%\end{proof}
%


\begin{defn}
Let $A \sub \VF^n$ and $f : A \fun \mdl P(\VF^m)$ be a definable function. For $a \in \VF^n$ and $L \sub \VF^m$, we say that $L$ is a \memph{limit set of $f$ at $a$}, written as $\lim_{A \rightarrow a} f \sub L$, if for every $\epsilon \in \Gamma$ there is a $\delta \in \Gamma$ such that if $c \in \go(a,\delta) \cap (A\setminus\{a\})$ then
\[
f(c)\sub\bigcup_{b\in L}\go(b,\epsilon).
\]
This is a local definition and imposes no global boundedness hypothesis. A finite limit set forces the union of the values $f(c)$ to be bounded on some punctured neighborhood of $a$. Whenever boundedness of the range is assumed, it means boundedness of $\bigcup_{c\in A}f(c)$.

A limit set $L$ of $f$ at $a$ is \memph{minimal} if no proper subset of $L$ is a limit set of $f$ at $a$.
\end{defn}

% BEGIN REFEREE INSERTION R038
\ResolvedRefereeRemark{038}{Clarification; revised}{Simplify the limit-set definition}{%
\textit{Original source line 1046.}\quad Say explicitly that $\bigcup_{a\in A}f(a)$ is bounded. The auxiliary subset $L'$ does no work: the displayed condition is equivalent to $f(c)\subseteq\bigcup_{b\in L}\go(b,\epsilon)$. The following definition of continuity is applied to finitary functions without a global boundedness hypothesis, so the present definition needs a compatible extension. One can use the same formula for arbitrary correspondences whenever a finite affine limit exists, or formulate limits in $(\mathbb P^1(\VF))^m$ from the outset; local boundedness near the point suffices for affine local uses.
}{The author approves the local limit formula without a global range bound, using all of $L$ in place of the redundant subfamily $L'$. A finite limit set supplies local boundedness on a punctured neighborhood. Any separately assumed range bound refers to the union of all values. Minimality and the puncture are unchanged.}
% END REFEREE INSERTION R038

Observe that if $\lim_{A \rightarrow a} f \sub L$ and $b \in L$ is not isolated in $L$ then actually $\lim_{ A \rightarrow  a} f \sub L \mi b$. So in a minimal limit set every element is isolated. Moreover, if a minimal limit set $L$ exists then its topological closure is unique:

\begin{lem}[{\cite[Lemma~9.2]{Yin:special:trans}}]
Let $L_1, L_2 \sub \VF^m$ be two minimal limit sets of $f$ at $a$. Then the closures of $L_1$, $L_2$ coincide.
\end{lem}


This lemma justifies writing $\lim_{ A \rightarrow a} f = L$ when $L$ is a closed (hence the unique) minimal limit set of $f$ at $a$, in particular, when $L$ is finite and minimal.

Let $A \sub \VF^n$ be a definable set with $\dim_{\VF}(A) = n$ and $f : A \fun \mdl P(\VF^m)$ a definable finitary function. We say that $f$ is \memph{continuous} if $\lim_{ A \rightarrow a} f = f(a)$ for all $a \in A$.

\begin{lem}\label{fini:conti}
Suppose $n>0$. Then $f$ is continuous almost everywhere.
\end{lem}
\begin{proof}
This can be reduced to the familiar case that $f$ is indeed a function into $\VF^m$ as follows. By Lemma~\ref{finite:VF:project:RV} and compactness, there are a definable set $U$ in $\RV$ and a definable function $g: A \times \VF^m \fun A \times U$ such that, for all $a \in A$, $\pr_A(g(a, b)) = a$  and $g \rest (a \times f(a))$ is injective. For each $t \in U$, let
\[
A_t = \set{a \in A \given (a,t) \in g(a \times f(a))}
\]
and $f_t \sub f$ be the obvious $t$-definable function $A_t \fun \VF^m$. The discontinuity locus $E$ of $f$ is definable: the limit condition is first-order, and minimality of the finite set $f(a)$ can be tested by removing one of its elements at a time. Suppose for contradiction that $\dim_{\VF}(E)=n$. By Lemma~\ref{full:dim:open:poly}, $E$ contains a nonempty open polydisc. Apply generic local constancy \cite[Lemma~8.9]{Yin:special:trans} to the finite label-set map
\[
a\efun g(a\times f(a))_{\RV}.
\]
After shrinking to a nonempty open polydisc $P\sub E$, this set is a fixed finite set $T$. If $T$ is empty, then $f$ is the constant empty correspondence on $P$ and is continuous there, a contradiction.

For each $t\in T$, the inverse branch $f_t$ is defined on all of $P$. Extend it by zero outside $A_t$ and apply Lemma~\ref{fun:almost:cont}. The finitely many branch-exception sets have union of $\VF$-dimension less than $n$, so choose $a\in P$ outside their union. All branches are continuous at $a$. Their finite union has limit set $f(a)$; this limit set is minimal because the distinct values at $a$ can be separated by sufficiently small discs and, since $n>0$, the punctured polydisc still accumulates at $a$. Thus $f$ is continuous at $a$, contradicting $P\sub E$. Hence $\dim_{\VF}(E)<n$.
\end{proof}

% BEGIN REFEREE INSERTION R039
\ResolvedRefereeRemark{039}{Gap; revised}{One continuous branch is not enough}{%
\textit{Original source line 1068.}\quad The argument finds one $t$ for which $A_t$ has full dimension and $f_t$ is generically continuous. That alone does not account for all branches of $f$. Apply generic local constancy to the finite $\RV$-label set $a\mapsto g(a\times f(a))_{\RV}$, then simultaneously apply generic continuity to its finitely many inverse branches on each neighborhood. To justify the global dimension conclusion, suppose the discontinuity locus had full dimension and choose an open polydisc in it where this label set is constant; the finite collection of branch-exception sets cannot fill that polydisc. The empty label set simply gives the constant empty correspondence and should be handled separately. This supplies the missing coverage argument.
}{The author approves the expanded proof and the restriction $n>0$. The original ending uses the dimension principle that an $\RV$-indexed definable union of sets of $\VF$-dimension less than $n$ still has dimension less than $n$. The expanded argument obtains the same coverage by making the finite label set locally constant and treating all branches together, including the empty-label case. The positive-arity hypothesis is needed for the unchanged punctured limit convention.}
% END REFEREE INSERTION R039

\begin{lem}[{\cite[Lemma~9.5]{Yin:special:trans}}]\label{lim:exists}
Suppose that $f : \MM \mi 0 \fun \mdl P (A)$ is a definable finitary function, where $A$ is a closed and bounded set in $\VF$. Then there is a definable finite set $L \sub A$ such that $\lim_{\MM \mi 0 \rightarrow 0} f = L$.
\end{lem}

We remark that the finite-limit theorem also holds in the coordinatewise projective line $(\mathbb P^1(\VF))^m$, with the valuation topology: if $A\subseteq\VF^m$ is not bounded, take its closure there and allow infinite coordinates in $L$. This is the projective version of the same cited theorem; no topological local compactness is assumed.


\yimu{revision of the next few lemmas}

\begin{lem}[Curve selection]\label{curve:select}
Suppose that $\mathbb S$ is $\VF$-generated. Let $A\sub\VF^n$ be definable and let $b$ be a point in its closure in the coordinatewise projective line $(\mathbb P^1(\VF))^n$, with the valuation topology. Then there is a $b$-definable continuous finitary function $g:\MM\setminus\{0\}\fun\mdl P(A)$ whose values are nonempty and such that
\[
\lim_{\MM\setminus\{0\}\to0}g=\{b\}.
\]
\end{lem}

% BEGIN REFEREE INSERTION R165
\ResolvedRefereeRemark{165}{Clarification; new}{Definability and nonempty values in curve selection}{%
\textit{Original source line 1081.}\quad State that $A$ is definable. In the proof require $g(c)\ne\varnothing$ for every $c$ (or at least on a punctured neighborhood of $0$): ``finitary'' in the earlier definition allows empty values, and $g(c)\subseteq\go(b,\vv(c))$ alone is also satisfied by the identically empty correspondence. The algebraic-closure construction, once its base hypothesis is corrected below, supplies finite nonempty values and hence the stronger conclusion $\lim g=\{b\}$.
}{The author approves explicitly taking $A\subseteq\VF^n$ definable and requiring finite nonempty values at every input. The construction gives the stronger singleton limit $\{b\}$. Projective charts handle infinite coordinates without changing the affine target set of the curve.}
% END REFEREE INSERTION R165
\begin{proof}
Use a coordinatewise projective chart around $b$ to reduce to the case $b\in\VF^n$: at an infinite coordinate use inversion and restrict $A$ to the chart. These changes are definable over $b$ and preserve convergence and continuity locally.

Naming the finite tuple $b$ and a nonzero $c\in\MM$ preserves $\VF$-generation and the standing nontriviality of the ground valuation. Hence $\acl(\mathbb S,b,c)$ is a model of $\ACVF(0,0)$. Since $b\in\overline A$, the definable set $A\cap\go(b,\vv(c))$ is nonempty, so elementarity gives a point of that intersection algebraic over $(\mathbb S,b,c)$. Compactness supplies a $b$-definable finitary function $g_0$ on $\MM\setminus\{0\}$ such that
\[
\varnothing\ne g_0(c)\sub A\cap\go(b,\vv(c))
\]
for every $c$. Thus $g_0$ converges to $\{b\}$, and nonempty values make this limit minimal.

By Lemma~\ref{fini:conti}, with its one-dimensional source, $g_0$ is continuous outside a finite set. Choose a $b$-definable clopen disc $D\sub\MM$ around $0$ such that $g_0$ is continuous on $D\setminus\{0\}$. Also $\acl(\mathbb S,b)$ is a model, so $A$ contains a finite nonempty $b$-definable set $E$. Define
\[
g(c)=\begin{cases}
g_0(c),&c\in D\setminus\{0\},\\
E,&c\in\MM\setminus D.
\end{cases}
\]
The two pieces are relatively clopen in $\MM\setminus\{0\}$, so $g$ is continuous, finite-valued and nonempty everywhere. Its limit at $0$ is still $\{b\}$. Finally transport the construction back through the projective chart. This avoids assuming that a definable value-group radius is the valuation of a definable scalar.
\end{proof}

% BEGIN REFEREE INSERTION R040
\ResolvedRefereeRemark{040}{Error}{Curve selection requires a restriction on the base}{%
\textit{Original source line 1088.}\quad The assertion is false for the arbitrary $\mathbb S$ allowed so far. Let $K$ be a nontrivially valued algebraically closed base field, add an $\RV$ parameter $t$ with $\vrv(t)=0$ transcendental over its residue field, and put $A=\{(x,y):y\ne0,\rv(x/y)=t\}$. Then $(0,0)\in\overline A$. A nonempty finite $\mathbb S$-definable choice in $A$ at an input $c\in K\cap(\MM\setminus\{0\})$ would have coordinates algebraic over $K,t,c$, hence in $K$ by Lemma~\ref{dcl:to:ac}; their ratio cannot have leading term $t$. Even though finitary values may be empty, $b\in\lim g$ forces their support to accumulate at $0$; by $C$-minimality that support contains a punctured neighborhood, whose defining radius lies in $\Gamma(K)$ since adjoining the residue parameter $t$ does not enlarge definable values. It therefore contains such an input $c\in K$. Thus the claimed curve does not exist over that base. Impose $\VF$-generation before this lemma (as already suggested in the human report), or allow explicitly specified extra $\VF$ representatives. Under $\VF$-generation, the added nonzero $c\in\MM$ makes the algebraic closure a model, and the displayed construction yields nonempty finite values converging to $b$ itself. For the last rescaling step, first shrink to a disc whose radius is the valuation of a definable scalar; the original definable radius need not be realized by such a scalar.
}{The author approves $\VF$-generation as a local hypothesis of this lemma, validating the algebraic-closure model argument. The continuous germ is extended by a constant finite nonempty set on a clopen complement, so no definable scalar realizing an arbitrary definable radius is assumed. Later geometric uses may temporarily enlarge the base; no global $\VF$-generation assumption is imposed here.}
% END REFEREE INSERTION R040

This is an analogue of the highly useful \omin-minimal curve selection.

\begin{rem}
Although the $\Gamma$-sort is essentially an \omin-minimal group, curve selection fails due to the presence of the top element $\infty$. For example, the set $\Gamma \times (0, \infty) \sub  \Gamma^2_\infty$ contains curves approaching the boundary point $(0, 0)$ or $(\infty, \infty)$ but no curves approaching the boundary point $(\infty, 0)$, because every definable one-parameter curve is piecewise rational-affine. Near a one-sided endpoint, one affine piece remains, and a positive affine coordinate cannot tend to $0$ while another coordinate tends to $\infty$.
\end{rem}

% BEGIN REFEREE INSERTION R166
\ResolvedRefereeRemark{166}{Typo; new}{Eventually affine curves in the value group}{%
\textit{Original source line 1094.}\quad Replace ``all curves are lines'' by ``every definable one-parameter curve is piecewise rational-affine''. Near a one-sided endpoint, one affine piece remains; a positive affine coordinate cannot tend to $0$ while another coordinate tends to $\infty$. This is the precise reason the stated boundary point cannot be approached.
}{The author approves the precise piecewise rational-affine formulation. An eventual affine piece gives the stated obstruction to approaching $(\infty,0)$; the example and the conclusion are unchanged.}
% END REFEREE INSERTION R166

%In addition, we can make this ``curve'' $f \circ g$ approach $b$ as fast as we wish. More precisely, if $l : \Gamma^+ \fun \Gamma^+$ is a definable increasing function then $g$ can be so constructed that $\vv(f(g(c))) > l(\vv(c))$ for all $c \in \MM \mi 0$.


\begin{lem}\label{cb:to:cb}
Let $A\subseteq\VF^n$ be definable, closed and bounded, and let $f:A\fun\VF^m$ be definable and continuous. Then $f(A)$ is also closed and bounded.
\end{lem}
\begin{proof}
The conclusion is geometric, so we may temporarily enlarge the parameter base to a small $\VF$-generated base over which all the data are definable. The standing nontriviality of the ground valuation is preserved. Over this base, and after naming any additional finite $\VF$ tuple, algebraic closure is a model of $\ACVF(0,0)$. Elementarity and compactness therefore give finite nonempty definable choices in the nonempty families below. No choice over the original parameter base is asserted.

First let $b\in\overline{f(A)}\setminus f(A)$, where this closure is affine, and name $b$. For every $c\in\MM\setminus\{0\}$ the set
\[
\{x\in A:f(x)\in\go(b,\vv(c))\}
\]
is nonempty. Applying the model argument over the enlarged base together with $(b,c)$, and then compactness, gives a definable finitary function $F:\MM\setminus\{0\}\fun\mdl P(A)$ such that every $F(c)$ is nonempty and $f(F(c))\subseteq\go(b,\vv(c))$. Lemma~\ref{lim:exists}, in its finite-dimensional form, gives a finite minimal limit set $L\subseteq A$. This set is nonempty, since $F$ has nonempty values arbitrarily close to $0$. For every $a\in L$, minimality implies that values of $F$ meet every neighborhood of $a$ for inputs arbitrarily close to $0$: otherwise $a$ could be removed from $L$. Continuity at $a$ and the displayed containment force $f(a)=b$, contradicting $b\notin f(A)$. Thus $f(A)$ is closed.

If $f(A)$ were unbounded, one coordinate $f_i$ would be unbounded. For every $c\in\MM\setminus\{0\}$ the definable set
\[
\{x\in A:\vv(f_i(x))<-\vv(c)\}
\]
would be nonempty. The same model argument and compactness give finite nonempty choices $F(c)$ inside these sets. Again Lemma~\ref{lim:exists} supplies a nonempty finite minimal limit set $L\subseteq A$. Choose $a\in L$. Continuity makes $f_i$ bounded on a neighborhood of $a$ in $A$, while minimality makes $F(c)$ meet that neighborhood for arbitrarily small nonzero $c$. The inequality $\vv(f_i(x))<-\vv(c)$ for all $x\in F(c)$ contradicts this local bound. Hence $f(A)$ is bounded as well.
\end{proof}

% BEGIN REFEREE INSERTION R041
\ResolvedRefereeRemark{041}{Gap}{Selecting and lifting the limiting curve}{%
\textit{Original source line 1104.}\quad If the image is unbounded but closed in affine space, its affine frontier need not contain a point outside it. Take the frontier in the coordinatewise projective line and allow an infinite coordinate. Also the existence of algebraic points in every fiber needs justification under the chosen parameter assumptions; for the curve pullback it is enough to include the curve parameter $c$ in the algebraic closure and apply the model argument there. Preserve a nonempty finite-valued parametrization throughout and apply the limit lemma in its multidimensional form.
}{The author approves treating affine missing boundary points and unbounded output separately. Finite nonempty choices are justified by elementarity and compactness over a temporary $\VF$-generated base. The finite-dimensional limit theorem and continuity then give both closedness and boundedness, without a global $\VF$-generation assumption or an unjustified original-base choice.}
% END REFEREE INSERTION R041



\begin{cor}\label{conti:homeo}
If the function $f$ in Lemma~\ref{cb:to:cb} is injective then  it is a homeomorphism from $A$ onto $f(A)$.
\end{cor}

\begin{lem}[The closed graph criterion for continuity]\label{closed:gra}
Let $A\subseteq\VF^n$ and $B\subseteq\VF^m$ be definable, with $B$ closed in affine space. Write $B^{\mathrm p}$ for the closure of $B$ in $(\mathbb P^1(\VF))^m$, with the valuation topology. A definable function $f:A\fun B$ is continuous if and only if its graph is closed in $A\times B^{\mathrm p}$. If $B$ is bounded, then $B^{\mathrm p}=B$, so this is the usual relative closed-graph condition.
\end{lem}
\begin{proof}
The graph of a continuous map into the Hausdorff space $B^{\mathrm p}$ is closed, proving the forward implication.

Conversely, suppose the graph is closed but $f$ is not continuous at $a\in A$. Choose a clopen projective polydisc $U$ around the finite point $f(a)$ such that
\[
D=\{x\in A:f(x)\notin U\}
\]
has $a$ in its closure. As in the preceding proof, temporarily enlarge the parameter base to a small $\VF$-generated base containing the parameters of these data. This changes neither continuity nor graph closure. By Lemma~\ref{curve:select}, there is a definable continuous finitary function $C:\MM\setminus\{0\}\fun\mdl P(D)$ with nonempty values and limit $\{a\}$.

Apply the projective version of Lemma~\ref{lim:exists} to the finite-valued function $c\efun f(C(c))$. It has a finite minimal limit set $L$, which is nonempty because all the values are nonempty. Every point of $L$ belongs to the closed set $B^{\mathrm p}\setminus U$: minimality makes the image meet each neighborhood of such a point for inputs arbitrarily close to $0$. Choose $b'\in L$. This same property, together with $C(c)\to\{a\}$, shows that $(a,b')$ belongs to the closure of the graph of $f$ in $A\times B^{\mathrm p}$. Closedness of the graph gives $b'=f(a)$, contrary to $b'\notin U$. This proves continuity without a global $\VF$-generation hypothesis.
\end{proof}

% BEGIN REFEREE INSERTION R042
\ResolvedRefereeRemark{042}{Gap; revised}{The compactified closed-graph criterion}{%
\textit{Original source line 1117.}\quad The parenthesis appears intended to impose the correct compactified-target hypothesis; make it explicit by taking $B$ to be a closed definable subset of $(\mathbb P^1(\VF))^m$, and take graph closure in $A\times B$. The example $f(0)=0$, $f(x)=1/x$ for $x\ne0$ has closed graph in affine $\VF^2$ but is discontinuous at $0$: this is a counterexample only to the unbounded affine-target version, not to the intended compactified statement. For the reverse implication, apply curve selection to the set of input points $x$ with $f(x)\notin\gb$, approaching $a$; apply the finite-limit lemma to the resulting finite-valued image in $B$. A limit point $b'$ lies outside the clopen neighborhood $\gb$, and $(a,b')$ lies in the graph closure. Also, ``projective line with the valuation topology'' is more precise than a literal topological one-point compactification; local compactness is not assumed here.
}{The author approves making the intended projective ambient space explicit: the affine target is $B$, and graph closure is in $A\times B^{\mathrm p}$, with the valuation topology. The reverse proof selects bad inputs first and then takes a finite projective limit of their images. Temporary base enlargement supplies curve selection without changing the geometric conclusion or imposing global $\VF$-generation.}
% END REFEREE INSERTION R042


\begin{lem}\label{lim:weak}
Let $f : A  \fun B$ be a definable bijection between two open subsets of $\VF$. Suppose that $B$ is bounded. Then $f$ is weakly concentric, that is, for every $a \in A$ there is a $b \in \VF$ such that, for every sufficiently large $\gamma \in \Gamma$, $f(\go(a, \gamma) \mi a) = \go(b, \delta) \mi b$ for some $\delta \in \Gamma$.
\end{lem}

Note that, compared to Lemma~\ref{open:pro}, the assumption of this lemma is stronger and its conclusion weaker, but the point is that we do not need a partition to achieve it.

\begin{proof}
Fix $a\in A$. After translating the input and restricting to a small punctured disc about $a$, Lemma~\ref{lim:exists}, applied with the closed bounded target $\overline B$, gives a finite minimal limit. Its cardinality is at most one by \cite[Lemma~9.4]{Yin:special:trans}. It is nonempty because the function has a value at every point of the punctured disc. Write this limit as $\{b\}$.

By shrinking the punctured disc, exclude the at most one input for which $f(x)=b$. The valuation calculation in the proof of \cite[Lemma~8.3]{Yin:special:trans}, applied to the translated germ $f(x)-b$, gives a rational-affine law
\[
\vv(f(x)-b)=\rho(\vv(x-a)),\qquad \rho(\gamma)=q\gamma+c,
\]
for all sufficiently small nonzero $x-a$, where $q\in\mathbb Q_{>0}$ and $c\in\Gamma$. Indeed, unary polynomials have eventually monomial valuation, the definable germ selects one of the finitely many resulting valuation laws, and convergence to $b$ forces positive slope.

By $C$-minimality, $B$ contains a punctured neighborhood of $b$. On this neighborhood the inverse of $f$ has a unique projective limit, using the projective finite-limit theorem and the same cardinality bound. That limit must be $a$, since $x\to a$ implies $f(x)\to b$ and the image avoids $b$ on a sufficiently small punctured disc. Apply the same valuation calculation to the inverse to obtain a positive rational-affine law $\sigma$ for its tail radii. Substituting $y=f(x)$ shows that $\sigma\circ\rho$ is the identity on a tail of $\Gamma$, so the two affine laws are inverse there.

Consequently, for every sufficiently large $\gamma$, with $\delta=\rho(\gamma)$, the image of the whole source thin annulus is exactly the whole target thin annulus:
\[
f\bigl(\gc(a,\gamma)\setminus\go(a,\gamma)\bigr)
 =\gc(b,\delta)\setminus\go(b,\delta).
\]
The forward inclusion is the law for $f$; the reverse inclusion follows by applying the inverse law to every point of the target annulus. Taking the unions over all radii greater than $\gamma$ proves
\[
f\bigl(\go(a,\gamma)\setminus\{a\}\bigr)
 =\go(b,\rho(\gamma))\setminus\{b\},
\]
as required. No assertion that $q=1$ is needed.
\end{proof}

% BEGIN REFEREE INSERTION R043
\ResolvedRefereeRemark{043}{Gap; revised}{The annulus equality needs a proof}{%
\textit{Original source line 1132.}\quad Specify that $A$ and $B$ each have one $\VF$-coordinate. The finite limit is a singleton $\{b\}$ by \cite[Lemma~9.4]{Yin:special:trans}; nonemptiness follows because the values are nonempty on a punctured neighborhood. The equality of whole thin annuli needs more explanation. Here is a repair. The valuation calculation in the proof of \cite[Lemma~8.3]{Yin:special:trans}, applied to $x\mapsto\vv(f(x)-b)$, gives
\[
 \vv(f(x)-b)=q\vv(x-a)+c\qquad(q\in\mathbb Q_{>0})
\]
on a sufficiently small punctured neighborhood: each unary polynomial has eventually monomial valuation, and the limit forces positive slope. Remove the at most one exceptional input with $f(x)=b$. By $C$-minimality, $B$ contains a punctured neighborhood of $b$. The inverse has a unique projective limit there, necessarily $a$, since points $x\to a$ give $f(x)\to b$. Apply the same valuation calculation to the inverse. The resulting two affine maps of tail radii are inverse, so every sufficiently small target shell consists exactly of images of its corresponding source shell. Taking their unions proves the asserted punctured-disc equality. Thus the lemma is repairable; an additional assertion that $q=1$ is unnecessary.
}{The arity objection is withdrawn: the explicit subsets $A,B\subseteq\VF$ are already unary, and no amendment is needed. The author agrees that the original proof is quick and that the remaining comments spell out details. The expanded proof makes the singleton limit and the inverse tail-radius laws explicit, giving equality of whole thin annuli. The statement and proof strategy are retained; this is an expository expansion, not a correction of a false claim.}
% END REFEREE INSERTION R043

\begin{defn}
Let $p : A \fun \Gamma$ be a definable function. We say that $p$ is an \memph{$\go$-partition} of $A$ if, for every $a \in A$, the function $p$ is constant on $\go(a, p(a)) \cap A$.
\end{defn}


This makes sense even when $A$ has no $\VF$-coordinates, in which case any definable function $A \fun \Gamma$ is an $\go$-partition of $A$.

\begin{lem}\label{vol:par:bounded}
Let $p$ be an $\go$-partition of $A$. Suppose that $A_{\RV}$ is doubly bounded and $A_t$ is closed and bounded for every $t \in A_{\RV}$. Then $p(A)$ is doubly bounded.
\end{lem}
\begin{proof}
First suppose that $A$ has no $\RV$-coordinates. The empty set is immediate. Choose $\eta\in\Gamma$ such that $A\subseteq\OO_\eta^n$. If $p(a)<\eta$ for some $a\in A$, then $A\subseteq\go(a,p(a))$, so the partition property makes $p$ constant on $A$. Otherwise $p(A)\subseteq[\eta,\infty)$. In either case $p(A)$ is bounded below.

Suppose for contradiction that $p(A)$ is unbounded above. We may temporarily enlarge the parameter base to a small $\VF$-generated base, as in Lemma~\ref{cb:to:cb}; the desired bound is a geometric conclusion. For every $c\in\MM\setminus\{0\}$ the definable set
\[
A_{\vv(c)}=\{x\in A:p(x)>\vv(c)\}
\]
is nonempty. Elementarity of the algebraic-closure model over this base together with $c$, followed by compactness, supplies a definable finitary function $F:\MM\setminus\{0\}\fun\mdl P(A)$ with finite nonempty values $F(c)\subseteq A_{\vv(c)}$. Continuity of this choice is not needed.

By Lemma~\ref{lim:exists}, $F$ has a finite minimal limit set $L\subseteq A$, nonempty because its values are nonempty arbitrarily close to $0$. Choose $e\in L$. Minimality implies that, for arbitrarily small nonzero $c$, $F(c)$ meets $\go(e,p(e))$. Choose such a $c$ with $\vv(c)>p(e)$ and $x\in F(c)\cap\go(e,p(e))$. The partition property at $e$ gives $p(x)=p(e)$, whereas the choice gives $p(x)>\vv(c)>p(e)$, a contradiction. This proves the upper bound.

For the general case, each nonempty fiber image $p(A_t)$ is bounded above and below by the case just proved, with $t$ temporarily named. Definable completeness in $\Gamma$ gives finite endpoints
\[
\alpha_t=\inf p(A_t),\qquad \beta_t=\sup p(A_t).
\]
Their defining greatest-lower-bound and least-upper-bound formulas make these functions definable over the original parameter base. Thus
\[
t\efun\gamma_t=\max\{|\alpha_t|,|\beta_t|,0\}
\]
is a definable $\Gamma$-valued function on $A_{\RV}$. By Lemma~\ref{db:to:db}, its image is doubly bounded. A common upper bound for the $\gamma_t$ bounds $p(A)$ in both directions, proving the lemma.
\end{proof}

% BEGIN REFEREE INSERTION R044
\ResolvedRefereeRemark{044}{Gap}{Bounds for an open-disc partition}{%
\textit{Original source line 1152.}\quad The lower bound has a short direct proof: choose $\eta$ with $A\subseteq\OO_\eta^n$. If $p(a)<\eta$ for some $a$, then $A\subseteq\go(a,p(a))$, forcing $p$ to be constant; otherwise $p(A)\subseteq[\eta,\infty)$. For the upper-bound argument, justify the nonempty finite choices in $A_{\vv(c)}$ using the corrected curve-selection/model argument, and use the nonempty finite limit supplied by the limit lemma. In the general case, take a definable bound for each fiber image $p(A_t)$ by its definable endpoints in $\Gamma$, then apply Lemma~\ref{db:to:db}; fiberwise boundedness by itself is not the uniform conclusion.
}{The author agrees that the original proof is quick and the comments spell out details. The expansion retains its lower-bound, limit-point contradiction and uniform-fiber-bound strategy. It gives the direct lower-bound calculation, justifies finite nonempty choices over a temporary base, and writes out the definable endpoint bound used with Lemma~\ref{db:to:db}. The statement is unchanged.}
% END REFEREE INSERTION R044

This lemma, in its various incarnations, is crucial for the good behavior of motivic Fourier transform (see \cite[\S~11]{hrushovski:kazhdan:integration:vf} and \cite{yin:hk:part:3}). For essentially the same reason, the  main construction of this paper depends heavily on  it (often in combination with Lemma~\ref{db:to:db}).



\begin{lem}\label{lem:open:fiberwise}
Let $A\subseteq \VF^n$ be a definable set. Suppose that, for each $a \in A_1$, the fiber $A_a \sub \VF^{n-1}$ is open (respectively, closed). Then there is a finite definable set $A' \subseteq A_1$ such that $A\cap ((A_1 \mi A')\times \VF^{n-1})$ is open (respectively, closed) in $(A_1 \mi A')\times \VF^{n-1}$.
\end{lem}
\begin{proof}
We prove the open case over an arbitrary definable unary base $I\subseteq\VF$, allowing empty fibers, with $A\subseteq I\times\VF^{n-1}$. The closed case then follows by taking the relative complement in $I\times\VF^{n-1}$. In the statement take $I=A_1$.

Put $X=I\times\VF^{n-1}$ and let
\[
S=A\cap\Cl_X(X\setminus A).
\]
Thus no point of $S$ is interior to $A$ relative to $X$. It is enough to show that $S_1$ is finite. For each $\gamma\in\Gamma$, set
\[
T_\gamma=\{(a,x)\in S:\go(x,\gamma)\subseteq A_a\}.
\]
Openness of every fiber gives
\[
S_1=\bigcup_{\gamma\in\Gamma}(T_\gamma)_1.
\]
If every $(T_\gamma)_1$ were finite, this definable $\Gamma$-indexed union would be finite as well: lift the index to $\RV$ and use the $\RV$-indexed $\VF$-dimension bound \cite[Lemma~4.3]{Yin:special:trans}. Consequently, if $S_1$ is infinite, some $(T_\gamma)_1$ contains an open disc $D$.

For this geometric argument we may temporarily enlarge the parameter base to a small $\VF$-generated base containing the data, as in Lemma~\ref{cb:to:cb}. For every $a\in D$, the nonempty definable fiber $(T_\gamma)_a$ has an algebraic point by elementarity of the algebraic-closure model over this base together with $a$. Compactness therefore gives a definable finitary function
\[
F:D\fun\mdl P(\VF^{n-1}),\qquad
\varnothing\ne F(a)\subseteq(T_\gamma)_a.
\]
By Lemma~\ref{fini:conti}, choose $a_0\in D$ at which $F$ is continuous, and put $E=F(a_0)$. After shrinking $D$ around $a_0$, continuity gives
\[
F(a)\subseteq\bigcup_{e\in E}\go(e,\gamma)\qquad(a\in D).
\]
Here $E$ is finite and nonempty. The finitely many definable sets
\[
D_e=\{a\in D:F(a)\cap\go(e,\gamma)\ne\varnothing\},\qquad e\in E,
\]
cover $D$. Hence one of them contains an open subdisc $D'$. Fix its index $e$. For each $a\in D'$, choose $x\in F(a)\cap\go(e,\gamma)$. Since $(a,x)\in T_\gamma$ and intersecting open polydiscs of the same radius are equal,
\[
\go(e,\gamma)=\go(x,\gamma)\subseteq A_a.
\]
Thus $D'\times\go(e,\gamma)$ is an open subset of $X$ contained in $A$ and meeting $S$, a contradiction. Therefore $S_1$ is finite. Away from $S_1$, every point of $A$ is interior relative to $X$, which proves the open case and, by relative complementation, the closed case.
\end{proof}

% BEGIN REFEREE INSERTION R045
\ResolvedRefereeRemark{045}{Gap; revised}{The fiberwise openness proof needs reconstruction}{%
\textit{Original source line 1173.}\quad The displayed proof contains genuine mistakes, but the conclusion has a repair. Retain $d=\dim_{\VF}(S)<n$ instead of replacing it by $n-1$. Define
\[
 \alpha(a,x)=\min\{\beta\geq0:\go(x,\beta)\subseteq A_a\}\quad ((a,x)\in S).
\]
This minimum exists: for a nonempty complement it is the maximum of $0$ and the supremum of the valuative distances from $x$ to that complement; for a full fiber it is $0$. The original unrestricted minimum fails for full fibers. If $S_1$ is infinite, choose a smooth algebraic chart of a stratum of $S$ whose projection to the first coordinate is nonconstant. In characteristic zero, a separating coordinate projection, including the first coordinate, is generically \'etale; a local branch over an open polydisc is therefore definable and continuous. Pull $\alpha$ back to this chart and shrink to an open polydisc where it is constant, say $\epsilon$, using generic local constancy of maps into $\Gamma$. Fix the remaining chart coordinates and shrink the first-coordinate disc $D$ until the resulting section $a\mapsto(a,s(a))\in S$ satisfies $s(a)\in\go(s(a_0),\epsilon)$. Then
\[
 D\times\go(s(a_0),\epsilon)\subseteq A,
\]
contradicting $(a_0,s(a_0))\in S$. This proves that $S_1$ is finite, including when $d<n-1$. The contradiction concerns interior points of $A$, not an open polydisc inside $S$. For the closed case and the next proposition, formulate the lemma over a fixed definable base $I\subseteq\VF$, allowing empty fibers; passage between $I$ and the actual projection costs only its finite unary boundary.
}{The author approves the independently reconstructed bad-locus proof. A definable family of valid radii and the $\RV$-indexed union principle supply one radius on a whole disc of base points. Finite nonempty choices and unary generic continuity give an open subset of $A$ meeting $S$, the required contradiction. No codimension-one assumption or minimum-radius choice is needed. The statement is retained, and working over a fixed unary base with empty fibers allowed justifies the closed case by relative complement.}
% END REFEREE INSERTION R045



\begin{prop}\label{prop:continuity:fiberwise}
Let $A$ be as in Lemma~\ref{lem:open:fiberwise} and $f : A \fun \VF$  a definable function such that
for each $a \in A_1$, the induced function $f_a : A_a \fun \VF$ is continuous.
Then there is a finite definable set $A' \sub A_1$ such that $f$ is continuous on
$\bigcup_{a\in A_1 \mi A'} a \times A_a$.
\end{prop}

For the closed-graph argument, the projective line $\mathbb P^1(\VF)$ is understood with the valuation topology, as in Lemma~\ref{closed:gra}.

\begin{proof}
Put $I=A_1$ and write $H_a$ for the graph of $f_a$. Let
\[
G_a=\Cl_{\VF^{n-1}\times\mathbb P^1(\VF)}(H_a),\qquad
G=\bigcup_{a\in I}\{a\}\times G_a.
\]
This family is definable: fiberwise closure is expressed uniformly by the valuative neighborhood formula in each affine chart. Every $G_a$ is closed in $\VF^{n-1}\times\mathbb P^1(\VF)$.

Apply the fixed-base version proved in Lemma~\ref{lem:open:fiberwise}, with empty fibers allowed, to $G$ in each of the two standard affine charts of $\mathbb P^1(\VF)$. Outside the union $I'$ of the resulting two finite definable exceptional sets, $G$ is closed in
\[
(I\setminus I')\times\VF^{n-1}\times\mathbb P^1(\VF),
\]
since closedness is local on this open chart cover. Fiberwise continuity gives
\[
H_a=G_a\cap\bigl(A_a\times\mathbb P^1(\VF)\bigr).
\]
Consequently, the graph of the restricted function is closed relative to
\[
\bigl(A\cap((I\setminus I')\times\VF^{n-1})\bigr)\times\mathbb P^1(\VF).
\]
The closed-graph criterion (Lemma~\ref{closed:gra}) proves the asserted continuity. This spells out the relative-closure step in the original implication.
\end{proof}

% BEGIN REFEREE INSERTION R046
\ResolvedRefereeRemark{046}{Gap; revised}{Relative graph closedness is insufficient for the cited lemma}{%
\textit{Original source line 1188.}\quad Continuity of $f_a$ makes its graph closed in $A_a\times\mathbb P^1(\VF)$, whereas the preceding closed-fiber lemma concerns ambient closed fibers. Define instead
\[
 G_a=\overline{\operatorname{graph}(f_a)}\subseteq\VF^{n-1}\times\mathbb P^1(\VF).
\]
Use the preceding lemma in its fixed-base form over $I=A_1$, allowing empty fibers, in each of the two affine charts of $\mathbb P^1$. Equivalently, also discard the finite unary boundaries of the chart projections when applying the lemma as currently stated. Outside the union of these finite exceptional sets, $G$ is relatively closed over $I$. Fiberwise continuity gives $\operatorname{graph}(f)=G\cap(A\times\mathbb P^1)$, so the graph is closed in $A\times\mathbb P^1$, and the compactified closed-graph criterion applies. This proves the proposition without an open/closed-fiber hypothesis on $A_a$; clarify the phrase `as in Lemma~\ref{lem:open:fiberwise}' accordingly.
}{The author approves expanding the original implication through its relative-closure step. Ambient closures of the fiber graphs form a definable closed-fiber family; Lemma~\ref{lem:open:fiberwise} applies in the two affine charts over the fixed unary base. Intersecting back with the domain gives relative graph closedness, so Lemma~\ref{closed:gra} gives continuity. The statement and its existing scope are retained. The target wording now explicitly uses the projective line with the valuation topology.}
% END REFEREE INSERTION R046

This is analogous to \cite[Corollary~6.2.4]{dries:1998} in the \omin-minimal environment and \cite[Lemma~2.38]{Yin:tcon:1.5} in the \T-convex environment.


\section{Hrushovski-Kazhdan style integral}\label{sec:HK:main}

\begin{defn}
\label{def:Kgroup}
The Grothendieck semigroup $\gsk \mdl C$ of a category $\mdl C$ endowed with a binary operation ``$\mi$'' and a binary relation ``$\subseteq$'' (subobject) is the free semigroup generated by the isomorphism classes of $\mdl C$, subject to the usual scissor relation
$[A \mi B] + [B] = [A]$, when $B\subseteq A$,
where $[A]$, $[B]$ denote the isomorphism classes of the objects $A$, $B$, and ``$\mi$'' and ``$\subseteq$'' are compatible with isomorphisms between objects, usually they are just set subtraction and set inclusion. Sometimes $\mdl C$ is also equipped with a binary operation (for example, cartesian product) that induces multiplication in $\gsk \mdl C$, in which case $\gsk \mdl C$ becomes a (commutative) semiring. The formal groupification of  $\gsk \mdl C$, which is then a ring, is denoted by $\ggk \mdl C$. %If a group $G$ acts on the objects of $\mdl C$ and the morphisms of $\mdl C$ are $G$-equivariant, that is, they commute with $G$-actions, then the corresponding $G$-equivariant Grothendieck ring is denoted by $\ggk^G \mdl C$. If $G = \lim_n G_n$ is profinite then we shall always impose the condition that a $G$-action factor through some $G_n$-action.

% BEGIN REFEREE INSERTION R047
\RefereeRemark{047}{Clarification; revised}{Commutative Grothendieck construction}{%
\textit{Original source line 1199.}\quad Optional precision: call this the free \emph{commutative monoid} on the isomorphism classes, subject to the scissor relations, with the empty object representing $0$. The final phrase ``which is then a ring'' can be read as referring to the preceding case where multiplication has been supplied; one can make that qualification explicit. This is a convention clarification, not a defect in the subsequent Grothendieck constructions.
}
% END REFEREE INSERTION R047
\end{defn}

\begin{defn}
An \memph{$\RV$-fiber} of a definable set $A$ is a set of the form $A_a$, where $a \in A_{\VF}$. The \memph{$\RV$-fiber dimension} of $A$ is the maximum of the $\RV$-dimensions of its $\RV$-fibers and is denoted by $\dim^{\fib}_{\RV}(A)$.
\end{defn}


\begin{lem}[{\cite[Lemma~4.13]{Yin:special:trans}}]\label{RV:fiber:dim:same}
Let $f : A \fun A'$ be a definable bijection. Then $\dim^{\fib}_{\RV}(A) = \dim^{\fib}_{\RV} (A')$.
\end{lem}

\begin{defn}[$\VF$-categories]\label{defn:VF:cat}
The objects of the category $\VF[k]$ are the definable sets of $\VF$-dimension less than or equal to $k$ and $\RV$-fiber dimension $0$ (that is, all the $\RV$-fibers are finite). Any definable bijection between two such objects is a morphism of $\VF[k]$. Set $\VF_* = \bigcup_k \VF[k]$.
\end{defn}

As soon as one considers adding volume forms to definable sets in $\VF$, the question of ambient dimension arises and, consequently, one has to take ``essential bijections'' as morphisms.

 \begin{defn}
\label{defn:VF:mu:G}
An object of the category $\mgVF[k]$ is a definable pair $(A, \omega)$, where $A \in \VF[k]$, $A_{\VF} \sub \VF^k$, and $\omega : A \fun \Gamma$ is a function, which is understood as a \memph{definable $\Gamma$-volume form} on $A$. A \memph{morphism} between two such objects $(A, \omega)$, $(B, \sigma)$ is a definable \memph{essential bijection} $F : A \fun B$, that is, a bijection that is defined outside definable subsets of $A$, $B$ of $\VF$-dimension less than $k$, such that, for almost all $x \in A$,
\begin{equation}\label{cond:mg}
\omega(x) = \sigma(F(x)) + \vv(\jcb_{\VF} F(x)).
\end{equation}
We also say that such an $F$ is \memph{$\Gamma$-measure-preserving}.
\end{defn}

\begin{nota}
In \cite{hrushovski:kazhdan:integration:vf}, the category $\mgVF[k]$ is denoted by ${\mu_{\Gamma}}{\VF[k]}$ to indicate that the volume forms take values in $\Gamma$ as opposed to $\RV$. Here the subscript ``$\Gamma$'' is dropped since  we will not consider $\RV$-volume forms at all.
\end{nota}

In the definition above and other similar ones below, for the cases $k =0$, the reader should interpret things such as $\VF^0$ and how they interact with other things in a natural way. For instance, $\VF^0$ may be treated as the empty tuple, the only definable set of $\VF$-dimension less than $0$ is the empty set, and $\jcb_{\VF}$ is always $1$ on sets that have no $\VF$-coordinates. So $(A, \omega) \in \mgVF[0]$ if and only if $A$ is a finite definable set in $\RV_\infty$.

Set $\mgVF[{\leq} k] = \coprod_{i \leq k} \mgVF[i]$ and $\mgVF[*] = \coprod_{k} \mVF[k]$; similarly for the other categories below (with or without volume forms).


\begin{rem}\label{mor:equi}
Let $F : (A, \omega) \fun (B, \sigma)$ be a $\mgVF[k]$-morphism. Our intention is that such an $F$ should identify the two objects, that is, an isomorphism. However, if $F$ is not defined everywhere in $A$ then obviously it does not admit an inverse. We remedy this by introducing the following congruence relation $\sim$ on $\mgVF[k]$. Let $G : (A, \omega) \fun (B, \sigma)$ be another $\mgVF[k]$-morphism. Then $F \sim G$ if $F(a) = G(a)$ for all $a \in A$  outside a definable subset  of $\VF$-dimension $< k$. The morphisms of the quotient category $\mgVF[k] /{\sim}$ have the form $[F]$, where $F$ is a $\mgVF[k]$-morphism. Clearly every $(\mgVF[k] / {\sim})$-morphism is an isomorphism and hence $\mgVF[k] / {\sim}$ is a groupoid. In fact, all the categories of definable sets we shall work with should be and are groupoids.

It is certainly more convenient to work with representatives than equivalence classes. In the discussion below, this quotient category $\mgVF[k] /{\sim}$ will almost never be needed except when it comes to forming the Grothendieck semigroup or, by abuse of terminology, when we speak of two objects of $\mgVF[k]$ being isomorphic.

% BEGIN REFEREE INSERTION R167
\RefereeRemark{167}{Typo; new}{Missing verb}{%
\textit{Original source line 1238.}\quad In the preceding paragraph, write ``that is, \emph{be} an isomorphism'' after ``should identify the two objects.''
}
% END REFEREE INSERTION R167
\end{rem}



\begin{defn}[$\RV$-categories]\label{defn:c:RV:cat}
The objects of the category $\RV[k]$ are the pairs $(U, f)$ with $U$ a set in $\RVV$ and $f : U \fun \RV^k$ a definable finite-to-one function. Given two such objects $(U, f)$, $(V, g)$, any definable bijection $F : U \fun V$ is a \memph{morphism} of $\RV[k]$.
\end{defn}

By Lemma~\ref{finite:VF:project:RV}, the two categories $\VF[0]$, $\RV[0]$ are actually equivalent, similarly for other such categories.

\begin{nota}\label{0coor}
We emphasize that if $(U, f)$ is an object of $\RV[k]$ then $f(U)$ is a subset of $\RV^k$ instead of $\RV^k_\infty$, while $\infty$ can occur in any coordinate of $U$. An object of  $\RV[*]$ of the form $(U, \id)$ is often just written as $U$.

More generally, if $f : U \fun \RV_\infty^k$ is a definable finite-to-one function then $(U, f)$ denote the obvious object of $\RV[{\leq} k]$. For example, the inclusion $\{\infty\} \fun \RV_{\infty}$ gives rise to an object of $\RV[0]$, the inclusion $\{(1, \infty)\} \fun \RV_{\infty}^2$ gives rise to an object of $\RV[1]$, and so on. Often $f$ will be a coordinate projection (every object in $\RV[*]$ is isomorphic to an object of this form). In that case, $(U, \pr_{\leq k})$ is simply denoted by $U_{\leq k}$ and its class in $\gsk \RV[k]$ by $[U]_{\leq k}$, etc.

% BEGIN REFEREE INSERTION R048
\RefereeRemark{048}{Typo}{Agreement}{%
\textit{Original source line 1252.}\quad Change the verb ``denote'' to ``denotes'' in the sentence about $(U,f)$.
}
% END REFEREE INSERTION R048
\end{nota}

\begin{defn}[$\RES$-categories]\label{defn:RES:cat}
The category $\RES[k]$ is the full subcategory of $\RV[k]$ such that $(U, f) \in \RES[k]$ if and only if $\vrv(U)$ is finite.
\end{defn}

\begin{nota}\label{mor:ftf}
Every $\RV[k]$-morphism $F: (U, f) \fun (V, g)$  induces a definable finite-to-finite correspondence $F^\dag \sub f(U) \times g(V)$, which may also be thought of as a finitary function $f(U) \fun \mdl P(g(V))$. Recall the operator $\jcb_{\Gamma}$ from Notation~\ref{jcbG:op}. For $ u \in U$, we abbreviate $\jcb_{\Gamma}  F^\dag(f(u), g \circ F(u))$ as $\jcb_{\Gamma}  F^\dag(u)$

% BEGIN REFEREE INSERTION R049
\RefereeRemark{049}{Typo}{Missing punctuation}{%
\textit{Original source line 1260.}\quad Add a full stop after the abbreviation $\jcb_{\Gamma}F^\dag(u)$.
}
% END REFEREE INSERTION R049
\end{nota}



\begin{defn}[$\RV$- and $\RES$-categories with $\Gamma$-volume forms]\label{defn:RV:mu}
An object of the category $\mgRV[k]$ is a definable triple $(U, f, \omega)$, where $(U, f)$ is an object of $\RV[k]$ and $\omega : U \fun \Gamma$ is a function, which is understood as a \memph{definable $\Gamma$-volume form} on $(U, f)$. A \memph{morphism} between two such objects $(U,f, \omega)$, $(V, g, \sigma)$ is an $\RV[k]$-morphism $F: (U, f) \fun (V, g)$ such that for every $ u \in U$,
\begin{equation}\label{cond:mg2}
\omega(u) = (\sigma \circ F)(u) + \jcb_{\Gamma}  F^\dag(u).
\end{equation}

The category $\mgRES[k]$ is the obvious full subcategory of $\mgRV[k]$.
\end{defn}



\begin{defn}[$\Gamma$-categories]\label{def:Ga:cat}
The objects of the category $\Gamma[k]$ are the finite disjoint unions of definable subsets of $\Gamma^k$. A definable bijection between  two such objects is a  morphism of $\Gamma[k]$ if and only if it is the $\vrv$-contraction of a definable bijection.


The category $\Gamma^{\fin}[k]$ is the full subcategory of $\Gamma[k]$ such that $I \in \Gamma^{\fin}[k]$ if and only if $I$ is finite.
\end{defn}

\begin{rem}
If $\gamma \in \Gamma$ is definable then it is in the divisible hull of $\Gamma(\mathbb{S})$. This does not mean that the definable set $\gamma^{\sharp} \sub \RV$ contains a definable point unless $\gamma \in \Gamma(\mathbb{S})$.
\end{rem}


\begin{rem}\label{GLZ:char}
By \cite[Remark~2.28]{Yin:int:expan:acvf}, if a definable function between two sets in $\Gamma$ is a $\vrv$-contraction then it is piecewise $\Z$-affine (with constant terms of the form $\vrv(t)$, where $t \in \RV$ is definable, that is, $t \in \RV(\mathbb{S})$). Moreover,  a definable bijection  between two objects of $\Gamma[k]$ is a  $\Gamma[k]$-morphism if and only if it is definably a piecewise $\mgl_k(\Z)$-transformation. The ``if'' direction is clear. For the ``only if'' direction, see \cite[Lemma~10.1]{hrushovski:kazhdan:integration:vf} or \cite[Lemma~2.29]{Yin:int:expan:acvf}.

% BEGIN REFEREE INSERTION R050
\ResolvedRefereeRemark{050}{Clarification}{Piecewise integral-affine behavior}{%
\textit{Original source line 1289.}\quad The first sentence needs ``piecewise $\mathbb Z$-affine,'' not globally ``$\mathbb Z$-linear.'' For example, $\gamma\mapsto |\gamma|$ is the $\vrv$-contraction of the definable map $t\mapsto t$ for $\vrv(t)\geq0$, $t\mapsto t^{-1}$ otherwise. Its two integral-linear formulas differ. Retain the specified restriction on the translation constants.
}{The author approves the piecewise $\Z$-affine formulation. The restriction on translation constants and the following piecewise $\mgl_k(\Z)$ characterization are retained.}
% END REFEREE INSERTION R050
\end{rem}


\begin{defn}[$\Gamma$-categories with volume forms]\label{def:Ga:cat:mu}
An object of the category $\mG[k]$ is a definable pair $(I, \omega)$, where $I \in \Gamma[k]$ and $\omega : I \fun \Gamma$ is a function. A $\mG[k]$-morphism between two objects $(I, \omega)$, $(J, \sigma)$ is a $\Gamma[k]$-morphism $F : I \fun J$ such that for all $\alpha  \in I$,
\[
\omega(\alpha) = \sigma(F(\alpha)) + \jcb_{\Gamma} F(\alpha).
\]

The category $\mG^{\fin}[k]$ is the obvious full subcategory of $\mG[k]$.
\end{defn}

\begin{nota}\label{comp:form}
For each $(U, f, \omega) \in \mgRV[k]$, write $\omega_f : U \fun \Gamma$ for the function given by $u \efun \Sigma(\vrv \circ f)(u) + \omega(u)$. Similarly, for $(I, \sigma) \in \mG[k]$, write $\sigma_I : I \fun \Gamma$ for the function given by $\gamma \efun \Sigma \gamma + \sigma(\gamma)$.
\end{nota}

%\begin{defn}[$\Gamma$-categories with counterweights]\label{def:Ga:wgt}
%The category ${\mu_0}{\Gamma[k]}$ is the full subcategory of $\mG[k]$ such that $(I, \omega) \in {\mu_0}{\Gamma[k]}$ if and only if $\omega(\alpha) + \Sigma \alpha = 0$ for all $\alpha \in I$.
%
%The category ${\mu_0}{\Gamma^c[k]}$ is the obvious full subcategory of ${\mu_0}{\Gamma[k]}$.
%\end{defn}

There is a natural map $\Gamma[*] \fun \RV[*]$ given by $I \efun (I^\sharp, \id)$ (see Notation~\ref{gamwhat}). This map induces a commutative diagram in the category of graded semirings:
\[
\bfig
 \Square(0,0)/>->`>->`>->`>->/<400>[{\gsk \Gamma^{\fin}[*]}`{\gsk \RES[*]}`{\gsk \Gamma[*]}`{\gsk \RV[*]}; ```]
\efig
\]
where all the arrows are monomorphisms. The map
\[
\gsk \RES[*] \times \gsk \Gamma[*] \fun \gsk \RV[*]
\]
determined by the assignment
\begin{equation}\label{bil:def}
([(U, f)], [I]) \efun [(U \times I^\sharp, f \times \id)]
\end{equation}
is well-defined and is clearly $\gsk \Gamma^{\fin}[*]$-bilinear. Hence it induces a $\gsk \Gamma^{\fin}[*]$-linear map
\begin{equation}\label{tensor:expre}
\Psi: \gsk \RES[*] \otimes_{\gsk \Gamma^{\fin}[*]} \gsk \Gamma[*] \fun \gsk \RV[*],
\end{equation}
which is a homomorphism of graded semirings. By the universal mapping property, groupifying a tensor product in the category of $\gsk \Gamma^{\fin}[*]$-semimodules is, up to isomorphism, the same as taking the corresponding tensor product in the category of $\ggk \Gamma^{\fin}[*]$-modules; the groupification of $\Psi$ is still denoted by $\Psi$. Similarly, there is a
$\gsk \mG^{\fin}[*]$-linear map
\begin{equation}\label{bdd:to:2bdd}
{\mu}\Psi: \gsk \mgRES[*] \otimes_{\gsk \mG^{\fin}[*]} \gsk \mG[*] \fun \gsk \mgRV[*].
\end{equation}
We shall abbreviate $\otimes_{\gsk \Gamma^{\fin}[*]}$, $\otimes_{\gsk \mG^{\fin}[*]}$ both as ``$\otimes$'' below when no confusion can arise.

\begin{prop}[{\cite[Corollary~10.3, Proposition~10.10(1)]{hrushovski:kazhdan:integration:vf}}]\label{red:D:iso}
$\Psi$ and ${\mu}\Psi$ are both isomorphisms of graded semirings.
\end{prop}

\begingroup\normalfont\small\color{blue}
\noindent\textbf{Formalization review T001: generators of the unmeasured kernel.}
Keep the original ground $\mathbb S$ fixed. Put
\[
 T=\ggk\RES[*]\otimes_{\ggk\Gamma^{\fin}[*]}\ggk\Gamma[*].
\]
The following is a mathematical presentation of $\ker\Psi$, rather than
an assertion that its complete proof is already formalized in Lean.
It concerns the unmeasured map only.

\smallskip
Take any $X\in\RES[r]$, any definable $\delta\in\Gamma$, and any
definable function $c:X\to\delta^\sharp$. No point of
$\delta^\sharp$ is required to be definable over $\mathbb S$ alone.
For any $U\in\Gamma[m]$, $m\geq1$, and coordinate vector $e_i$, put
\[
 g_{X,c,U,i}=[X]_r\otimes
       \bigl([U]_m-[U+\delta e_i]_m\bigr).
\]
Then the kernel has the homogeneous generating presentation
\[
 \boxed{\quad\ker\Psi=
   \bigl(g_{X,c,U,i}:X,c,U,i\text{ as above}\bigr)\subseteq T.\quad}
\]
This is a potentially redundant family; no finite or minimal generating
set is claimed. Equivalently, for each $X$ use every $\delta$ for which
$\delta^\sharp$ has a definable point \emph{depending on $x\in X$}.
One can make $c$ a coordinate projection by replacing $X$ with its graph,
without changing its RES isomorphism class or its degree.

\smallskip
\noindent\textbf{The generating bijections.}
Each generator maps to zero because the original-ground definable map
\[
 (x,u_1,\ldots,u_m)\longmapsto
 (x,u_1,\ldots,c(x)u_i,\ldots,u_m)
\]
is an RV bijection from $X\times U^\sharp$ to
$X\times(U+\delta e_i)^\sharp$. Its inverse divides by $c(x)$.
The finite labels need not be preserved by an unmeasured RV morphism.
If $\delta$ already has an RV representative definable over the original
ground, the displayed difference is already zero in $T$.

\smallskip
\noindent\textbf{Why these generators suffice.}
Normalize Gamma factors by integral unimodular changes of coordinates,
and transfer their fixed Gamma coordinates into RES by the existing
finite-Gamma balancing relation. Given an RV bijection of product
unions, work at the largest remaining free Gamma arity $m$. On product
neighbourhoods of full-rank types, the local generic-fibre calculation
gives the form
\[
 F(x,u)=\bigl(h(x),c(x)u^M\bigr),\qquad
 M\in\operatorname{GL}_m(\mathbb Z),\qquad
 \vrv(c(x))=\boldsymbol\delta,
\]
where $h$ is a RES bijection and the value tuple
$\boldsymbol\delta$ is constant and definable over $\mathbb S$.
The coefficients $c(x)$ are retained as functions of $x$.
The image is exactly
$h(X_0)\times(MU_0+\boldsymbol\delta)^\sharp$.
The unimodular part is already allowed in $\Gamma[*]$; its translation
part is removed one coordinate at a time by the displayed generators.
Compactness and finite product refinement match finitely many disjoint
such rectangles. The remainder has lower Gamma rank on both sides, so
induction applies. At rank zero, every Gamma factor has been absorbed
into RES and the equality is a RES isomorphism. This argument uses no
cancellation and makes no same-ground choice of the coefficients.

\smallskip
For the positive semiring, replace ``ideal generated by the differences''
with ``congruence generated by the pairs''
\[
 [X]_r\otimes[U]_m\ \sim\
 [X]_r\otimes[U+\delta e_i]_m.
\]
Groupification then gives the stated ideal, without assuming that the
positive monoid is cancellative. In particular, an arbitrary kernel
element can first be written $a-b$; equality after groupification is a
positive equality after adding a common RV class, which can be lifted
to a positive tensor by surjectivity.

\smallskip
\noindent\textbf{A single uniform annihilator.}
In the graded tensor put
\[
 \tau=1\otimes[\{0\}]_1=[\mathbb G_m]_1\otimes1.
 \qquad\text{Then}\qquad \tau\ker\Psi=0.
\]
Indeed $c$ gives $X\times0^\sharp\simeq X\times\delta^\sharp$.
Balance the latter degree-one coefficient to the Gamma side and use the
integral shear
$(\alpha,\gamma)\mapsto(\alpha,\gamma+\alpha e_i)$ on the slice
$\alpha=\delta$. Balance back to obtain
$\tau g_{X,c,U,i}=0$. Thus $\Psi[\tau^{-1}]$ is an isomorphism.
Here $\tau$ has degree one and cannot be cancelled in the original
ring. The stronger equality $\ker\Psi=\operatorname{Ann}(\tau)$ is
not asserted.

\smallskip
The previously found nonzero element is the special generator
\[
 [\{r:r^5=\rv(t)\}]_0\otimes
 \bigl([(0,\infty)]_1-[(\tfrac15,\infty)]_1\bigr),
 \qquad c(r)=r,
\]
over $\mathbb S=\mathbb C((t))$, $\vv(t)=1$.
The detailed mathematical proof, including the local rectangle and
lower-rank arguments, is recorded in the section's verification files.
The complete uniform compactness and scissor argument remains to be
exported in Lean. No HK theorem is taken as a premise, no statement
about the measured kernel is inferred, and the working manuscript has
not been changed.
\par\endgroup

Note that, however, \cite[Proposition~10.10(2)]{hrushovski:kazhdan:integration:vf} does not hold. One of the main motivations of this paper is to construct an alternative to be used in applications.

\begin{nota}\label{drop:0}
For simplicity, we often drop the constant volume form $0$ from the notation. For instance, if $\bm U$ is an object of $\RV[*]$ then it may also denote the  object $(\bm U, 0)$ of $\mgRV[*]$ with the constant volume form $0$.
\end{nota}

\begin{nota}
Let $[1] \in \gsk \RES[1]$ be the class of the singleton $\{1\}$. The class of the singleton $\{1\}$ in $\gsk \RES[0]$ is the multiplicative identity of $\gsk \RES[*]$ and hence is simply denoted by $1$. We have $[\RV^{\circ\circ}_\infty] = [\RV^{\circ\circ}] + 1$ in $\gsk \RV[{\leq} 1]$.

Let $\isp$ be the (nonhomogeneous) semiring congruence relation on $\gsk \RV[*]$ generated by the pair $([1], [\RV^{\circ\circ}_\infty])$. Let
\[
\bm P = [1]-[\RV^{\ccirc}]  \in \ggk \RV[1].
\]
The corresponding principal ideal of $\ggk \RV[*]$ is thus generated by the element $\bm P - 1$.


Similarly, let $\mgisp$ be the  semiring congruence relation on $\gsk \mgRV[*]$ generated by the pair $([1], [\RV^{\circ\circ}])$, which is homogeneous, and the corresponding principal ideal of $\ggk \mgRV[*]$ is generated by the element $\bm P$.
\end{nota}

% BEGIN REFEREE INSERTION R051
\ResolvedRefereeRemark{051}{Typo}{Spelling of homogeneous}{%
\textit{Original source line 1358.}\quad In this notation, replace ``nonhomogenous'' and ``homogenous'' by ``nonhomogeneous'' and ``homogeneous,'' respectively.
}{The author approves correcting the spellings to nonhomogeneous and homogeneous. The congruences and their generators are unchanged.}
% END REFEREE INSERTION R051

\begin{nota}\label{RVcat:tag}
For each  $\bm U = (U, f) \in  \RV[k]$, let $U_f$ be the set $\bigcup \set{f(u)^\sharp \times u \given u \in U}$. Let $\mathbb{L}_{\leq k}: \RV[{\leq}k] \fun \VF[k]$ be the map  given by $\bm U \efun U_f$. Set $\bb L = \bigcup_k \bb L_{\leq k}$.

Let $\mgL_k: \mgRV[k] \fun \mgVF[k]$ be the map given by $(\bm U, \omega) \efun (\bb L \bm U, \bb L \omega)$, where $\bb L \omega$ is the obvious function on $\bb L \bm U$ induced by $\omega$. Set $\mgL = \bigoplus_k \mgL_{k}$.
\end{nota}

Suppose that $\mathbb S$ is $(\VF,\Gamma)$-generated. The map $\bb L$ induces a surjective semiring homomorphism $\gsk \RV[*] \fun \gsk \VF_*$ and the map $\mgL$ induces a surjective graded semiring homomorphism $\gsk \mgRV[*] \fun \gsk \mgVF[*]$, see \cite[\S~4, \S~6]{hrushovski:kazhdan:integration:vf} or \cite[Corollaries~7.7, 10.6]{Yin:special:trans}; we use the same notation for these homomorphisms as well as their groupifications. We have

% BEGIN REFEREE INSERTION R053
\ResolvedRefereeRemark{053}{Error; revised}{Scope of the parameter hypothesis}{%
\textit{Original source line 1366.}\quad State here the $(\VF,\Gamma)$-generation hypothesis used in Theorem~\ref{main:prop:HK}, before asserting that lifting induces these homomorphisms. Without a parameter restriction, even well-definedness can fail. For example, let the parameters be generated by a trivially valued algebraically closed field $C$ and a named $\tau\in\RV$ of positive value. The singleton objects $\{\tau\}$ and $\{1\}$ are isomorphic in $\RV[1]$. But $\tau^\sharp$ has no definable $\VF$ point: Lemma~\ref{dcl:to:ac} places every such point in $\acl(C)\cap\VF=C$. In contrast, $1+\MM$ contains the definable point $1$, so their lifts cannot be definably bijective. Moving the stated hypothesis forward removes this obstruction.
}{The author approves stating the $(\VF,\Gamma)$-generation hypothesis before both homomorphism assertions. The lift on individual objects needs no such hypothesis. Under the standing nontrivial-ground-value-group assumption, take instead a small algebraically closed nontrivially valued subfield $F$ and name $\tau\in\RV$ with $\vrv(\tau)\notin\Gamma(F)$. Lemma~\ref{dcl:to:ac} puts every definable $\VF$ point over the generated structure in $F$. Thus $\tau^\sharp$ has no definable point, whereas $1^\sharp$ contains $1$, although the singleton objects are isomorphic in $\RV[1]$.}
% END REFEREE INSERTION R053
\[
\bb L ([1]) = [1 + \MM] = [\MM] \dand \bb L ([\RV^{\circ\circ}])=[\MM \mi 0],
\]
and hence $\bb L (\bm P - 1)=0$. Similarly, $\mgL (\bm P)=0$ since, as we have seen above, $\MM$ and $\MM \mi 0$ are in essential bijection. It so happens that these relations are the only ones needed to describe the kernels of $\bb L$ and $\mgL$.


\begin{thm}\label{main:prop:HK}
Suppose that $\mathbb{S}$ is $(\VF, \Gamma)$-generated. For each $k \geq 0$ there is a canonical isomorphism of semigroups
\[
\int_{+} : \gsk  \VF[k] \fun \gsk  \RV[{\leq}k] /  \isp
\]
such that $ \int_{+} [A] = [\bm U]/  \isp$ if and only if $[A] = [{\bb L}{\bm U}]$. These isomorphisms are  compatible with the inductive systems. Passing to  the colimit of the groupifications, we obtain a canonical isomorphism of rings
\[
\int : \ggk \VF_* \fun \ggk  \RV[*] /  (\bm P - 1).
\]
Similarly, for each $k \geq 0$ there is a canonical isomorphism of semigroups
\[
\int_{+}^\mu : \gsk  \mgVF[k] \fun \gsk  \mgRV[k] /  \mgisp
\]
such that $\int_{+}^{\mu} [\bm A] = [\bm U]/  \mgisp$  if and only if $[\bm A] = [{\mgL}{\bm U}]$.

% BEGIN REFEREE INSERTION R054
\ResolvedRefereeRemark{054}{Typo}{Measured integral notation}{%
\textit{Original source line 1386.}\quad The integral in this characterization should be $\int_+^\mu$, consistently with the immediately preceding display.
}{The author approves adding the measured superscript $\mu$ in this characterization. The measured congruence and the asserted equality of lift classes are unchanged.}
% END REFEREE INSERTION R054
Taking the direct sum of the groupifications yields a canonical isomorphism of graded rings
\[
\int^\mu : \ggk \mgVF[*] \fun \ggk  \mgRV[*] /  (\bm P).
\]
\end{thm}

This is a combination of two main theorems, Theorems~8.8 and 8.29, of \cite{hrushovski:kazhdan:integration:vf}.

\section{Uniform retraction to $\RES$}
\label{sec:retract}

\begin{defn}[Dimension-free $\RES$-categories]\label{df:res}
The objects of the category $\RES$ is obtained from $\RES[*]$ by forgetting the function $f$ in the pair $(U, f)$. Any definable bijection between two such objects is a morphism of $\RES$.

% BEGIN REFEREE INSERTION R055
\RefereeRemark{055}{Typo}{Agreement}{%
\textit{Original source line 1399.}\quad Change the verb ``is'' to ``are'' in the sentence beginning ``The objects of the category $\RES$.''
}
% END REFEREE INSERTION R055
\end{defn}

%\begin{nota}
%Denote the half interval $(0, \infty) \sub \Gamma$ simply by $\bm H$. We can associate two Euler characteristics $\chi_{g}$, $\chi_{b}$ with the $\Gamma$-sort, which are distinguished by $\chi_{g}(\bm H) = -1$ and $\chi_{b}(\bm H) = 0$. They agree on doubly bounded definable sets. Both of them shall be denoted simply by $\chi$ when no distinction is needed. See \cite[Section 9]{hrushovski:kazhdan:integration:vf} for details, where $\chi_g$ is denoted by $\chi$ and $\chi_b$ by $\chi_c$.
%\end{nota}

\begin{rem}
As  the  $\Gamma$-sort is an \omin-minimal group, there is  the \memph{geometric} \omin-minimal Euler characteristic  $\chi_g : \ggk \Gamma[*] \fun \Z$. This is roughly constructed as follows. Every definable set $X \sub \Gamma^k$ is definably bijective to a disjoint union of open cubes $\prod^k_{i= 1}(\alpha_i,\beta_i)$, where $\alpha_i, \beta_i \in \Gamma \cup \{-\infty, +\infty\}$. One sets $\chi_g(\prod^k_{i= 1}(\alpha_i,\beta_i)) = (-1)^k$ and then defines $\chi_g(X)$ by additivity. This does not depend on the chosen partition of $X$ and is invariant under definable bijections. It can also be shown that $\chi_g(X \cap [-\gamma, \gamma]^k)$ stabilizes as $\gamma$ approaches $+\infty$. The induced map, denoted by $\chi_b: \ggk \Gamma[*] \fun \Z$,  is the so-called \memph{bounded} \omin-minimal Euler characteristic. Obviously the two maps $\chi_g$, $\chi_b$  coincide on doubly bounded sets. On the other hand,  $\chi_g((0,+\infty))=-1$ whereas $\chi_b((0,+\infty))=0$. Both of them shall be denoted simply by $\chi$ when no distinction is needed.

% BEGIN REFEREE INSERTION R056
\RefereeRemark{056}{Error}{Cell dimensions in the Euler characteristic}{%
\textit{Original source line 1407.}\quad The dimension of a cell need not equal the ambient $k$: a singleton in $\Gamma^k$ already contradicts the stated decomposition into open $k$-cubes. The clean correction is to take a finite cell decomposition $X=\coprod_j C_j$ and define $\chi_g(X)=\sum_j(-1)^{\dim C_j}$. If a rectilinearization description is retained, use variable dimensions $d_j\leq k$, including $d_j=0$.
}
% END REFEREE INSERTION R056

See \cite[\S~9]{hrushovski:kazhdan:integration:vf} for more details, where $\chi_g$ is denoted by $\chi$ and $\chi_b$ by $\chi_c$.
\end{rem}


\begin{nota}\label{res:note}
Let $[\A] \in \ggk \RES$ denote the class of the affine line over the residue field. The class of the multiplicative torus $[\A] - 1$ over the residue field is written as $[\G_m]$; note that the multiplicative identity $1$ of $\ggk \RES$ is indeed the class $[1]$, but $[1]$ is not the multiplicative identity of the graded ring $\ggk \RES[*]$.

% BEGIN REFEREE INSERTION R057
\RefereeRemark{057}{Clarification}{Define the homogeneous affine-line class}{%
\textit{Original source line 1414.}\quad The graded and ungraded uses of $[\mathbb A]$ still need an explicit convention. In $\ggk\RES[*]$ set $\mathbb L_1=[\mathbb G_m]+[1]\in\ggk\RES[1]$, whereas in $\ggk\RES$ set $\mathbb L=[\mathbb G_m]+1$. Specify that every graded localization by $[\mathbb A]$ below means localization by $\mathbb L_1$, and that $!I[*]$ is generated by homogeneous degree-one differences. This is not merely cosmetic: $[\mathbb G_m]+1$ in the graded ring is inhomogeneous.
}
% END REFEREE INSERTION R057

Let $!I$ be the ideal of $\ggk \RES$ generated by the elements $[\gamma^\sharp] - [\G_m]$, where $\gamma \in \Gamma$ is definable. The quotient ring $\ggk \RES / !I$ is written as $\sggk \RES$. The ideal $!I[*]$ of $\ggk \RES[*]$ and the (graded) quotient ring $\sggk \RES[*] = \ggk \RES[*] / !I[*]$ are constructed in the same way.

The ideal of $\ggk \mgRES[*]$ generated by the  elements
\[
[(\G_m, \alpha + \beta)] - [(\alpha^\sharp, \beta)] \in \ggk \mgRES[1],
\]
where $\alpha, \beta \in \Gamma$ are definable, is denoted by $!\mu I[*]$, and the (graded) quotient ring $\ggk \mgRES[*] / !\mu I[*]$ by $\sggk \mgRES[*]$.
\end{nota}


The quotient maps from ``$\ggk$'' to ``$\sggk$'' will all be denoted by $\iota$, and will often be omitted when it is clear in the context that the latter one is meant. We will also use the same notation for elements when passing from the former to the latter.

Note that the category $\mgRES[*]$ and the ring $\sggk \mgRES[*]$  will be dismissed after Proposition~\ref{prop:eu:retr:k} as inadequate for our purpose, and a new construction is introduced thereafter.

\subsection{Without volume forms}\label{subs:uni:retr:novol}

The discussion here more or less follows \cite[\S~2.5]{hru:loe:lef} (also see \cite[Theorem~10.5]{hrushovski:kazhdan:integration:vf}).

For $0 \leq l \leq n$, consider the maps
\[
\epsilon^{n,l}_g : \ggk \RES[n{-}l] \otimes \ggk \Gamma[l] \fun \sggk \RES[n]
\]
given by
\begin{equation}\label{kill:gamma}
x \otimes y \efun \chi_g(y) x [\G_m]^l.
\end{equation}
Note that here we are forced to pass from ``${\ggk}$'' to ``$\sggk$'' in the target by the tensor $\otimes_{\gsk \Gamma^{\fin}[*]}$, for otherwise $\epsilon^{k,l}_g$ is not well-defined. Then $\bigoplus_{0 \leq l \leq n} \epsilon^{n,l}_g$ and $\Psi^{-1}$ induce a homomorphism

% BEGIN REFEREE INSERTION R058
\RefereeRemark{058}{Clarification}{Tensor products of homogeneous pieces}{%
\textit{Original source line 1442.}\quad Strictly speaking, an individual homogeneous piece such as $\ggk\RES[n-l]$ is not a module over the whole graded ring $\ggk\Gamma^{\fin}[*]$. Define the displayed component maps first on $\ggk\RES[n-l]\otimes_{\mathbb Z}\ggk\Gamma[l]$, then show that their sum descends through the homogeneous balancing relations in the full graded tensor product. The quotient $!I[*]$ gives exactly the required compatibility.
}
% END REFEREE INSERTION R058
\[
\ggk \RV[n] \fun \sggk \RES[n].
\]
Composing this with the homomorphism
\[
\epsilon^{n,m}_g : \sggk \RES[n]  \fun \sggk \RES[m] : x \efun x [\A]^{m-n}
\]
and taking the direct sum of the resulting homomorphisms over all $n$ with $n \leq m$, we obtain a homomorphism
\[
\epsilon^m_g : \ggk \RV[{\leq} m] \fun \sggk \RES[m].
\]
These homomorphisms $\epsilon^m_g$, $m = 0, 1, \ldots$, are compatible in the obvious sense. Consequently, they give rise to a homomorphism
\[
\bb E_g : \ggk \RV[*] \fun \sggk \RES[*][[\A]^{-1}]
\]
whose range is precisely the zeroth piece $(\sggk \RES[*][[\A]^{-1}])_0$ of the $\Z$-graded ring on the righthand side. We have
\begin{equation}\label{eg:isp:van}
\bb E_g([\RV^{\ccirc}] + 1 - [1]) = - [\G_m][\A]^{-1} + 1 - [\A]^{-1} = 0.
\end{equation}

% BEGIN REFEREE INSERTION R059
\RefereeRemark{059}{Error}{Missing degree-one point factor}{%
\textit{Original source line 1461.}\quad In the graded target, $\mathbb E_g([1])=[1][\mathbb A]^{-1}$, not $[\mathbb A]^{-1}$. Thus the correct computation is $-[\mathbb G_m][\mathbb A]^{-1}+1-[1][\mathbb A]^{-1}=0$, using the homogeneous identity $[\mathbb A]=[\mathbb G_m]+[1]$. As written, the final term has degree $-1$, while the first two have degree $0$. The formula $\mathbb E_g([1])=[\mathbb A]^{-1}$ becomes correct only after forgetting the grading.
}
% END REFEREE INSERTION R059
Thus $\bb E_g$ induces an eponymous homomorphism
\[
\bb E_g : \ggk \RV[*] / (\bm P - 1) \fun \sggk \RES[*][[\A]^{-1}].
\]

There is a parallel construction, where  we replace $\chi_g$ with $\chi_b$ and $[\A]^{-1}$ with $[1]^{-1}$. The latter replacement is needed so to make $(\bm P - 1)$ vanish as in (\ref{eg:isp:van}). The resulting homomorphism is
\begin{equation}\label{EB}
\bb E_b : \ggk \RV[*] / (\bm P - 1) \fun \sggk \RES[*][[1]^{-1}],
\end{equation}
whose range is precisely the zeroth piece  $(\sggk \RES[*][[1]^{-1}])_0$.

\begin{rem}\label{eb:no:eg}
The zeroth graded piece $(\sggk \RES[*][[\A]^{-1}])_0$ of $\sggk \RES[*][[\A]^{-1}]$ is canonically isomorphic to the colimit  of the groups $\sggk \RES[n]$ with the transition morphisms $\epsilon^{n,m}_g $, which is actually what appears in the construction above. So there is an injective homomorphism from $(\sggk \RES[*][[\A]^{-1}])_0$ to the similar-looking ring $\sggk \RES[[\A]^{-1}]$, induced by the assignment $x \efun x[\A]^{-n}$ for $x \in \sggk \RES[n]$. Note that this homomorphism is not surjective.

% BEGIN REFEREE INSERTION R060
\RefereeRemark{060}{Error; revised}{The asserted injection is false}{%
\textit{Original source line 1474.}\quad The natural homomorphism in this sentence need not be injective, already over an untwisted residue field $\mathbb C$. To distinguish the two constructions precisely, let $R_n$ be the Grothendieck group of complex constructible sets of dimension at most $n$, with scissor relations imposed \emph{within that category}, and set $R=\bigoplus_{n\geq0}R_n$. This is the graded $\RES$ ring in the untwisted case; $R_n$ is not by definition the $n$th dimension-filtration subgroup of the ordinary Grothendieck ring. Let $\mathbb L_1=[\mathbb A^1]_1$. Ito--Miura--Okawa--Ueda give smooth projective threefolds $X,Y$ that are not stably birational but satisfy $([X]-[Y])\mathbb L=0$ in the ordinary ring (\href{https://arxiv.org/pdf/1606.04210}{Theorems~1.1 and~1.3}). Thus $z=([X]_3-[Y]_3)\mathbb L_1^{-3}$ maps to zero in $K_0(\mathrm{Var}_{\mathbb C})[\mathbb L^{-1}]$. Yet $z\neq0$: for every $m\geq0$, the additive invariant on $R_{3+m}$ recording its top-dimensional birational classes distinguishes $X\times\mathbb A^m$ and $Y\times\mathbb A^m$. Hence no homogeneous power of $\mathbb L_1$ kills the numerator. To fit the paper's parameter conventions, take the full algebraically closed Puiseux field over $\mathbb C$ as the parameter field. Every definable value then has a named RV lift, so all $!I$ generators already vanish and $\RES[n]$ identifies with the indicated constructible-set category. Replace ``injective homomorphism'' by ``natural homomorphism.'' The comparison of $\mathbb E_g$ with $\mathbb E_b$ uses only that natural map, and the point-localized identification in the next paragraph is unaffected.
}
% END REFEREE INSERTION R060

Similarly, there is an injective homomorphism $(\sggk \RES[*][[1]^{-1}])_0 \fun \sggk \RES[[1]^{-1}] \cong \sggk \RES$, and this is indeed an isomorphism.

If the codomain of $\bb E_b$ is changed to $\sggk \RES[[\A]^{-1}]$ in the obvious way then it may be compared with $\bb E_g$. They are different since
\begin{equation}\label{eb:neq:eg:counter}
\bb E_b([1]) = 1 \dand \bb E_g([1]) = \bb E_g([\RV^{\circ\circ}]) + 1 = -[\G_m][\A]^{-1} + 1 = [\A]^{-1}.
\end{equation}
We can equalize them by forcing $[\A] = 1$ (hence $\bb E_{g}(x \otimes y) = \bb E_{b}(x \otimes y) = 0$ if $y \notin \ggk \Gamma[0]$). The (complex version of the) resulting homomorphism is indeed the one constructed in \cite[(2.5.7)]{hru:loe:lef}.
\end{rem}

The homomorphism $\bb E_b$ will be used in constructing motivic Milnor fiber below, but not $\bb E_g$. It is not that  $\bb E_g$ would not work; it is just that, conceptually, $\bb E_b$ is a better choice,  see Remarks~\ref{eb:not:eg} and \ref{rem-comp-milnor-fib}.

% BEGIN REFEREE INSERTION R168
\RefereeRemark{168}{Typo; new}{Missing article}{%
\textit{Original source line 1485.}\quad Write ``used in constructing \emph{the} motivic Milnor fiber below.''
}
% END REFEREE INSERTION R168


%\begin{rem}\label{bbE:dim}
%In retrospect, since ``dimensions'' are lost in $\sggk \RES[[\A]^{-1}]$, in computing $\bb E_g([\bm U])$ for $\bm V = (V, g) \in \RV[*]$, we may treat $x$ in (\ref{kill:gamma}) simply as an element in $\ggk \RES$, similarly for $\bb E_b$.
%\end{rem}

\subsection{With volume forms}
Let $\fn(\Gamma, \Z)$ be the set of functions $f: \Gamma \fun \Z$ such that the range of $f$ is finite and $f^{-1}(m)$ is definable for every $m \in \Z$. Addition in $\fn(\Gamma, \Z)$ is defined in the obvious way. Multiplication in $\fn(\Gamma, \Z)$ is given by the convolution product as follows. For $f, g \in \fn(\Gamma, \Z)$ and $\gamma \in \Gamma$, let $h_\gamma$ be the function $\Gamma \fun \Z$ given by $\alpha \efun f(\alpha)g(\gamma - \alpha)$. The range of $h_\gamma$ is obviously finite. For each $m \in \Z$, $h_\gamma^{-1}(m)$ is a finite disjoint union of sets of the form
\[
f^{-1}(m') \cap (\gamma - g^{-1}(m'')),
\]
where $m' m'' = m$, and hence is $\gamma$-definable. Let
\[
 h_\gamma^* = \sum_{m} m \chi(h_\gamma^{-1}(m)) \in \Z
\]
and $f * g : \Gamma \fun \Z$ be the function  given by $\gamma \efun h_\gamma^*$. It is not hard to see that, by \omin-minimality in the $\Gamma$-sort, $f * g = g * f$ and it belongs to $\fn(\Gamma, \Z)$. Thus $\fn(\Gamma, \Z)$ is a commutative ring.

% BEGIN REFEREE INSERTION R061
\RefereeRemark{061}{Clarification; revised}{Associativity of convolution}{%
\textit{Original source line 1501.}\quad Optional proof detail: add that $p_0$ is the multiplicative identity and that associativity follows by checking the additive generators in Notation~\ref{gam:equ:cons}. For example, writing $rq_\beta=\eta r$, the only nontrivial $r,q,q$ check reduces to $\eta^2=-\eta$, which holds for both $\eta=-1$ and $\eta=0$. The ring assertion is correct for both choices of $\chi$.
}
% END REFEREE INSERTION R061

\begin{nota}\label{gam:equ:cons}
In $\fn(\Gamma, \Z)$, let $r$ be the constant function $1$ and, for each definable element $\gamma \in \Gamma$, let $p_{\gamma}$,  $q_{\gamma}$ be the characteristic functions of $\gamma$, $(\gamma, \infty)$, respectively. By \omin-minimality, $\fn(\Gamma, \Z)$ is generated as a $\Z$-module by elements of the forms $r$, $p_\gamma$, $q_\gamma$. We have the following equalities:
\[
r p_\beta   = r, \quad p_\alpha p_\beta = p_{\alpha+ \beta}, \quad p_\alpha q_\beta  = q_{\alpha+ \beta}, \quad q_\alpha q_\beta = - q_{\alpha+ \beta}.
\]
In addition,
\begin{itemize}
  \item if $\chi = \chi_{g}$ then $r^2 = r q_\beta = - r$,
  \item if $\chi = \chi_{b}$ then $r^2 = r$ and $r q_\beta = 0$.
\end{itemize}
\end{nota}

\begin{rem}\label{int:Gam}
For each $\bm I = (I, \mu) \in \mG[k]$, let $\mu_I : I \fun \Gamma$ be the function given by $\gamma \efun \mu(\gamma) + \Sigma\gamma$ and $\lambda_{\bm I} : \Gamma \fun \Z$  the function given by $\gamma \efun \chi(\mu^{-1}_I(\gamma))$. By \omin-minimality, the range of $\lambda_{\bm I}$ is finite and, for every $m \in \Z$, $\lambda_{\bm I}^{-1}(m)$ is definable. So $\lambda_{\bm I} \in \fn(\Gamma, \Z)$. If $F :(I, \mu) \fun (J, \sigma)$ is a $\mG[k]$-morphism then $F$ restricts to a bijection from $\mu^{-1}_I(\gamma)$ to $\sigma_J^{-1}(\gamma)$ for every $\gamma \in \Gamma$ and hence $\lambda_{\bm I} = \lambda_{\bm J}$ and we may write $\lambda_{\bm I}$ as $\lambda_{[\bm I]}$. Let $\bm J = (J, \sigma) \in \mG[l]$. If $k=l$ then clearly $\lambda_{[\bm I]} + \lambda_{[\bm J]} = \lambda_{[\bm I]+[\bm J]}$. An easy computation shows that, for every $\gamma \in \Gamma$,
\[
(\lambda_{[\bm I]} * \lambda_{[\bm J]})(\gamma) = \chi( (\mu \sigma)^{-1}_{I \times J}(\gamma)),
\]
where $\mu \sigma : I \times J \fun \Gamma$ is the function given by $(\alpha, \beta) \efun \mu(\alpha) + \sigma(\beta)$, and hence $\lambda_{[\bm I]} * \lambda_{[\bm J]} = \lambda_{[\bm I][\bm J]}$. Let $\Z\Gamma$ be the group ring of $\Gamma(\mathbb{S})$ over $\Z$, which is viewed as the subring of $\fn(\Gamma, \Z)$ generated, as a $\Z$-module, by $p_{\gamma}$. We have $\ggk \mG[0] \cong \Z\Gamma$. It follows that the map

% BEGIN REFEREE INSERTION R062
\RefereeRemark{062}{Error}{The group ring uses definable value-group constants}{%
\textit{Original source line 1520.}\quad The degree-zero coefficient group is the group of \emph{definable} elements of $\Gamma$, not necessarily $\Gamma(\mathbb S)$. With the present conventions, put $\Gamma_{\mathrm{def}}=\Gamma\cap\dcl^{\mathrm{eq}}(\mathbb S)$ and replace $\mathbb Z\Gamma$ throughout this construction by $\mathbb Z[\Gamma_{\mathrm{def}}]$; in the usual valued-field-generated setting $\Gamma_{\mathrm{def}}=\mathbb Q\otimes\Gamma(\mathbb S)$. For example, over $\mathbb C((t))$, the singleton object of $\mu\Gamma[0]$ with constant volume form $\operatorname{val}(t)/2$ is allowed and maps to $p_{1/2}$, although $\Gamma(\mathbb S)=\mathbb Z$. Thus the stated target misses legitimate degree-zero classes. This distinction is already acknowledged in the remark following Definition~\ref{def:Ga:cat}.
}
% END REFEREE INSERTION R062
\[
\lambda : \ggk \mG[*] \fun \Z\Gamma \oplus X \fn(\Gamma, \Z)[X]
\]
determined by the assignment
\[
[\bm I] \efun \lambda_{[\bm I]}X^k, \quad \bm I \in \mG[k],
\]
is a graded ring homomorphism. Of course, there are two such homomorphisms $\lambda_g$ and $\lambda_b$, corresponding to the two choices $\chi = \chi_{g}$ and $\chi = \chi_{b}$.
\end{rem}


%recall Notation~\ref{comp:form}



We claim that the assignments
\begin{equation}\label{Gam:ass}
p_\gamma X^k \efun  [\G_m]^{k-1}[(\G_m, \gamma)], \quad q_\gamma X^k \efun  [\G_m]^{k-1}[(\{1\}, \gamma)], \quad r X^k \efun 0
\end{equation}
induce a graded ring homomorphism
\[
\psi : \Z\Gamma \oplus X\fn(\Gamma, \Z)[X] \fun \sggk \mgRES[*] / ([\A]).
\]
Note that $\Z\Gamma$ may be regarded as a subring of $\ggk \mgRES[0]$ by identifying $p_\gamma$ with $[(\infty, \gamma)]$, which is the intended interpretation of the first assignment for $k = 0$. For the claim, it suffices to check that $\psi$ obeys the equational constraints for the generators above in Notation~\ref{gam:equ:cons}, which are all straightforward (the equality $\psi(q_\alpha) \psi( q_\beta) = \psi(- q_{\alpha + \beta})$ holds since $[\A] = [\G_m] +[1] = 0$ is forced).

% BEGIN REFEREE INSERTION R063
\RefereeRemark{063}{Typo}{Restore the degrees in the multiplicativity check}{%
\textit{Original source line 1544.}\quad The symbols $q_\alpha$ alone are not degree-zero elements of the displayed domain. State the check as $\psi(q_\alpha X^k)\psi(q_\beta X^l)=\psi(-q_{\alpha+\beta}X^{k+l})$ for $k,l\geq1$. Also state explicitly that the second and third assignments in~\eqref{Gam:ass} apply only for $k\geq1$; the separate interpretation at $k=0$ applies to $p_\gamma$.
}
% END REFEREE INSERTION R063


Let $\gh$ denote the element $[(\RV^{\ccirc},  -\vrv)] \in \ggk \mgRV[1]$.

\begin{prop}\label{prop:eu:retr:k}
There are two graded ring homomorphisms
\[
  \mgE_{g}, \mgE_{b}: \ggk \mgRV[*] \fun   \sggk \mgRES[*] / ([\A])
\]
such that
\begin{itemize}
  \item $\bm P \in \ggk \mgRV[1]$ vanishes under both of them,
  \item for all $x \in \ggk \mgRES[*]$, $\mgE_{g}(x) = \mgE_{b}(x) = x/([\A])$,
  \item they are distinguished by $\mgE_{g}(\gh) = -[1]/([\A])$ and $\mgE_{b}(\gh) = 0$.
\end{itemize}
\end{prop}

% BEGIN REFEREE INSERTION R064
\RefereeRemark{064}{Error}{Sign of the geometric retraction on h}{%
\textit{Original source line 1560.}\quad The construction gives the opposite sign in the third bullet. For $\bm I=((0,\infty),-\mathrm{id})$, the corrected weight $\mu_I$ is identically zero, so $\lambda_g([\bm I])=-p_0X$. Consequently $\mu\mathbb E_g(\mathfrak h)=\psi(-p_0X)=-[\mathbb G_m]=+[1]$ modulo $([\mathbb A])$. The bounded value is indeed zero. Replace $-[1]/([\mathbb A])$ by $+[1]/([\mathbb A])$. This confirms the sign concern in the supplied human referee report; repeating the sign printed in the cited HK theorem does not fix its incompatibility with the present definitions.
}
% END REFEREE INSERTION R064
\begin{proof}
The product map
\[
(\iota /([\A])) \times (\psi \circ \lambda) : \ggk \mgRES[*] \times  \ggk \mG[*] \fun    \sggk \mgRES[*] /([\A])
\]
is $\ggk \mG^{\fin}[*]$-bilinear (because of  the ideal $!\mu I[*]$ in the target). We have
\[
{\mu}\Psi^{-1}(\bm P) = [1]-[(0, \infty)]  \dand {\mu}\Psi^{-1}(\gh) = [((0, \infty), -\id)].
\]
The existence of the desired homomorphisms follow from a straightforward computation.

% BEGIN REFEREE INSERTION R065
\RefereeRemark{065}{Typo}{Agreement}{%
\textit{Original source line 1570.}\quad Change the final verb ``follow'' to ``follows,'' since the subject is ``The existence.''
}
% END REFEREE INSERTION R065
\end{proof}

The discussion above is more or less an expanded version of the proof of \cite[Theorem~10.11]{hrushovski:kazhdan:integration:vf}.

\begin{rem}
Alternatively, for $\chi = \chi_{g}$, we may replace $r X^k \efun 0$ in (\ref{Gam:ass}) with $r X^k \efun [\G_m]^{k-1}[1]$, that is, sending $r X^k$ and $q_0 X^k$ to the same element.  Then $\psi$ is still a graded ring homomorphism provided that we enlarge the ideal $!\mu I[*]$ by adding the elements $[1] - [(\{1\}, \gamma)] \in \ggk \mgRES[1]$. Consequently, for instance, we can construct a ring homomorphism
\begin{equation}\label{res:forget}
\ggk \mgRV[*] \fun \sggk \mgRES[*] / ([\A]) \to^{\phi} \sggk \RES[*] / ([\A]),
\end{equation}
where the homomorphism $\phi$ is induced by the obvious forgetful functor.

% BEGIN REFEREE INSERTION R066
\RefereeRemark{066}{Clarification; revised}{Name the enlarged quotient}{%
\textit{Original source line 1580.}\quad Optional notation clarification: the preceding sentence already stipulates enlargement of $!\mu I[*]$. In the alternative construction, name the enlarged ideal $J$ and write the middle term of~\eqref{res:forget} as $\ggk\mgRES[*]/(J+([\A]))$, so that the old and enlarged quotients are visibly distinguished. Read with that stipulated enlargement, the construction is consistent. The original assignment $rX^k\mapsto0$ needs no such additional relations.
}
% END REFEREE INSERTION R066

We shall not, however, pursue this further, since taking the quotient by the ideal $([\A])$ is detrimental to our purpose here. To avoid it, we need to work with the categories of doubly bounded objects.
\end{rem}

\begin{defn}[Doubly bounded $\RV$-categories]\label{db:RV}
An object $(U, f) \in \RV[k]$ is \memph{doubly bounded} if $U$ is. Denote the full subcategory of $\RV[k]$ of doubly bounded objects by $\RV^{\db}[k]$, and similarly for the categories $\mgRV^{\db}[k]$, $\mG^{\db}[k]$.
\end{defn}




\begin{rem}\label{tensor:bd}
By Lemma~\ref{db:to:db}, if $(U, f, \omega) \in \mgRV^{\db}[k]$ then $f(U)$, $\omega(U)$ are both doubly bounded as well. For essentially the same reason, ${\mu} \Psi$ restricts to an isomorphism
\[
{\mu} \Psi^{\db}: \gsk \mgRES[*] \otimes \gsk \mG^{\db}[*] \fun \gsk \mgRV^{\db}[*].
\]
\end{rem}

\begin{nota}\label{db:generator}
In $\fn(\Gamma, \Z)$, for each definable element $\gamma \in \Gamma \mi 0$, let $o_{\gamma}$ be the characteristic function of $(0, \gamma)$; also set $o_0 = -p_0$. As a $\Z$-module, $\fn^{\db}(\Gamma, \Z)$ is generated by elements of the forms $p_\gamma$, $o_\gamma$. Of course $o_{\gamma} = q_0 - q_\gamma - p_\gamma$ if $\gamma> 0$, etc. An easy  computation shows that

% BEGIN REFEREE INSERTION R067
\RefereeRemark{067}{Error}{Negative interval generators and bounded support}{%
\textit{Original source line 1600.}\quad For $\gamma<0$, the interval $(0,\gamma)$ is empty, so the definition as written cannot yield the claimed generators or quotient: the relation $o_\gamma+p_0$ would even kill the unit $p_0$. Define $o_\gamma=\mathbf1_{(\gamma,0)}$ when $\gamma<0$, retaining $o_\gamma=\mathbf1_{(0,\gamma)}$ when $\gamma>0$ and $o_0=-p_0$. Also define $\fn^{\db}(\Gamma,\mathbb Z)$ explicitly as the subgroup of functions whose nonzero locus is doubly bounded. These conventions make the subsequent formulas meaningful.
}
% END REFEREE INSERTION R067
\begin{itemize}
  \item $p_\alpha o_\beta$ equals $o_{\alpha + \beta} - o_\alpha - p_\alpha$ or $- o_{\alpha + \beta} + o_\alpha - p_{\alpha + \beta}$ or $ o_{\alpha + \beta} + o_\alpha + p_{0}$,
  \item $o_\alpha o_\beta$ equals $- o_{\alpha + \beta}$ or $- o_{\alpha} - o_\beta - p_{0}$.
\end{itemize}
\end{nota}

% BEGIN REFEREE INSERTION R068
\RefereeRemark{068}{Clarification}{Specify which interval identity applies}{%
\textit{Original source line 1605.}\quad After correcting the negative-interval convention, the alternatives can be made unambiguous as follows. Set $a=\min(\alpha,\alpha+\beta)$ and $b=\max(\alpha,\alpha+\beta)$. For $\beta\neq0$, $p_\alpha o_\beta$ is the characteristic function of $(a,b)$: it is $o_b-o_a-p_a$ if $0\leq a<b$, $o_a-o_b-p_b$ if $a<b\leq0$, and $o_a+o_b+p_0$ if $a<0<b$. For $\beta=0$ it equals $-p_\alpha$. Likewise $o_\alpha o_\beta=-o_{\alpha+\beta}$ when $\alpha,\beta\geq0$ or $\alpha,\beta\leq0$, and equals $-o_\alpha-o_\beta-p_0$ when $\alpha$ and $\beta$ have opposite strict signs.
}
% END REFEREE INSERTION R068

\begin{nota}\label{db:gam:fn}
The homomorphism
\begin{equation*}
\lambda : \ggk \mG^{\db}[*] \fun \Z\Gamma \oplus X \fn^{\db}(\Gamma, \Z)[X]
\end{equation*}
is constructed as before (there is only one such homomorphism now since  $\chi_{g}$, $\chi_{b}$ agree on doubly bounded sets).
\end{nota}


\begin{nota}\label{pgammabis}
Recall that $\RV^{\ccirc}_\gamma = \rv(\MM_{\gamma} \mi 0)$ (Notation~\ref{defn:disc}). For each definable $\gamma \in \Gamma^+_0$, let
\[
\bm P_\gamma = [\RV^{\ccirc} \mi \RV^{\ccirc}_{\gamma}] + [\{t_\gamma\}] - [1] \in \ggk \RV^{\db}[1],
\]
where $t_\gamma \in \gamma^\sharp$ is any definable point.  It also stands for the corresponding element  in $\ggk \mRV^{\db}[1]$ (with the constant volume form $0$). Note that, although $[\{t_\gamma\}] = [1]$ in $\ggk \RV^{\db}[1]$,  $[\{t_\gamma\}] \neq [1]$ in $\ggk \mRV^{\db}[1]$  unless $\gamma = 0$.

% BEGIN REFEREE INSERTION R069
\RefereeRemark{069}{Error}{The radius must admit a definable RV lift}{%
\textit{Original source line 1621.}\quad Not every definable $\gamma$ admits the definable point $t_\gamma$ required here. Restrict the definition to nonnegative $\gamma$ for which $\gamma^\sharp$ contains an $\mathbb S$-definable point, equivalently $\gamma\in\vrv(\RV(\dcl\mathbb S))$. In the later valued-field-generated setting this is the value group of the parameter field, rather than its divisible hull. For example, over $\mathbb C((t))$, $\gamma=1/2$ is definable but there is no definable RV element of that value. Make the same restriction in the generating set of $(\bm P_\Gamma)$. The later specialization to integral radii should refer back to this restriction.
}
% END REFEREE INSERTION R069

Clearly  $\bm P_\gamma$ does not depend on the choice of $t_\gamma \in \gamma^\sharp$. The ideal of $\ggk \mRV^{\db}[*]$   generated by the elements $\bm P_\gamma$ is denoted by $(\bm P_\Gamma)$. The images of $(\bm P_\Gamma)$ under the obvious forgetful homomorphisms
\[
\ggk  \mgRV^{\db}[*] \fun \ggk \RV[*], \quad \ggk  \mgRV^{\db}[*] \fun \ggk  \mgRV[*]
\]
are contained in $(\bm P - 1)$, $(\bm P)$, respectively,

% BEGIN REFEREE INSERTION R070
\RefereeRemark{070}{Typo}{Closing punctuation}{%
\textit{Original source line 1627.}\quad End this sentence with a full stop, not a comma.
}
% END REFEREE INSERTION R070
\end{nota}

\begin{rem}\label{db:gam:res}
Let $!P$ denote the ideal of $\fn^{\db}(\Gamma, \Z)$ generated by the elements $p_\gamma - p_0$, $o_\gamma + p_0$. By the computation in Notation~\ref{db:generator}, $\fn^{\db}(\Gamma, \Z) / !P \cong \Z$; the quotient homomorphism is denoted by $\chi$ because a simple computation shows that if $\bm I = (I, \mu) \in \mG^{\db}[*]$ then indeed $\lambda_{\bm I} / !P = \chi(I)$. Now it is straightforward to check that the assignments
\begin{equation}\label{Gam:ass:bd}
p_\gamma X^k \efun  [\G_m]^{k}, \quad o_\gamma X^k \efun  - [\G_m]^{k}
\end{equation}
induce a graded ring homomorphism
\[
\psi^{\db} : \Z\Gamma \oplus X\fn^{\db}(\Gamma, \Z)[X] \fun \sggk \RES[*].
\]
In a nutshell, the  homomorphism $\psi^{\db} \circ \lambda$ is a ``forgetful'' map  given by $[(I, \mu)] \efun \chi(I)[\G_m]^{k}$, $(I, \mu) \in \mG^{\db}[k]$.
\end{rem}


The forgetful homomorphism $\phi : \ggk \mgRES[*] \fun \sggk \RES[*]$ is similar to the one in  (\ref{res:forget}).

\begin{prop}\label{prop:eu:retr:k:db}
There is a graded ring homomorphism
\[
  \mgE^{\db}: \ggk \mgRV^{\db}[*] \fun \sggk \RES[*]
\]
such that $(\bm P_\Gamma)$ vanishes and, for all $x \in \ggk \mgRES[*]$, $\mgE^{\db}(x) = \phi(x)$.
\end{prop}
\begin{proof}
The homomorphism $\phi \times (\psi^{\db} \circ \lambda)$ is still $\ggk \mG^{\fin}[*]$-bilinear. We have
\[
({\mu} \Psi^{\db})^{-1}(\bm P_\gamma) = [(0, \gamma)] + [\{\gamma\}] + [\{t_\gamma\}] - [1],
\]
which vanishes under $\phi \otimes (\psi^{\db} \circ \lambda)$.

% BEGIN REFEREE INSERTION R071
\RefereeRemark{071}{Error}{The zero-radius case in the proof}{%
\textit{Original source line 1657.}\quad The displayed decomposition and its vanishing computation apply only for $\gamma>0$. At $\gamma=0$, $\bm P_0=0$, but the displayed right side leaves the nonzero term $[\{0\}]$ in the $\Gamma[1]$ factor, whose image is $[\mathbb G_m]$. Handle $\gamma=0$ separately, and for $\gamma>0$ write the calculation explicitly: $\chi((0,\gamma))[\mathbb G_m]+\chi(\{\gamma\})[\mathbb G_m]+[1]-[1]=-[\mathbb G_m]+[\mathbb G_m]=0$. The artificial convention $o_0=-p_0$ concerns functions; it does not change the actual class of the empty interval $(0,0)$.
}
% END REFEREE INSERTION R071
\end{proof}

The composition of $\mgE^{\db}$ with the obvious forgetful homomorphism $\sggk \RES[*] \fun \sggk \RES$ is denoted by $\bb E^{\diamond}$.  By Remark~\ref{db:gam:res}, the diagram commutes:
\begin{equation}\label{edb:eb:com}
\bfig
 \Vtriangle(0,0)/->`->`->/<400, 400>[{\ggk \mgRV^{\db}[*]}`{\ggk \RV[*]}`\sggk \RES; `\bb E^{\diamond}`\bb E_b]
\efig
\end{equation}

% BEGIN REFEREE INSERTION R072
\RefereeRemark{072}{Clarification; revised}{Identify the map labelled E\_b}{%
\textit{Original source line 1665.}\quad Optional clarification of the diagram's shorthand: the right arrow is the quotient map $\ggk\RV[*]\to\ggk\RV[*]/(\bm P-1)$, followed by~\eqref{EB} and the degree-zero identification $(\sggk\RES[*][[1]^{-1}])_0\cong\sggk\RES$ from Remark~\ref{eb:no:eg}. This is a valid composite. Its point-localized identification is unaffected by the affine-localized injectivity issue noted above.
}
% END REFEREE INSERTION R072
which may serve as an alternative and more direct (but less comprehensive) construction of $\bb E^{\diamond}$.

\begin{rem}\label{eb:not:eg}
Changing $\sggk \RES$ to $\sggk \RES[[\A]^{-1}]$ in (\ref{edb:eb:com}), we see that  $\bb E_g$ does not commute with $\bb E^{\diamond}$: by (\ref{eb:neq:eg:counter}),
\[
\bb E_g([1]) = [\A]^{-1} \quad \text{whereas} \quad \bb E^{\diamond}([1]) = \bb E_b([1]) = 1.
\]
This is but an example of a general phenomenon, see Remark~\ref{eg:eb:factor}.


We can certainly redefine $\bb E^{\diamond}$ as the composition of $\mgE^{\db}$ with  the homomorphism $\sggk \RES[*] \fun \sggk \RES[[\A]^{-1}]$ implicit in Remark~\ref{eb:no:eg} and thereby make it compatible with $\bb E_g$, but as we have mentioned above, $\bb E_b$ is ultimately a better choice; see Remark~\ref{rem-comp-milnor-fib}.
\end{rem}

\section{Interlude: Reduced cross-section and Fubini theorems}\label{csndag}

A \memph{cross-section} of $\Gamma$ is a group homomorphism $\csn : \Gamma \fun \VF^{\times}$ such that $\vv \circ \csn = \id$. The corresponding \memph{reduced cross-section} of $\Gamma$ is the function
\[
\rcsn = \rv \circ \csn : \Gamma \fun \RV^{\times}.
\]
These are usually augmented by $\csn(\infty) = 0$ and $\rcsn(\infty) = \infty$.

We can add a reduced cross-section $\rcsn$ to the language $\lan{}{RV}{}$,  denoted by $\lan{}{RV}{}^\dag$ the extension, and consider the corresponding integration theory; this has been worked out in \cite{Yin:int:expan:acvf}. We shall, however, only need a few facts about definable sets in $\RV$ in this setting. Strictly speaking, this extra ingredient is not really needed, but it does help simplify certain computations.

% BEGIN REFEREE INSERTION R073
\RefereeRemark{073}{Typo}{Verb form}{%
\textit{Original source line 1687.}\quad Replace ``denoted by $\lan{}{RV}{}^\dag$ the extension'' by ``denote the extension by $\lan{}{RV}{}^\dag$''. This correction from the human report remains outstanding.
}
% END REFEREE INSERTION R073

In this section we work in an $\lan{}{RV}{}^\dag$-expansion $\bb U^\dag$ of $\bb U$. Definability, if unqualified, is interpreted accordingly.

% BEGIN REFEREE INSERTION R074
\RefereeRemark{074}{Clarification; revised}{Parameter base for the expansion}{%
\textit{Original source line 1689.}\quad Specify separately the original parameter base $\mathbb S$ and the base used for the $\lan{}{RV}{}^\dag$ expansion. For the original-language definability assertions below, one needs the appropriate closure under the reduced section. For example, if $a\in\RV(\mathbb S)$ has value $\alpha$, then $\{a/\rcsn(\alpha)\}$ is expanded-language definable and has value $0$, but need not be original-language definable over $\mathbb S$ if the section introduces a new residue coefficient. Closing the original base under the section would also change its Galois group. The intended two-base convention should therefore be explicit.
}
% END REFEREE INSERTION R074

\begin{defn}
A \memph{$\Gamma$-partition} of a definable set $A$ in $\RV$ is a definable function $p : A \fun \Gamma_{\infty}^l$ such that, for all $\gamma \in \Gamma_{\infty}^l$, $\vrv(A_{\gamma})$ is a singleton and is $\rcsn(\gamma)$-$\lan{}{RV}{}$-definable. If $p$ is a $\Gamma$-partition of $A$ then the \memph{$\RV^\dag$-dimension} of $p$, denoted by $\dim_{\RV^\dag} (p)$, is the number $\max \set{\dim_{\RV} (A_{\gamma}) \given \gamma \in \Gamma_{\infty}^{l} }$.
\end{defn}

% BEGIN REFEREE INSERTION R075
\RefereeRemark{075}{Typo; revised}{Definition of a Gamma-partition}{%
\textit{Original source line 1693.}\quad Restrict the condition to $\gamma\in p(A)$, since fibers outside the image are empty. The original-language definability condition should apply to $A_\gamma$ itself, not merely to $\vrv(A_\gamma)$: write ``$A_\gamma$ is $\mathbb S\cup\{\rcsn(\gamma)\}$-$\lan{}{RV}{}$-definable and $\vrv(A_\gamma)$ is a singleton.'' This is the fiber whose original $\RV$-dimension is used in the next sentence.
}
% END REFEREE INSERTION R075

\begin{lem}[{\cite[Lemma~3.2]{Yin:int:expan:acvf}}]\label{G:part:dim:inv}
If $p_1$, $p_2$ are $\Gamma$-partitions of $A$ then  $\dim_{\RV^\dag} (p_1) = \dim_{\RV^\dag} (p_2)$.
\end{lem}

So the $\RV^\dag$-dimension $\dim_{\RV}(A)$ of a definable set $A$ in  $\RV$ may be defined as the $\RV^\dag$-dimension of any $\Gamma$-partition of $A$. It may also be shown that there is a definable finite-to-one function $f : A \fun \RV^k \times \Gamma_{\infty}^l$ if and only if there is a definable function $f : A \fun \RV^k$ such that all fibers of $f$ are of $\RV^\dag$-dimension $0$ if and only if $\dim_{\RV^\dag}(A) \leq k$.

% BEGIN REFEREE INSERTION R076
\RefereeRemark{076}{Typo; revised}{Dimension notation}{%
\textit{Original source line 1699.}\quad The first $\dim_{\RV}(A)$ should be $\dim_{\RV^\dag}(A)$, as noted in the human report. A reference to the existence of $\Gamma$-partitions would also make the definition self-contained.
}
% END REFEREE INSERTION R076

\begin{defn}[$\RV^\dag$-categories]\label{defn:f:RV:cat}
The objects of the category $\RV^\dag[k]$ are the pairs $(U, f)$ with $U$ a set in $\RV_\infty$ and $f : U \fun \RV^k$ a definable function such that $\dim_{\RV^\dag}(U_t) = 0$ for all $t \in \RV^k$. For two such objects $(U, f)$, $(V, g)$, any definable bijection $F : U \fun V$ is a \memph{morphism} of $\RV^\dag[k]$.
\end{defn}

The categories $\RES^\dag[k]$, $\RES^\dag$ are formulated similar to $\RES[k]$, $\RES$.

% BEGIN REFEREE INSERTION R077
\RefereeRemark{077}{Typo}{Adverb}{%
\textit{Original source line 1705.}\quad Replace ``are formulated similar to'' by ``are formulated similarly to''.
}
% END REFEREE INSERTION R077


\begin{defn}\label{twistback}
The \memph{twistback} function $\tbk : \RV \fun \K$ is given by $u \efun u / \rcsn(\vrv(u))$, where $\infty / \infty = 0$. For any set $U \sub \RV^n_\infty$ and $\gamma \in \Gamma_{\infty}^n$, the set $\tbk(U_{\gamma}) \sub \K^n$ is called the \memph{$\gamma$-twistback} of $U$. If $\tbk(U_{\gamma}) = \tbk(U_{\gamma'})$ for all $\gamma, \gamma' \in \vrv(U)$ then $U$ is called a \memph{twistoid}, in which case we simply write $\tbk(U)$ for the unique twistback.
\end{defn}

% BEGIN REFEREE INSERTION R078
\RefereeRemark{078}{Typo}{Domain of twistback}{%
\textit{Original source line 1710.}\quad The displayed domain should be $\RV_\infty$, since the definition includes the value at $\infty$ and is immediately applied to subsets of $\RV_\infty^n$. This correction from the human report remains outstanding.
}
% END REFEREE INSERTION R078

\begin{lem}[{\cite[Corollaries~2.23, 3.4]{Yin:int:expan:acvf}}]\label{gam:red:lrv}
Every definable set in $\Gamma$ is $\lan{}{RV}{}$-definable and every definable set in $\RV$ with $\vrv(U)$ finite is $\rcsn(\vrv(U))$-$\lan{}{RV}{}$-definable.
\end{lem}

% BEGIN REFEREE INSERTION R079
\RefereeRemark{079}{Typo}{Unnamed set in the lemma}{%
\textit{Original source line 1714.}\quad Write ``every definable set $U$ in $\RV$ with $\vrv(U)$ finite''. Currently the symbol $U$ has not been introduced in this statement. The original-language definability claims also require the parameter convention requested above.
}
% END REFEREE INSERTION R079

\begin{lem}[{\cite[Lemma~3.3]{Yin:int:expan:acvf}}]\label{shift:K:csn:def}
Let $U \sub \RV^n$ be a definable set and $\vrv(U) = D$. Then there is a definable finite partition $(D_i)_i$ of $D$ such that each $U_i = U \cap D_i^\sharp$ is a definable twistoid and the corresponding twistback is $\lan{}{RV}{}$-definable.
\end{lem}


A definable finite partition $(U_i)_i$ of $U$ is called a \memph{twistoid decomposition} of $U$ if every $U_i$ is a twistoid. For instance,  the partition $(U_i)_i$ of $U$ induced by the partition $(D_i)_i$ of $D$ above is a \memph{$\Gamma$-cohesive} twistoid decomposition, which, by Lemma~\ref{gam:red:lrv}, is  $\lan{}{RV}{}$-definable if $U$ is.

\begin{ter}\label{twist:rv:ob}
For any object $\bm U = (U, f) \in \RV^\dag[*]$, by a ``twistoid decomposition of $\bm U$'' we actually mean a twistoid decomposition $(V_i)_i$ of the graph of the function $f$. In that case, clearly $(U_i)_i = (\pr_U(V_i))_i$ is a twistoid decomposition of $U$. For convenience, we shall refer to $(U_i)_i$ as a twistoid decomposition of $\bm U$.
\end{ter}


Let  $(U_i)_i$ be a twistoid decomposition of $\bm U = (U, f) \in \RV^\dag[*]$. Set
\begin{equation}\label{ebdag}
 \bb E^\dag_b(\bm U) = \sum_i \chi_b(\vrv(U_i))[\tbk(U_i)] \in \ggk \RES^\dag.
\end{equation}

\begin{prop}[{\cite[Propositions~3.21, 3.30]{Yin:int:expan:acvf}}]
The assignment (\ref{ebdag}) does not depend on the choice of the twistoid decomposition  and is invariant on isomorphism classes. The resulting map
\[
\bb E^\dag_b : \ggk \RV^\dag[*] \fun \ggk \RES^\dag
\]
is a ring homomorphism that vanishes on the ideal $(\bm P - 1)$.
\end{prop}

We consider $\RV[*]$ as a subcategory of $\RV^\dag[*]$ and denote the induced homomorphism between the Grothendieck rings by
\[
\Lambda : \ggk\RV[*] \fun \ggk\RV^\dag[*].
\]


In the current environment, the transition from ``$\ggk \RES^\dag$'' to ``$\sggk \RES^\dag$'' is superfluous since we already have $[\gamma^\sharp] = [\G_m]$ in $\ggk \RES^\dag$ for all definable $\gamma \in \Gamma$ (every element in  $\gsk \RES^\dag$ is represented by a definable set in $\K$). So there is a natural homomorphism  $\sggk \RES \fun \ggk \RES^\dag$, which is also denoted by $\Lambda$. It then follows from  Proposition~\ref{red:D:iso}, more precisely, the correspondence (\ref{bil:def}), that we have a commutative diagram
\begin{equation}\label{dag:red}
\bfig
\Square(0,0)/->`->`->`->/<400>[{\ggk \RV[*]}`{\sggk \RES}`{\ggk \RV^\dag[*]}`{\ggk \RES^\dag}; \bb E_b`\Lambda`\Lambda`{\bb E^\dag_b}]
\efig
\end{equation}


\begin{rem}\label{dag:g}
We may replace $\chi_b$ with $\chi_g$ and $[\tbk(U_i)]$ with $[\tbk(U_i)][\A]^{-n}$ (assuming $\bm U \in \RV^\dag[n]$) in (\ref{ebdag}) and thereby obtain a ring homomorphism
\[
\bb E^\dag_g : \ggk \RV^\dag[*] \fun \ggk \RES^\dag[[\A]^{-1}]
\]
that also vanishes on the ideal $(\bm P - 1)$. The diagram (\ref{dag:red}) still commutes if $\bb E_b$, $\bb E^\dag_b$ are replaced by $\bb E_g$, $\bb E^\dag_g$.
\end{rem}

One main advantage of $\bb E^\dag_b$ over $\bb E_b$ is that it makes computations much easier, essentially because  we can use twistoid decompositions given by Lemma~\ref{shift:K:csn:def} instead of ``tensor decompositions'' given by Proposition~\ref{red:D:iso}. To make (\ref{ebdag}) work for $\bb E_b$ as intended, though, we need a special kind  of twistoid decompositions:

\begin{defn}\label{bipo:defn}
Suppose that $(U_i)_i$ is an $\lan{}{RV}{}$-definable twistoid decomposition of $U$.
If, for each $i$, there is an $\lan{}{RV}{}$-definable bijection $U_i \fun V_i \times I_i^\sharp$, where $\vrv(V_i)$ is a singleton and $I_i \in \Gamma[*]$, such that its $\vrv$-contraction is a bijection  and its graph is a twistoid as well, then we say that $(U_i)_i$ is  \memph{bipolar}. Naturally $U$ is a \memph{bipolar twistoid} if it admits a trivial bipolar twistoid decomposition.
\end{defn}


\begin{rem}\label{bipo:desc}
The point of Definition~\ref{bipo:defn} is this: Suppose that $\bm U = (U, f) \in \RV[n]$ and  $(U_i)_i$ is a bipolar twistoid decomposition of $U$, as witnessed by bijections $U_i \fun V_i \times I_i^\sharp$ with $I_i \in \Gamma[l_i]$. By dimension theory, for each $i$, there is an $\lan{}{RV}{}$-definable finite-to-one function $g_i: V_i \fun \RV^{n-l_i}$, that is, $(V_i, g_i) \in \RES[n{-}l_i]$. Let $\gamma_i$ be a definable element in $\vrv(U_i)$. Since $U_{i,\gamma_i}$ is $\lan{}{RV}{}$-definably bijective to $V_i \times \alpha_i^{\sharp}$ for some definable $\alpha_i \in I_i$, we see that, in $\sggk \RES$
\[
[U_{i,\gamma_i}] = [V_i] [\alpha_i^{\sharp}] = [V_i] [\G_m]^{l_i}.
\]
So,  by the construction of $\bb E_b$ and bipolarity,
\begin{equation}\label{Eb:twist:com}
\bb E_b([\bm U]) = \sum_i \chi_b(I_i)[V_i] [\G_m]^{l_i} = \sum_i \chi_b(\vrv(U_i))[U_{i,\gamma_i}].
\end{equation}
Observe that this does not really follow without some stronger condition such as bipolarity.
\end{rem}

% BEGIN REFEREE INSERTION R169
\RefereeRemark{169}{Clarification; new}{Scope of the sampled value fibers}{%
\textit{Original source line 1778.}\quad The second expression presupposes that each $\vrv(U_i)$ contains a definable point. This is automatic over the later Puiseux base, but not over every base allowed here: over a trivially valued algebraically closed base, $U=\RV^{\ccirc}$ is bipolar and $\vrv(U)=(0,\infty)$ contains no definable point. The first, tensor-product expression is still valid (and is $0$ in this example). State the sampled-fiber formula subject to the existence of the chosen $\gamma_i$, or add the nontrivial definable value-group hypothesis when invoking it. This is a scope issue, not a counterexample to $\mathbb E_b$.
}
% END REFEREE INSERTION R169

Obviously a $\Gamma$-cohesive twistoid decomposition of a bipolar twistoid is also bipolar. There is indeed an abundant supply of bipolar twistoid decompositions:

\begin{lem}\label{coh:decom}
Every $\lan{}{RV}{}$-definable set $U \sub \RV^n$ admits a $\Gamma$-cohesive twistoid decomposition that is also bipolar.
\end{lem}
\begin{proof}
This follows from \cite[Lemmas~3.21, 3.25]{hrushovski:kazhdan:integration:vf}. In more detail, the proof of \cite[Lemmas~3.21]{hrushovski:kazhdan:integration:vf} shows that the definable finite partition $(D_i)_i$ given by Lemma~\ref{shift:K:csn:def} can be refined so as to make the following condition hold:  each $U_i = U \cap D_i^\sharp$ is of the form $\set{t \in D_i^\sharp \given N_it \in W_i}$, where $N_i$ is a $k \times n$ matrix over $\Z$ and $W_i$ is a definable subset of $\alpha_i^\sharp$ for some $\alpha_i \in \Gamma^k$. Next, there exists a matrix $M_i \in \mgl_n(\Z)$ such that $N_iM_i$ is in lower echelon form (in general $M_i$ is not a product of ``standard'' column operations since $\Z$ is not a field, but  such an invertible matrix exists over any  principal ideal domain). Observe that if $N_iM_i$ does not have zero columns then $D_i$ must be a singleton. At any rate, the set $M_i^{-1}U_i$ must be of the form $\pr_{\leq m}(M_i^{-1}U_i) \times I_i^\sharp \sub \RV^n$, where $n-m$ is the number of zero columns in $N_iM_i$ and $\vrv(\pr_{\leq m}(M_i^{-1}U_i))$ is a singleton.

% BEGIN REFEREE INSERTION R080
\RefereeRemark{080}{Error}{Integer column operations}{%
\textit{Original source line 1789.}\quad The parenthetical assertion about $\mathbb Z$ is incorrect: every matrix in $\mathrm{GL}_n(\mathbb Z)$ is a product of elementary integer column additions, column interchanges, and sign changes. One may simply invoke column Hermite reduction to obtain $N_iM_i=(C_i\;0)$ with $C_i$ of full column rank. Then $C_i\vrv(v)=\alpha_i$ determines the first coordinates uniquely, which supplies the asserted product decomposition. Also, the cited item 3.25 of \cite{hrushovski:kazhdan:integration:vf} is a corollary, and ``Lemmas~3.21'' should be ``Lemma~3.21''.
}
% END REFEREE INSERTION R080
\end{proof}



\subsection{Tropical motivic Fubini theorem}\label{jp:fubini}

In this subsection we work in the field $\puC$ of complex Puiseux series, considered as an elementary substructure of $\bb U$, and take the parameter space  to be the substructure $\mathbb{S} = \C \dpar{ t }$. The value group $\Gamma$ is identified with $\Q$. As an elementary substructure of $\bb U^\dag$, $\puC$ is denoted by  $\tilde\C^{\dag}$. We reiterate that the reduced cross-section $\rcsn$ is not something of central interest, it is just an (implicit) computational tool.

% BEGIN REFEREE INSERTION R081
\RefereeRemark{081}{Clarification; revised}{Normalization and language in the Puiseux setting}{%
\textit{Original source line 1796.}\quad For the Puiseux computation, explicitly normalize $\rcsn(1)=\rv(t)$ and choose compatible roots. Retain the distinction already introduced above: unadorned categories use the original language over $\mathbb C((t))$, while daggered categories use the expansion. A normalized reduced section makes the compatible roots of $t$ definable in the expanded language, so this distinction is necessary to retain monodromy.
}
% END REFEREE INSERTION R081

\begin{rem}
We may think of the profinite group $\hat {\mu}$ as the Galois group $\gal(\puC / \C \dpar t)$, since they are canonically isomorphic. For each element $\xi = (\xi_n)_n \in \hat \mu$, the assignment $n \efun \xi_n t^{1/n}$ indeed induces a reduced cross-section $\rcsn_\xi : \Q \fun \RV$, and the map given by $\xi \efun \rcsn_\xi$ is indeed a bijection between $\hat \mu$ and the set $\Omega$ of reduced cross-sections $\rcsn : \Q \fun \RV$ with $\rcsn(1) = \rv(t)$; in other words, $\hat \mu$ acts freely and transitively on $\Omega$ via multiplication in the obvious way.

% BEGIN REFEREE INSERTION R082
\RefereeRemark{082}{Typo}{Formula for a reduced section}{%
\textit{Original source line 1799.}\quad The assignment should be $1/n\mapsto\rv(\xi_n t^{1/n})$; the present $n\mapsto\xi_nt^{1/n}$ has the wrong input and takes values in $\VF$, rather than $\RV$. Equivalently, write $\rcsn_\xi(a/n)=\rv(\xi_n^a t^{a/n})$ and mention the compatibility that makes this independent of the fraction.
}
% END REFEREE INSERTION R082

Here is a slightly different perspective on $\hat \mu$. By the structural theory of valued fields (in particular, \cite[Lemma~5.3.2]{engler:prestel:2005}), an element $\sigma \in \gal(\puC / \C \dpar t)$ is in the ramification subgroup if and only if it fixes $\RV$ pointwise. But it can be easily checked that every $\sigma \in \gal(\puC / \C \dpar t)$ moves some element of $\RV$ unless $\sigma = \id$. So $\gal(\puC / \C \dpar t) \cong \hat \mu$ may be identified with $\aut(\RV / \RV(\C \dpar t))$, the group of  model-theoretic automorphisms, where $\RV(\C \dpar t)$ is equal to the subgroup generated by $\rv(t)$ over $\K^\times$.
\end{rem}

\begin{nota}\label{bipo:gdv}
The presence of $\rcsn$ induces an intrinsic isomorphism $\RV \cong {\K^\times} \oplus \Q$. Consequently, $\ggk \RES^\dag \cong \ggk \var_{\C}$. Let $\phi : \gmv \fun \ggk \var_{\C}$ be the obvious forgetful homomorphism.  Using the twistback function $\tbk$ as in  \cite[\S~4.3]{hru:loe:lef}, we construct an isomorphism
\begin{equation}
\Theta : \sggk \RES \fun \gmv
\end{equation}
such that $\phi \circ \Theta = \Lambda$. By the discussion in Remark~\ref{bipo:desc}, if $U$ is a bipolar twistoid then the class $[U_\gamma] \in \sggk \RES$ is invariant for definable $\gamma \in \vrv(U)$, and hence it makes sense to denote the class $\Theta([U_\gamma]) \in \gmv$ by $[\tbk(U)]^{\hat \mu}$. In general, for any $\bm U = (U, f) \in \RV[*]$, we may write
\begin{equation}\label{twist:C}
  (\Theta \circ \bb E_b)([\bm U]) = \sum_i \chi_b(\vrv(U_i))[\tbk(U_{i})]^{\hat \mu} \in \gmv,
\end{equation}
where $(U_i)_i$ is any bipolar twistoid decomposition of $U$.
\end{nota}


\begin{nota}\label{mov:vol}
The \memph{motivic volume operator} is the composition
\[
  \Vol : \ggk \VF_* \to^{\int} \ggk \RV[*] / (\bm P - 1)    \to^{\bb E_b} \sggk \RES \to^{\Theta} \gmv.
\]
\end{nota}


For each $\gamma \in \Q^n$, let $A_\gamma \in \VF_*$ and $\bm U_\gamma \in \RV[*]$ such that $\int [A_\gamma] = [\bm U_\gamma] / (\bm P - 1)$,  which we think of  as holding at the semiring level. Suppose that $A = \bigcup_{\gamma \in \Q^n} \gamma^\sharp  \times A_\gamma \in \VF_*$, that is, $A$ is definable, for instance,  there could be a definable set $B \sub \VF^n$ such that $A_\gamma = B \cap \gamma^\sharp$ for all $\gamma \in \Q^n$. For ease of notation, we take each $\bm U_\gamma$ to be simply a definable set $U_\gamma$ in $\RV$ and assume, by compactness, $U = \bigcup_{\gamma \in \Q^n} \gamma^\sharp  \times U_\gamma$ is definable as well. By Lemma~\ref{coh:decom} and compactness again, there is a definable finite partition $(D_i)_i$ of $\vrv(U)$ such that

% BEGIN REFEREE INSERTION R083
\RefereeRemark{083}{Clarification; revised}{Uniform representatives and compactness}{%
\textit{Original source line 1825.}\quad Make this compactness step precise by first interpreting the family in a saturated elementary extension. Apply Theorem~\ref{main:prop:HK} over $\mathbb S\cup\{\gamma\}$ there, uniformize the representatives together with their finite bijection witnesses, and then restrict to the Puiseux model. Compactness cannot be applied just to an arbitrary list of representatives chosen at rational points. Also distinguish the valued-field fibers $\vv^{-1}(\gamma)$ used in $A$ from the leading-term fibers $\vrv^{-1}(\gamma)$ used in $U$.
}
% END REFEREE INSERTION R083
\begin{itemize}
  \item every set $U \cap D_{i}^\sharp$ is a bipolar twistoid,
  \item the sets $E_i = \pr_{\leq n}(D_i)$ form a partition of $\Q^n$,
  \item for all $\gamma, \gamma' \in E_i$, $\chi_b(D_{i,\gamma}) = \chi_b(D_{i,\gamma'})$.
\end{itemize}

% BEGIN REFEREE INSERTION R084
\RefereeRemark{084}{Error}{Projected pieces need not form a partition}{%
\textit{Original source line 1830.}\quad The second bullet cannot in general hold together with the first. For example, use the constant family $U_\gamma=\{1,2\rv(t),3\rv(t)\}$. In $U=\bigcup_\gamma\gamma^\sharp\times U_\gamma$, the twistbacks at internal values $0$ and $1$ have respectively one and two points in the last coordinate. Thus they must lie in different twistoid pieces, whose projections to the base overlap at every $\gamma$. Use a separately indexed base partition $(E_a)_a$, subordinate to the projections of all $D_i$, such that every contribution indexed by $i$ is constant on each $E_a$. Then define $v_a$ by summing all those contributions. Include a base piece for empty fibers when necessary.
}
% END REFEREE INSERTION R084
It follows  that the map $\Q^n \fun \ggk^{\hat \mu} \var_{\C}$ given by
\begin{equation}\label{trop:fub:cons}
\gamma \efun \Vol([A_\gamma]) = \sum_{i} \chi_b(D_{i,\gamma}) [\tbk (U_\gamma \cap D^\sharp_{i,\gamma})]^{\hat \mu}
\end{equation}
is constant over each $E_i$, which we denote by $v_i$, and hence
\[
 \Vol([A]) = \sum_i \chi_b(E_i)v_i[\G_m]^{n}.
\]
This  is essentially the content of the tropical motivic Fubini theorem  of \cite{Nic:Pay:trop:fub}.

% BEGIN REFEREE INSERTION R085
\RefereeRemark{085}{Clarification; revised}{The extra torus factor}{%
\textit{Original source line 1839.}\quad This correctly formulated target is an inflated-family identity: $A=\bigcup_\gamma\vv^{-1}(\gamma)\times A_\gamma$. It has an additional factor $[\mathbb G_m]^n$ compared with the usual tropical Fubini identity for a set $B$ and its tropical fibers. State this scope explicitly. To deduce $\operatorname{Vol}(B)=\int\operatorname{Vol}(B\cap\vv^{-1}(\gamma))\,d\chi_b$, a compatible relative decomposition is still needed; cancellation of $[\mathbb G_m]^n$ in the Grothendieck ring is not justified. Compare \cite[Theorem~3.1.3]{Nic:Pay:trop:fub}.
}
% END REFEREE INSERTION R085


\begin{rem}
The definable sets $E_i$ may be viewed as rational polyhedrons in $\R^n$, that is, intersections of closed half spaces in $\R^n$ defined by linear equations with rational coefficients and (definable) constant terms. In general they are not cones, in other words, some of the constant terms may be nonzero (every element in $\Q$ is definable). This issue is moot, however, if the parameter space is $\mathbb{S} = \C$, since in that case the only definable element in $\Q$ is $0$. This is needed for the applications in the motivic Donaldson-Thomas theory in  \cite{Nic:Pay:trop:fub}.
\end{rem}

% BEGIN REFEREE INSERTION R086
\RefereeRemark{086}{Error}{Polyhedra versus constructible polyhedral sets}{%
\textit{Original source line 1844.}\quad A definable set in $\Gamma^n$ need not be an intersection of closed half-spaces; $(0,1)$ is already a counterexample. Replace this by ``finite Boolean combinations of rational polyhedra'', or refine to relatively open rational polyhedral cells. Also say inequalities, rather than equations, when defining half-spaces. Over the trivially valued base the defining linear conditions are homogeneous, but the resulting pieces can be relatively open cones or Boolean combinations of cones. This is the convention used in \cite[\S~2.2]{Nic:Pay:trop:fub}.
}
% END REFEREE INSERTION R086


\subsection{Fubini theorem}
Let $A, A' \in \VF_*$ and $f : A \fun B$, $f': A' \fun B$ be definable functions. Write the fibers $f^{-1}(b)$, $b \in B$, as $A_b$, similarly for $f'^{-1}(b)$. The purpose of this subsection is to prove the following result.

\begin{thm}\label{fubini}
If, for all $b \in B$, there are $b$-definable objects $\bm U_{b}, \bm U'_{b} \in \RV[*]$ such that $[A_b] = [\bb L \bm U_{b}]$, $[A'_b] = [\bb L \bm U'_{b}]$, and $\bb E_b ([\bm U_b]) = \bb E_b (  [\bm U'_b])$, then $\bb E_b (\int [A]) = \bb E_b ( \int [A'])$.
\end{thm}

% BEGIN REFEREE INSERTION R087
\RefereeRemark{087}{Clarification; revised}{Witnesses for the lifting equalities}{%
\textit{Original source line 1852.}\quad Specify that $[A_b]=[\mathbb L\bm U_b]$ and its primed analogue hold at the semiring level, equivalently through $b$-definable bijections, as required by the replacement of total spaces in the proof. If ring equalities are intended instead, include and uniformize their stabilization objects before cancelling them at the end. The saturated interpretation of $b$ is already covered by the standing convention, and the coefficient base is clarified in Remark~\ref{rem:volb}. Also make explicit that the ground base satisfies the $(\VF,\Gamma)$-generation hypothesis of Theorem~\ref{main:prop:HK}, which is needed for the displayed conclusion involving $\int$.
}
% END REFEREE INSERTION R087

\begin{rem}\label{rem:volb}
Suppose that $B$ is a set in $\VF$, or more generally, for every $b \in B$, the substructure $\bb S \la b \ra$ generated by it over $\bb S$ is $\VF$-generated. The latter condition holds if, for instance, there is a definable bijection $B \fun B'$ with $B'$ a set in $\VF$. Then the hypothesis of the theorem can be simply reformulated as $\bb E_b (\int [A_b]) = \bb E_b ( \int [A'_b])$. This equality is construed over  $\bb S \la b \ra$, which makes sense since $\bb S \la b \ra$ is $\VF$-generated (recall the provisions of Theorem~\ref{main:prop:HK}).

As in \S~\ref{jp:fubini}, we are mainly interested in applying this result in $\puC$ with $\bb S = \C \dpar{ t }$, in which case $\bb S\la b \ra = \C \dpar{ t }(b)$ is an algebraic extension of $\C \dpar{ t }$ and the conclusion of the theorem  may take the form $\Vol([A]) = \Vol([A'])$; however,  it is clearly not enough to test the hypothesis just inside $\puC$.

Take $\bb S = \C \dpar{ t }$  and $b \in B$ in $\bb U$. Denote the algebraic closure of the residue field $\K(\C \dpar{ t }( b ))$ by $\C_b$. Let $r_b$  be the $\Q$-rank of the value group $\Gamma(\C \dpar{ t }( b ))$. Then, working in the algebraic closure of $\C \dpar{ t }( b )$ (this is indeed the smallest model of $\ACVF(0, 0)$ that contains $\C \dpar{ t }( b )$), we can construct, similar to $\Theta$, an isomorphism
\[
\Theta_b:  \sggk \RES \fun \ggk^{{\hat \mu}^{r_b}}\var_{\C_b}.
\]

% BEGIN REFEREE INSERTION R088
\RefereeRemark{088}{Error; revised}{Arithmetic residue descent is missing}{%
\textit{Original source line 1862.}\quad The proposed twistback/base-change comparison does not retain residue-field descent for arbitrary $b$. Take a unit $b$ with residue $s$ transcendental over $\mathbb C$, endowed with the Gauss valuation. Then $k(\mathbb C((t))(b))=\mathbb C(s)$ and its value group is $\mathbb Z$. The residue object $\{z:z^2=s\}$ has a nontrivial arithmetic Galois action over $\mathbb C(s)$, but its base change to $\mathbb C_b$ is two points with trivial tame-inertia action. Their classes are distinguished before base change by the arithmetic Galois representation on cohomology. Thus the proposed comparison with only $\hat\mu^{r_b}$-action is not an isomorphism. Keep $\sggk\RES$ over $\mathbb S\langle b\rangle$, or construct a target retaining residue-field descent as well as inertia. Taking an algebraic closure as the ambient model does not enlarge the definability base.
}
% END REFEREE INSERTION R088
Setting $\Vol_b = \Theta_b \circ \bb E_b \circ \int$, the hypothesis of the theorem can then be equivalently reformulated as  $\Vol_b([A_b]) = \Vol_b([A'_b])$ for all $b \in B$. Since this construction is not really needed, we omit the details. Only note that if $\C \dpar{ t }( b )$ is an unramified extension of $\C \dpar{ t }$ then $\Theta_b$ lands in $\ggk^{{\hat \mu}}\var_{\C_b}$ and hence it is essentially the same as $\Theta$. It will also be clear from the proof that one can test the hypothesis only for those $b \in B$ that generate such unramified extensions (because ultimately the task can be reduced to checking a relation among certain objects of $\RES$ whose projection into $\Gamma$ is in $\Z = \Gamma(\bb S)$).

% BEGIN REFEREE INSERTION R089
\RefereeRemark{089}{Error}{The unramified-only test is false as stated}{%
\textit{Original source line 1863.}\quad The last assertion has a simple counterexample in the general definable setting of the theorem. Put $B=\{b:b^2=t\}$, $A=B$, $f=\mathrm{id}_B$, and $A'=\varnothing$. Every point of $B$ generates a ramified quadratic extension of $\mathbb C((t))$, so the proposed test only at unramified points is vacuous. Nevertheless $\operatorname{Vol}(A)$ is the nonzero class of the two-element monodromy set, whose ordinary Euler characteristic is $2$. Delete this assertion, or state and prove a different test on a specifically constructed auxiliary residue base. It cannot be imposed directly on the original points $b\in B$.
}
% END REFEREE INSERTION R089
\end{rem}



The following fact shall be needed:

\begin{lem}\label{special:term:constant:disc}
Let $F_i(x) = F_i(x_1, \ldots, x_n)$ be a finite list of polynomials with coefficients in $\VF(\bb S)$. Let $A \sub \VF^n$ be a definable set.  Then there is a definable bijection $T : A \fun A^\flat$ such that $A^\flat$ is an $\RV$-pullback and  every composite function $F_i \circ T^{-1}$ is $\rv$-contractible, that is, for each $i$ there is a commutative diagram
\[
\bfig
  \square(0,0)/`->`->`->/<1000,400>[A^\flat`\VF`\rv(A^\flat)`\RV_\infty;
  `\rv`\rv`(F_i \circ  T^{-1})_{\downarrow}]
  \morphism(0,400)<500,0>[A^\flat`A; T^{-1}]
  \morphism(500,400)<500,0>[A`\VF; F_i]
 \efig
\]
\end{lem}
This lemma follows from \cite[Theorem~5.5]{Yin:special:trans}, cited as Lemma~\ref{special:bi:term:constant} below, where $T$ may be taken to be a so-called special bijection, as in Definition~\ref{defn:special:bijection} but much simpler (we shall not need this here, though, because no volume forms are involved).

\begin{proof}[Proof of Theorem~\ref{fubini}]
By partitioning $B$ if necessary, it is enough to prove the claim either when $f$, $f'$ are both surjective or when $f$ is surjective and $A'$ is empty. We will  deal with the first case, the second being similar but simpler (just skip all the steps that involve $A'$ in the argument below).

%Write  $\int [A_b] = [\bm U_{b}] /  (\bm P - 1)$ and $\int [A'_b] = [\bm U'_{b}] /  (\bm P - 1)$.

By  compactness, there is a definable finite partition $(B_i)_i$ of $B$ such that, for each $i$, the objects $\bm U_{b}$,  $\bm U'_{b}$ are defined uniformly over $b \in B_i$. Let $\xi_i(x, y)$, $\xi'_i(x, y')$ be quantifier-free formulas such that $\xi_i(b, y)$, $\xi'_i(b, y')$  define the objects $\bm U_{b}$, $\bm U'_{b}$ for each $b \in B_i$, respectively.

Suppose that $B$ is a set in $\VF$. Applying Lemma~\ref{special:term:constant:disc} to the top terms $F(x)$ of $\xi_i(x, y)$, $\xi'_i(x, y')$  (recall Convention~\ref{topterm}), we obtain a definable bijection $T_i : B_i \fun B_i^\flat$ as described. So, to show the claim, we may replace $B$ with any $B_i^\flat$, and thereby assume that $B$ is an $\RV$-pullback and for each $b \in B$, the objects $\bm U_{b}$, $\bm U'_{b}$ depend on $\rv(b)$ instead of $b$ itself; henceforth $\bm U_{b}$, $\bm U'_{b}$ are denoted by $U_{\rv(b)}$, $U'_{\rv(b)}$ (as above, for ease of notation, we take them to be simply definable sets in $\RV$). In fact, by Theorem~\ref{main:prop:HK}, setting $S = \rv(B)$, we may just assume that $A = \bigcup_{t \in S} \bb L U_{t} \times t^\sharp$ and $f$ is the obvious coordinate projection. Similarly  $A' = \bigcup_{t \in S} \bb L U'_{t} \times t^\sharp$.


Of course, in general, $B$ has $\RV$-coordinates. But the same arrangement can be made by first applying the argument above to each fiber $B_t$, $t \in  \pr_{\RV}(B)$, and then compactness.
Without loss of generality, we just take $A$, $A'$, $B$ to be sets in $\VF$.

% BEGIN REFEREE INSERTION R090
\RefereeRemark{090}{Gap}{RV coordinates cannot simply be discarded}{%
\textit{Original source line 1894.}\quad The fiberwise argument retains the $\RV$-parameter $t$; compactness makes the resulting maps uniform but does not remove that coordinate. In particular an infinite $\RV$-set does not admit a definable injection into a power of $\VF$ (see \cite[Lemma~3.41]{hrushovski:kazhdan:integration:vf}). Thus the asserted reduction to $A,A',B$ all lying in $\VF$ needs a separate argument. A possible repair is to carry the $\RV$-coordinates and the finite-to-one structure maps throughout the proof, checking that all total spaces remain objects of the relevant categories; otherwise restrict the theorem to a $\VF$ base and prove the extra coding needed for its total spaces.
}
% END REFEREE INSERTION R090

Write $U = \rv(A)$, $D = \vv(A)$, $U '= \rv(A')$, $D' = \vv(A')$, and $E = \vv(B)$. Since  $\Gamma$ is stably embedded, upon further partitioning of $S$, we may assume that $\vrv(U_{t})$ depends on $\vrv(t)$ instead of $t \in S$ itself, and hence may be written as $D_{\vrv(t)}$; similarly, we may write $\vrv(U'_{t})$ as $D'_{\vrv(t)}$.

% BEGIN REFEREE INSERTION R091
\RefereeRemark{091}{Clarification; revised}{Justification of the value-parameter dependence}{%
\textit{Original source line 1896.}\quad The finite partition asserted here is compatible with the usual parameter-control argument, but stable embeddedness alone is not its justification. Use $\dcl(\mathbb S,t)\cap\Gamma=\operatorname{span}_{\mathbb Q}(\dcl(\mathbb S)\cap\Gamma,\vrv(t))$, together with the piecewise rational-affine description of definable value-group data. After finitely partitioning $S$, this gives the required dependence only on $\vrv(t)$. This is the precise replacement for the justification requested in the human report; it does not contradict the assertion after partitioning.
}
% END REFEREE INSERTION R091

For each $t \in S$, let $(D_{\vrv(t),i})_i$ be a definable finite partition of $D_{\vrv(t)}$ as given by Lemma~\ref{coh:decom}, that is, $(U_{t,i})_i = (U_t \cap D_{\vrv(t),i}^\sharp)_i$ is a bipolar and $\Gamma$-cohesive twistoid decomposition of $U_t$. Again, the notation  $(D_{\vrv(t),i})_i$ makes sense because the partition only depends on $\vrv(t)$; in fact, we may even assume so  for the $\vrv$-contractions of the witness bijections as required by Definition~\ref{bipo:defn}. Then, by compactness and \omin-minimality, we are further reduced to the case that $(D_{\gamma,i})_i$ is uniformly defined over $\gamma \in E$, and $(\chi_b (D_{\gamma,i}))_i$ does not depend on $\gamma \in E$ at all and hence may be simply written as $(e_i)_i$. Let $D_i = \bigcup_{\gamma \in E} D_{i,\gamma} \times \gamma$. So $\chi_b(D_i) = e_i \chi_b(E)$. All this is just a more elaborate version of what we did before (\ref{trop:fub:cons}).

% BEGIN REFEREE INSERTION R092
\RefereeRemark{092}{Typo; revised}{Order of the subscripts}{%
\textit{Original source line 1898.}\quad Use $D_{\gamma,i}$ consistently: the definition of $D_i$ changes it to $D_{i,\gamma}$, and the primed paragraph makes the same switch. In the final refinement paragraph below, replace $(D'_{i'})_i'$ by $(D'_{i'})_{i'}$.
}
% END REFEREE INSERTION R092

Similarly we obtain a bipolar and $\Gamma$-cohesive twistoid decomposition $(U'_{t,i'})_{i'} = (U'_t \cap {D'}_{\vrv(t),i'}^\sharp)_{i'}$ of $U'_t$, uniformly defined over $\gamma\in E$, such that $e_{i'}=\chi_b (D'_{\gamma,i'})$ do not depend on $\gamma \in E$ and hence $\chi_b(D'_{i'}) = e_{i'} \chi_b(E)$, where $D'_{i'}=\bigcup_{\gamma \in E} D'_{i',\gamma} \times \gamma$.

Applying Lemma~\ref{coh:decom} again to $U$, $U'$ and refining $(D_{i})_i$, $(D'_{i'})_{i'}$ if necessary, we may now assume that $(U_i)_i = (U \cap D^\sharp_i)_i$, $(U'_{i'})_{i'} = (U' \cap {D'}^\sharp_{i'})_{i'}$ are bipolar and $\Gamma$-cohesive twistoid decompositions of $U$, $U'$, respectively. They induce two twistoid decompositions of $S$ (not necessarily bipolar, which we do not need). Refining $(D_{i})_i$, $(D'_{i'})_i'$ one last time, we may assume that both induce the same decomposition of $S$ after all and hence, replacing $S$ with one of the pieces,  $S$ itself is a twistoid.

% BEGIN REFEREE INSERTION R093
\RefereeRemark{093}{Clarification; revised}{Explain the common refinement}{%
\textit{Original source line 1902.}\quad The common refinement can be justified more explicitly. First refine the value-group base so that the relevant projections of the $U_i$ to $S$ have the same value support. Their common Boolean refinement then consists of twistoids. Restrict each $U_i$ to a selected base piece; it remains a twistoid, with a fixed twistback projecting onto that base twistback. Apply Lemma~\ref{coh:decom} to these restrictions using $\Gamma$-cohesive refinements. Their projections to the selected base are now $\Gamma$-cohesive, so a final value-group refinement makes each surviving projection surjective. On that final refinement, also re-establish constancy of the fiber Euler characteristics $\chi_b(D_{\gamma,i})$. This supplies the compatibility used in the next paragraph; the earlier objection should not be read as an obstruction to the existence of the refinement.
}
% END REFEREE INSERTION R093

\begin{rem}
The witness bijection for $U_i$ as a bipolar twistoid is in general unrelated to those for $U_{t, i}$, $t \in S$, since parameterized bipolar twistoids over a base may not even form a single  bipolar twistoid. But our argument will not require this.
\end{rem}

Let $\gamma_i \in \vrv(U_i)$, $\gamma_{i'} \in \vrv(U'_{i'})$ be definable elements such that $\pr_E(\gamma_{i}) = \pr_E(\gamma_{i'})$ for all $i$, $i'$. Set $\dot U_i = U_{i, \gamma_i}$ and $\dot U'_{i'} = U'_{i', \gamma_{i'}}$. So $\pr_S(\dot U_{i}) = \pr_S(\dot U'_{i'})$ for all $i$, $i'$, which we denote by $\dot S$.

%Set
%\[
%\bar U_i = \tbk(U_i) = \tbk(\dot U_i), \quad \bar U'_{i'} = \tbk(U'_{i'}) = \tbk(\dot U'_{i'}), \quad \bar S = \tbk(S) = \tbk(\dot S).
%\]
%By Lemma~\ref{gam:red:lrv}, these are definable sets (constructible sets) in the residue field $\K$.

%$\ggk^{\hat \mu}_+ \RES$, it is equivalent to showing, in  $\ggk^{\hat \mu}_+ \RES$,
%\[
%\sum_i e'_i [\dot U_{i}] + [V] = \sum_j e''_j [\dot U_{j}] + [V] \quad \text{for some } V \in \RES.
%\]
%This may be further reduced to the condition:

Denote by $I$, $I'$ the sets of indices $i$, $i'$. Set $J = I \uplus I'$. For each $j\in J$, let $V_j=\dot U_j$ if $j\in I$ and $V_j=\dot U'_j$ otherwise. Set
\[
J_1=\set{i\in I \given e_i>0} \cup \set{i'\in I' \given e_{i'}<0},
\]
\[
J_2=\set{i\in I \given e_i<0}\cup \set{i'\in I'  \given e_{i'}>0},
\]
and for $j\in J$, $\tilde e_j = \abs{e_j}$.

By (\ref{Eb:twist:com}), we have $\bb E_b (\int [A] ) = \chi_b(E) \sum_{i\in I} e_i [\dot U_i]$ and $\bb E_b ( \int [A'] ) = \chi_b(E) \sum_{i'\in I'} e_{i'} [\dot U'_{i'}]$. Thus, to prove the theorem, it is enough to show that $\sum_{j\in J_1} \tilde e_j [V_j]=\sum_{j\in J_2} \tilde e_j [V_j]$ in $\sggk \RES$.  Since both sides can be understood in $\gsk \RES$, it is enough to show that the following condition holds.

\begin{test}\label{test}
There are objects $W_1, W_2, X \in \RES$ such that
\begin{itemize}
  \item $[W_1] = \sum_{j\in J_1} \tilde e_j [V_j] + [X]$ and $[W_2] = \sum_{j\in J_2} \tilde e_j [V_j] + [X]$ in $\gsk \RES$,
  \item $([W_1], [W_2])$ is in the semiring congruence relation $\mdl E_{sc}$ on $\gsk \RES$ generated by pairs of the form $([\gamma^\sharp \times Z], [\G_{m}^n \times Z])$, where $Z \in \RES$ and $\gamma \in \Gamma^n$ is definable (we allow $n = 0$, which gives the pairs $(Z, Z)$).
\end{itemize}
\end{test}

What the second item means is that  if we close off the set $\mdl A$ of these pairs $([\gamma^\sharp \times Z], [\G_{m}^n \times Z])$ by applying inductively the six conditions, namely identity (since each stage needs to be accumulative), the three equivalence conditions (symmetry, reflexivity, and transitivity), and the two semiring congruence conditions (if $([X], [Y])$ is already in at some stage then, for any $[Z] \in \gsk \RES$, $([X \uplus Z], [Y \uplus Z])$ and  $([X \times Z], [Y \times Z])$ are thrown in at the next stage), then $([W_1], [W_2])$ appears at some stage. Actually it is clear that reflexivity is redundant. We show that, beyond $\mdl A$, product is also redundant, as follows.

\begin{rem}
There are in general more than one way to inductively construct an element $([X], [Y])$ in $\mdl E_{sc}$. Each inductive construction of $([X], [Y])$ may be represented by a finite graph $T$, partitioned into levels (numbered inversely), such that
\begin{itemize}
  \item every vertex is an element in $\mdl E_{sc}$, in particular, there is only one vertex in the top level, which is $([X], [Y])$ itself, and the bottom vertices are all elements in $\mdl A$,
  \item each vertex above level $0$ has $1$ or $2$ (the latter only for transitivity) labeled edges connecting it with vertices in the level below, with the labels indicating how it is constructed from these latter vertices,
  \item each vertex below the top level is connected via an edge with a vertex in the level above (this is for efficiency, that is, no redundant vertices for the construction).
\end{itemize}
We call $T$ a \memph{construction graph} of $([X], [Y])$, and the number of levels in it its \memph{height}. Observe that, for any construction graph $T$ of $([X], [Y])$ and any $[Z] \in \gsk \RES$, if we replace every vertex $([X'], [Y'])$ of $T$ with $([X' \times Z], [Y' \times Z])$ then the resulting graph is a construction graph for $([X \times Z], [Y \times Z])$. Let $\rho([X], [Y])$ be the least natural number such that there is a construction graph for $([X], [Y])$ of height $\rho([X], [Y])$; let $T$ be such a construction graph. We claim that $T$ may be chosen so that  none of its edges is a product. This is vacuous if $\rho([X], [Y]) = 0$ and is clear if $\rho([X], [Y]) = 1$. For $\rho([X], [Y]) > 1$, note that if $([X'], [Y'])$, $([X''], [Y''])$ are the (possibly identical) vertices in level $\rho([X], [Y]) - 1$ of $T$ then
\[
\max(\rho([X'], [Y']), \rho([X''], [Y''])) = \rho([X], [Y]) - 1.
\]
Also, if the edge going into $([X], [Y])$ is a product then, by what we have just observed, there would be a construction graph of $([X], [Y])$ of height $\rho([X], [Y]) - 1$, contradicting the choice of $\rho([X], [Y])$. So the desired conclusion follows from a routine induction on $\rho([X], [Y])$.
\end{rem}

% BEGIN REFEREE INSERTION R094
\RefereeRemark{094}{Typo; revised}{Height convention}{%
\textit{Original source line 1955.}\quad Define height as the maximum level number, equivalently the number of levels minus one, to agree with the induction and the indices $\rho-1$ below. The product-free construction argument is otherwise valid. At the start of the remark, replace ``There are in general more than one way'' by ``There is in general more than one way''. 
}
% END REFEREE INSERTION R094


We need this seemingly obscure fact so as to be able to apply compactness uniformly over a base below.

\begin{rem}
The semiring congruence relation $\mdl E_{sc}$ may also be constructed ``without the brackets,'' that is, it may be  expressed as an equivalence relation $\mdl E^*_{sc}$ on the objects of $\RES$ that is generated from the set of definable pairs of the form $(\gamma^\sharp \times Z, \G^n_{m} \times Z')$, where $Z \cong Z'$ in $\RES$, via the six conditions. However, if $T$ is a construction graph of  $([X], [Y])$ then simply deleting the brackets from the vertices of $T$ does not give us a graph that witnesses $(X, Y) \in \mdl E^*_{sc}$. There are two issues. The first is transitivity, that is, if $([X], [Y])$ and $([Y'], [Z])$ are in  $\mdl E_{sc}$ with $[Y] = [Y']$ then of course so is $([X], [Z])$, but to conclude $(X, Z) \in \mdl E^*_{sc}$ from $(X, Y), (Y', Z) \in  \mdl E^*_{sc}$ we must add a middle step $(Y, Y')$, justified  by $Y \cong Y'$ in $\RES$, and apply transitivity twice. The second is similar: we can now translate $T$ into a graph that witnesses $(X', Y') \in \mdl E^*_{sc}$, but only with $[X'] = [X]$ and $[Y'] = [Y]$. But then, adding at most two levels on top, we arrive at a \memph{construction graph} $T^*$ of $(X, Y) \in \mdl E^*_{sc}$. Moreover, if $T$ is product-free as given above then $T^*$ is also  product-free, in the obvious sense.
\end{rem}
The reason that we  opt for working with $T^*$ instead of $T$ is that it can be described by  an $\lan{}{RV}{}$-formula and thereby makes itself amenable to compactness.


Finally, by the assumption, we have $\sum_{j\in J_1} \tilde e_j [V_{j,t}] =  \sum_{j\in J_2} \tilde e_j [V_{j,t}] $ in $\sggk \RES$ for every $t \in \dot S$. For each $t \in \dot S$, we can choose $t$-definable objects $W_{1,t}$, $W_{2,t}$, $X_t$ of $\RES$ with respect to which Test~\ref{test} holds, as witnessed by a product-free construction graph $T^*_t$. By compactness, there is a definable finite partition $(\dot S_i)_i$ of $\dot S$ such that, for each $i$, the graph $T^*_t$ is uniformly defined over $t \in \dot S_i$ by an $\lan{}{RV}{}$-formula; without loss of generality, we may assume this partition is trivial. Moreover, if level $0$ of $T^*_t$ consists of definable pairs of the forms $(\gamma^\sharp \times Z_t, \G^n_{m} \times Z'_t)$ then $Z_t$, $Z'_t$ in general depend on $t$, but $\gamma^\sharp$ does not, and neither do other $\Gamma$-data over $t$, such as $\vrv(Z_t)$, $\vrv(\dot U'_t)$, and $\vrv(V_t)$ (of course they do depend on $\vrv(\dot S)$). Let

% BEGIN REFEREE INSERTION R095
\RefereeRemark{095}{Gap; revised}{Descent from VF parameters to RV parameters}{%
\textit{Original source line 1966.}\quad The earlier application of Lemma~\ref{special:term:constant:disc} makes the objects depend only on $t=\rv(b)$, but it does not yet make the equality witnesses $t$-definable. One way to repair this is to uniformize the full finite witness data over $b$ before choosing $T$: include the fiberwise bipolar decompositions, their sampled value fibers, and the product-free congruence graphs constructed above. Apply the lemma to every valued-field top term occurring in those formulas, not just the object formulas. Then the witnesses themselves depend on $\rv(b)$, and the required relation over $t$ follows. Finally, refine the base finitely so that all the auxiliary $\Gamma$-data are constant on each piece. In the list of such data, $\dot U'_t$ and $V_t$ have not been defined; use the actual indexed objects, for example $V_{j,t}$.
}
% END REFEREE INSERTION R095
\[
W_1 = \bigcup_{t \in \dot S} W_{1,t}, \quad W_2 = \bigcup_{t \in \dot S} W_{2,t}, \quad X = \bigcup_{t \in \dot S} X_t.
\]
Thus, and this is why we have taken effort to find product-free construction graphs above, the union $\bigcup_{t \in \dot S} T^*_t$ is indeed a construction graph of $(W_1, W_2)$. So Test~\ref{test} holds with respect to $W_1$, $W_2$, and $X$. So $\bb E_b (\int [A]) = \bb E_b ( \int [A'])$.
\end{proof}



\begin{cor}
In $\puC$ with $\bb S = \C \dpar{ t }$, if $f$ is a morphism of varieties over $\C$ with $\Vol_b([A_b]) = 0$ for all $b \in B$ then $\Vol([A]) = 0$.
\end{cor}

\section{Special covariant bijections}\label{sec-pscb}


From here on, we also assume that the substructure $\mathbb{S}$ of parameters is $\VF$-generated. So every definable disc contains a definable point (Lemma~\ref{algebraic:balls:definable:centers}). Another simple fact we shall occasionally use is that if $\gamma \in \Gamma$ is definable then there is a $\gamma' \geq \gamma$ such that $\gamma'^\sharp$ contains a definable point.

\subsection{Invariance and covariance}\label{sec:inv:cov}

\begin{defn}\label{defn:binv}
For $\gamma \in \Gamma$, let $\pi_\gamma: \VF \fun \VF / \MM_{\gamma}$ be the quotient  map. If $\gamma = (\gamma_1, \ldots, \gamma_n) \in \Gamma^n$ then $\pi_{\gamma}$ denotes the product of the maps $\pi_{\gamma_i}$.

Let $\alpha \in \Gamma^n$ and $\beta \in \Gamma^m$. We say that a function $f :A \fun B$ with $A \sub \VF^n$ and $B \sub \VF^m$  is \memph{$(\alpha, \beta)$-covariant} if it $(\pi_{\alpha}, \pi_{\beta})$-contracts to a function $f_\downarrow : \pi_{\alpha}(A) \fun \pi_{\beta}(B)$, that is, $\pi_{\beta} \circ f = f_\downarrow \circ \pi_{\alpha}$. For simplicity, we shall often suppress parameters in notation and refer to $(\alpha, \beta)$-covariant functions as $(\alpha, -)$-covariant or $(-,\beta)$-covariant or just  covariant functions.

A set  in $\VF$ is \memph{$\alpha$-invariant} if its characteristic function is $(\alpha, 0)$-covariant.


More generally, for sets $A$, $B$ with $\RV$-coordinates, a function $f : A \fun B$ is \memph{covariant} if every one of its $\VF$-fibers    is covariant (recall Terminology~\ref{rvfiber}). Similarly, a set is \memph{invariant} if every one of its $\VF$-fibers is invariant.
\end{defn}


\begin{defn}\label{prop:pseu}
A definable function $f : A \fun B$ is \memph{proper covariant} if
\begin{itemize}
  \item the sets $A_{\VF}$, $f(A)_{\VF}$ are bounded and the sets $A_{\RV}$, $f(A)_{\RV}$ are  doubly bounded,
  \item for each $\VF$-fiber $f_{t}$ of $f$ there is a $t$-definable tuple $(\alpha_t, \beta_t) \in \Gamma$ such that $f_t$ is $(\alpha_{t}, \beta_{t})$-covariant, $\dom(f_{t})$ is $\alpha_t$-invariant, and $\ran(f_{t})$ is $\beta_t$-invariant.
\end{itemize}
%Accordingly, an invariant set $A$ is \memph{proper} if its projection into the $\RV$-coordinates is proper covariant, in particular, $A_{\VF}$ is bounded and $A_{\RV}$ is doubly bounded.
A definable set $A$ is \memph{proper invariant}   if $A_{\VF}$ is bounded, $A_{\RV}$ is doubly bounded, and for each $\VF$-fiber $A_t$ of $A$ there is a $t$-definable tuple $\alpha_t\in \Gamma$ such that $A_t$ is $\alpha_t$-invariant.
\end{defn}

Note that covariance makes sense even when the domain or the range of the function $f$ has no $\VF$-coordinates, in which case the contraction $f_\downarrow$ would be the empty function. In particular, if $B$  is a subset of $\RV^m$ or even $\mdl P(\RV^m)$ then $f$ being proper covariant simply means that every $f^{-1}(x)$ is proper invariant. Then $A$ being proper invariant may also be construed as the condition that its projection into the $\RV$-coordinates is proper covariant.

% BEGIN REFEREE INSERTION R098
\RefereeRemark{098}{Error}{Fiberwise invariance does not supply global bounds}{%
\textit{Original source line 2008.}\quad The asserted characterization drops the global boundedness conditions in Definition~\ref{prop:pseu}. For example, let $A=1^\sharp\times\RV$ and let $f:A\to\RV$ be the second projection. Every fiber $f^{-1}(t)$ is proper invariant (over $t$), but $A_{\RV}=\RV$ is not doubly bounded, so $f$ is not proper covariant. Retain the boundedness requirements explicitly, or formulate this as a characterization under those requirements. Also, when the target has no $\VF$-coordinates the relevant contraction is a map to the singleton empty tuple; it is not the empty function unless its domain is empty. The extension to powerset-valued functions should specify separately which boundedness condition is imposed on their RV-valued relation.
}
% END REFEREE INSERTION R098

\begin{exam}
A proper invariant set  is not required to be doubly bounded. However, if $\bm U \in \RV^{\db}[*]$ then $\bb L \bm U$ is indeed a doubly bounded proper invariant set.
\end{exam}

\begin{rem}
In defining proper covariance above, if  $A$ is closed then  there is no need to demand that $f(A)_{\RV}$ be doubly bounded, because it is implied by the other conditions, as follows. There is an $\go$-partition $p: A \fun \Gamma$ such that, for each open polydisc $\gb$ in question, that is, for each open polydisc $\gb \sub A$ of the form $\go(a, p(a))$, $f(\gb)$ lies in the range of some $\VF$-fiber of $f$. By Lemma~\ref{vol:par:bounded}, $p(A)$ is doubly bounded. So the definable surjection $A \fun \vrv(f(A)_{\RV})$ induced by $f$ satisfies the assumption of Lemma~\ref{invar:db} below. So  $f(A)_{\RV}$ is doubly bounded.
\end{rem}

\begin{lem}\label{inv:crit}
Let $A$ be a definable set such that $A_{\VF}$ is bounded and $A_{\RV}$ is doubly bounded. Then $A$ is proper invariant if and only if $A$ is clopen.
\end{lem}
\begin{proof}
The ``only if'' direction is clear. For the ``if'' direction, since every $\VF$-fiber of $A$ is open, there is an obvious  $\go$-partition $p: A \fun \Gamma$. Since every $\VF$-fiber of $A$ is also closed, by Lemma~\ref{vol:par:bounded}, $p(A)$ is doubly bounded. The claim follows.
\end{proof}


\begin{lem}\label{int:inv}
If $A$, $B$ are proper invariant subsets of $\VF^n \times \RV^m$ then $A \cap B$, $\RVH(A) \mi A$, and $B \mi A$ are also proper invariant.
\end{lem}
\begin{proof}
Clearly $(A \cap B)_{\VF}$ is bounded, $(A \cap B)_{\RV}$ is doubly bounded, and every $\VF$-fiber of $A \cap B$ is clopen; similarly for the other cases. So this follows from Lemma~\ref{inv:crit}.

% BEGIN REFEREE INSERTION R099
\RefereeRemark{099}{Clarification}{Explain the RV-hull case}{%
\textit{Original source line 2030.}\quad The assertion for $\RVH(A)\setminus A$ additionally needs the fact that $\RVH(A)$ is proper invariant. This follows directly: if a fiber $A_t$ is $\delta$-invariant, then its RV-hull is also $\delta$-invariant (in each coordinate an RV-disc of radius at least $\delta$ meeting $A_t$ is contained in a $\delta$-ball already in $A_t$; larger discs are themselves $\delta$-invariant). Boundedness is preserved by taking the RV-hull. Stating this fact explicitly also justifies the later applications to $\RVH(A)$.
}
% END REFEREE INSERTION R099
\end{proof}

\begin{lem}\label{conti:covar}
Let $f : A \fun B$ be a definable continuous surjection between two proper invariant sets such that  every $\dom(f_t)$ is open. If every $\ran(f_t)$ is clopen then $f$ is proper covariant. If $f$ is bijective then the same holds under the weaker assumption that   every $\ran(f_t)$ is open.
\end{lem}
\begin{proof}
For every $(r, s) \in f_{\RV}$ and every $b \in \ran(f_{r,s})$, let
\[
q(b,r,s) = \min\set{\beta \in \Gamma \given \go((b,s), \beta) \sub \ran(f_{r,s})}.
\]
The definable function $q$ is clearly an $\go$-partition of its domain, which, by the assumption, is closed. By Lemma~\ref{vol:par:bounded}, $\ran(q)$ is doubly bounded by, say, $\beta \in \Gamma^+$. By continuity, there is another $\go$-partition $p: A \fun \Gamma$ such that, for every open polydisc $\gb \sub A$ in question, $f(\gb)$ lies in an open polydisc of radius $\beta$. By Lemmas~\ref{inv:crit} and \ref{vol:par:bounded} again,  $p(A)$ is doubly bounded as well. It follows that $f$ is  proper covariant.

For the second claim, since $f$ is bijective, for each $(b, s) \in B$ there is only one  $f_{r,s}$ with $(b,r, s) \in \ran(f_{r,s})$. So the parameter $r$ is redundant in the definition of $q$, and  $q$ is  simply an $\go$-partition of $B$. The rest of the argument is the same.

% BEGIN REFEREE INSERTION R100
\RefereeRemark{100}{Clarification}{Make the partition radii construction explicit}{%
\textit{Original source line 2043.}\quad In the first paragraph take $p$ small enough in geometric size that its balls also lie in a single $\dom(f_{r,s})$; continuity of the separate $f_{r,s}$ alone does not express that requirement. In the bijective case the unique $r=r(b,s)$ is not literally a redundant coordinate: it is a definable function of $(b,s)$, locally constant because the sets $\ran(f_{r,s})$ are open. Hence $q$ descends to an $\go$-partition of the closed set $B$, and Lemma~\ref{vol:par:bounded} applies. This is the missing justification for the last paragraph.
}
% END REFEREE INSERTION R100
\end{proof}


\begin{lem}\label{conv:unif}
If $f : A \fun B$ is a proper covariant surjection then both $A$ and $B$ are proper invariant. If, in addition, $f$ is continuous and bijective then $f$ is a proper covariant homeomorphism, which just means that $f^{-1}$ is also proper covariant and continuous. Moreover, in that case, if $A'$ is a proper invariant subset of $A$ then $f \rest A'$ is a proper covariant homeomorphism as well.
\end{lem}
\begin{proof}
Since $f_{\RV}$ is  doubly bounded, by Lemma~\ref{db:to:db}, for all sufficiently large  $\epsilon \in \Gamma$, every $\VF$-fiber $f_{t}$ of $f$ is $(\epsilon, \beta_{t})$-covariant (but not necessarily $(\epsilon, \epsilon)$-covariant) and $\dom(f_{t})$,  $\ran(f_{t})$ are both $\epsilon$-invariant. The first claim follows.

For the second claim, since $\dom(f_{t})$ is clopen, by  Corollary~\ref{conti:homeo}, every $f_t$ is a homeomorphism. Since  $\ran(f_{t})$ is open, it follows from the continuity of $f^{-1}$, as in the proof of Lemma~\ref{conti:covar}, that every  $f^{-1}_t$ is  $(\delta, \epsilon)$-covariant  for all sufficiently large $\delta > \epsilon$.

For the third claim, clearly every $\dom(f_t) \cap A'$ is   clopen and, since $f_t$ is a homeomorphism, every $\ran(f_t) \cap f(A')$ is  clopen too. This is sufficient for the proof of (the first claim of) Lemma~\ref{conti:covar} to go through. By the second claim, then, $f \rest A'$ is a proper covariant homeomorphism.
\end{proof}

\begin{lem}\label{pr:co:home:com}
If $f : A \fun B$ and $g : B  \fun C$ are proper covariant homeomorphisms then $g \circ f$ is also a proper covariant homeomorphism.
\end{lem}
\begin{proof}
By Lemma~\ref{conv:unif}, we just need to show that $g \circ f$ is proper covariant. Then, by Lemma~\ref{conti:covar}, we are reduced to showing that the domain and range of every $\VF$-fiber of $g \circ f$ are open. Arguing as in the proofs of Lemmas~\ref{conti:covar} and \ref{conv:unif}, we see that there are $\epsilon, \epsilon', \epsilon'' \in \Gamma$, each one  sufficiently larger than the next, such that every $\VF$-fiber of $f$ is $(\epsilon, \epsilon')$-covariant and every $\VF$-fiber of $g$ is $(\epsilon', \epsilon'')$-covariant. It follows that the domain of every $\VF$-fiber of $g \circ f$ is $\epsilon$-invariant and hence is clopen. Since the same holds for $(g \circ f)^{-1}$, the range of every $\VF$-fiber of $g \circ f$ is clopen as well.
\end{proof}

\begin{conv}\label{conv:can}
The trivial-looking operation we describe here is actually quite crucial for understanding the discussion below, especially those passages that involve special covariant bijections (defined below). For each definable set $A$, let
\[
\can(A) = \set{(a, \rv(a), t) \given (a, t) \in A \tand a \in A_{\VF}}.
\]
The obvious definable bijection $\can : A \fun \can(A)$ is  referred to as the \memph{regularization} of $A$.

Observe that, for a proper invariant set $A$, $\can(A)$ is proper invariant if and only if $A_{\VF}$ is  doubly bounded. Going forward, when $A$ is doubly bounded and proper invariant, it is often more convenient or even necessary to substitute $\can(A)$ for $A$. Whether such a substitution has been performed (tacitly) or not should be clear in the context (or rather, it is always performed).
\end{conv}
%\begin{defn}
%Let $A$ be a definable set. A \memph{regularization} map is a map of the form
%\[
%x=(x_1, \ldots , x_n) \in A \efun (x,\rv(x_i)) \in A \times \RV,
%\]
%where the $i$th coordinate of $A$ is assumed to be a $\VF$-variable.
%\end{defn}

\subsection{Proper covariant bijections that are special}

\begin{defn}\label{defn:special:bijection}
Let $A$ be a doubly bounded proper invariant set  or, more generally, a finite disjoint union of such sets,  regularized in the sense of  Convention~\ref{conv:can}. Suppose that the first coordinate of $A$ is a $\VF$-coordinate (of course nothing is special about the first $\VF$-coordinate, we choose it simply for ease of notation).

Let $C \sub \RVH(A)$ be an $\RV$-pullback. Let
\[
\lambda: \pr_{>1}(C \cap A) \fun \VF
\]
be a definable continuous function whose graph is contained in $C$, that is, for each $\RV$-polydisc $\gp \sub C$, $\lambda$ restricts to a function
\[
\lambda_{\gp} : \pr_{>1}(\gp \cap A) \fun \pr_{1}(\gp).
\]
Suppose that there is an $\rv(\gp)$-definable tuple $(\alpha_{\gp}, \beta_{\gp}) \in \Gamma$ such that $\gp \cap A$ is $(\alpha_{\gp}, \beta_{\gp})$-invariant and $\lambda_{\gp}$ is  $(\beta_{\gp}, \alpha_{\gp})$-covariant. Choose a definable $\gamma \in \Gamma^+_0$  such that there is a definable point $t \in \gamma^\sharp$ and for all $\gp$,
\[
\gamma_{\gp} \coloneqq \rad(\pr_{1}(\gp)) + \gamma \geq \alpha_{\gp};
\]
the existence of such a $\gamma$ is guaranteed by Lemma~\ref{db:to:db} and the assumption that $\mathbb{S}$ is $\VF$-generated. For each $\gp$,  set
\[
t_{\gp} = t \cdot \rv(\pr_{1}(\gp)) \in \RV.
\]
Then the \memph{centripetal transformation $\eta$ on $A$ with respect to $\lambda$} is given by
\begin{equation*}\tag{CT}\label{centri}
\begin{cases}
  \eta (a, x) = (a - \lambda(x), x), & \text{if } (a, x) \in \gp \cap A \text{ and } \pi_{\gamma_{\gp}}(a) \neq \pi_{\gamma_{\gp}}(\lambda(x)),\\
  \eta (a, x) = (t_\gp, x), & \text{if } (a, x) \in \gp \cap A \text{ and } \pi_{\gamma_{\gp}}(a) = \pi_{\gamma_{\gp}}(\lambda(x)),\\
  \eta = \id, & \text{on } A \mi C.
\end{cases}
\end{equation*}
Observe that $\eta(A)$ is doubly bounded. The function $\lambda$ is called the \memph{focus} of $\eta$, the $\RV$-pullback $C$  the \memph{locus} of $\lambda$ (or $\eta$), and the pair $(\gamma, t)$   the \memph{aperture} of $\lambda$ (or $\eta$). Note that if $(\gamma, t)$ is the aperture of $\lambda$ then every other pair $(\gamma', t')$
of this form with $\gamma' \geq \gamma$ could be an aperture of $\lambda$ as well, so the aperture of $\lambda$ must be given as a part of $\lambda$ itself. Actually, all the data above should be regarded as part of $\lambda$, including the tuples $(\alpha_{\gp}, \beta_{\gp})$.

A \memph{special covariant transformation} $T$ on $A$ is an alternating composition of centripetal transformations and regularizations. The \memph{length} of such a special covariant transformation $T$, denoted by $\lh(T)$, is the number of centripetal transformations in it.

Choose a definable point $c \in t^\sharp$ (again, this is possible since $\mathbb{S}$ is assumed to be $\VF$-generated). If $(a, x) \in \gp \cap A$ and $\pi_{\gamma_{\gp}}(a) = \pi_{\gamma_{\gp}}(\lambda(x))$  then $(\lambda(x), x) \in \gp \cap A$. Thus  the second clause of (\ref{centri}) may be changed to
\[
\eta (a, x) = (a - \lambda(x) + \lambda(x)c, x).
\]
If so, the images from the first two clauses of (\ref{centri}) may  overlap, but we take their disjoint union and thereby always assume that the resulting function $\eta^\flat$ is injective. In so doing, a special covariant transformation may be \memph{lifted} to a \memph{special covariant bijection} $T^\flat$ on $A$. This of course does depend on the choice of the definable point $c \in t^\sharp$.
The image of $T^\flat$ is often denoted by $A^{\flat}$.
\end{defn}

The formulation of a special covariant bijection here is considerably more complicated than  that of a special bijection in \cite[Definition~5.1]{Yin:special:trans}. A special covariant bijection is a special bijection if all of the apertures are $(\infty, 0)$ in the obvious sense or the second clause of (\ref{centri}) does not occur. The extra generality is needed to achieve better control in transforming one proper invariant set into another.

\begin{rem}\label{bigger:aper}
If we change the aperture $(\gamma, t)$ of $\eta$ to $(\gamma', t')$ with $\gamma' > \gamma$ and write the resulting  centripetal transformation as $\eta'$ then  every open polydisc $\gb \sub \eta(A) \mi \eta'(A)$ of radius $\alpha_\gp$, where $\gp$ is the $\RV$-polydisc that contains $\eta^{-1}(\gb)$, has an extra $\RV$-coordinate $t_\gp$ that is contracted from open discs of radius $\gamma_\gp$ in the first $\VF$-coordinate of $A$. Each of these open polydiscs has a counterpart in $\eta'(A) \mi \eta(A)$, in which $t_\gp$ is replaced by  $(\MM_{\gamma_\gp} \mi \MM_{\gamma'_\gp}) \cup t'_{\gp}$. On the other hand, $\eta(A) = \eta'(A)$ if and only if the second clause of (\ref{centri}) does not occur.
\end{rem}

\begin{rem}\label{spec:addendum}
Suppose that $A$ is proper $\delta$-invariant but is not doubly bounded. Then, by definition, there  are no special covariant transformations on $A$. We can prepare $A$, though, as follows. Choose a definable point $t_\delta \in \delta^\sharp$. For each open polydisc $\gb_1 \times \ldots \times \gb_n \times t \sub A$ of radius $\delta$, if $\gb_i= \MM_\delta$ then replace it with $t_\delta$. The resulting definable set, or rather, the resulting finite disjoint union of definable sets, is doubly bounded and remains $\delta$-invariant. We may and do regard this operation as a centripetal transformation with respect to the constant focus map $0$.
\end{rem}

% BEGIN REFEREE INSERTION R101
\RefereeRemark{101}{Gap; revised}{Choose an attainable radius and spell out the lift}{%
\textit{Original source line 2132.}\quad A definable $\delta\in\Gamma$ need not admit a definable $t_\delta\in\delta^\sharp$; for example, over $\C((t))$ one may take $\delta=\vv(t)/2$. The opening paragraph only guarantees attainability after increasing the radius. First replace $\delta$ by an attainable $\delta'\geq\delta$, which preserves invariance. Distinguish the contracted preparation, which replaces $\MM_{\delta'}$ by the RV-point $t_{\delta'}$, from its bijective lift $x\mapsto x+c$, where $\rv(c)=t_{\delta'}$; its image is $t_{\delta'}^\sharp$. Carry disjoint-union tags to distinguish overlaps with the other pieces. This translation is the additional preparation rule at the zero-centered piece; it is not obtained by substituting the focus $0$ into the expression in Definition~\ref{defn:special:bijection} $\lambda(x)c$. This distinction is needed when Proposition~\ref{simplex:with:hole:rvproduct} calls the preparation a special covariant bijection.
}
% END REFEREE INSERTION R101

\begin{lem}\label{spec:inve}
Special covariant bijections  are  partially differentiable proper covariant homeomorphisms with the constant Jacobian $1$.
\end{lem}
\begin{proof}
In the notation of Definition~\ref{defn:special:bijection}, since $\pr_{>1}(C \cap A)$ is clopen (implied by proper invariance and Lemma~\ref{inv:crit}) and  $\lambda$ is  continuous, arguing as in the proof of Lemma~\ref{conti:covar}, we see that for all sufficiently large $\epsilon \in \Gamma$ there is an even larger $\delta \in \Gamma$ such that the image under $\lambda$ of any open polydisc of radius $\delta$ contained in $\pr_{>1}(C \cap A)$ is contained in an open disc of radius $\epsilon$. Thus, for all $\RV$-polydisc $\gp \sub C$, the image under the lifted centripetal transformation $\eta^\flat$ of any open polydisc of radius $(\gamma_{\gp}, \delta)$ contained in $\gp  \cap A$ is again an open polydisc of radius $(\gamma_{\gp}, \delta)$. This means that $\eta^\flat(A)$ is proper invariant and, moreover, the domain and range of every $\VF$-fiber of $\eta^\flat$ are clopen. More generally, since the $\RV$-coordinates are carried along by centripetal transformations, it is not hard to see that this actually holds  for all  special covariant bijections. Now the claim follows from  Lemmas~\ref{conti:covar} and \ref{conv:unif} (that the Jacobian is $1$ everywhere is obvious).
\end{proof}

% BEGIN REFEREE INSERTION R102
\RefereeRemark{102}{Error}{The everywhere Jacobian assertion fails under composition}{%
\textit{Original source line 2139.}\quad The constant-Jacobian assertion is false for the stated class of continuous covariant foci: partial derivatives do not obey the chain rule at every point. Here is a counterexample within Definition~\ref{defn:special:bijection}, using only its first clause. Choose a definable $d\ne0$ with $\delta=\vv(d)>0$, put $\epsilon=2\delta$, and let $A=(1+\MM_\epsilon)^3$, with coordinates $(a,x,y)$ and the stipulated constant regularization tags. Define
\[
h(x,y)=\begin{cases}-x,&\vv(x)<\vv(y),\\0,&\vv(x)\geq\vv(y).\end{cases}
\]
The function $h$ is continuous and $1$-Lipschitz for the maximum norm, and all its partial derivatives exist; at $(0,0)$ they are $(-1,0)$. Apply three centripetal transformations, first in $x$, then in $y$, then in $a$, with respective focus maps
\[
2-d-a,\qquad 2-d-a,\qquad 1-d+h(x-d,y-d).
\]
Here the $x,y$ in the last focus are the already transformed coordinates. The successive domains are product polydiscs with centers $(1,1,1)$, $(1,d,1)$, and $(1,d,d)$ and radius $\epsilon$. All three foci lie in the target coordinate's RV-disc, are continuous and $(\epsilon,\epsilon)$-covariant, and admit aperture $(\epsilon,\rv(d^2))$. Every transformed difference has valuation $\delta<\epsilon$, so the second clause of (\ref{centri}) never occurs. Thus these are permitted special covariant bijections.

Writing $(u,v,w)=(a-1,x-1,y-1)$, their composite, after subtracting its constant output center $(d,d,d)$, is
\[
(u,v,w)\longmapsto
\bigl(u-h(v+u,w+u),\ v+u,\ w+u\bigr).
\]
Since $h(t,t)=0$, $h(t,0)=-t$, and $h(0,t)=0$, its partial-Jacobian matrix at the origin is
\[
\begin{pmatrix}1&1&0\\1&1&0\\1&0&1\end{pmatrix},
\]
whose determinant is $0$, although every individual transformation has determinant $1$ everywhere. This does not contradict the almost-everywhere change-of-variables formula. One repair is to state Jacobian $1$ only almost everywhere and check all later uses; another is to require total differentiability of the foci, with a corresponding strengthened existence proof. An almost-everywhere repair is not automatic for the standard-contraction setup below, which explicitly demands constant Jacobian $1$. The relevant chain rule in \cite[Lemma~9.7]{Yin:special:trans} requires total differentiability; almost-everywhere partial differentiability by itself cannot justify the present claim.
}
% END REFEREE INSERTION R102


So calling them ``special proper covariant bijections'' may seem more appropriate, but the verbosity is somewhat confusing.


\begin{rem}\label{spec:is:iso}
Let $(A,\omega)$, $(B, \sigma)$ be objects of $\mgVF[k]$ and $F : A \fun B$ a special covariant bijection. Then  $F$ represents a $\mgVF[k]$-isomorphism if $\omega(x) = \sigma(F(x))$ for all $x \in A$ outside a definable subset of $\VF$-dimension less than $k$ (recall Remark~\ref{mor:equi}).
\end{rem}



\begin{lem}[{\cite[Theorem~5.5]{Yin:special:trans}}]\label{special:bi:term:constant}
Suppose that $A$ is a definable set in $\VF$ and $f : A \fun \mdl P(\RV^m)$ is a definable function. Then there is a special bijection $T$ on $A$ such that $A^\flat$ is an $\RV$-pullback and the function $f \circ T^{-1}$ is $\rv$-contractible.
\end{lem}

% BEGIN REFEREE INSERTION R103
\RefereeRemark{103}{Clarification}{Explain the extension of the cited preparation theorem}{%
\textit{Original source line 2153.}\quad The cited theorem is a polynomial-term preparation statement, whereas this lemma allows an arbitrary definable relation $f\subseteq A\times\RV^m$. Add the short deduction: take a quantifier-free defining formula for that relation, simultaneously prepare its finitely many occurring VF-polynomials (also those defining $A$), and observe that the truth value of the resulting formula is constant on each transformed RV-polydisc for every RV-parameter. This gives precisely the asserted powerset-valued contraction. See the quantifier-free-formula argument in \cite[Theorem~5.5 and Corollary~5.6]{Yin:special:trans}.
}
% END REFEREE INSERTION R103

\begin{rem}\label{special:pseudo}
The proof of Lemma~\ref{special:bi:term:constant} in \cite{Yin:special:trans} actually shows that for every definable set $A$, there is a special bijection $T$ on  $\RVH(A)$  such that $A^{\flat}$ is an $\RV$-pullback. In the present context, for reasons that will become clear, we would like to extend this result, using special covariant bijections  on proper invariant sets $A$. This is not guaranteed by Lemma~\ref{special:bi:term:constant} since the focus maps in $T$ are not required to be (suitably) covariant within each $\RV$-polydisc; except when $A_{\VF} \sub \VF$, in which case  the covariance requirement is half vacuous and it is easy to see how to turn $T$ into a special covariant bijection whose components all have the same aperture. The technique underlying such a modification could be generalized so as to be applicable in other similar situations, see  Terminology~\ref{FMT} below.
\end{rem}

\begin{ter}\label{FMT}
Often in a construction there is obstructing data that is contained in finitely many distinct definable open discs $\ga_i$, all of the same radius, and is presented in the form of a definable function $f : A \fun \VF$ with $A$ a set in $\VF$ and $\ga_i \sub \pr_1(A)$. Then the \memph{FMT procedure} may be applied to circumvent the difficulties, which is simply this: choose a definable point $a_i$ in each $\ga_i$ and replace $f$ with another function $f' : A \fun \VF$ given by $(a,b) \efun f(a_i, b)$ if $a \in \ga_i$ and $(a,b) \efun f(a,b)$ otherwise, now the construction over each $\ga_i$ depends on $a_i$ and hence no longer varies.
\end{ter}

% BEGIN REFEREE INSERTION R104
\RefereeRemark{104}{Gap}{FMT needs a product domain}{%
\textit{Original source line 2161.}\quad The replacement $f(a,b)\mapsto f(a_i,b)$ is not defined from the hypotheses as written: $(a,b)\in A$ and $a,a_i\in\ga_i$ do not imply $(a_i,b)\in A$. Require that $A\cap(\ga_i\times\VF^{n-1})=\ga_i\times A_{a_i}$ (with RV-coordinates carried along), or explicitly require the corresponding domain inclusion for each replacement. The later uses of FMT must verify this condition for the current focus-map domain, not just for the original $A$.
}
% END REFEREE INSERTION R104




\begin{prop}\label{simplex:with:hole:rvproduct}
Let $H$ be a proper invariant $\RV$-pullback, not necessarily doubly bounded, and $(A_i)_i$ a definable finite partition of $H$ such that each $A_i$ is proper invariant. Then there is a special covariant bijection $T$ on $H$ such that
\begin{itemize}
  \item every $T \rest A_i$ is indeed a special covariant bijection,
  \item every $A_i^{\flat} \sub H^{\flat}$  is a doubly bounded $\RV$-pullback.
%  \item $f \circ T^{-1}$ is $\rv$-contractible.
\end{itemize}
\end{prop}
\begin{proof}
To begin with, by Remark~\ref{spec:addendum}, we may assume that $H$ is doubly bounded. It is equivalent and less cumbersome to construct a special covariant transformation such that its lift is as required (see the last paragraph of this proof). To that end, we proceed by induction on $n$, where $H_{\VF} \sub \VF^n$. The base case $n=1$ follows from Lemma~\ref{special:bi:term:constant} and the discussion in Remark~\ref{special:pseudo}.


For the inductive step, let $A_{i1} = \pr_1(A_i)$ and $H_1 = \pr_1(H)$. By Lemma~\ref{int:inv}, we may assume that $(A_{i1})_i$ form a partition of $H_1$. For each $a \in H_1$, the inductive hypothesis gives an $a$-definable special covariant transformation $T_a$ on $H_a$ such that every $T_a \rest A_{ia}$ is a special covariant transformation and every $T_a(A_{ia}) \sub T_a(H_a)$ is an $\RV$-pullback. Our goal then is to fuse together these transformations $T_a$ so to obtain one special covariant transformation on $H$ as desired. This is in general not possible without first modifying $H_1$ in a suitable way, which constitutes the bulk of the work below.

% BEGIN REFEREE INSERTION R105
\RefereeRemark{105}{Clarification}{Refine the projected partition explicitly}{%
\textit{Original source line 2178.}\quad The projections $\pr_1(A_i)$ need not be disjoint. Replace them by the Boolean atoms of the finite family $(\pr_1(A_i))_i$, and restrict each original piece over these atoms. Several nonempty pieces can still have the same first-coordinate projection; group them over that atom instead of treating the projected family as a disjoint partition indexed by the original $i$. The subsequent fiberwise induction can then be stated without changing the original partition problem.
}
% END REFEREE INSERTION R105


Let $U_{ak}$ enumerate the loci of the components of $T_{a}$, $\lambda_{ak}$ the corresponding continuous focus maps, and $(\gamma_{ak}, t_{ak})$ their apertures; for each $\RV$-polydisc $\gq \sub U_{ak}$, the map $\lambda_{ak\gq}$ is $(\alpha_{ak\gq}, \beta_{ak\gq})$-covariant. By compactness, there is a definable set $V \sub \VF \times \RV^l$ such that $\pr_1(V) = H_1$ and, for each $a \in H_1$, the set $V_a$ contains the following $\RV$-data of $T_a$:
\begin{itemize}
  \item $\rv(T_{a} \rest A_{ia})$, $\rv(T_{a})$, and $\rv(U_{ak})$,
  \item the $\VF$-coordinates targeted by the focus maps $\lambda_{ak}$,
  \item the $a$-definable apertures $(\gamma_{ak}, t_{ak})$,
   \item the $(a, \rv(\gq))$-definable tuples $(\alpha_{ak\gq}, \beta_{ak\gq})$;
\end{itemize}
the set  $\rv(T_{a})$ is determined by other data in this list and hence is redundant, but we add it in anyway for clarity. Note that $V$ is  doubly bounded but not necessarily proper invariant.

% BEGIN REFEREE INSERTION R106
\RefereeRemark{106}{Gap; revised}{Bounded fiber data need not be uniformly bounded}{%
\textit{Original source line 2188.}\quad Compactness of the choices does not establish that $V$ is doubly bounded. Even on $H_1=1^\sharp$, the definable function $\gamma(a)=\vv(a-1)$ for $a\ne1$, with $\gamma(1)=0$, is unbounded. These can actually be aperture data: on the second coordinate of $H=(1^\sharp)^2$, take the constant focus $1$, the aperture $(\gamma(a),\rv(a-1))$ for $a\ne1$, and $(0,1)$ for $a=1$. All the separate transformations are permitted, but their collected aperture data are unbounded. This does not show that a bounded choice is impossible; it shows that the asserted bound needs a choice argument. A possible repair is to postpone it until the finitely many exceptional discs have been frozen by FMT, and then prove that all the data factor through a doubly bounded RV-base. Definable fiberwise bounds and Lemma~\ref{db:to:db} can then give a uniform bound.
}
% END REFEREE INSERTION R106

Let $\phi(x, y)$ be a quantifier-free formula that defines $V$ and $G_{i}(x)$ enumerate its top terms. By Lemma~\ref{special:bi:term:constant}, there is a special bijection $R : H_1 \fun H_1^\flat$ such that each $A_{i1}^\flat \sub H_1^\flat$ is an $\RV$-pullback and every function $G_{i} \circ R^{-1}$ is $\rv$-contractible. This implies that,  for every $\RV$-polydisc $\gp \sub H_1^\flat$, the $\RV$-data $V_a$ is constant over $a \in R^{-1}(\gp)$. Also, since $H_1$ is an $\RV$-pullback of $\RV$-fiber dimension $0$, by \cmin-minimality and Lemma~\ref{RV:fiber:dim:same}, each focus map in $R$ consists of only finitely many points.

Fix a sufficiently large definable $\delta \in \Gamma$ such that every $A_i$ is proper $\delta$-invariant.  Then there is an algebraic set of open discs $\ga_j \sub H_1$ of radius $\delta$ such that the restriction of $R$ to $H_1 \mi \bigcup_j \ga_j$ is actually a special covariant bijection --- the reason simply being that, after deleting all the discs $\ga_j$, each focus map in $R$ lies outside the set in question. Each $\ga_j$ contains an $\code {\ga_j}$-definable point $a_j$ and, for all $a, a' \in \ga_j$ and every $A_i$, we have $A_{ia} = A_{ia'}$ (because  $A_i$ is $\delta$-invariant). It follows that, over each $\ga_j$, we can use the same  transformation $T_{a_j}$ to achieve the desired effect. This is the first instance where FMT is applied, namely to $T$, and then $R$ is adjusted so that every $\ga_j$ is mapped to the same $\RV$-disc $t^\sharp_{\delta}$, where $t_{\delta} \in \delta^\sharp$ is a chosen definable point.

%\begin{ter}\label{FMT}
%This operation just performed --- modifying a construction over finitely many definable open discs with a chosen definable point in each of them --- will be applied many more times below (even after this proof) and, to reduce verbosity, we shall refer to it by the acronym FMT. Which construction it is being applied to should be clear in context.
%\end{ter}

Therefore, we may assume that $R$ is a special covariant bijection whose components all have the same  aperture $(\delta, t_{\delta})$.


By compactness, there is  a definable finite partition of $H_1$ such that, over each piece, the  focus maps $\lambda_{ak}$ are uniformly defined by formulas $\lambda_k(a, y, z)$. By \cmin-minimality, there are only finitely many open discs $\ga_i \sub H_1$ of radius $\delta$ that are split by this partition. Thus, by FMT again, we may assume that the partition is indeed trivial. Since $R$ induces a special covariant bijection on $H$, we may actually assume that $R$ is trivial as well.

% BEGIN REFEREE INSERTION R107
\RefereeRemark{107}{Clarification; revised}{Explain uniformization after FMT}{%
\textit{Original source line 2201.}\quad Here ``trivial'' presumably means that one global definition of the focus maps suffices. Explain the finite gluing step: FMT makes the formula-choice pieces clopen, but does not by itself ensure that an operation on one piece extends by the identity as a permitted centripetal transformation, whose locus must be an RV-pullback. Include the finite formula-choice partition among the sets prepared by $R$, pad shorter transformations by identities, and glue the focus definitions over the resulting RV-pullback pieces, checking their loci. This is the intended uniformization. Similarly, taking $R$ to be trivial means replacing $H$ by its transformed image and transporting the entire partition and all fiberwise data along $R$.
}
% END REFEREE INSERTION R107


Over each $t^\sharp \sub H_1$, we can now write $U_{ak}$ as $U_{tk}$, $\alpha_{ak\gq}$ as $\alpha_{k\gq}$ (the first $\RV$-coordinate of $\gq$ is actually $t$), and so on. We are almost ready to fuse together the transformations $T_a$ over $a \in t^\sharp$. The remaining problem  is that, for any $a, a' \in t^\sharp$, although the two focus maps $\lambda_{a1\gq}$, $\lambda_{a'1\gq}$ are both $(\alpha_{1\gq}, \beta_{1\gq})$-covariant,  the images of the same open polydisc of radius $\alpha_{1\gq}$ may lie in two distinct open discs of radius $\beta_{1\gq}$. To solve this problem, consider an open polydisc $\gp$ of radius $\alpha_{1\gq}$ that is contained in $\dom(\lambda_{a1\gq})$ for some (hence all) $a \in t^\sharp$. For each $b \in \gp$, let $\lambda_{1b} : t^\sharp \fun \VF$ be the function defined by $\lambda_1(x, b, z)$. By \cmin-minimality and Corollary~\ref{part:rv:cons}, there are a $b$-definable finite set $C_b \sub \VF$ and, for any $a \in t^\sharp \mi C_b$, an open disc $\ga_a \sub t^\sharp \mi C_b$ around $a$ such that $\lambda_{1b}(\ga_a)$ lies in an open disc of radius $\beta_{1\gq}$. Since for any other $a' \in \ga_a$,  $\lambda_{a'1\gq}(\gp)$ also lies in an open disc of radius $\beta_{1\gq}$, we see that $\lambda_1(\ga_a \times \gp)$ lies in an open disc of radius $\beta_{1\gq}$, where $\lambda_1$ stands for the function defined by $\lambda_1(x, y, z)$. Therefore, we may assume that the finite set $C_b$ is actually $\code \gp$-definable. But then, by Lemma~\ref{ball:to:ac}, it is even definable, that is, it does not depend on $\gp$ (or rather it works for all $\gp$). By compactness and FMT, we may assume that there is an $\go$-partition $p: \dom(\lambda_1) \fun \Gamma$ such that, for each open polydisc $\gb$ in question,  $\lambda_1(\gb)$ lies in a disc of radius $\beta_{1\gq}$, where $\gq$ is related to $\gb$ in the obvious way. By Lemma~\ref{vol:par:bounded}, the image of $p$ is doubly bounded and, by Lemma~\ref{db:to:db}, there is a definable $\gamma_1 \in \Gamma$ with $\gamma_1 \geq \gamma_{t1}$ for all $t \in \rv(H_1)$. We may choose such a $\gamma_1$ together with a definable point $t_1 \in \gamma^\sharp_1$ (see the opening paragraph of this section). All this means that $\lambda_1$ can serve as the focus map of a centripetal transformation $T_1$ on $H$ of aperture $(\gamma_{1}, t_{1})$. To ensure that $\lambda_1$ is continuous, we can appeal to Proposition~\ref{prop:continuity:fiberwise} and, once again, FMT.

% BEGIN REFEREE INSERTION R108
\RefereeRemark{108}{Clarification; revised}{Use a canonical exceptional set before eliminating ball parameters}{%
\textit{Original source line 2204.}\quad Use a canonical exceptional set instead of an arbitrary chosen $C_b$. With the other RV labels $z$ fixed, let $E_{\gp,z}$ consist of the $a$ at which no neighborhood $\ga$ has $\lambda_1(\ga\times\gp)$ contained in one radius-$\beta_{1\gq}$ disc. Covariance on $\gp$ shows that $E_{\gp,z}\subseteq C_b$ for every $b\in\gp$, so this is a finite set definable from the polydisc code and the RV labels. Lemma~\ref{ball:to:ac} removes the polydisc code from algebraic dependence; Lemma~\ref{dcl:to:ac} then removes the remaining RV labels. The union of the uniformly definable family $(E_{\gp,z})$ is contained in $\acl(\mathbb S)$, and hence is finite by compactness. This supplies the common finite exceptional set needed for FMT. Algebraicity of each separate chosen $C_b$ would not on its own make that choice independent of $\gp$.
}
% END REFEREE INSERTION R108

At this point the proof would be complete if we could repeat the procedure above for $\lambda_2$, and so on. We still have a small issue, namely some part of the locus $U_{a2}$ may have disappeared because the aperture of  $\lambda_{a1}$ is bumped up to $(\gamma_{1}, t_{1})$; see Remark~\ref{bigger:aper}. It is not hard to see that the inductive hypothesis may be applied to the $\RV$-pullback $W$ contained in $T_1(H)$ that corresponds to the missing locus, ignoring the coordinate of $W$ that was targeted by $T_1$, because the missing locus must have one less $\VF$-coordinate and its intersection with each $T_1(A_i)$ is proper invariant.
\end{proof}

% BEGIN REFEREE INSERTION R109
\RefereeRemark{109}{Gap}{Complete the aperture-change induction}{%
\textit{Original source line 2207.}\quad The final paragraph leaves the induction unfinished: the text does not construct the replacement transformations on the part created when an aperture increases, nor explain why the remaining loci still define valid transformations. The useful precise assertion would be that the affected part $W$ is a pullback under a projection forgetting the newly released VF-coordinate, that every $W\cap T_1(A_i)$ is a pullback under the same projection, and that the projected sets are proper invariant. Then apply the $(n-1)$-dimensional induction to that projected partition, lift its transformations while carrying the released coordinate along, and verify compatibility with the remaining loci. Give the explicit projection/factorization and the finite induction on the remaining components. Merely saying that the old missing locus has one fewer VF-coordinate does not prove these facts for $W\subseteq T_1(H)$.
}
% END REFEREE INSERTION R109

This proposition can indeed be applied to any proper invariant set $A$ against its $\RV$-hull $\RVH(A)$, without specifying a partition of $\RVH(A)$ as required, because, by  Lemma~\ref{int:inv},   $\RVH(A) \mi A$ is proper invariant too.

\begin{rem}\label{cons:extra}
As in the statement of Lemma~\ref{special:bi:term:constant} above, if $f : H \fun \mdl P(\RV^m)$ is a proper covariant function then we can demand as a clause of the conclusion of Proposition~\ref{simplex:with:hole:rvproduct} that the function $f \circ T^{-1}$ be $\rv$-contractible, which would follow if it is simply included as a clause of the inductive hypothesis in the proof.
\end{rem}

Let $A$ be a definable set and $\omega : A \fun \Gamma$ a definable function. We say that the definable pair $\bm A = (A, \omega)$ is \memph{proper invariant} if $A$ is proper invariant and $\omega$ is proper covariant, that is, every fiber $A_\gamma$ is proper invariant.

\begin{lem}\label{L:measure:surjective}
Suppose that $A$ is an $\RV$-pullback and  $(A_i)_i$ is a definable finite partition of $A$ such that each $\bm A_i = (A_i, \omega \rest A_i)$ is proper invariant. Then there is a special covariant bijection $T$ as given in Proposition~\ref{simplex:with:hole:rvproduct} such that $\omega \circ T^{-1}$ is $\rv$-contractible. In particular, if $\bm A_i \in \mVF[k]$ then $T$ induces a morphism  $\bm A_i \fun \mL(\bm U_i, \pi_i)$, where $(\bm U_i, \pi_i) \in \mRV^{\db}[k]$.
\end{lem}
\begin{proof}
We view $\omega$ as a definable function $A \fun \mdl P(\RV)$. By Proposition~\ref{simplex:with:hole:rvproduct}, there is a special covariant bijection $T$ on $A$ such that, for each $i$, $T \rest A_{i}$ is a special covariant bijection and $A_{i}^\flat$ is an $\RV$-pullback. Moreover, applying Remark~\ref{cons:extra}, we get that each function $\sigma_i \coloneqq (\omega \circ T^{-1}) \rest A_{i}^\flat$ is constant on every $\RV$-polydisc in $A_{i}^\flat$ and hence induces a function $\pi_i : \rv(A_{i}^\flat) \fun \Gamma$. Now, if $\bm A_i \in \mVF[k]$ then, by Lemma~\ref{RV:fiber:dim:same}, $\bm U_i\coloneqq \rv(A_{i}^\flat)_{\leq k} \in \RV^{\db}[k]$ (recall Notation~\ref{0coor}). We see that $(\bm U_i, \pi_i)$ is as required because, by Remark~\ref{spec:is:iso}, $\bm A_i$ is isomorphic to $(\mathbb{L} \bm U_i, \sigma_i)$,
\end{proof}

% BEGIN REFEREE INSERTION R110
\RefereeRemark{110}{Typo}{Sentence-ending punctuation}{%
\textit{Original source line 2222.}\quad Replace the final comma after ``$(\mathbb L\bm U_i,\sigma_i)$'' by a full stop.
}
% END REFEREE INSERTION R110

As above, this lemma may be applied to any proper invariant object $(A, \omega) \in \mVF[k]$, since  the object $(\RVH(A) \mi A, 0) \in \mVF[k]$  is also proper invariant.


% BEGIN REFEREE INSERTION R170
\RefereeRemark{170}{Error; new}{The RV-hull can create infinite RV-fibers}{%
\textit{Original source line 2225.}\quad The auxiliary complement need not belong to $\mVF[k]$. For a definable $\delta>0$, put $U=\{t\in\RV:0<\vrv(t)<\delta\}$ and
\[
 A=\{(x,t)\in(1+\MM)\times U:t=\rv(x-1)\}.
\]
With constant volume form $0$, this is a proper $\delta$-invariant object of $\mVF[1]$: each RV-fiber over $x$ has at most one point. Its regularization merely adds the constant tag $\rv(x)=1$. However,
\[
 \RVH(A)=(1+\MM)\times U,
\]
and the fiber of $\RVH(A)\setminus A$ over $x=1$ is the infinite set $U$. Thus the complement is not in $\VF[1]$. The application of the lemma is nevertheless repairable: delete ``$\in\mVF[k]$'' for this auxiliary piece. The lemma permits arbitrary proper invariant pieces, and only its ``in particular'' conclusion requires the selected original piece $(A,\omega)$ to lie in $\mVF[k]$.
}
% END REFEREE INSERTION R170


\begin{lem}\label{invar:db}
Let $A$ be a definable set with $A_{\VF}$  bounded and $A_{\RV}$  doubly bounded. Let $f : A \fun \Gamma$ be a definable function  that is constant on $A \cap \ga$ for all open polydiscs $\ga$ of radius $\gamma \in \Gamma$. Then $f(A)$ is doubly bounded.
\end{lem}
\begin{proof}
Observe that, without loss of generality, we can replace $A$ with its $\gamma$-tubular neighborhood and thereby  assume that $(A,f)$ is indeed proper invariant. By Lemma~\ref{L:measure:surjective}, $f$ factors through a doubly bounded definable set in $\RV$. So, by Lemma~\ref{db:to:db},  $f(A)$ must be doubly bounded.
\end{proof}

\begin{cor}
If $(A, \omega)$ is proper invariant then $\omega(A)$ is doubly bounded.
\end{cor}
\begin{proof}
We have seen this many times above: as a consequence of Lemma~\ref{vol:par:bounded}, for all sufficiently large $\beta \in \Gamma$, every fiber $A_\gamma$ of $\omega$ is $\beta$-invariant. So the claim is immediate by Lemma~\ref{invar:db}.
\end{proof}

\subsection{Standard contractions}

\begin{defn}\label{rela:unary}
Let $A \sub \VF^{n} \times \RV^{m}$, $B \sub \VF^{n} \times \RV^{m'}$ be definable sets. A definable bijection $G : A \fun B$ is \memph{relatively unary}, or more precisely, \memph{relatively unary in the $(i, j)$-coordinate}, where $i, j \in [n]$, if $(\pr_{\tilde{j}} \circ G)(x) = \pr_{\tilde{i}}(x)$ for all $x \in A$. If, in addition, $G \rest A_a$ is also a special covariant bijection for every $a \in \pr_{\tilde{i}} (A)$ then $G$ is \memph{relatively special covariant} in the $(i, j)$-coordinate.

Let $\bm U = (U, f)$, $\bm V = (V, g)$ be objects of $\RV[k]$. A morphism $F : \bm U \fun \bm V$ is \memph{relatively unary in the $(i, j)$-coordinate}, where $i,j \in [k]$, if $(\pr_{\tilde{j}} \circ g \circ F)(u) = (\pr_{\tilde{i}} \circ f)(u)$ for all $u \in U$.
\end{defn}

Alternatively, for simplicity, we may take $i = j$ in this definition and stipulate that the condition in question holds up to coordinate permutations.

Since identity functions are relatively unary in any coordinate, if a  bijection  is piecewise (over a definable finite partition) a composition of relatively unary bijections then it is indeed a composition of relatively unary bijections.

\begin{exam}
Every  special covariant bijection $T$ of length $1$ is relatively special covariant, but clearly not vice versa (to begin with, the relevant data is not necessarily uniform across the fibers in each $\RV$-polydisc).
\end{exam}

Let $A \sub \VF^{n} \times \RV^{m}$ be a proper invariant set. Let $i \in [n]$ and $T_i : A \fun A^\flat$ be a proper covariant bijection that is
\begin{itemize}
  \item relatively unary in the $(i, i)$-coordinate (up to coordinate permutations),
  \item partially differentiable with the constant Jacobian $1$.
\end{itemize}
Moreover, the $\RV$-coordinates of $A$ are passed on to $A^\flat$ via $T_i$ in the obvious way (the simplest example is that $T_i$ is the regularization of $A$). Adjusting $T_i$ by way of Remark~\ref{spec:addendum} if necessary, we take $A$ and $A^\flat$ to be both doubly bounded. By Lemma~\ref{conv:unif}, $A^\flat$ is proper invariant as well, and $T_i$ is a proper covariant homeomorphism. Suppose that, for all $a \in A_{\tilde i}$ and all $t \in (A_{a})_{\RV}$, the image $T_i(A_{a, t}) \sub A_{a, t}^\flat$ of the fiber $A_{a, t} \sub A$ is an $\RV$-pullback.  Let

% BEGIN REFEREE INSERTION R111
\RefereeRemark{111}{Error}{Continuity is missing from the standard-contraction hypotheses}{%
\textit{Original source line 2263.}\quad Lemma~\ref{conv:unif} only gives a homeomorphism if continuity has already been assumed. Proper covariance plus everywhere partial differentiability and Jacobian $1$ does not imply continuity. For example, put
\[
q(x,y)=\begin{cases}x/y,&y\ne0\ \text{and}\ \vv(x)\geq\vv(y),\\0,&\text{otherwise}.\end{cases}
\]
This bounded function has every partial derivative everywhere, but $q(t,t)=1$ for $t\ne0$ and $q(0,0)=0$. For a definable $0\ne d\in\MM$, the map
\[
T(a,x,y)=\bigl(a+dq(x-1,y-1),x,y\bigr)
\]
on $A=(1+\MM)^3$ is a relatively unary proper covariant bijection $A\to A$, with constant partial Jacobian $1$ and RV-pullback fibers, yet is discontinuous at $(1,1,1)$. All regularization tags are constant, so forgetting or retaining them does not change this example. Add ``continuous'' to the hypotheses on $T_i$ before invoking Lemma~\ref{conv:unif}.
}
% END REFEREE INSERTION R111
\[
 A_{-i} = \bigcup_{a \in A_{\tilde i}} a \times (A_{a}^\flat)_{\RV} \sub \VF^{n-1} \times \RV^{m_i},
\]
which is doubly bounded and proper invariant, and  we think of it as obtained from $A^\flat$ by ``$\rv$-contracting'' the $\VF$-coordinate in question.  Also, under Convention~\ref{conv:can}, $A^\flat$ can be ``fully recovered'' from  $A_{-i}$. Let $\hat T_i : A \fun A_{-i}$ be the function induced by $T_i$.

For any $j \in [n{-}1]$, we repeat the above setup on $A_{-i}$ with respect to the $j$th $\VF$-coordinate, and thereby obtain a doubly bounded proper invariant set $A_{-j} \sub \VF^{n-2} \times \RV^{m_j}$ and a function $\hat T_{j} : A_{-i} \fun A_{-j}$ that is induced by a relatively unary, partially differentiable, proper covariant homeomorphism $T_j : A^\flat \fun A^{\flat\flat}$ (the $\VF$-coordinate in $A^\flat$ that gets $\rv$-contracted in $A_{-i}$ is simply carried along by $T_j$). Continuing thus,  a sequence of such homeomorphisms  $T_{\sigma(1)}, \ldots, T_{\sigma(n)}$ and a corresponding function $\hat T_{\sigma} : A \fun \RV^{l}$ result, where $\sigma$ is the permutation of $[n]$ in question. The composition $T_{\sigma(n)} \circ \ldots \circ T_{\sigma(1)}$, which is referred to as the \memph{lift} of $\hat T_{\sigma}$, is denoted by $T_{\sigma}$. By Lemma~\ref{pr:co:home:com}, $T_{\sigma}$ is a proper covariant homeomorphism. Also note that $T^{-1}_{\sigma}$ is $\rv$-contractible.


\begin{defn}\label{defn:standard:contraction}
The function $\hat T_{\sigma}$, or the set $\hat T_{\sigma}(A)$, is called a \memph{standard contraction} of  $A$.
\end{defn}



\begin{lem}\label{exit:stand}
Every proper invariant set $A \sub \VF^{n} \times \RV^{m}$  admits a standard contraction with respect to any permutation $\sigma$ of $[n]$.
\end{lem}
\begin{proof}
We first remark that, although the lift of a standard contraction is in effect a more general (or rather less stringent) notion than that of a special covariant bijection, for existence, it is enough to exhibit a special covariant bijection that is also the lift of a standard contraction. To that end, the proof of Proposition~\ref{simplex:with:hole:rvproduct} almost works. The  issue is that a standard contraction is required to transform fiberwise a $\VF$-coordinate into an $\RV$-pullback in one fell swoop and subsequently carry the result in the $\RV$-coordinates, whereas in a special covariant bijection one is allowed to come back to modify the same $\VF$-coordinate many times over. We can of course simply add a clause in the inductive hypothesis to fulfill the new requirement. In that case, the special covariant bijection induced by $R$ there --- let us call it $R^*$ here --- should also be made to comply. This needs work: it is clear how to transform fiberwise the first $\VF$-coordinate into an $\RV$-pullback on top of $R^*$, but continuity is not automatic. Fortunately, it is only for convenience that  $R^*$ is carried out in the beginning, that is, we may equivalently graft the $\RV$-coordinates of $\ran(R)$ onto $A$ directly in the obvious way (strictly speaking this operation is not allowed in  a special covariant bijection, and that is why we cannot really opt for it in the proof of Proposition~\ref{simplex:with:hole:rvproduct}) and carry out $R^*$ in the end. With this modification, the same construction yields a standard contraction as desired.
\end{proof}

% BEGIN REFEREE INSERTION R112
\RefereeRemark{112}{Gap}{Check the deferred transformation and fixed order}{%
\textit{Original source line 2283.}\quad To complete this existence argument, explicitly define the grafting map $a\mapsto(a,\tau(a))$, where $\tau$ records the RV-coordinates produced by $R$, and show that its fibers are uniformly invariant, so grafting is a proper covariant homeomorphism. Then specify the transported relatively unary maps while $R^*$ is deferred, and verify that each already contracted VF-coordinate is subsequently only carried along. These are the conditions defining a standard contraction; the last sentence currently asserts rather than establishes them. This proof also depends on resolving the aperture-change induction above.
}
% END REFEREE INSERTION R112

\begin{rem}\label{stand:vol}
Suppose that there is a $k \in 0 \cup [m]$ such that  $(A_{a})_{\leq k} \in \RV^{\db}[k]$  for  every $a \in A_{\VF}$. If $k=0$ then this just means $A \in \VF_*$. By Lemma~\ref{RV:fiber:dim:same}, $\hat T_{\sigma}(A)_{\leq n+k}$ is an object of $\RV^{\db}[n{+}k]$. Let $\bm A = (A, \omega)$ be a proper invariant pair such that the proper covariant function $\omega : A \fun \Gamma$   induces a function $\omega_{\hat T_{\sigma}}$ on $\hat T_{\sigma}(A)$ via $T_{\sigma}$, which just means that $\omega \circ T^{-1}_{\sigma}$ is invariant on $\RV$-polydiscs. Then the function $\hat T_{\sigma}$, or the object
\[
\hat T_{\sigma}(\bm A)_{\leq n+k} = (\hat T_{\sigma}(A), \omega_{\hat T_{\sigma}})_{\leq n+k} \in \mRV^{\db}[n{+}k],
\]
is understood as  a \memph{standard contraction} of $\bm A$. Using Lemma~\ref{exit:stand} instead of Proposition~\ref{simplex:with:hole:rvproduct} in the proof of Lemma~\ref{L:measure:surjective}, we see that such a $\hat T_{\sigma}$ does exist. Moreover, if we set $A^\sharp = \bigcup_{a \in A_{\VF}} a \times \bb L (A_a)_{\leq k}$, which is a  proper invariant object of $\VF_*$, and denote the  proper covariant function on $A^\sharp$ induced by  $\omega$ still by $\omega$, then  $T_{\sigma}$  indeed induces a $\mVF[n{+}k]$-morphism
\[
(A^\sharp, \omega) \fun {\mL}\hat T_{\sigma}(\bm A)_{\leq n+k}
\]
that can be written as a composition of relatively unary proper covariant homeomorphisms.

Here $k$ is dubbed the \memph{head start} of $\hat T_{\sigma}$, which is usually implicit in the sequel. In fact,  it never has to be anything other than $0$ except in Lemma~\ref{isp:VF:fiberwise:contract}, and can be circumvented even there. This seemingly needless gadget only serves to make the discussion more streamlined: if $A \in \VF_*$ then the intermediate steps of a standard contraction of $A$ may or may not result in objects of $\VF_*$ and hence the definition cannot be formulated entirely within $\VF_*$.
\end{rem}


%So $\bm A^\sharp \coloneqq (A^\sharp, \omega)$ is an object of $\mVF^{\diamond}[n{+}k]$.

\subsection{Lifting from $\RV$ to $\VF$}






\begin{lem}\label{bijection:partitioned:unary}
Every morphism in $\RV[k]$ can be written as a composition of relatively unary morphisms, similarly in $\mRV[k]$, $\RV^{\db}[k]$, and $\mRV^{\db}[k]$.
\end{lem}
\begin{proof}
It is enough to show this piecewise, and the situation can be  simplified accordingly step by step. Let $F : (U, f) \fun (V, g)$ be an $\RV[k]$-morphism. The idea is to adjust  $f(u)$ towards $g(F(u))$, simultaneously for all $u \in U$, one coordinate at a time, and the issue is how to stay in $\RV[k]$, that is, how to make the adjusted $f$ remain finite-to-one along the way. To that end, we proceed by induction on $n = \dim_{\RV}(U)$ and, without loss of generality, may take $F = \id$. The base case $n=0$ is straightforward and is left to the reader. For the inductive step, we may assume that both $\pr_{\leq n} \rest f(U)$ and $\pr_{\leq n} \rest g(V)$ are finite-to-one. Let $f' = \pr_{\leq n} \circ f$ and $g' = \pr_{\leq n} \circ g$. Since, as in the base case, the two obvious $\RV[k]$-morphisms
\[
(U,f) \fun (U \times 1, f' \times 1), \quad (U \times 1, g' \times 1) \fun (U, g)
\]
are compositions of relatively unary morphisms, it is enough to show that $(U, f') \fun (U, g')$ is such a morphism as well. So we may assume $n=k$.

Write $\dot U = f(U)$ and $\ddot U = g(U)$. For each $t \in \dot U_{< n}$ and each $i \in [n]$, let
$F^\dag_{t, i} : \dot U_t \fun \mdl P(\ddot U_i)$
be the definable finitary function induced by $F^\dag$ (recall Notation~\ref{mor:ftf}). Then, by Lemma~\ref{finite:fib:def}, there is a $t$-algebraic subset $B_{t, i} \sub \ddot U_i \sub \RV$ such that every fiber of $F^\dag_{t, i}$ over an element in $\ddot U_i \mi B_{t, i}$ is finite; denote by $W_{t, i} \sub \dot U$ the union of such fibers. Indeed, by the compactness argument thereof,  $B_{t, i}$ can be defined uniformly. Let $U_{t, i}$ be the set of those $u \in f^{-1}( t \times W_{t, i})$ such that $g(u)_i \notin B_{t, i}$, and $f'_{t, i} : U_{t, i} \fun \RV^n$ the function given by $u \efun (t, g(u)_i)$, which is finite-to-one. Let $U_i = \bigcup_{t}  U_{t, i}$ and $f'_i = \bigcup_{t}  f'_{t, i}$. Then $(U_i, f \rest U_i) \fun (U_i, f'_i)$ is a relatively unary $\RV[n]$-morphism. By applying the inductive hypothesis fiberwise to $(U_i, f'_i) \fun (U_i, g \rest U_i)$ in the obvious sense, we may actually assume $U_i = \0$, for every $i \in [n]$. Now, since $F^\dag(\dot U_t) \sub \prod_{i \in [n]} B_{t, i}$, which is finite, we must have $\dim_{\RV}(U) < n$. So the claim follows from the inductive hypothesis.

If $F$ is indeed a $\mRV[k]$-morphism then it is straightforward to equip the object in each intermediate step with a volume form so that the map in question  becomes a $\mRV[k]$-morphism. The other cases are rather similar.

% BEGIN REFEREE INSERTION R114
\RefereeRemark{114}{Clarification}{Give the intermediate volume form}{%
\textit{Original source line 2323.}\quad For an intermediate finite-to-one coordinate map $h:U\to\RV^k$, the required volume form is explicitly
\[
\omega_h(u)=\omega(u)+\Sigma\vrv(f(u))-\Sigma\vrv(h(u)).
\]
Thus $\omega_h+\Sigma\vrv\circ h=\omega+\Sigma\vrv\circ f$, and the volume identities telescope through the composition. If $f'$ has only $n<k$ coordinates, write $f'\times\mathbf1_{k-n}$ in the displayed intermediate objects to make the ambient degree unambiguous.
}
% END REFEREE INSERTION R114
\end{proof}


\begin{lem}\label{RVlift}
Let $U$, $V$ be sets in $\RV$ and $f : U \fun \mdl P(V)$ a definable finitary function. Then there is a definable finitary function $f^\sharp : U^\sharp \fun \mdl P(V^\sharp)$ that $\rv$-contracts to $f$. Moreover, for every $(u, v) \in f$, $f^\sharp_{u, v} \coloneqq (u, v)^\sharp \cap f^\sharp$ is a partially differentiable function $u^\sharp \fun v^\sharp$.
\end{lem}

Similar to Definition~\ref{defn:special:bijection}, this is another place where we need the property that every algebraic $\RV$-disc contains an algebraic point, and the easiest way to achieve
this is to assume that the substructure $\bb S$ is $\VF$-generated, which validates the repeated use of Lemma~\ref{algebraic:balls:definable:centers} below.

\begin{proof}
We proceed by induction on $\dim_{\RV}(U) = n$. For the base case $n=0$, $U$ is finite and hence, for every $u \in U$, the $\RV$-polydisc $u^\sharp$ contains a $u$-definable point, and similarly for every element in $f(u) \sub V$. So lifting $f$ as desired is trivial. For the inductive step, we may assume that $\pr_{\leq n} \rest U$ is finite-to-one. Let $U' = \pr_{\leq n}(U)$ and $f' : U' \fun \mdl P(V)$ be the definable finitary function induced by $f$. Observe that if $f'$ can be lifted as desired then  $f$ can be lifted (again trivially) as desired as well. So, without loss of generality,  $U \sub \RV^n$. By Lemma~\ref{algebraic:balls:definable:centers} and compactness, there is indeed a definable finitary function $f^\sharp : U^\sharp \fun \mdl P(V^\sharp)$ that $\rv$-contracts to $f$. By Lemma~\ref{diff:almost:every}, every $f^\sharp_{u, v}$ is partially differentiable almost everywhere. By Lemma~\ref{RV:bou:dim}, the $\RV$-boundary of the differential locus is of $\RV$-dimension less than $n$. The lemma follows.
\end{proof}

% BEGIN REFEREE INSERTION R115
\RefereeRemark{115}{Gap}{Supply the single-valued lift and exceptional-locus induction}{%
\textit{Original source line 2336.}\quad Spell out two steps hidden here. For $a\in U^\sharp$, centers are chosen over the full VF-tuple $a$, not over an arbitrary nonalgebraic $u\in U$: the RV-elements in $f(\rv(a))$ are algebraic over $a$, so the center lemma and compactness give a finite $a$-definable family meeting their RV-polydiscs. To obtain one output in each RV-polydisc, average the finitely many chosen points in that polydisc (coordinatewise); residue characteristic $0$ ensures the average remains in it. Finally remove the RV-hull of the nondifferentiability locus, prove that its RV-parameter set has dimension $<n$, and replace the lift over that whole set by the inductively constructed one. The dimension bound on a boundary alone is not the construction of an everywhere partially differentiable lift.
}
% END REFEREE INSERTION R115

Let $F: \bm U \fun \bm V$ be an $\RV[k]$-morphism, where $\bm U = (U, f)$ and $\bm V = (V, g)$. If $F^{\sharp} : \bb L \bm U \fun \bb L \bm V$ is a definable bijection that $\rv$-contracts to the function $U_f \fun V_g$ induced by $F$ (recall Notation~\ref{RVcat:tag}) then it is called a \memph{lift} of $F$. We shall also think of the finite-to-finite correspondence $F^\dag$ as the $\rv$-contraction of such a lift.

%Assume that the substructure $\mathbb{S}$ is $(\VF, \Gamma)$-generated for the remainder of this section.
%\arthur{We can remove this sentence since $\mathbb{S}$ is assumed $\VF$-generated since the beginning of the section.}

%\begin{lem}[{\cite[Proposition~6.1, Lemma~6.2]{hrushovski:kazhdan:integration:vf}, \cite[Corollary~7.7]{Yin:special:trans}}]\label{L:semigroup:hom}
% $F$ admits a lift.
%\end{lem}

\begin{lem}[{\cite[Lemma~6.3]{hrushovski:kazhdan:integration:vf}}]\label{iso:lifted:vol}
There is a lift $F^{\sharp}$ of $F$ such that
\begin{equation}\label{jac:mat}
\vv(\jcb_{\VF} F^{\sharp}(a,  u)) = \jcb_{\Gamma} F^\dag(u), \quad \text{for all } (a, u) \in \bb L \bm U.
\end{equation}
\end{lem}

See \cite[Corollary~7.7]{Yin:special:trans} for an alternative proof (the weaker qualifier ``for almost all $(a, u) \in \bb L \bm U$'' there may be upgraded to ``for all $(a, u) \in \bb L \bm U$'' here by appealing to Lemma~\ref{RV:bou:dim} at suitable places).

\begin{rem}\label{rvdisc:stre}
Let $t, s \in \RV$ such that the $\RV$-discs $t^\sharp$, $s^\sharp$ contain definable points. It is easy to see that for any definable $c \in \VF^\times$ with $\vv(c) = \vrv(s/t)$, there is a definable bijection $f : t^\sharp \fun s^\sharp$ such that $\ddx f = c$.
\end{rem}


\begin{lem}\label{L:measure:class:lift}
Every $\mRV[k]$-morphism $F : (\bm U, \omega) \fun (\bm V, \sigma)$ admits a lift
\[
F^\sharp: {\mu}\bb{L}(\bm U, \omega) \fun {\mu}\bb{L}(\bm V, \sigma)
\]
that is piecewise a composition of partially differentiable relatively unary $\mVF[k]$-morphisms. Moreover, if $F$ is a $\mRV^{\db}[k]$-morphism then every component of the composition is a  proper covariant homeomorphism and hence so is $F^\sharp$.
\end{lem}
\begin{proof}
We do induction on $\dim_{\RV}(U) = n$. For the base case $n=0$, we may assume that $U$ is just a singleton and $F$ is relatively unary. Since  every $\RV$-disc $t^\sharp$ involved contains a $t$-definable point, it is easy to lift $F$  as desired by applying Remark~\ref{rvdisc:stre} in the coordinate in question.

For the inductive step, again, we may assume that $f' = \pr_{\leq n} \circ f$, $g' = \pr_{\leq n} \circ g$ are finite-to-one. Write $\bm U' = (U, f')$ and $\bm V' = (V, g')$. By Lemma~\ref{RVlift}, there is a definable partially differentiable finitary function that $\rv$-contracts to the finitary function $f(U)_{\leq n} \fun \mdl P(f(U)_{> n})$. This means that the obvious $\mRV[k]$-morphism $(\bm U, \omega) \fun (\bm U' \times 1, \omega')$ can be lifted as desired by applying Remark~\ref{rvdisc:stre} fiberwise as in the base case above, and similarly for the obvious $\mRV[k]$-morphism $(\bm V' \times 1, \sigma') \fun (\bm V, \sigma)$. Thus we are reduced to the case $n = k$. By Lemma~\ref{bijection:partitioned:unary}, we may assume that $F$ is  relative unary. Then, applying Lemma~\ref{iso:lifted:vol} fiberwise in the obvious way, we can lift $F$ to a relatively unary bijection. By Lemmas~\ref{diff:almost:every} and~\ref{RV:bou:dim}, the lift is partially differentiable outside a definable subset of $\RV$-dimension less than $k$, and the restriction  is indeed a $\mVF[k]$-morphism due to the condition (\ref{jac:mat}). So we are done by  the inductive hypothesis.

The second claim is a consequence of Lemmas~\ref{conti:covar}, \ref{conv:unif}, and \ref{pr:co:home:com}.

% BEGIN REFEREE INSERTION R116
\RefereeRemark{116}{Gap; revised}{The continuous lifting conclusion needs an additional check}{%
\textit{Original source line 2373.}\quad The preceding construction establishes partial differentiability, whereas the cited covariance/homeomorphism lemmas require continuity as an input. Partial differentiability does not imply continuity, as the example at the standard-contraction setup shows. In the generic lifting step also remove the RV-hull of the noncontinuity locus, using Lemma~\ref{fun:almost:cont} together with Lemma~\ref{RV:bou:dim}; formulate the induction so that it constructs continuous branches as well. For the finitary center functions, do this uniformly for all the finitely many RV-indexed branches: the projection of their bad locus to VF still has dimension $<n$, so its whole RV-hull is indexed by an RV set of dimension $<n$. The resulting pieces are whole RV-pullbacks, hence clopen, and their continuous lifts glue. Check these center functions when forming the triangular corrections as well. Also change ``relative unary'' in the preceding paragraph to ``relatively unary''. This is a proposed completion of the proof, not a counterexample to the lifting statement.
}
% END REFEREE INSERTION R116
\end{proof}


\begin{defn}\label{defn:diamond}
The subcategory $\mVF^{\diamond}[k]$ of $\mVF[k]$ consists of the proper invariant objects and the morphisms that are compositions of relatively unary proper covariant homeomorphisms.
\end{defn}

\begin{rem}\label{dia:cat:well}
Obviously the composition law holds in  $\mgVF^{\diamond}[k]$ and hence it is indeed a category. Moreover, every morphism in it is an honest  bijection, as opposed to merely an essential bijection, and is in effect required to admit an inverse.  Thus, $\mgVF^{\diamond}[k]$ is already a groupoid and there is no need to pass to a quotient category as in Remark~\ref{mor:equi}. On the other hand, Proposition~\ref{simplex:with:hole:rvproduct}, Lemma~\ref{spec:inve}, and Remark~\ref{stand:vol} give various constructions of nontrivial morphisms.
\end{rem}

\begin{rem}
The main reason that we need the lift in Lemma~\ref{L:measure:class:lift} to be in that particular form is that, as it will become clear, we are forced to work with explicit compositions of relatively unary proper covariant morphisms, due to the failure of generalizing Lemma~\ref{simul:special:dim:1} below to higher dimensions; see \cite[\S~1]{Yin:int:acvf} for further explanation.

Every definable bijection between definable sets with the same number of $\VF$-coordinates is a composition of relatively unary bijections, see \cite[Lemma~5.6]{Yin:int:acvf}.   It is certainly more satisfying to simply take proper covariant homeomorphisms as morphisms in $\mVF^{\diamond}[k]$ and then show that they are all indeed compositions of relatively unary proper covariant homeomorphisms. We suspect this is not true, although it does not seem straightforward to come up with a counterexample.
\end{rem}




\begin{cor}\label{L:sur:c}
The lifting maps ${\mu}\bb L_{k}$ induce surjective homomorphisms, which are often simply denoted by ${\mu}\bb L$, between the Grothendieck semigroups
\[
 \gsk \mRV^{\db}[k] \epi \gsk \mVF^{\diamond}[k].
\]
\end{cor}
\begin{proof}
By Lemma~\ref{L:measure:class:lift}, every $\mRV^{\db}[k]$-morphism can be lifted to a $\mVF^{\diamond}[k]$-morphism. Thus, ${\mu}\bb L_{k}$ induces a map on the isomorphism classes, which is easily seen to be a semigroup homomorphism. By Lemma~\ref{L:measure:surjective} (see the remark thereafter) or Remark~\ref{stand:vol}, this  homomorphism is surjective.
\end{proof}

\begin{lem}\label{simul:special:dim:1}
Let $f : (A, \omega) \fun (B, \sigma)$ be a $\mgVF^{\diamond}[1]$-morphism.  Then there are special covariant bijections $T_A : A \fun A^{\flat}$, $T_B : B \fun B^{\flat}$ such that $A^{\flat}$, $B^{\flat}$ are doubly bounded $\RV$-pullbacks and the functions $f^{\flat}$, $f^{\flat}_{\downarrow}$ are bijective in
the diagram
\[
\bfig
  \square(0,0)/->`->`->`->/<600,400>[A`A^{\flat}`B`B^{\flat};
  T_A`f``T_B]
 \square(600,0)/->`->`->`->/<600,400>[A^{\flat}`\rv(A^{\flat})`B^{\flat} `\rv(B^{\flat});  \rv`f^{\flat}`f^{\flat}_{\downarrow}`\rv]
 \efig
\]
where the second square commutes (but not necessarily the first one, although it fails at only finitely many points). Moreover, $\omega' = \omega \circ T_A^{-1}$ and $\sigma' = \sigma \circ T_B^{-1}$ are both $\rv$-contractible, and  $f^{\flat}_{\downarrow}$ is a $\mRV^{\db}[1]$-morphism $(\rv(A^{\flat}), \omega')_{1} \fun (\rv(B^{\flat}), \sigma')_1$ (recall Notation~\ref{0coor}).
\end{lem}
\begin{proof}
This is a variation of \cite[Lemma~5.2]{Yin:int:acvf} adapted to the current environment and can be established by modifying two special bijections $T'_A$, $T'_B$ as constructed in the proof thereof as follows.

By Lemma~\ref{L:measure:surjective}, we may assume that $A$, $B$ are doubly bounded $\RV$-pullbacks and $\omega$, $\sigma$ are already $\rv$-contractible. By Lemmas~\ref{dcl:to:ac} and \ref{RV:fiber:dim:same}, each focus map in $T'_A$ consists of an algebraic set of focus points, similarly for the focus maps in $T'_B$. We think of these focus points as contained in $A$, $B$. Each one of the focus maps may produce $0$ in the $\VF$-coordinate in the resulting set, but this is not allowed in a special covariant bijection. To remedy  this, we choose pairwise disjoint open discs of sufficiently large definable radii such that each of them contains exactly one focus point and the volume forms $\omega$, $\sigma$ are constant on them. By the construction in the proof of \cite[Lemma~5.2]{Yin:int:acvf}, $f$ has the disc-to-disc property outside these open discs. Inside these discs, by Lemma~\ref{lim:weak}, the radii may be chosen so that $f$ is weakly concentric at the focus points. In particular, since $f^\flat$ $\rv$-contracts to a bijection,  these open discs can be paired up via the induced bijection  in question. At this point it is not hard to see how to modify $T'_A$, $T'_B$ according to the second clause of (\ref{centri}). The first claim follows. For the second claim, since the volume forms $\omega'$, $\sigma'$ are $\rv$-contractible, we can use the disc-to-disc property (or weak concentricity for the exceptional discs) to gauge the difference of the radii of the $\RV$-discs in question against the constant valuation of the differentials, and this shows that $\omega'$, $\sigma'$  do match under $f^{\flat}_{\downarrow}$ in the required sense.
\end{proof}

% BEGIN REFEREE INSERTION R117
\RefereeRemark{117}{Gap}{Record the radius and volume compatibility in the unary matching}{%
\textit{Original source line 2420.}\quad The final volume comparison should be written explicitly. For each pair of matched RV-discs $\gp,\gq$, prove
\[
\omega'(\gp)-\sigma'(\gq)
 =\rad(\gq)-\rad(\gp)
 =\vv(f'(a))
\]
at a generic point $a$ of the corresponding original source disc. The first equality is exactly the required $\Gamma$-Jacobian identity. For the exceptional discs this requires choosing paired aperture radii compatibly with the weakly concentric image, rather than choosing the radii independently. Also distinguish the already constructed conjugation using $T'_A,T'_B$ from the desired modified $f^\flat$ in the sentence ``since $f^\flat$ RV-contracts to a bijection''; as written this uses the intended conclusion before defining the modification. Give the finite-point adjustment explicitly and verify that it does not change the constant volume forms.
}
% END REFEREE INSERTION R117

\section{Integrating proper invariant sets}
\label{section:int}

\subsection{Annular blowups}
Recall that we write $\RV$ to mean the $\RV$-sort without the  element $\infty$, and $\RV_\infty$ otherwise, although quite often the difference is immaterial and it does not matter which set $\RV$ stands for. In the definition below, this difference does matter, in particular, $\RV^{\ccirc}$ denotes the set $\rv(\MM \mi 0)$.

\begin{nota}\label{punc:RVd}
Recall from Notation~\ref{defn:disc} that the set $\rv(\MM_{\gamma} \mi 0)$ is denoted by $\RV^{\ccirc}_\gamma$. For $\gamma \in \Gamma^+$ and $t \in \gamma^\sharp$, write
\[
\RV^{\ccirc}(t)_\gamma = (\RV^{\ccirc} \mi \RV^{\ccirc}_\gamma) \uplus t;
\]
if $\gamma = \infty$ then $\RV^{\ccirc}(t)_\gamma = \RV^{\ccirc}_\infty$ and if  $\gamma = 0$ then $\RV^{\ccirc}(1)_\gamma = 1 \in \RV$.

% BEGIN REFEREE INSERTION R118
\RefereeRemark{118}{Typo; revised}{Zero aperture not fully defined}{%
\textit{Original source line 2433.}\quad The next definition permits any definable $t$ with $\vrv(t)=0$, whereas the zero-aperture convention here specifies only $t=1$. Either set $\RV^{\ccirc}(t)_0=t$ for every such $t$, or require $t=1$ at aperture zero. The singleton notation itself is legitimate under Notation~\ref{indexing}.
}
% END REFEREE INSERTION R118
\end{nota}

% be a definable function and $\tau_*: U \fun$ the function given by
%\[
%t \efun \vrv(\tau(t)) - \vrv(f_j(t)).
%\]

\begin{defn}\label{defn:blowup:coa}
Let $\bm U = (U, f, \omega)$ be an object of $\mRV[k]$, where $k > 0$, such that the function $\pr_{\tilde j} \rest f(U)$ is finite-to-one for some $j \in [k]$. Write $f = (f_1, \ldots, f_k)$. Let $t \in \RV_\infty$ be a definable element with $\tau = \vrv(t)$ nonnegative.  An \memph{elementary annular blowup} of $\bm U$ in the $j$th coordinate of \memph{aperture} $(\tau, t)$ is the triple $\bm U^{\flat} = (U^{\flat}, f^{\flat}, \omega^\flat)$, where $U^{\flat} = U \times \RV^{\ccirc}(t)_{\tau}$ and for every $(u, s) \in U^{\flat}$,
\[
f^{\flat}_{i}(u, s) = f_{i}(u) \text{ for } i \neq j, \quad f^{\flat}_{j}(u, s) = s f_{j}(u), \quad \omega^\flat(u, s) = \omega(u).
\]

% BEGIN REFEREE INSERTION R119
\RefereeRemark{119}{Error}{Infinite aperture leaves the stated category}{%
\textit{Original source line 2445.}\quad Restrict the annular construction used below to $t\in\RV$ and finite $\tau\geq0$. As written, $t=\infty$ is allowed, but $\RV^{\ccirc}(\infty)_\infty=\rv(\MM)$ contains $\infty$; consequently $f^\flat$ takes an $\infty$-value and is not a map into $\RV^k$, as required in Definition~\ref{defn:c:RV:cat}. Its nonzero part also has unbounded valuation, contradicting the next assertion about $\mRV^{\db}[k]$. The infinite-aperture comparison with ordinary blowups should be stated separately, with its different grading convention.
}
% END REFEREE INSERTION R119

Let $\bm V = (V, g, \sigma)$ be another object of $\mRV[k]$ and $C_i \sub V$ finitely many pairwise disjoint definable sets. Each triple
\[
\bm C_i = (C_i, g \rest C_i, \sigma \rest C_i) \in \mRV[k]
\]
is referred to as a \memph{subobject} of $\bm V$. Suppose that $F_i : \bm U_i \fun \bm C_i$ is a $\mRV[k]$-morphism and $\bm U^{\flat}_i$ is an elementary annular blowup of $\bm U_i$. Let $C = \bigcup_i C_i$, $\bm C = \bigcup_i \bm C_i$,  and $F = \biguplus_i F_i$. Then the object
$
 (\bm V \mi \bm C) \uplus \biguplus_i \bm U_i^{\flat}
$
is  an \memph{annular blowup of $\bm V$ via $F$}, denoted by $\bm V^{\flat}_F$. The subscript $F$ is usually dropped. The object $\bm C$ (or the set $C$) is referred to as the \memph{locus} of $\bm V^{\flat}_F$.

An \memph{annular blowup of length $n$} is a composition of $n$ annular blowups.
\end{defn}

\begin{rem}
An elementary annular blowup could be an elementary blowup (see \cite[Definition~6.1]{Yin:int:acvf}) if the aperture is allowed to be $(\infty, \infty)$. If there is an elementary annular blowup of $\bm U$ then, \textit{a posteriori}, $\dim_{\RV}(U) < k$. For any coordinate of $f(U)$, there is at most one elementary blowup of $\bm U$, whereas there could be many annular elementary blowups of $\bm U$. We should have included the coordinate that is blown up as a part of the data. However, in context, either this is
clear or it does not need to be spelled out, and we shall suppress mention of it for ease of notation.
\end{rem}

\begin{rem}
Clearly if $\bm U \in \mRV^{\db}[k]$  then $\bm U^{\flat} \in \mRV^{\db}[k]$ too. We shall only consider annular blowups  in $\mRV^{\db}[k]$ (then, for simplicity, they will just be referred to as blowups). This notion is dual to that of a  special covariant bijection. As a matter of fact, to make this duality precise is essentially what is left to do for the rest of our main construction.
\end{rem}



For the next few lemmas, let $\bm U = (U, f, \omega)$ and $\bm V = (V, g, \sigma)$ be objects of $\mRV^{\db}[k]$.

%and the aperture $\alpha$ of $\bm U^{\flat}$ is no greater than the aperture $\beta$ of  $\bm V^{\flat}$

\begin{lem}\label{elementary:blowups:preserves:iso:vol}
Suppose that $[\bm U] = [\bm V]$ and $\bm U^{\flat}$, $\bm V^{\flat}$ are elementary blowups of $\bm U$, $\bm V$ in the $k$th coordinate. Then there are  blowups $\bm U^{\flat\flat}$, $\bm V^{\flat\flat}$ of $\bm U^{\flat}$, $\bm V^{\flat}$ of length $1$ such that $[\bm U^{\flat\flat}] = [\bm V^{\flat\flat}]$.
\end{lem}
\begin{proof}
Let $(\rho, r)$, $(\tau, t)$ be the apertures of  $\bm U^{\flat}$, $\bm V^{\flat}$, respectively. If $\rho = \tau$ then clearly any morphism $\bm U \fun \bm V$ induces a morphism $\bm U^{\flat} \fun \bm V^{\flat}$. So we may assume $\rho' \coloneqq \tau - \rho > 0$. Choose a definable element $r' \in \rho'^\sharp$. Let $W = U \times r$ and $\bm W$ be the corresponding subobject of $\bm U^{\flat}$. Then it is easy to see that the blowup $\bm U^{\flat\flat}$  of $\bm U^{\flat}$ with locus $\bm W$ and aperture $(\rho', r')$ is isomorphic to $\bm V^{\flat}$.

% BEGIN REFEREE INSERTION R120
\RefereeRemark{120}{Clarification; revised}{Explicit common refinement}{%
\textit{Original source line 2479.}\quad For completeness, the stated equal lengths can be obtained by taking an empty-locus blowup on the $\bm V^\flat$ side; the definition already permits this identity operation. If $\rho=\tau$, take it on both sides. When $\rho<\tau$, take $r'=t/r$. Multiply the second-stage annular parameter by $r$, send its distinguished endpoint to $t$, and use the original isomorphism on $U$. The total valuation in \eqref{cond:mg2} is then preserved. This makes the short common-refinement construction explicit.
}
% END REFEREE INSERTION R120
\end{proof}

\begin{cor}\label{eblowup:same:loci:iso}
Let $F : \bm U \fun \bm V$ be a morphism and $\bm U^{\flat}$, $\bm V^{\flat}$ blowups of  $\bm U$, $\bm V$ of length $1$. Then there are blowups $\bm U^{\flat\flat}$, $\bm V^{\flat\flat}$ of $\bm U^{\flat}$, $\bm V^{\flat}$ of length $1$ such that $[\bm U^{\flat\flat}] = [\bm V^{\flat\flat}]$.
\end{cor}
\begin{proof}
Let  $C$, $D$ be the loci of $\bm U^{\flat}$, $\bm V^{\flat}$, respectively. By Lemma~\ref{elementary:blowups:preserves:iso:vol}, we can make the claim hold by restricting the loci of $\bm U^{\flat}$, $\bm V^{\flat}$ to $C \cap F^{-1}(D)$, $F(C) \cap D$. But in the meantime we can also blow up $\bm U^{\flat}$, $\bm V^{\flat}$ at the loci $F^{-1}(D) \mi C$, $F(C) \mi D$ using the apertures induced by those on $D \mi F(C)$, $C \mi F^{-1}(D)$, respectively. Then the resulting compounded blowups of $\bm U^{\flat}$, $\bm V^{\flat}$ of length $1$ are as desired.
\end{proof}


This corollary also holds in $\RV[k]$, which is essentially the only thing that the otherwise formal proof of \cite[Lemma~6.5]{Yin:special:trans} depends on. Thus, the same proof yields the following analogue in the present context:

% BEGIN REFEREE INSERTION R121
\RefereeRemark{121}{Typo}{Wrong cited paper}{%
\textit{Original source line 2490.}\quad Replace \texttt{Yin:special:trans} by \texttt{Yin:int:acvf} in this citation: the common-refinement result is Lemma~6.5 of \emph{Integration in algebraically closed valued fields} (2011). Section~6 of \emph{Special transformations} concerns elimination of field quantifiers. This issue from the supplied human report is still present.
}
% END REFEREE INSERTION R121

\begin{lem}\label{blowup:equi:class:coa}
If $[\bm U] = [\bm V]$ and $\bm U_1$, $\bm V_1$ are blowups of $\bm U$, $\bm V$ of lengths $m$, $n$, respectively, then there are blowups $\bm U_2$, $\bm V_2$ of $\bm U_1$, $\bm V_1$ of lengths $n$, $m$, respectively, such that $[\bm U_2] = [\bm V_2]$.
\end{lem}

\begin{cor}\label{blowup:equi:class}
Suppose that $[\bm U] = [\bm U']$ and $[\bm V] = [\bm V']$. If there are isomorphic blowups of $\bm U$, $\bm V$ then there are isomorphic blowups of $\bm U'$, $\bm V'$.
\end{cor}

\begin{defn}
Let $\mispdb[k]$ be the set of pairs $(\bm U, \bm V)$ of objects of $\mRV^{\db}[k]$ such that there exist isomorphic blowups $\bm U^{\flat}$, $\bm V^{\flat}$. Let $\mispdb[*] = \bigoplus_{k} \mispdb[k]$.
\end{defn}

We will just write $\mispdb$ for all these sets when there is no danger of confusion. By Corollary~\ref{blowup:equi:class}, they may be regarded as binary relations on isomorphism classes.


\begin{lem}\label{isp:congruence:vol}
$\mispdb[k]$ is a semigroup congruence relation and $\mispdb[*]$ is a semiring congruence relation.
\end{lem}
\begin{proof}
This is an easy consequence of Lemma~\ref{blowup:equi:class:coa} (the proof of \cite[Lemma~6.8]{Yin:special:trans}, which uses \cite[Lemma~6.5]{Yin:special:trans}, contains more details).

% BEGIN REFEREE INSERTION R122
\RefereeRemark{122}{Typo}{Same bibliographic correction}{%
\textit{Original source line 2511.}\quad Both references in this sentence should again use \texttt{Yin:int:acvf}; its Lemmas~6.8 and~6.5 are the congruence and common-refinement results being used.
}
% END REFEREE INSERTION R122
\end{proof}



Let  $T$ be a special covariant bijection on $\bb L (U,f)$. Denote the object $(\mathbb{L} (U,f)^\flat_{\RV})_{\leq k} \in \RV^{\db}[k]$ by $(U_{T})_{\leq k}$ and  the object $(U_{T}, \omega_T)_{\leq k} \in \mRV^{\db}[k]$ by $\bm U_{T}$, where $\omega_T : U_{T} \fun \RV$ is the volume form induced by $\omega$.

% BEGIN REFEREE INSERTION R123
\RefereeRemark{123}{Typo}{Wrong codomain for the volume form}{%
\textit{Original source line 2516.}\quad Replace $\omega_T:U_T\longrightarrow\RV$ by $\omega_T:U_T\longrightarrow\Gamma$. Throughout this paper the volume forms are $\Gamma$-valued; see Definition~\ref{defn:RV:mu}.
}
% END REFEREE INSERTION R123


\begin{lem}\label{special:to:blowup:coa}
The object $\bm U_T$ is isomorphic to a blowup of $\bm U$ of the same length as $T$.
\end{lem}
\begin{proof}
By induction on the length $\lh (T)$ of $T$ and Lemma~\ref{blowup:equi:class:coa}, this is immediately reduced to the case $\lh (T) = 1$. Without loss of generality, we may assume that the locus of $T$ is indeed $\mathbb{L} (U,f)$. Then it is clear how to construct an (elementary) blowup of $\bm U$ as desired, using the aperture provided by $T$.

% BEGIN REFEREE INSERTION R171
\RefereeRemark{171}{Clarification; new}{Why the proposed elementary blowup is allowed}{%
\textit{Original source line 2523.}\quad The required finite-to-one projection condition has a short verification worth inserting. Suppose the focus is in coordinate $j$, fix $r\in f(U)_{\tilde j}$ and a representative $a\in r^\sharp$, and vary $u$ with $f_{\tilde j}(u)=r$. The focus gives a definable function $u\mapsto\lambda(a,u)$ from an $\RV$-set into $\VF$, so its image is finite. Since $\rv(\lambda(a,u))=f_j(u)$, there are only finitely many possible $f_j(u)$ over $r$. Thus $\pr_{\tilde j}\rest f(U)$ is finite-to-one, exactly as required by Definition~\ref{defn:blowup:coa}. The annular parameter is then the translated leading term divided by $f_j(u)$, with the distinguished endpoint retaining its extra tag.
}
% END REFEREE INSERTION R171
\end{proof}



\begin{lem}\label{blowup:same:RV:coa}
Let $\bm U^{\flat}$ be a blowup of $\bm U$. Then $\mL \bm U^{\flat}$ is isomorphic to $\mL \bm U$ in $\mVF^{\diamond}[k]$.
\end{lem}
\begin{proof}
By induction on the length $l$ of $\bm U^{\flat}$, this is immediately reduced to the case $l=1$. We may assume that $\bm U^{\flat}$ is an elementary blowup in the first coordinate. Let $(\tau, t)$ be the aperture of $\bm U^{\flat}$. Fix a definable point $c \in t^\sharp$. For $u \in U$ and $a \in  f(u)_{\tilde 1}^\sharp$, every $\RV$-disc in $f(u)_{1}^\sharp$ contains an $a$-algebraic point $b_a$; moreover, by Lemma~\ref{RVlift}, these points may be chosen uniformly via a partially differentiable finitary function. Thus, there is a continuous (fiberwise) additive translation, with respect to the points $b_a$ and $b_a - c b_a$, between the two sets in question, which is a $\mVF^{\diamond}[k]$-morphism by Lemmas~\ref{conti:covar} and \ref{conv:unif}.

% BEGIN REFEREE INSERTION R124
\RefereeRemark{124}{Clarification; revised}{Specify the lift and its tags}{%
\textit{Original source line 2532.}\quad Write the chosen center as $b=b(a,u)\in f_1(u)^\sharp$; the phrase ``every $\RV$-disc in $f(u)_1^\sharp$'' should simply refer to this disc. The lift can be written explicitly: send $x$ to $x-b$ when $\vv(x-b)\leq\vv(b)+\tau$, and to $x-b+cb$ otherwise. The latter branch fills the extra copy of $(tf_1(u))^\sharp$; retain the two tags as already stipulated in Definition~\ref{defn:special:bijection}. On each branch this is a triangular translation of Jacobian $1$. To justify the asserted continuity, use the strengthened version of Lemma~\ref{RVlift} with continuous as well as partially differentiable center branches (see the remark at that lemma). This supplies the input needed for Lemmas~\ref{conti:covar} and~\ref{conv:unif}.
}
% END REFEREE INSERTION R124
\end{proof}



\begin{lem}\label{kernel:dim:1:coa}
Let $[\bm A] = [\bm B]$ in  $\gsk \mVF^{\diamond}[1]$ and $\bm U, \bm V \in \mRV^{\db}[1]$ be standard contractions of $\bm A$, $\bm B$, respectively. Then $([\bm U], [\bm V]) \in \mispdb$.
\end{lem}
\begin{proof}
This is immediate by Lemmas~\ref{simul:special:dim:1} (applied to $\mL \bm U$, $\mL \bm V$) and \ref{special:to:blowup:coa}.
\end{proof}

\begin{lem}\label{isp:VF:fiberwise:contract}
Let $A'$, $A''$ be definable sets with $A'_{\VF} = A''_{\VF} = A \sub \VF^n$ and $\omega'$, $\omega''$ definable functions from $A'$, $A''$ into $\Gamma$, respectively. Write $\bm A' = (A', \omega')$ and $\bm A'' = (A'', \omega'')$. Suppose that $\bm A'$, $\bm A''$ are proper invariant and there is a $k \in \N$ such that for  every $a \in A$,
\begin{equation}\label{condi:fiber}
([(A'_a, \omega')]_{\leq k}, [(A''_a, \omega'')]_{\leq k}) \in \mispdb.
\end{equation}
Let $\hat T_{\sigma}$, $\hat R_{\sigma}$ be standard contractions of $\bm A'$, $\bm A''$, respectively. Then
\[
([\hat T_{\sigma}(\bm A')]_{\leq n+k}, [\hat R_{\sigma}(\bm A'')]_{\leq n+k}) \in \mispdb.
\]
\end{lem}

Note that condition (\ref{condi:fiber}) makes sense only  over the substructure $\mathbb{S} \la a \ra$.

\begin{proof}
By induction on $n$, this is immediately reduced to the case $n=1$. So assume $A \sub \VF$. Using Proposition~\ref{simplex:with:hole:rvproduct} and FMT, we can construct a special covariant bijection $F : A \fun A^\sharp$, as in the proof of \cite[Lemma~6.14]{Yin:int:acvf}, such that for all $\RV$-polydisc $\gp \sub A^{\sharp}$ and all $a_1, a_2 \in F^{-1}(\gp)$,
\[
\omega' \rest A'_{a_1} = \omega' \rest A'_{a_2} \dand \omega'' \rest A''_{a_1} = \omega'' \rest A''_{a_2}.
\]
Therefore, using Lemma~\ref{kernel:dim:1:coa} in place of \cite[Corollary~6.11]{Yin:int:acvf}, that proof goes through here with virtually no changes, and the additional computations involving $\jcb_{\Gamma}$ are all straightforward. (Also consult the proof of \cite[Lemma~5.36]{Yin:tcon:I}, which is in a different environment but is formally the same and is  written down with more details.)

% BEGIN REFEREE INSERTION R125
\RefereeRemark{125}{Gap; revised}{Verify the bounded version of witness constancy}{%
\textit{Original source line 2562.}\quad The displayed equalities, read as equalities of functions, already imply equality of the underlying $\RV$-fibers. The cited proof of \cite[Lemma~6.14]{Yin:int:acvf} additionally makes the blowup witnesses constant on each contracted disc. Please explain how its third condition is maintained here: after FMT, the loci, finite apertures, and comparison isomorphisms must be given by data independent of $a\in F^{-1}(\gp)$. This is what lets the fiberwise congruences descend to the contracted $\RV$-parameters. The old proof supplies the strategy, but the proper-covariant, finite-aperture implementation is the point requiring verification.
}
% END REFEREE INSERTION R125
\end{proof}


\begin{cor}\label{contraction:same:perm:isp}
Let $\bm A', \bm A'' \in \mVF^{\diamond}[n]$ and suppose that there is a morphism $F : \bm A' \fun \bm A''$ that is relatively unary in the $i$th $\VF$-coordinate. Then for any permutation $\sigma$ of $[n]$ with $\sigma(1) = i$ and any standard contractions $\hat T_{\sigma}$, $\hat R_{\sigma}$ of $\bm A'$, $\bm A''$,
\[
([\hat T_{\sigma}(\bm A')]_{\leq n}, [\hat R_{\sigma}(\bm A'')]_{\leq n}) \in \mispdb.
\]
\end{cor}
\begin{proof}
Note that, since $F$ is differentiable outside a definable subset of $\VF$-dimension less than $n$, it may not induce a morphism $\bm A'_a \fun \bm A''_a$ for every $a \in \pr_{\tilde i}(A)$. Nevertheless, for all sufficiently large $\delta \in \Gamma^+$, both $A'$ and $\hat T_{\sigma(1)}(A')$ are $\delta$-invariant. This means that for all $a \in \pr_{\tilde i}(A)$ there is a $b \in \pr_{\tilde i}(A)$ such that they are contained in the same open polydisc of radius $\delta$, and $F$ indeed induces a morphism $\bm A'_b \fun \bm A''_b$, and $\hat T_{\sigma(1)}(A')_a = \hat T_{\sigma(1)}(A')_b$. Thus, the claim follows immediately from Lemmas~\ref{kernel:dim:1:coa} and \ref{isp:VF:fiberwise:contract}.

% BEGIN REFEREE INSERTION R126
\RefereeRemark{126}{Gap; revised}{Use a common radius on both sides}{%
\textit{Original source line 2573.}\quad Here $A$ has not been defined; use the common base $\pr_{\tilde i}(A')=\pr_{\tilde i}(A'')$. Choose $\delta$ for both first-stage contracted pairs, including their induced volume forms, and for $\bm A'$ and $\bm A''$. Let the bad set include failures of differentiability and of the measure identity. The parameters with a positive-dimensional bad unary fiber have dimension at most $n-2$, so every open base polydisc contains a good $b$. The two contracted fibers at $a$ and $b$ then agree on both sides, and the unary kernel lemma applies at $b$. These additions supply the claimed reduction; invariance of the first side alone is insufficient.
}
% END REFEREE INSERTION R126
\end{proof}

\subsection{$2$-cells}

\begin{defn}[$\vv$-affine and $\rv$-affine]\label{rvaffine}
Let $\ga$ be an open disc and $f : \ga \fun \VF$ an injection.
We say that $f$ is \memph{$\vv$-affine} if there is a (necessarily unique) $\gamma \in \Gamma$, called the \memph{shift} of $f$, such that, for all $a, a' \in \ga$,
\[
\vv(f(a) - f(a')) = \gamma + \vv(a - a').
\]
We say that $f$ is \memph{$\rv$-affine} if there is a (necessarily unique) $t \in \RV$, called the \memph{slope} of $f$, such that, for all $a, a' \in \ga$,
\[
\rv(f(a) - f(a')) = t \rv(a - a').
\]
\end{defn}

\begin{defn}\label{defn:balance}
Let $A \sub \VF^2$ be a definable set such that $\ga_1 \coloneqq \pr_1(A)$ and $\ga_2 \coloneqq \pr_2(A)$ are both open discs. Let $f : \ga_1 \fun \ga_2$ be a definable bijection that has the disc-to-disc property.  We say that $f$ is \memph{balanced in $A$} if $f$ is actually $\rv$-affine and there are $t_1, t_2 \in \RV_\infty$, called the \memph{paradigms} of $f$, such that, for every $a \in \ga_1$,
\[
A_a = t_2^\sharp + f(a) \dand f^{-1}(A_a) = a - t_1^\sharp.
\]
\end{defn}

\begin{rem}\label{2cell:linear}
Suppose that $f$ is balanced in $A$ with paradigms $t_1$, $t_2$.
If one of the paradigms is $\infty$ then the other one must be $\infty$. In this case $A$ is just the (graph of the) bijection $f$ itself.

Assume that $t_1$, $t_2$ are not $\infty$. Let $\gB_1$, $\gB_2$ be the sets of closed subdiscs of $\ga_1$, $\ga_2$ of radii $\vrv(t_1)$, $\vrv(t_2)$, respectively. Let $a_1 \in \gb_1 \in \gB_1$ and $\go_1$ be the maximal open subdisc of $\gb_1$ containing $a_1$. Let $\gb_2 \in \gB_2$ be the smallest closed disc containing the open disc $\go_2 \coloneqq A_{a_1}$. Then, for all $a_2 \in \go_2$, we have
\[
\go_2 = t_2^\sharp + f(\go_1) = A_{a_1} \dand A_{a_2} = f^{-1}(\go_2) + t_1^\sharp = \go_1.
\]
This internal symmetry of $A$ is illustrated by the following diagram:
\[
\bfig
  \dtriangle(0,0)|amb|/.``<-/<600,250>[\go_1`f^{-1}(\go_2)`\go_2; \pm t_1^\sharp`\times`f^{-1}]
  \ptriangle(600,0)|amb|/->``./<600,250>[\go_1`f(\go_1)`\go_2; f`` \pm t_2^\sharp]
 \efig
\]
Since $f$ is $\rv$-affine, we see that its slope must be $-t_2/t_1$.

Let $\tor(\code {\gb_1})$, $\tor(\code {\gb_1})$ denote the sets of the maximal open subdiscs of $\gb_1$, $\gb_1$. These may be viewed as $\K$-torsors and as such  are equipped with much of the structure of $\K$. Then the set $A \cap (\gb_1 \times \gb_2)$ may be thought of as the ``line'' in $\tor(\code {\gb_1}) \times \tor(\code {\gb_2})$ given by the equation

% BEGIN REFEREE INSERTION R127
\RefereeRemark{127}{Typo}{Repeated ball index}{%
\textit{Original source line 2614.}\quad In the opening sentence, the second $\tor(\code{\gb_1})$ should be $\tor(\code{\gb_2})$, and the second $\gb_1$ should likewise be $\gb_2$. The same typographical correction was requested in the supplied human report.
}
% END REFEREE INSERTION R127
\[
x_2 = - \tfrac{t_2}{t_1}(x_1 - \code{\go_1}) + (\code{\go_2} - t_2).
\]

% BEGIN REFEREE INSERTION R128
\RefereeRemark{128}{Error}{Incorrect intercept of the torsor line}{%
\textit{Original source line 2617.}\quad The line describing $A$ must pass through $(\code{\go_1},\code{\go_2})$. The displayed formula instead gives $\code{\go_2}-t_2$ at $x_1=\code{\go_1}$. Remove the term $-t_2$: the required equation is $x_2=-(t_2/t_1)(x_1-\code{\go_1})+\code{\go_2}$. The printed equation describes the induced graph of $f$, before translation by $t_2$. For example, when $f=\id$ and $A=\{(x,y):\rv(y-x)=t_2\}$ on a suitable pair of discs, the required residue line is $y=x+t_2$, not the diagonal.
}
% END REFEREE INSERTION R128
Thus, by Lemma~\ref{simul:special:dim:1}, the obvious bijection between $A_1 \times t_2^\sharp$ and $t_1^\sharp \times A_2$ is the lift of an $\RV^{\db}[2]$-morphism modulo special covariant bijections; see Lemma~\ref{2:unit:contracted} below for details.
\end{rem}

\begin{defn}\label{def:units}
We say that a set $A$ is a \memph{$1$-cell} if it is either an open disc contained in a single $\RV$-disc or a point in $\VF$. We say
that $A$ is a \memph{$2$-cell} if
\begin{itemize}
 \item $A$ is a subset of $\VF^2$ contained in a single $\RV$-polydisc and $A_1$ is a $1$-cell,
 \item there is a function $\epsilon : A_1  \fun \VF$ and a $t \in \RV$ such that, for every $a \in A_1$, $A_a = t^\sharp + \epsilon(a)$,
 \item one of the following three possibilities occurs:
  \begin{itemize}
   \item $\epsilon$ is constant,
   \item $\epsilon$ is injective, has the disc-to-disc property, and $\rad(\epsilon(A_1)) \geq \vrv(t)$,
   \item $\epsilon$ is balanced in $A$.
  \end{itemize}
\end{itemize}
The function $\epsilon$ is called the \memph{positioning function} of $A$ and the element $t$ the \memph{paradigm} of $A$.

More generally, a set $A$ with exactly one $\VF$-coordinate is a \memph{$1$-cell} if, for each $t \in A_{\RV}$, $A_t$ is a $1$-cell in the above sense; the parameterized version of the notion of a $2$-cell is formulated in the same way.
\end{defn}

\begin{lem}\label{decom:into:2:units}
Let  $A \sub \VF^2$ be a proper invariant set. Then there is a standard contraction $\hat T_\sigma$ of $A$ such that $A_s$ is a $2$-cell for every $s \in \hat T_\sigma(A)$.
\end{lem}
\begin{proof}
This is a variation of  \cite[Lemma~4.8]{Yin:int:acvf}. The proof of the latter proceeds by constructing a positioning function $\epsilon_{(t,s)}$ in each $\VF$-fiber, which heavily relies on the use of standard contractions in the preceding auxiliary results, namely \cite[Lemmas~4.2, 4.4]{Yin:int:acvf}.  It still goes through here, provided that FMT is applied at suitable places to cut out finitely many exceptional open discs in the first coordinate at which the desired properties of $\epsilon_{(t,s)}$ do not hold, as we have done above, say, in the proof of Proposition~\ref{simplex:with:hole:rvproduct}.

% BEGIN REFEREE INSERTION R130
\RefereeRemark{130}{Gap}{The cited cell lemma has a weaker affine condition}{%
\textit{Original source line 2643.}\quad The cited proof produces a $\vv$-affine positioning function, whereas Definition~\ref{defn:balance} now requires an $\rv$-affine one. The implication is not automatic: for definable $\delta>0$, the bijection of $\MM$ given by $x\mapsto2x$ when $\vv(x)\geq\delta$ and $x\mapsto x$ otherwise preserves all valuations of differences but has no single $\RV$-slope. Please supply the refinement or replacement of the positioning functions that gives the stronger condition, and verify that the fibers, paradigms, and proper covariance are preserved. Identifying the exceptional discs and explaining the FMT modification is necessary here; the older lemma plus the phrase \emph{at suitable places} does not supply this new step. An economical alternative, if only the $\Gamma$-Jacobian computation is used, is to retain the older $\vv$-affine definition and replace the assertions about a single $\RV$-slope and a torsor line by the corresponding valuation-shift statement.
}
% END REFEREE INSERTION R130
\end{proof}

We also remark that Lemma~\ref{decom:into:2:units} holds fiberwise for proper invariant sets $A \sub \VF^n$ with $n \geq 2$, that is, there is a standard contraction $\hat T_\sigma$ of $A$ such that for every $(a, s) \in \hat T_{\sigma(2)}(A)$, $\hat T_{\sigma(2)}^{-1}(a, s)$ is of the form $a \times C$, where $C$ is a $2$-cell. This follows from Lemma~\ref{decom:into:2:units} and the argument in the proof of  Lemma~\ref{exit:stand} (Lemma~\ref{decom:into:2:units} serves as the base case and hence, in the inductive step, we can assume that $\hat T_{\sigma(2)}$ is already as desired fiberwise).

% BEGIN REFEREE INSERTION R131
\RefereeRemark{131}{Clarification; revised}{Cumulative-contraction notation}{%
\textit{Original source line 2646.}\quad For clarity, specify that $\hat T_{\sigma(2)}$ here denotes the cumulative contraction of the first two coordinates. With the stage-map notation of the definition, write $P=\hat T_{\sigma(2)}\circ\hat T_{\sigma(1)}$ and use $P(A)$ and $P^{-1}(a,s)$. This is only a notation clarification.
}
% END REFEREE INSERTION R131


For the next two lemmas, let $12$, $21$ denote the permutations of $[2]$ and $\bm A = (A, \omega) \in \mVF^\diamond[2]$.


\begin{lem}\label{2:unit:contracted}
Suppose that $A \sub \VF^2$ is a $2$-cell and $\omega$ is constant. Then there are standard contractions $\hat T_{12}$, $\hat R_{21}$ of $\bm A$ such that $[\hat T_{12}(\bm A)]_{\leq 2} = [\hat R_{21}(\bm A)]_{\leq 2}$.
\end{lem}
\begin{proof}
All we need to do is to check that the maps constructed in the proof of  \cite[Lemma~5.7]{Yin:int:acvf} are indeed $\mRV^{\db}[2]$-morphisms. By inspection of that proof, we see that there are two cases: $A$ is a product of two open discs or the positioning function $\epsilon$ in question is balanced in $A$ with nonzero paradigms $t_1$, $t_2$. The first case is obvious since we can simply use the identity map. In the second case,  a morphism between the standard contractions can be easily constructed using  $\epsilon$  (we could also cite Lemma~\ref{simul:special:dim:1}, but the situation here is much simpler), and  the requirement on $\jcb_{\Gamma}$ is satisfied since the slope of the $\rv$-affine function $\epsilon$ is $-t_2/t_1$ (see the last paragraph of Remark~\ref{2cell:linear} for further explanation).

% BEGIN REFEREE INSERTION R132
\RefereeRemark{132}{Clarification; revised}{Explicit balanced-cell comparison}{%
\textit{Original source line 2656.}\quad Here is an explicit completion of the Jacobian check. Put $\gamma=\vrv(t_2)-\vrv(t_1)$. For order $12$, first use $(x,y)\mapsto(x-\epsilon^{-1}(y),y)$, whose image is $t_1^\sharp\times A_2$; for order $21$, use $(x,y)\mapsto(x,y-\epsilon(x))$, whose image is $A_1\times t_2^\sharp$. The translated coordinates have the fixed leading terms $t_1,t_2$, so finite apertures chosen strictly beyond their radii make the exceptional CT clause unnecessary on these sets. Apply Lemma~\ref{simul:special:dim:1} to $\epsilon$ with constant volume forms $0$ on $A_1$ and $-\gamma$ on $A_2$. It gives compatible unary contractions and an $\RV$-map $\psi$ with $\jcb_\Gamma\psi=\gamma$. Keeping all auxiliary tags, the map $(u,t_2)\mapsto(t_1,\psi(u))$ between the product contractions has $\Gamma$-Jacobian $\gamma+\vrv(t_1)-\vrv(t_2)=0$. This proves the desired measure comparison, conditional on that unary lemma. The same calculation uses only the valuation-affine shift, so it also supports the weaker-definition alternative mentioned above.
}
% END REFEREE INSERTION R132
\end{proof}

\begin{lem}\label{subset:partitioned:2:unit:contracted}
There are a morphism $\bm A \fun \bm A^*$, relatively unary in both coordinates, and two standard contractions $\hat T_{12}$, $\hat R_{21}$ of $\bm A^*$ such that $[\hat T_{12}(\bm A^*)]_{\leq 2} = [\hat R_{21}(\bm A^*)]_{\leq 2}$.
\end{lem}
\begin{proof}
This follows from  Lemmas~\ref{decom:into:2:units}, \ref{2:unit:contracted}, and compactness (see the proof of \cite[Corollary~5.8]{Yin:int:acvf} for a bit more details).

% BEGIN REFEREE INSERTION R133
\RefereeRemark{133}{Gap; revised}{Apply the available volume-form preparation}{%
\textit{Original source line 2663.}\quad Please spell out the preparation of $\omega$ and the definition of $\bm A^*$. Lemma~\ref{2:unit:contracted} assumes a constant form on each cell. The preparation tools in Remark~\ref{cons:extra}, Lemma~\ref{L:measure:surjective}, and Remark~\ref{stand:vol} are relevant: arrange proper covariant $\RV$-data through which $\omega$ factors, and carry these data through the cell construction. Then $A^*$ is the graph of the resulting cell-parameter map, with the transported volume form. Verify that these parameters have proper invariant fibers and doubly bounded range, so the parameter-adjoining map and its inverse are proper covariant. This is the additional check needed to turn the cited ordinary cell partition into a morphism in the present category.
}
% END REFEREE INSERTION R133
\end{proof}

\begin{lem}\label{contraction:perm:pair:isp}
Let $\bm A = (A, \omega) \in \mVF^\diamond[n]$. Let $i, j \in [n]$ with $i \neq j$ and $\sigma_1$, $\sigma_2$ be  permutations of $[n]$ such that
\[
\sigma_1(1) = \sigma_2(2) = i, \quad \sigma_1(2) = \sigma_2(1) = j, \quad \sigma_1
\rest \set{3, \ldots, n} = \sigma_2 \rest \set{3, \ldots, n}.
\]
Then, for any standard contractions $\hat T_{\sigma_1}$, $\hat T_{\sigma_2}$ of $\bm A$,
\[
([\hat T_{\sigma_1}(\bm A)]_{\leq n}, [\hat T_{\sigma_2}(\bm A)]_{\leq n}) \in \mispdb.
\]
\end{lem}
\begin{proof}
Let $ij$, $ji$ denote the permutations of $\{i, j\}$ and $E = [n] \mi \{i, j\}$. By compactness and
Lemma~\ref{isp:VF:fiberwise:contract}, it is enough to show
that for  all $a \in A_E$ and all standard
contractions $\hat T_{ij}$, $\hat T_{ji}$ of $\bm A_a = (A_a, \omega \rest A_a)$,
\[
([\hat T_{ij}(\bm A_a)]_{\leq 2}, [\hat T_{ji}(\bm A_a)]_{\leq 2})
\in \mispdb.
\]
By Corollary~\ref{contraction:same:perm:isp} and Lemma~\ref{isp:VF:fiberwise:contract}, it is enough to find a morphism $\bm A_a \fun \bm B$, relatively unary in both coordinates, and  standard contractions $\hat R_{ij}$, $\hat R_{ji}$ of $\bm B$ such that $[\hat R_{ij}(\bm B)]_{\leq 2} = [\hat R_{ji}(\bm B)]_{\leq 2}$. This is just Lemma~\ref{subset:partitioned:2:unit:contracted}.
\end{proof}

\subsection{The kernel of $\mL$ and the bounded integral}

The following proposition is the culmination of the preceding technicalities, it identifies the
congruence relation  $\mispdb$ with that induced by
$\mL$.


\begin{prop}\label{kernel:L}
For $\bm U, \bm V \in \mRV^{\db}[k]$,
\[
[\mu\bb L \bm U] = [\mu\bb L \bm V] \text{ in }\gsk \mVF^\diamond[k] \quad \text{if and only if} \quad ([\bm U], [\bm V]) \in \mispdb.
\]
\end{prop}
\begin{proof}
The ``if'' direction simply follows from Lemma~\ref{blowup:same:RV:coa} and
Corollary~\ref{L:sur:c}.

For the ``only if'' direction, we proceed by induction on $k$. The base case $k = 1$ is of course Lemma~\ref{kernel:dim:1:coa}. For the inductive step, let
\[
\mu\bb L \bm U = \bm B_1 \to^{G_1} \bm B_2 \cdots \bm B_l \to^{G_l} \bm B_{l+1} = \mu\bb L \bm V
\]
be a sequence of relatively unary $\mVF^\diamond[k]$-morphisms, which exist by definition. For each $j \leq l - 2$, we can choose five standard contractions
\[
[\bm U_j]_{\leq k}, \quad [\bm U_{j+1}]_{\leq k}, \quad [\bm U'_{j+1}]_{\leq k}, \quad [\bm U''_{j+1}]_{\leq k}, \quad [\bm U_{j+2}]_{\leq k}
\]
of $\bm B_j$, $\bm B_{j+1}$, $\bm B_{j+1}$, $\bm B_{j+1}$, $\bm B_{j+2}$ with the permutations $\sigma_{j}$, $\sigma_{j+1}$, $\sigma'_{j+1}$, $\sigma''_{j+1}$, $\sigma_{j+2}$ of $[k]$, respectively, such that
\begin{itemize}
  \item $\sigma_{j+1}(1)$ and $\sigma_{j+1}(2)$ are the $\VF$-coordinates targeted by $G_{j}$ and $G_{j+1}$, respectively,
  \item $\sigma''_{j+1}(1)$ and $\sigma''_{j+1}(2)$ are the $\VF$-coordinates targeted by $G_{j+1}$ and $G_{j+2}$, respectively,
  \item $\sigma_{j} = \sigma_{j+1}$, $\sigma''_{j+1} =  \sigma_{j+2}$,  and $\sigma'_{j+1}(1) = \sigma''_{j+1}(1)$,
  \item the relation between $\sigma_{j+1}$ and $\sigma'_{j+1}$ is as described in Lemma~\ref{contraction:perm:pair:isp}.
\end{itemize}

% BEGIN REFEREE INSERTION R134
\RefereeRemark{134}{Error}{Incomplete chain and impossible permutation choices}{%
\textit{Original source line 2720.}\quad Two bookkeeping repairs are required. First, consecutive relatively unary maps can target the same coordinate; then the first or second bullet would assign the same value to positions $1$ and $2$ of a permutation. Compose such consecutive maps before choosing the permutations. Second, $j\leq l-2$ only connects contractions through $\bm B_l$, not through the last arrow to $\bm B_{l+1}$, and for $l=1,2$ there are no comparisons at all. Handle the short chains and the last arrow explicitly using Corollary~\ref{contraction:same:perm:isp}, and compare the endpoint contractions with the tautological contractions of $\mL\bm U$ and $\mL\bm V$. Alternatively, first prove independence of the complete contraction order and then apply that corollary successively to all $G_j$. Also record the immediate degree-zero case, since the proposition and the next theorem include $k=0$.
}
% END REFEREE INSERTION R134
By Corollary~\ref{contraction:same:perm:isp} and Lemma~\ref{contraction:perm:pair:isp}, all the adjacent pairs of these standard contractions are indeed $\mispdb$-congruent, except $([\bm U'_{j+1}]_{\leq k}, [\bm U''_{j+1}]_{\leq k})$.  Since we can choose $[\bm U'_{j+1}]_{\leq k}$, $[\bm U''_{j+1}]_{\leq k}$ so that they start with the same contraction in the first targeted $\VF$-coordinate of $\bm B_{j+1}$, the resulting sets from this step are the same. Thus, applying  the inductive hypothesis in each fiber over the just contracted coordinate, we see that this last pair is also $\mispdb$-congruent. This completes the ``only if'' direction.

% BEGIN REFEREE INSERTION R135
\RefereeRemark{135}{Gap; revised}{Induction after an RV contraction requires descent}{%
\textit{Original source line 2721.}\quad Make the parameter base and the final gluing explicit. The new parameter after the first contraction lies in $\RV$, and adjoining it need not leave $\mathbb S$ $\VF$-generated; the intermediate contracted object can also have positive head start. Work instead with the lifted intermediate object, fix a $\VF$-representative $a$ of its first leading term, and apply induction to the remaining fibers over $\mathbb S\langle a\rangle$. These fibers still have finite $\RV$-fibers and degree $k-1$. To descend, form the two families $\bigcup_a a\times U_{\rv(a)}$ and $\bigcup_a a\times V_{\rv(a)}$ of the final $\RV$-outputs (with the first leading coordinate omitted from each output). Lemma~\ref{isp:VF:fiberwise:contract}, with one $\VF$-coordinate and head start $k-1$, then glues their fiberwise congruences; tautologically contracting $a$ recovers the two full outputs. This repair depends on the uniform-witness step discussed at that lemma.
}
% END REFEREE INSERTION R135
\end{proof}

\begin{rem}\label{thekernel}
Proposition~\ref{kernel:L} shows that  the kernel  of $\mL$ in $\gsk \mRV^{\db}[*]$ is generated by the pairs $([1], [\RV^{\ccirc}(t)_\gamma])$, where $\gamma \in \Gamma$ and $t \in \gamma^\sharp$ are definable (Notation~\ref{punc:RVd}), and hence  the corresponding ideal of the graded ring $\ggk \mRV^{\db}[*]$  is $(\bm P_\Gamma)$ (Notation~\ref{pgammabis}).

% BEGIN REFEREE INSERTION R136
\RefereeRemark{136}{Clarification}{Show the generator calculation and restrict the aperture}{%
\textit{Original source line 2725.}\quad Here $\gamma$ must be nonnegative, as in Notation~\ref{punc:RVd}. The passage from arbitrary elementary blowups to these degree-one generators has a short useful verification. For a blowup in coordinate $j$, put $\bm W=(U,f_{\tilde j},\omega+\vrv(f_j))\in\mRV^{\db}[k-1]$. The measure condition gives $\bm U\simeq\bm W\times(1,0)$ and $\bm U^\flat\simeq\bm W\times(\RV^{\ccirc}(t)_\gamma,0)$. Hence $[\bm U^\flat]-[\bm U]=[\bm W]\bm P_\gamma$. This proves the claimed ideal description by additivity and iteration, once Proposition~\ref{kernel:L} is established.
}
% END REFEREE INSERTION R136
\end{rem}

\begin{thm}\label{main:prop}
For each $k \geq 0$ there is a canonical isomorphism of Grothendieck semigroups
\[
  \int^{\diamond}_+ : \gsk \mVF^\diamond[k] \fun \gsk  \mRV^{\db}[k] /  \mispdb
\]
such that $\int_{+}^{\diamond} [\bm A] = [\bm U]/  \mispdb$ if and only if $[\bm A] = [{\mL}\bm U]$. Taking the direct sum of the groupifications
yields a canonical isomorphism of graded rings
\[
  \int^{\diamond} : \ggk \mVF^\diamond[*] \fun \ggk  \mRV^{\db}[*] /  (\bm P_\Gamma).
\]
\end{thm}
\begin{proof}
This is immediate by Corollary~\ref{L:sur:c} and Proposition~\ref{kernel:L}.
\end{proof}


\begin{cor}\label{db:to:novol}
The diagram (\ref{diag-diamond-interpol}) commutes.
\end{cor}
\begin{proof}
Note that, in (\ref{diag-diamond-interpol}),  the first horizontal arrow  is induced by the obvious forgetful functor (it exists
because the morphisms in $\mgVF^\diamond[*]$ are bijections instead of essential bijections, see Remark~\ref{dia:cat:well}), the second horizontal arrow is induced by the subcategory relation, and the reason that the bottom row exists is given in Notation~\ref{pgammabis}. That the upper half of the diagram commutes is then simply a consequence of the fact that image of $[\bm A]$ under any of the three integrals  can be represented by a standard contraction of $\bm A$. The rest of it is just a combination of (\ref{edb:eb:com}) and Proposition~\ref{prop:eu:retr:k}.

% BEGIN REFEREE INSERTION R137
\RefereeRemark{137}{Typo}{Proof refers to a deleted half of the diagram}{%
\textit{Original source line 2749.}\quad The current Figure~\ref{diag-diamond-interpol} has only the two integral squares; the additional row involving retractions is commented out in the source. Replace \emph{the upper half of the diagram} by \emph{both squares} and delete the final sentence about \emph{the rest}. The standard-contraction argument already proves commutativity of the entire displayed diagram. Also insert \emph{the} in \emph{the fact that the image}. This obsolete-diagram issue from the human report remains unresolved.
}
% END REFEREE INSERTION R137
\end{proof}


\begin{rem}\label{eg:eb:factor}
We may define the motivic volume operator in Notation~\ref{mov:vol} using $\bb E_g$ instead of $\bb E_b$ (of course $[\A]$ will have to be inverted, see Remark~\ref{eb:no:eg}).
For proper invariant objects, this only changes the image by a factor. To see this, first note that  Theorem~\ref{main:prop} (or rather Corollary~\ref{L:sur:c}) shows that every proper invariant object $A \in \VF_*$ with $\dim_{\VF}(A) = n$ is isomorphic to an object of the form $\bb L \bm U$ with $(\bm U, 0) \in  \mgRV^{\db}[n]$. Therefore, in light of the isomorphism ${\mu}\Psi$ in (\ref{bdd:to:2bdd}) and the defining conditions of $\bb E_g$, $\bb E_b$ in \S~\ref{subs:uni:retr:novol}, we have
\begin{equation}\label{off:fac:A}
\bb E_g \Big (\int [A] \Big ) = \bb E_b \Big (\int [A] \Big )[\A]^{-n}.
\end{equation}
\end{rem}

\section{Motivic Milnor fiber}\label{sec:spec:pui}

We begin with a very brief discussion on specialization to henselian subfields. Let $\mathbb{M}$ be a $\VF$-generated  substructure. Recall that $\mathbb{S}$ is a part of the language and hence all other substructures contain it. If $X \sub \VF^n  \times \RV^m$ is a definable (and hence quantifier-free definable) set then the \memph{trace} of $X$ in $\mathbb{M}$, denoted by $X(\mathbb{M})$, is the set of $\mathbb{M}$-rational points of $X$, that is,
\[
X(\mathbb{M}) = X \cap (\VF(\mathbb{M})^n  \times \RV(\mathbb{M})^m).
\]
Such a trace is also called a \memph{constructible set} in $\mathbb{M}$ since it is indeed quantifier-free definable in $\mathbb{M}$. Note that, however, if $f : X \fun \Gamma$ is a definable function then the image $f(X(\mathbb{M}))$  is not necessarily a set in $\Gamma(\mathbb{M})$, but rather a set in the divisible hull $\Q \otimes \Gamma(\mathbb{M})$  of $\Gamma(\mathbb{M})$. For instance, if $\mathbb{M} = \C \dpar t$ then $\Gamma(\mathbb{M}) = \Z$ and hence $\gamma \in \Gamma$ is definable if and only if $\gamma \in \Q \otimes \Gamma(\mathbb{M}) = \Q$. On the other hand, if $X$ is a set in $\Gamma$ and $f$ is a piecewise $\mgl_k(\Z)$-transformation on $X$ with constant $\Gamma(\mathbb{M})$-terms  then $f(X(\mathbb{M}))$ is of course a set in $\Gamma(\mathbb{M})$; this is the situation in the $\Gamma$-categories (see Remark~\ref{GLZ:char}).


%\begin{rem}\label{lift:hen:bij}
%Suppose that $X$, $Y$ are definable sets and $g : X(\mathbb{M}) \fun Y(\mathbb{M})$ is a constructible bijection in $\mathbb{M}$. Obviously there may or may not be a definable bijection between $X$ and $Y$, let alone one whose trace equals $g$. But the quantifier-free formula that defines $g$ in $\mathbb{M}$ also gives rise to a definable bijection $f : X' \fun  Y'$ with $X' \sub X$ and $Y' \sub Y$ such that its trace $f'(\mathbb{M}): X'(\mathbb{M}) \fun  Y'(\mathbb{M})$ indeed equals $g$.
%\end{rem}


If $\mathbb{M}$ is definably closed and $\Gamma(\mathbb{M})$ is nontrivial, or equivalently (see Lemma~\ref{cut:to:hensel:substru}), the valued field $(\VF(\mathbb{M}), \OO(\mathbb{M}))$  is henselian, then $\mathbb{M}$ is \memph{functionally closed}, that is,  for any definable set $X$ and any definable function $f$ on $X$, the image $f(X(\mathbb{M}))$ is a set that is definable in $\mathbb{M}$ (which then, ex post facto, is  constructible in $\bb M$). This is all we need to deduce the results below.

% BEGIN REFEREE INSERTION R138
\RefereeRemark{138}{Error; revised}{Henselian definability is not quantifier-free constructibility}{%
\textit{Original source line 2775.}\quad The parenthetical assertion is false, as already observed in the supplied human report. For $\mathbb M=\C\dpar t$ and $f(x)=x^2$, the image $f(\VF(\mathbb M))$ is the set of squares. It is definable in $\mathbb M$, but is neither the trace of the ambient image $f(\VF)=\VF$ nor quantifier-free definable in the stated language: on nonzero elements it detects parity of the valuation. Henselianity gives closure under ambient definable functions with values in the home sorts, not quantifier elimination in the henselian substructure. Delete the constructibility claim. The home-sort restriction is already implicit in Terminology~\ref{nota:sub}; the preceding paragraph correctly treats $\Gamma$-valued functions separately. The later $\RV$-bijection arguments need only the valid statement that a definable bijection and its inverse preserve $\RV(\mathbb M)$.
}
% END REFEREE INSERTION R138


%
%\begin{defn}[$\mathbb{M}$-specialized $\RV$-categories]
%An object of the category $\mgRV[k]_{\mathbb{M}}$ is a triple of the form $\bm U(\mdl M) = (U(\mdl M), f(\mdl M),  \omega \rest U(\mdl M))$, where $\bm U = (U, f, \omega) \in \mgRV[k]$. Any constructible function of the form $F(\mdl M) : \bm U(\mdl M) \fun \bm V(\mdl M)$, where $F : \bm U \fun \bm V$ is a $\mgRV[k]$-morphism, is a morphism of $\mgRV[k]_{\mdl M}$.
%\end{defn}
%
%%The categories $\mgRES[k]_{\mdl M}$, $\mG[k]_{\mdl M}$, etc., are formulated analogously.
%
%We call $\gsk \mgRV[*]_{\mdl M}$,  etc.,  \memph{$\mdl M$-constructible Grothendieck semirings} associated with $\mdl M$.
%
%
%Suppose that $[\bm A] = [\bm C]$, $[\bm B] = [\bm D]$ in $\gsk \mgRV[*]$ and there is a $\gsk \mgRV[*]_{\mdl M}$-morphism $\bm A(\mdl M) \fun \bm B(\mdl M)$. By Remark~\ref{lift:hen:bij}, there is a $\mgRV[*]$-morphism $F : \bm A' \fun \bm B'$ with $\bm A' \sub \bm A$, $\bm B' \sub \bm B$, $\bm A(\mdl M) = \bm A'(\mdl M)$, and $\bm B(\mdl M) = \bm B'(\mdl M)$, and hence there is a $\mgRV[*]$-morphism $G : \bm C' \fun \bm D'$ with $\bm C' \sub \bm C$, $\bm D' \sub \bm D$, $\bm C(\mdl M) = \bm C'(\mdl M)$, and $\bm D(\mdl M) = \bm D'(\mdl M)$. Thus the following binary relation is well-defined and is indeed a semiring congruence relation:
%\[
%\mathfrak{R}_{\mdl M} = \{ ([\bm A], [\bm B]) \in (\gsk \mgRV[*])^2 : [\bm A(\mdl M)] = [\bm B(\mdl M)] \text{ in } \gsk \mgRV[*]_{\mdl M}\}.
%\]
%Obviously $\gsk \mgRV[*]_{\mdl M} \cong \gsk \mgRV[*] / \gR_{\mdl M}$.


%It is also routine to check that the diagram commutes:
%\begin{equation}\label{dia:tensor:hen}
%\bfig
% \Square(0,0)/->`->`->`->/<400>[{\gsk \mgRES[*] \otimes \gsk \mG[*]}`{\gsk \mgRES[*]_{\mdl M} \otimes \gsk \mG[*]_{\mdl M}}`{\gsk \mgRV[*]}`{\gsk \mgRV[*]_{\mdl M}}; -/{\gR_{\mdl M}} \otimes - / {\gG_{\mdl M}}`\mu \Psi`{\mu \Psi_{\mdl M}}`-/{\gR_{\mdl M}}]
%\efig
%\end{equation}
%where the tensor product on the righthand side is taken over $\gsk \mG^{\fin}[*]_{\mdl M}$ and the isomorphism .
%Similar constructions are available for the other categories and their doubly bounded versions.To simplify the notation, we shall write $\mgRV[*]_{\C \dpar{ t^{1/m} }}$ as $\mgRV[*]_m$.


\subsection{Piecewise retraction to $\RES$}\label{piece:res}

As in \S~\ref{jp:fubini}, from here on, we work in the field $\puC$ of complex Puiseux series (with an implicit reduced cross-section) and take  $\mathbb{S} = \C \dpar{ t }$. The value group $\Gamma$ is identified with $\Q$. We also consider the family of henselian subfields $\C \dpar{ t^{1/m} }$,  $m \in \Z^+$, regarded as substructures. The value group $\Gamma(\C \dpar{ t^{1/m} })$ is identified with $m^{-1} \Z$.



\begin{nota}\label{loc:rings}
For any ring $R$, let $R[T^{\Q}]$ denote the ring of Puiseux polynomials over $R$, that is, the group ring of $\Q$ over $R$.

Let $\gmv[[\A]^{-1}][T,T^{-1}]$ be the obvious subring of $\gmv[[\A]^{-1}][T^{\Q}]$. The canonical image of $\gmv [T^{\Q}]$ in  $\gmv[[\A]^{-1}][T^{\Q}]$ is just denoted as such. The assignment $T \efun [\A]$ determines a ring homomorphism
\begin{equation}\label{pui:laurent}
  \bm \eta : \gmv[[\A]^{-1}][T,T^{-1}] \fun \gmv[[\A]^{-1}].
\end{equation}
\end{nota}

%(recall Terminology~\ref{twist:rv:ob})

Recall Notations~\ref{comp:form} and \ref{bipo:gdv}. Let $\bm U = (U, f, \omega) \in \mgRV^{\db}[k]$. In particular, if $f=(f_1,\dots,f_n)$, we have introduced the shorthand
\[
\omega_{f}: x\in U \efun \omega(x)+\sum_{i=1}^n\vrv(f_i(x)).
\]
The image $\omega_f(U_\gamma)$ is finite for every $\gamma \in \vrv(U)$, because the $\K$-sort and the $\Gamma$-sort are orthogonal to each other. It follows that there is a bipolar twistoid decomposition $(U_i)_i$ of $U$ such that each restriction $\omega_f \rest U_i$ $\vrv$-contracts to a function $\sigma_i : I_i = \vrv(U_i) \fun \Gamma$. Write $I_{i,m}$ for the set $I_i(m^{-1}\Z)$ of  $m^{-1}\Z$-rational points of $I_i$. We assign to $\bm U$ the expression
\begin{equation}\label{hm:ass}
  h_m(\bm U) = \sum_i [\tbk(U_i)]^{\hat \mu} \sum_{\gamma \in I_{i,m}} T^{-m \sigma_i(\gamma)},
\end{equation}
which is  a finite sum and hence  belongs to $\gmv[[\A]^{-1}][T^{\Q}]$.

If $\bm U$ is already of the form $(V, g, \omega') \times (I, \sigma) \in \mgRES[i] \times \mG^{\db}[j]$, where $\vrv(g)$ is a singleton, $\vrv(V) \in m^{-1}\Z$, and $\omega'_g$ is constant, then clearly
\[
h_m(\bm U) = \Theta([V])T^{-m\omega'_g} \cdot [\G_m]^{j} \sum_{\gamma \in I_{m}} T^{-m (\sigma(\gamma) + \Sigma \gamma )};
\]
note that if $\vrv(V) \notin m^{-1}\Z$ then simply  $h_m(\bm U) = 0$. In light of the isomorphism ${\mu} \Psi^{\db}$ in Remark \ref{tensor:bd}, we may as well take this as the definition of $h_m$, which is essentially how it is done in \cite[\S~8.2]{hru:loe:lef}. Computationally, though,  (\ref{hm:ass}) is more effective.

% BEGIN REFEREE INSERTION R139
\RefereeRemark{139}{Clarification}{Partition the RES factor by its valuation}{%
\textit{Original source line 2836.}\quad The condition $\vrv(V)\in m^{-1}\Z$ is not well-typed: $\vrv(V)$ is a finite set of tuples, and the hypothesis that $\vrv(g)$ is a singleton does not make $\vrv(V)$ a singleton. First partition $V$ by its finitely many valuation tuples, say $\vrv(V)=\{\delta\}$ on each piece. Then use $\delta\in(m^{-1}\Z)^r$ in the displayed formula and assign zero to the other pieces. This also resolves the corresponding outstanding point in the human report.
}
% END REFEREE INSERTION R139


The ``piecewise retraction'' formula (\ref{hm:ass}) is indeed quite similar to the ``uniform retraction'' formula (\ref{twist:C}). The difference is just that the coefficient $\chi_b(I_i)$ in the latter is replaced in the former by a formal ``counting'' sum $\sum_{\gamma \in I_{i,m}} T^{-m \sigma_i(\gamma)}$ that also takes the volume form into account. As mentioned in Remark~\ref{bipo:desc}, that the ``uniform retraction'' formula does not depend on the choice of the bipolar twistoid decomposition is simply a consequence of the construction of $\bb E_b$ itself. The same holds for the ``piecewise retraction'' formula, but we do need an argument.

\begin{lem}\label{muhat:section}
The  assignment (\ref{hm:ass}) does not depend on the choice of the bipolar twistoid decomposition and is invariant on isomorphism classes.
\end{lem}
\begin{proof}
Let $\bm V = (V, g, \pi) \in \mgRV^{\db}[k]$  and $f : U \fun V$ be a $\mgRV^{\db}[k]$-morphism. Let $D = (U_i)_i$, $E = (V_i)_i$ be bipolar twistoid decompositions of $U$, $V$ satisfying the condition above. We need to show that $h_m(\bm U)$, which depends on $D$, and $h_m(\bm V)$, which depends on $E$, are  equal.

% BEGIN REFEREE INSERTION R140
\RefereeRemark{140}{Typo}{The letter f has two incompatible meanings}{%
\textit{Original source line 2845.}\quad Rename the isomorphism $f:U\fun V$ to $F$ throughout this proof. The symbol $f$ already denotes the structural map in $\bm U=(U,f,\omega)$; expressions such as $\omega_f$ must continue to refer to that structural map.
}
% END REFEREE INSERTION R140
This is clear if $\bm U = \bm V$,  and $E$ is trivial (so $U$ is already a bipolar twistoid and $\omega_f$ already $\vrv$-contracts to a function on $\vrv(U)$), and $D$ is $\Gamma$-cohesive or $\vrv(U_i) = \vrv(U_j)$ for all $i$, $j$. The case that $\bm U = \bm V$ and $D$ is a refinement of $E$  follows easily from this since there is a refinement $(U_{ij})_{ij}$ of $D$ such that $(U_{ij})_{j}$ is a $\Gamma$-cohesive twistoid decomposition of $U_i$  and $\vrv(U_{ij}) = \vrv(U_{i'j})$ for all $i$, $i'$. Finally,  if $\omega_f$, $\pi_g$ both $\vrv$-contract to a function and $f$ $\vrv$-contracts to a bijection then, by Lemma~\ref{coh:decom}, we may assume that $U$, $V$ are already bipolar twistoids. In that case the desired equality follows because $\C \dpar{ t^{1/m} }$ is functionally closed.

For the general case, we first remark that the image $\vrv(f(U_\gamma))$ is finite for every $\gamma \in \vrv(U)$. It then follows from Lemmas~\ref{shift:K:csn:def} and~\ref{coh:decom} that there is a twistoid decomposition $(f_i)_i$ of $f$ such that every $f_i$ is a $\mgRV^{\db}[k]$-morphism as in the last special case considered above and, moreover, is compatible with $D$, $E$ in the obvious sense, in other words, the domains and ranges of these $f_i$ induce bipolar refinements of $D$, $E$.  The result follows.
\end{proof}

% BEGIN REFEREE INSERTION R141
\RefereeRemark{141}{Clarification; revised}{Explain lattice and weight preservation}{%
\textit{Original source line 2849.}\quad In the special case with a bijective valuation contraction $\bar F$, spell out the two identities needed for the counting argument: $\bar F(I\cap(m^{-1}\Z)^r)=J\cap(m^{-1}\Z)^s$ and $\omega_f=\pi_g\circ F$. The first follows from closure of $\RV(\C\dpar{t^{1/m}})$ under $F$ and $F^{-1}$; the second follows from the $\Gamma$-Jacobian condition. The induced maps on each residue fiber are equivariant, so their classes agree under $\Theta$. This supplies the currently omitted bridge from functional closure to equality of the weighted, equivariant sums.
}
% END REFEREE INSERTION R141

Thus, $h_m$ may be viewed as a map on $\gsk \mgRV^{\db}[k]$. It is now routine to check that  we have in effect constructed a ring homomorphism
\[
\bm h_m : \ggk \mgRV^{\db}[*] \fun  \ggk^{\hat \mu} \var_{\C}[[\A]^{-1}][T^{\Q}].
\]

Let $\ggk^\natural_m \mgRV^{\db}[*]$ denote the subring $(\bm h_m)^{-1}(\gmv[[\A]^{-1}][T, T^{-1}])$ of $\ggk \mgRV^{\db}[*]$.

\begin{lem}\label{PG:van}
The homomorphism
\[
\bm \eta \circ \bm h_m : \ggk^\natural_m \mgRV^{\db}[*] \fun \gmv[[\A]^{-1}]
\]
vanishes on $(\bm P_{\Gamma})$.
\end{lem}
Note that the ideal $(\bm P_\Gamma)$ of $\ggk \mgRV^{\db}[*]$ in Notation~\ref{pgammabis} is now generated by the elements $\bm P_\gamma$ with $\gamma \in \Z$ (because the point $t_\gamma$ there needs to be definable, which is possible only if $\gamma \in \Z$ in the current setting).

% BEGIN REFEREE INSERTION R142
\RefereeRemark{142}{Typo}{The generator radius is nonnegative}{%
\textit{Original source line 2865.}\quad Replace $\gamma\in\Z$ here and at the beginning of the proof by $\gamma\in\Z_{\geq0}$, or by $\gamma\in\Z_{>0}$ after discarding the zero generator. Notation~\ref{pgammabis} defines $\bm P_\gamma$ only for nonnegative radii; the subsequent sum from $1$ to $m\gamma$ uses this restriction.
}
% END REFEREE INSERTION R142

\begin{proof}
For $\gamma \in \Z$, the image of $[\RV^{\circ \circ} \mi \RV^{\circ \circ}_{\gamma}] + [\{t_\gamma\}]$ under $\bm h_m$ in $\gmv[[\A]^{-1}][T, T^{-1}]$ is
\[
 ([\A] - 1) \sum_{i = 1}^{m\gamma} T^{-i} + T^{-m\gamma},
\]
which, after passing to $\gmv[[\A]^{-1}]$ via $\bm \eta$, becomes $1 = \bm \eta (\bm h_m([1]))$.
\end{proof}

% BEGIN REFEREE INSERTION R172
\RefereeRemark{172}{Clarification; new}{Specify the ideal inside the natural subring}{%
\textit{Original source line 2873.}\quad The ideal $J=(\bm P_\Gamma)$ was defined in the full ring $B=\ggk\mgRV^{\db}[*]$, whereas $\bm\eta\bm h_m$ has domain $B_m^\natural$. Interpret ``vanishes on $J$'' as vanishing on $J\cap B_m^\natural$, and use this intersection in the quotient notation below. The displayed generator computation extends to that intersection as follows. The polynomials $p_\gamma=\bm h_m(\bm P_\gamma)$ have integer exponents. Decompose the Puiseux polynomial ring by exponent classes in $\Q/\Z$, and let $\operatorname{pr}_0$ retain its integer-exponent part. If $x=\sum_j a_j\bm P_{\gamma_j}\in B_m^\natural$, then
\[
 \bm h_m(x)=\sum_j\operatorname{pr}_0(\bm h_m(a_j))p_{\gamma_j}.
\]
Applying $\bm\eta$ gives zero by the displayed calculation. This supplies the small algebraic step needed when the coefficients $a_j$ themselves lie outside $B_m^\natural$.
}
% END REFEREE INSERTION R172

\begin{rem}\label{dia:to:db:dag}
If $\bm U = (U, f, l) \in \mgRV^{\db}[*]$ with $l \in m^{-1}\Z$ (constant volume form) then the exponents in (\ref{hm:ass}) are all integers and hence $[\bm U] \in \ggk^\natural_m \mgRV^{\db}[*]$. Actually we shall only need the case $l = 0$.

The ring $\bigcap_{m \in \Z^+} \ggk^\natural_m \mgRV^{\db}[*]$ is denoted by $\ggk^\natural \mgRV^{\db}[*]$.

If $\bm A = (A, l) \in \mgVF^\diamond[*]$ with $l$  constant then $\int^\diamond [\bm A]$ may be expressed as $[(U, f, l)]/ (\bm P_{\Gamma})$, and hence if $l \in \Z$ then $\int^\diamond [\bm A]$ belongs to $\ggk^\natural \mgRV^{\db}[*]/ (\bm P_{\Gamma})$. In that case, by Lemma~\ref{PG:van}, for every $m \in \Z^+$, the expression $(\bm \eta \circ \bm h_m \circ \int^\diamond)([\bm A])$ designates a unique element  in $\gmv[[\A]^{-1}]$.
\end{rem}


Denote by $\RES_{m}$ the full subcategory of $\RES$ such that $U \in \RES_{m}$ if and only if every $\gamma \in \vrv(U)$ is a tuple in $m^{-1} \Z$, or equivalently, $U \in \RES_{m}$ if and only if the action on $U$ of the kernel of the canonical projection $\hat \mu \fun \mu_m$ is trivial.


Let $\beta = (\beta_1, \ldots, \beta_n) \in (m^{-1} \Z)^n$ and $A \sub \OO^n \times \RV^l$ be a proper $\beta$-invariant definable set such that $\pvf \rest A$ is finite-to-one. Then there is a set
\[
A[m;\beta] \sub \prod_i \C[t^{1/m}] / t^{\beta_i + 1/m} \times \RV(\C \dpar{t^{1/m}})^l
\]
such that, for every $t \in \RV(\C \dpar{t^{1/m}})^l$, the $\VF$-fiber $A_t(\C \dpar{t^{1/m}})$ is, under the obvious quotient map, the preimage of the fiber $A[m;\beta]_t$. Note that $\pvf \rest A[m;\beta]$ is still finite-to-one, for otherwise, by the $\beta$-invariance of $A$, $\pvf \rest A$ would fail to be finite-to-one. We can think of $\C[t^{1/m}] / t^{\beta_i + 1/m}$ as the object $\prod_{0 \leq \gamma \leq \beta_i} \gamma^\sharp$ of $\RES_{m}$. Then, since every element in the value group $\Q$ is definable and $\vrv(A[m;\beta])$ is doubly bounded,  it follows that  $A[m;\beta]$ is indeed a finite disjoint union of definable sets in $\RV$ (no extra parameters other than those in $\mathbb{S} = \C \dpar{ t }$ are needed) and hence, as such, is an object of $\RES_{m}$ (see \cite[\S~4.2]{hru:loe:lef} for more details).

% BEGIN REFEREE INSERTION R143
\RefereeRemark{143}{Error; revised}{Truncated coefficient spaces need zero coefficients}{%
\textit{Original source line 2891.}\quad Take $\beta_i\geq0$, and replace the asserted coefficient space by
\[
 \prod_{\substack{\gamma\in m^{-1}\Z\\0\leq\gamma\leq\beta_i}}
       (\gamma^\sharp\cup\{0_\gamma\}).
\]
Here $0_\gamma$ denotes the zero of the one-dimensional residue vector space. By the manuscript's convention $\gamma^\sharp=\vrv^{-1}(\gamma)$ excludes zero, so the displayed product in the text allows no vanishing coefficients and is not the truncated polynomial space. The lattice restriction on the index is also necessary to make the product finite. Interpret the definition of $\RES_m$ accordingly when a coordinate is zero. More explicitly, in the preceding definition of $\RES_m$, allow coordinates equal to $\infty=\rv(0)$ in addition to coordinates with valuation in $m^{-1}\Z$: the kernel of $\hat\mu\to\mu_m$ fixes these zeros too.
}
% END REFEREE INSERTION R143

\begin{lem}[{\cite[Lemma~4.2.1]{hru:loe:lef}}]
Let $\beta' \in (m^{-1} \Z)^n$ with $\beta_i \leq \beta'_i$ for all $i$. Then
\begin{equation}\label{arc:com}
  [A[m;\beta']] = [A[m;\beta]][\A]^{m \Sigma (\beta' - \beta)} \in \ggk \RES_{m}.
\end{equation}
\end{lem}

% BEGIN REFEREE INSERTION R144
\RefereeRemark{144}{Error; revised}{The truncation relation requires the affine-action quotient}{%
\textit{Original source line 2898.}\quad The equality generally fails in the raw ring $\ggk\RES_m$. Take $A=\OO$, $m=2$, $\beta=0$ and $\beta'=1/2$. Then $A[2;0]$ is a trivial affine line, whereas $A[2;1/2]$ is its product with a line carrying the sign action of $\mu_2$. The $\mu_2$-fixed loci of the asserted left and right sides have classes $[\A]$ and $[\A]^2$, respectively, distinguished for example by their Hodge polynomials. The equality is valid after imposing $[\gamma^\sharp]=[\G_m]$: put it in $\sggk\RES_m$, and define $\tilde A[m]$ in its localization below. In the proof of the cited \href{https://arxiv.org/pdf/1111.1954}{Hrushovski--Loeser Lemma~4.2.1}, averaging a section makes that section equivariant, but does not turn the remaining linear fiber action into the trivial action; this is precisely why the quotient is needed.
}
% END REFEREE INSERTION R144

Thus the element $[A[m;\beta]][\A]^{- m\sum \beta }$ in $\ggk \RES_{m}[[\A]^{-1}]$ does not depend on $\beta$; we denote it by $\tilde A[m]$.

The localization of $\Theta$ at $[\A]$ is still denoted by $\Theta$.

\begin{lem}\label{inv:dir:com}
Suppose that $(A, 0) \in \mgVF^\diamond[n]$. Then, in $\gmv[[\A]^{-1}]$,
\[
(\bm \eta \circ \bm h_m) \Big( \int^\diamond [(A, 0)] \Big) = \Theta(\tilde A[m]).
\]
\end{lem}

% BEGIN REFEREE INSERTION R145
\RefereeRemark{145}{Clarification}{The right-hand side has only been defined inside O}{%
\textit{Original source line 2909.}\quad The definition of $\tilde A[m]$ above assumes $A\subseteq\OO^n\times\RV^l$, whereas this lemma allows an arbitrary proper invariant object, whose field projection need only be bounded. Either add the assumption $A\subseteq\OO^n\times\RV^l$ (which suffices for the application), or first extend the definition using finite Laurent-polynomial truncations beginning at a common lower valuation bound. A scaling into $\OO^n$ must also account for its Jacobian and cannot simply retain the zero volume form.
}
% END REFEREE INSERTION R145
\begin{proof}
By Lemma~\ref{L:measure:surjective}, there exists a special covariant bijection $F : (A, 0) \fun (\bb L \bm U, 0)$, where $\bm U \in \RV^{\db}[*]$. Observe that we can choose apertures $(\gamma_i, t_i)$ with $\gamma_i \in m^{-1} \Z$ so that $F$ induces a definable bijection between $A[m;\beta]$ and  $\bb L \bm U[m;\beta]$ for all sufficiently large $\beta$; it is for this purpose that each $\gamma_i$ needs to lie in $m^{-1} \Z$, due to the second clause of (\ref{centri}) and the mandatory regularizations in Definition~\ref{defn:special:bijection}. Thus, we may assume that $A$ is already of the form $\bb L \bm U$. Let $D$ be a $\Gamma$-cohesive bipolar twistoid decomposition of $\bm U$. Since both sides respect finite disjoint union of proper $\beta$-invariant definable sets, we may further assume that $D$ is actually trivial or even $\vrv(\bm U)$ is a singleton. Then the equality follows from a simple computation.
\end{proof}

% BEGIN REFEREE INSERTION R146
\RefereeRemark{146}{Gap; revised}{The singleton reduction needs the fixed-m trace}{%
\textit{Original source line 2912.}\quad Finite additivity cannot reduce an arbitrary positive-dimensional $\vrv(\bm U)$ to a singleton. For the fixed integer $m$, first restrict to the finite set $\vrv(U)\cap(m^{-1}\Z)^r$: the remaining pieces contribute neither to $h_m$ nor to the $\C\dpar{t^{1/m}}$-truncation. On a remaining fiber, let $\alpha_i=\vrv(f_i(u))$, after refining so these are constant, and take $\beta_i\geq\alpha_i$. There are $m(\beta_i-\alpha_i)$ free coefficients following the prescribed leading term in coordinate $i$. Thus its normalized class is
\[
 \Theta([U_\gamma])\,[\A]^{-m\sum_i\alpha_i},
\]
which is precisely its contribution to $\bm\eta\bm h_m$. This calculation uses the affine-action quotient noted above. For definability, the apertures can be chosen in $\Z_{\geq0}$, a cofinal subset of $m^{-1}\Z$, so the required points $t_i$ remain definable over $\C\dpar t$. This integrality observation does not settle the preceding truncation-compatibility step; see the separate remark on that step.
}
% END REFEREE INSERTION R146

% BEGIN REFEREE INSERTION R173
\RefereeRemark{173}{Gap; new}{Integral apertures do not ensure equal-cutoff truncation compatibility}{%
\textit{Original source line 2912.}\quad The first reduction needs a separate argument. Choosing integral apertures does not by itself make an allowed special covariant bijection induce equal-cutoff truncations for every sufficiently large vector $\beta$. For example, take $m=1$ and
\[
 A=(1+t+\MM_2)\times(t+\MM_3),\qquad \lambda(y)=t^{-1}y.
\]
The focus satisfies the required $(3,2)$-covariance. With aperture $(2,\rv(t^2))$, every point uses the first clause of (CT), because $\vv(x-\lambda(y))=1$. Following this shear by two constant-focus transformations yields a permitted preparation
\[
 F(x,y)=(x-t^{-1}y-t+t^2,\ y-t+t^3)
\]
onto the $\RV$-pullback $(t^2+\MM_2)\times(t^3+\MM_3)$; the additional foci are $t-t^2$ and $t-t^3$, with sufficiently large integral apertures. All regularization labels are constant on these sets. For every integer $N>3$, the inputs $(1+t,t)$ and $(1+t,t+t^{N+1})$ have the same $(N,N)$-truncation, but their first output coordinates differ by $-t^N$. Thus this $F$ induces no such truncation map for arbitrarily large diagonal cutoffs. Compatible anisotropic cutoffs do exist for this example, namely $\beta_y\geq\beta_x+1$. Prove invariance of the normalized truncation class using compatible cutoffs for each centripetal step and its inverse, or strengthen the preparation lemma to provide the needed compatibility. This example identifies a missing step in the proof; it does not disprove the lemma's coefficient identity or rule out a better choice of $F$.
}
% END REFEREE INSERTION R173

\subsubsection{Bounded version}
\label{sec:hm:bounded}
Although not needed, we can  extend the construction above to the category $\mgRV^{\bdd}[*]$; this category is formulated as $\mgRV^{\db}[*]$, but with doubly bounded objects replaced by bounded ones.

Let $\gmv[[\A]^{-1}][T^{\Q}]_{\loc}$ be the localization of $\gmv[[\A]^{-1}][T^{\Q}]$ by the multiplicative family generated by the elements $1 - T^{-i}$, $i \in \Z^+$. The ring $\gmv[[\A]^{-1}][T, T^{-1}]_{\loc}$ is the localization of $\gmv[[\A]^{-1}][T, T^{-1}]$ by the same multiplicative family. Similarly one defines the ring $\gmv[[\A]^{-1}]_{\loc}$ with respect to the multiplicative family generated by the elements $1 - [\A]^{-i}$, $i \in \Z^+$. As in (\ref{pui:laurent}), the assignment the $T \efun [\A]$ determines a ring homomorphism

% BEGIN REFEREE INSERTION R147
\RefereeRemark{147}{Typo}{Duplicated article}{%
\textit{Original source line 2918.}\quad Replace ``the assignment the $T\efun[\A]$'' by ``the assignment $T\efun[\A]$''. The symbol $T$ here is the weight variable used in $h_m$; using a different letter from the later zeta variable would also make the two operations easier to distinguish.
}
% END REFEREE INSERTION R147
\begin{equation}\label{pui:laurent:loc}
  \bm \eta_{\loc} : \gmv[[\A]^{-1}][T, T^{-1}]_{\loc} \fun \gmv[[\A]^{-1}]_{\loc}.
\end{equation}


We assign to each $\bm U= (U, f, \omega) \in \mgRV^{\bdd}[k]$ the same expression as in (\ref{hm:ass}), which is now a formal Laurent series in $T^{-1/l}$ for some integer $l > 0$, because each $\sigma_i$ is piecewise $\Q$-linear  and its graph is bounded. By \cite[Lemma~8.2.1]{hru:loe:lef}, if $l =1$ then the series belongs to $\gmv[[\A]^{-1}][T, T^{-1}]_{\loc}$ and hence,  in general, it belongs to $\gmv[[\A]^{-1}][T^{\Q}]_{\loc}$ (consider any isomorphism of $\puC$ that fixes $\C$ and sends $1/l$ to $1$).  So the assignment  yields a ring homomorphism $\bm h_{\loc,m}$ and consequently a subring $\ggk^\natural_m \mgRV^{\bdd}[*]$.

% BEGIN REFEREE INSERTION R148
\RefereeRemark{148}{Error}{Bounded support does not imply summability}{%
\textit{Original source line 2924.}\quad This extension is not defined on the declared category. Take $\bm U=(\RV^{\circ\circ},\mathrm{id},-\vrv)$. It is bounded in the stated sense, but $\omega_f=0$, and the proposed expression is
\[
 ([\A]-1)\sum_{i>0}1,
\]
which is not a formal Laurent series. Even assuming the graph of the weight is bounded below would not make its level sets finite. One needs a genuine summability condition, for example that the weight $\omega_f$ is bounded below and its sublevel sets on the relevant valuation lattice are finite, together with a proof that the resulting class of objects is compatible with the chosen morphisms and quotient. The cited Hrushovski--Loeser Lemma~8.2.1 concerns the sum-of-coordinates weight on a set bounded below and does not justify arbitrary $\omega$. Since this subsection is expressly unused, deleting this extension and the corresponding arrows is a safe alternative. Finally, clarify that ``sends $1/l$ to $1$ refers to the induced map on the value group, or write $t^{1/l}\mapsto t$: it cannot refer to field scalars in $\C$. To clear rational exponents without changing the parameter field, simply substitute $S=T^{1/l}$ and use that $1-S^{-a}$ divides $1-S^{-al}=1-T^{-a}$.
}
% END REFEREE INSERTION R148


\begin{lem}
\label{lem:hm:bounded}
The homomorphism
\[
\bm \eta_{\loc} \circ \bm h_{\loc,m} : \ggk^\natural_m \mgRV^{\bdd}[*] \fun \gmv[[\A]^{-1}]_{\loc}
\]
 vanishes on $(\bm P)$.
\end{lem}
\begin{proof}
The image of $[\RV^{\ccirc}]$ in $\gmv[[\A]^{-1}][T, T^{-1}]_{\loc}$ is
\[
 ([\A] - 1) \sum_{i > 0} T^{-i} = ([\A] - 1) / (T-1)
\]
and hence, after passing to $\gmv[[\A]^{-1}]_{\loc}$ via $\bm \eta_{\loc}$, it becomes $1 = \bm \eta_{\loc}( \bm h_{\loc,m}([1]))$.
\end{proof}




\subsection{Constructing motivic Milnor fiber}\label{sec:milnor}

Finally, we proceed to replicate the construction in \cite[\S~8]{hru:loe:lef} so to recover the motivic zeta function with coefficients in $\gmv[[\A]^{-1}]$ (this is the purpose of Lemma~\ref{inv:dir:com}) and the corresponding motivic Milnor fiber.

% BEGIN REFEREE INSERTION R149
\RefereeRemark{149}{Typo}{Missing as}{%
\textit{Original source line 2948.}\quad Replace ``so to recover'' by ``so as to recover''. The subsection title would read more naturally as ``Constructing the motivic Milnor fiber''. 
}
% END REFEREE INSERTION R149

Let $X$ be a smooth connected complex algebraic variety of dimension $d$ and $f : X \fun \A^1$ a nonconstant regular function. Suppose that $0 \in X(\C) \cap f^{-1}(0)$ is singular. For $m \geq 1$, we consider the set of truncated arcs

% BEGIN REFEREE INSERTION R150
\RefereeRemark{150}{Clarification}{Critical point of f, not singular point of X}{%
\textit{Original source line 2950.}\quad Since $X$ is assumed smooth, the phrase ``$0\in X(\C)\cap f^{-1}(0)$ is singular'' needs qualification. Write ``let $0\in f^{-1}(0)$ be a critical point of $f$, that is, $df_0=0$''. This avoids identifying a critical point with a singular point of the ambient variety, and handles a nonreduced zero fiber without ambiguity.
}
% END REFEREE INSERTION R150
\begin{equation}\label{f:m:arc}
\mdl X_{m} = \set{ \varphi \in X(\C[t] / (t^{m+1})) \given f(\varphi) = t^m \mod t^{m+1}},
\end{equation}

% BEGIN REFEREE INSERTION R151
\RefereeRemark{151}{Error}{The local arc space still omits its base point}{%
\textit{Original source line 2953.}\quad Add $\varphi(0)=0$ to the defining conditions. As written, $\mdl X_m$ includes jets based at every point of $f^{-1}(0)$, whereas $\mdl X_f$ is the tube at the chosen point. For example, with $X=\A^1$ and $f(x)=x^2(x-1)$, the chosen point $0$ is critical and there is no local jet of order $m=1$ there, but the displayed definition includes the jet $1+t$ based at $1$. Thus equation~\eqref{z:coeff:change} is false with the present definition. This correction was already requested in the human report and remains outstanding.
}
% END REFEREE INSERTION R151
which may be viewed in a natural way as the set of closed points of an algebraic variety over $\C$ that carries a natural $\mu_{m}$-action, and the so-called \memph{nonarchimedean Milnor fiber}
\[
\mdl X_f = \set{ x \in X(\MM) \given \rv(f(x)) = \rv(t)}.
\]
Note that $\mdl X_f$ may be constructed in an affine neighborhood of $0$ and hence is indeed a (quantifier-free) definable set. Moreover, it is $\beta$-invariant for every $\beta \geq 1$. Therefore, $\mdl X_f$ may be regarded as an object of $\mgVF^\diamond[*]$ equipped with the constant volume form $0$. The  \memph{motivic zeta function} of $f$ is then the power series

% BEGIN REFEREE INSERTION R152
\RefereeRemark{152}{Gap}{Use an etale chart to obtain a proper invariant object}{%
\textit{Original source line 2958.}\quad For a general smooth $X$, $X(\MM)$ is not literal scheme notation, since $\MM$ is not a unital ring; define it as $\operatorname{red}^{-1}(0)\subseteq X(\OO)$. More substantially, embedding an affine neighborhood into $\A^N$ does not make this tube open in $\VF^N$ when $N>d$, so it does not establish proper invariance as defined earlier. Choose an \'etale coordinate map $q:(X,0)\to(\A^d,0)$. Henselian lifting identifies the tube at $0$ with $\MM^d$; in these coordinates the function is defined by an \'etale algebraic germ and is $1$-Lipschitz on the tube, which gives the asserted $\beta$-invariance for $\beta\geq1$. Use the volume form $q^*(dx_1\wedge\cdots\wedge dx_d)$, whose chart transitions have Jacobian of valuation zero on the tube. State this reduction before assigning the constant volume form $0$.
}
% END REFEREE INSERTION R152
\begin{equation}\label{def:zeta}
Z_f(T) = \sum_{m \in \Z^+} \Big(\bm \eta \circ \bm h_m \circ \int^{\diamond} \Big)([\mdl X_f]) T^m \in \gmv[[\A]^{-1}] \dbra T.
\end{equation}
By Remark~\ref{dia:to:db:dag} and Lemma~\ref{inv:dir:com}, we have, in $\gmv[[\A]^{-1}]$,
\begin{equation}\label{z:coeff:change}
 \Big(\bm \eta \circ \bm h_m \circ \int^{\diamond} \Big)([\mdl X_f]) = \Theta(\tilde{\mdl X_f}[m]) = [\mdl X_{m}][\A]^{-md},
\end{equation}
where the second equality follows if  we change the variable $t$ to $t^{1/m}$ in (\ref{f:m:arc}) (see the discussion in \cite[\S~6.2]{hru:loe:lef} for  detail). This shows that the coefficients of $Z_f(T)$ may also be written as $[\mdl X_{m}][\A]^{-md}$, which is how they are usually defined.

% BEGIN REFEREE INSERTION R153
\RefereeRemark{153}{Clarification}{Give the equivariant coefficient comparison}{%
\textit{Original source line 2966.}\quad After correcting $\varphi(0)=0$, put $s=t^{1/m}$ and choose $\beta=(1,\ldots,1)$. The truncation is modulo $s^{m+1}$ and $\rv(f(x))=\rv(t)$ becomes $f(x)\equiv s^m\pmod{s^{m+1}}$. Hence $\mdl X_f[m;\beta]$ is exactly the corrected $m$-jet variety. The action $s\mapsto\zeta s$, $\zeta\in\mu_m$, is the usual action on those jets, and the normalization is $[\A]^{-m\sum_i\beta_i}=[\A]^{-md}$. Including these sentences proves both the factor and equivariance without leaving the crucial comparison to a general section reference.
}
% END REFEREE INSERTION R153

\begin{rem}
One of the key steps of the main construction in \cite[\S~8]{hru:loe:lef} is to show (\ref{z:coeff:change}), see \cite[Proposition~8.3.1]{hru:loe:lef}, but with $\int^\diamond$  replaced by $\int$. This, and a few other statements in \cite[\S~8]{hru:loe:lef}, rely on a tensor decomposition of the ring $\ggk \mathrm{vol}\RV[*]$, as claimed in \cite[Proposition~10.10(2)]{hrushovski:kazhdan:integration:vf}. We have pointed out above that, unfortunately, this result does not hold.
\end{rem}

\begin{nota}\label{dag:loc}
Let $\gmv[[\A]^{-1}][T]_{\dag}$ be the localization of $\gmv[[\A]^{-1}][T]$ by the multiplicative family generated by the elements $1 - [\A]^{a}T^b$, where $a \in \Z$ and $b \in \Z^+$; it may be regarded as a subring of  $\gmv[[\A]^{-1}] \dbra T$.
\end{nota}

It is known that $Z_f(T)$ belongs to $\gmv[[\A]^{-1}][T]_{\dag}$ and, letting ``$T$ go to infinity'' as described in \cite[\S~8.4]{hru:loe:lef}, we get a limit
\[
\mathscr S_f \coloneqq - \lim_{T \limplies \infty} Z_f(T) \in \gmv[[\A]^{-1}],
\]
which is understood as the \memph{motivic Milnor fiber} attached to $f$.

\begin{rem}
Recall the homomorphism $\bb E^{\diamond}$ from (\ref{edb:eb:com}).
By Lemma~\ref{coh:decom} and Notation~\ref{bipo:gdv}, we can construct the homomorphism
\[
\Theta \circ \bb E^{\diamond} : \ggk \mgRV^{\db}[*] \fun \gmv
\]
simply using the expression in (\ref{twist:C}) with respect to  bipolar twistoid decompositions.
\end{rem}

\begin{thm}\label{direct:mil}
$\mathscr S_f = (\Theta \circ \bb E^{\diamond} \circ \int^{\diamond})([\mdl X_f])$.
\end{thm}
\begin{proof}
Let $[\bm U] = [(U, f, l)] \in \ggk^\natural \mgRV^{\db}[*]$ with $l \in \Z$  and consider the zeta function
\begin{equation}\label{U:zeta}
Z([\bm U])(T) = \sum_{m \in \Z^+} (\bm \eta \circ \bm h_m)([\bm U]) T^m \in \gmv[[\A]^{-1}]\dbra T.
\end{equation}
If $[\bm U] / (\bm P_\Gamma) = \int^{\diamond}([\mdl X_f])$ then this is $Z_f(T)$. Thus it is enough to show that  $Z([\bm U])(T)$ belongs to $\gmv[[\A]^{-1}][T]_{\dag}$, and $\lim_{T \limplies \infty} Z([\bm U])(T)$ exists and equals $- (\Theta \circ \bb E^{\diamond})([\bm U])$.

Without loss of generality, we may assume that $U$ is already a bipolar twistoid and the function $l_f: U \fun \Gamma$ (Notation~\ref{comp:form}) $\vrv$-contracts to a function $\sigma : I = \vrv(U) \fun \Gamma$.  Set $v = [\tbk(U)]^{\hat \mu}$. Then
\[
(\Theta \circ \bb E^{\diamond})([\bm U]) = \chi(I) v.
\]
Let $Z(v)(T) = \sum_{m \in \Z^+} v T^m$. This is an element in $\gmv[[\A]^{-1}][T]_{\dag}$ and $\lim_{T \limplies \infty}Z(v)(T) = -v$. Write  $I_{m} = I(m^{-1}\Z)$ and let
\[
Z(\sigma)(T) = \sum_{m \in \Z^+} \sum_{\gamma \in I_{m}} [\A]^{-m \sigma(\gamma)} T^m.
\]
By Remark~\ref{GLZ:char}, $\sigma$ is $\Z$-linear. So, by \cite[Proposition~8.5.2]{hru:loe:lef}, $Z(\sigma)(T)$ is in $\gmv[[\A]^{-1}][T]_{\dag}$ and $\lim_{T \limplies \infty}Z(\sigma)(T) = -\chi(I)$. Since $Z([\bm U])(T)$ is the Hadamard product of $Z(v)(T)$ and $Z(\sigma)(T)$, by \cite[Lemma~8.4.1]{hru:loe:lef}, $\lim_{T \limplies \infty}Z([\bm U])(T) = - \chi(I) v$.

% BEGIN REFEREE INSERTION R154
\RefereeRemark{154}{Typo}{The cited result is a lemma}{%
\textit{Original source line 3009.}\quad The reference is \emph{Lemma}~8.5.2 of Hrushovski--Loeser, not Proposition~8.5.2. Its hypothesis is piecewise integral-affine weight, including an integral constant term. Make this explicit here: $l_f$ is the valuation of the definable $\RV$-map $u\mapsto\rv(t)^l\prod_j f_j(u)$, so Remark~\ref{GLZ:char} gives that hypothesis after finite refinement. Since $I$ is doubly bounded, the lemma applies directly. In this particular proof the Hadamard-product step can also be replaced by the simpler identity $Z([\bm U])=v\,Z(\sigma)$.
}
% END REFEREE INSERTION R154
\end{proof}


\begin{rem}
The construction above no longer needs to go through this additional localization process ``$\loc$'' employed throughout \cite[\S~8]{hru:loe:lef}.
\end{rem}


\begin{rem}
Using $\bm \eta_{\loc} \circ \bm h_{\loc,m}$ instead of $\bm \eta \circ \bm h_m$, we can define the zeta function $Z([\bm U])(T)$ for every $[\bm U] \in \ggk^\natural \mgRV^{\bdd}[*]$, and the proof of Theorem~\ref{direct:mil} shows that $\lim_{T \limplies \infty} Z([\bm U])(T)$ exists and belongs to $\gmv[[\A]^{-1}]_{\loc}$. This and  Lemma \ref{lem:hm:bounded} gives the right half of the diagram (\ref{diag-diamond-zeta-function}). However,  we are not  claiming  that $\lim_{T \limplies \infty} Z([\bm U])(T)$ equals $(\Theta \circ \bb E_b)( [\bm U])$, which is reflected in the fact that there is no commutative relation between the left and right columns in (\ref{diag-diamond-zeta-function}). In fact, they are clearly not equal, as seen for example by considering $\bm U = \RV^{\circ\circ}$ (then the first value is $-1$ and the second one is $0$).  Such a failure explains the need to introduce an interpolant, such as the composite map $- \lim \circ Z \circ \int^\diamond$, in (\ref{diag-diamond-zeta-function}), for which Theorem \ref{direct:mil} may be formulated and proven.

% BEGIN REFEREE INSERTION R155
\RefereeRemark{155}{Gap; revised}{The bounded zeta claim has no established domain}{%
\textit{Original source line 3019.}\quad The claim for every element of $\ggk^\natural\mgRV^{\bdd}[*]$ cannot follow from the proof of Theorem~\ref{direct:mil}: the proposed $\bm h_{\loc,m}$ is not defined on all bounded objects, as the counterexample in \S~\ref{sec:hm:bounded} shows, and that proof uses a doubly bounded valuation set with an integral-affine weight. After repairing the summability domain, give a separate uniform rationality and finite-limit argument, or omit this remark and the right-hand part of diagram~\eqref{diag-diamond-zeta-function}. The numerical example in the present paragraph is correct: $\lim Z([\RV^{\circ\circ}])=-1$ whereas $\Theta\bb E_b([\RV^{\circ\circ}])=0$. With the diagram's $-\lim$ convention the former value is $+1$, which still distinguishes the maps.
}
% END REFEREE INSERTION R155
\end{rem}


\begin{rem}\label{rem-comp-milnor-fib}
The expression $(\Theta \circ \bb E^{\diamond})([\bm U])$, as an element in $\gmv[[\A]^{-1}]$, does not actually involve  $[\A]^{-1}$. Still, since the coefficients of $Z([\bm U])(T)$ does involve $[\A]^{-1}$ and it is a fact that the natural homomorphism $\gmv \fun \gmv[[\A]^{-1}]$ is not injective, we cannot really take  the motivic Milnor fiber $\mathscr S_f$ of $f$ in $\gmv$, at least not if $\mathscr S_f$ is viewed as something obtained through  $Z_f(T)$. It is this point of view that forces us to work with an integral whose target only involves doubly bounded sets in $\RV$, namely $\int^{\diamond}$, instead of $\int$, so as to facilitate the computation of the coefficients of $Z_f(T)$, and consequently with the nonarchimedean Milnor fiber $\mdl X_f$, which is proper invariant, instead of, perhaps, the more obvious definable set

% BEGIN REFEREE INSERTION R156
\RefereeRemark{156}{Typo}{Agreement and a missing source for noninjectivity}{%
\textit{Original source line 3024.}\quad Replace ``the coefficients \ldots{} does involve'' by ``the coefficients \ldots{} do involve''. For the noninjectivity assertion, cite Borisov, \emph{The class of the affine line is a zero divisor in the Grothendieck ring}, \href{https://arxiv.org/abs/1412.6194}{arXiv:1412.6194}. His nonequivariant example remains nonzero after giving it the trivial $\hat\mu$-action, because forgetting the action is a left inverse to the trivial-action inclusion.
}
% END REFEREE INSERTION R156
\[
\mdl X_t = \set{x \in X(\MM) \given f(x) = t},
\]
which is not proper invariant. This set $\mdl X_t$ may be understood as containing the $\puC$-rational points of  the analytic Milnor fiber introduced in \cite{nicaise:sebag:motivic:serre} (see Remark~\ref{rem-comp-milnor-fib-bis} below for more) and does play a role in \cite{hru:loe:lef} (but not in this paper).


On the other hand, in light of Theorem~\ref{direct:mil}, we can forego the zeta function point of view and recover $\mathscr S_f$ directly as $(\Theta \circ \bb E^{\diamond} \circ \int^{\diamond})([\mdl X_f])$. In that case there is truly no need to invert $[\A]$. In fact, we can also recover $\mathscr S_f$  directly as  $\Vol([\mdl X_f])$ (Notation~\ref{mov:vol}), the result is the same because the left half of the diagram (\ref{diag-diamond-interpol}) commutes (Corollary~\ref{db:to:novol}).
\end{rem}


\begin{rem}\label{rem-comp-milnor-fib-bis}
For any $t'\in \rv(t)^\sharp$, there exists an immediate automorphism $\sigma$ of $\puC$ over $\C$ with $\sigma(t') = t$, where ``immediate'' just means that $\sigma$ fixes $\RV(\puC)$ pointwise. Therefore, $\int[\mdl X_{t'}]=\int[\mdl X_t]$. Since $\mdl X_f = \bigcup_{t'\in \rv(t)^\sharp} \mdl X_{t'}$, it follows from compactness  that
\begin{equation*}
\int [\mdl X_f] = \int [\rv(t)^\sharp] \int [\mdl X_t] =  [1]\int [\mdl X_t].
\end{equation*}

% BEGIN REFEREE INSERTION R157
\RefereeRemark{157}{Gap}{Uniformity and parameter descent in the last compactness argument}{%
\textit{Original source line 3039.}\quad The immediate-automorphism observation is useful, but the displayed product formula needs a uniform argument over the fixed parameter field. An automorphism sending $t'$ to $t$ fixes $\C$ and $\RV$, but generally does not fix $\mathbb S=\C\dpar t$. Explain how the fiber equalities are represented by finite special-bijection/blowup witnesses over the appropriate base, how these witnesses vary definably with $t'$, and why they are available for parameters in a sufficiently saturated extension before applying compactness. Pointwise equality only for the Puiseux-series points does not by itself give that compactness conclusion. One possible route is to carry out the construction first over the henselization of $\C(t)$, then prove the required uniformity and base-change statement. The parallel sketch in Hrushovski--Loeser, Corollary~4.2.3, is relevant but does not replace the missing details requested in the human report. This is a proof gap, not a counterexample to the final motivic-volume identity.
}
% END REFEREE INSERTION R157
Since $\bb E_b([1]) = 1 $, this and Remark \ref{rem-comp-milnor-fib} show that
\[
\mathscr S_f  = \Vol([\mdl X_f]) = \Vol([\mdl X_t]).
\]


This is first shown in \cite{nicaise_ps_tropical}, but without taking the $\hat \mu$-actions into account, and also in \cite{Nic:Pay:trop:fub, forey_virtual_2017}. The arguments there all rely on resolution of singularities as well as other algebro-geometric machineries.

% BEGIN REFEREE INSERTION R158
\RefereeRemark{158}{Typo}{Uncountable noun}{%
\textit{Original source line 3046.}\quad Replace ``other algebro-geometric machineries'' by ``other algebro-geometric machinery''. 
}
% END REFEREE INSERTION R158
\end{rem}



\providecommand{\bysame}{\leavevmode\hbox to3em{\hrulefill}\thinspace}
\providecommand{\MR}{\relax\ifhmode\unskip\space\fi MR }
% \MRhref is called by the amsart/book/proc definition of \MR.
\providecommand{\MRhref}[2]{%
  \href{http://www.ams.org/mathscinet-getitem?mr=#1}{#2}
}
\providecommand{\href}[2]{#2}
\begin{thebibliography}{1}

\bibitem{denef_loser_motivic_igusa}
Jan~Denef and Fran\c cois~Loeser,
\newblock Motivic {{Igusa}} zeta functions,
\newblock {\em Journal of Algebraic Geometry}, \textbf{7} (1998), no.~3, 505--537.

\bibitem{denefloeser:arc}
\bysame, \emph{{Geometry on arc spaces of algebraic varieties}}, {3rd European
  congress of mathematics (ECM), Barcelona, Spain, July 10--14, 2000. Volume
  I}, Basel: Birkh\"auser, 2001, pp.~327--348.

\bibitem{den:loe:2002}
\bysame, \emph{Lefschetz numbers of iterates of the monodromy and truncated
  arcs}, Topology \textbf{41} (2002), no.~5, 1031--1040.

\bibitem{dries:1998}
Lou van~den Dries, \emph{Tame topology and \omin-minimal structures}, LMS Lecture Note
  Series, vol. 248, Cambridge University Press, Cambridge, UK, 1998.

\bibitem{engler:prestel:2005}
Antonio~J. Engler and Alexander Prestel, \emph{Valued fields}, Springer-Verlag,
  Berlin, 2005.

\bibitem{FY:Mot-Int-Mil-Fib}
Goulwen Fichou and Yimu Yin, {Motivic integration and Milnor fiber}, arXiv:1810.03561, 2018.

% BEGIN REFEREE INSERTION R159
\RefereeRemark{159}{Clarification}{Update the published Fichou--Yin reference}{%
\textit{Original source line 3083.}\quad This paper has appeared as G. Fichou and Y. Yin, \emph{Motivic integration and Milnor fiber}, J. Eur. Math. Soc. \textbf{24} (2022), 1617--1678, DOI \href{https://doi.org/10.4171/JEMS/1123}{10.4171/JEMS/1123}. Update the preprint-only citation if this bibliography is intended to be current.
}
% END REFEREE INSERTION R159


\bibitem{forey_virtual_2017}
Arthur Forey, {Virtual rigid motives of semi-algebraic sets}, {\em Selecta Mathematica (New Series)}, \textbf{25} (2019),  no.~1, Art. 6.


\bibitem{hrushovski:kazhdan:integration:vf}
Ehud Hrushovski and David Kazhdan, {Integration in valued fields}, in
  {\em Algebraic geometry and number theory}, Progr. Math., vol. 253, Birkh\"{a}user,
  Boston, MA, 2006, pp.~261--405.

\bibitem{hru:loe:lef}
Ehud Hrushovski and Fran{\c{c}}ois Loeser, {Monodromy and the {L}efschetz
  fixed point formula}, {\em Annales Scientifiques de l'\'Ecole Normale Sup\'erieure} (4) \textbf{48} (2015),
  no.~2, 313--349.

\bibitem{loeser2019motivic}
Fran{\c{c}}ois Loeser and Dimitri Wyss, {Motivic integration on the Hitchin fibration}, arXiv:1912.11638, 2019

% BEGIN REFEREE INSERTION R160
\RefereeRemark{160}{Clarification}{Update the published Loeser--Wyss reference}{%
\textit{Original source line 3101.}\quad Replace the preprint-only entry by F. Loeser and D. Wyss, \emph{Motivic integration on the Hitchin fibration}, Algebraic Geometry \textbf{8} (2021), 196--230, DOI \href{https://doi.org/10.14231/AG-2021-004}{10.14231/AG-2021-004}. The current entry also lacks its final full stop.
}
% END REFEREE INSERTION R160


\bibitem{Nic:Pay:trop:fub}
Johannes Nicaise and Sam Payne, {A tropical motivic {F}ubini theorem with
  applications to {D}onaldson-{T}homas theory}, {\em Duke Mathematical Journal} \textbf{168} (2019), no.~10, 1843–1886.

\bibitem{nicaise_ps_tropical}
Johannes Nicaise, Sam Payne, and Franziska Schroeter,
{Tropical refined curve counting via motivic integration}, {\em Geometry and Topology} \textbf{22} (2018),  no.~6, 3175--3234.

\bibitem{nicaise:sebag:motivic:serre}
Johannes Nicaise and Julien Sebag,
{Motivic Serre invariants, ramification, and the analytic Milnor fiber}, {\em Inventiones Mathematicae} \textbf{168}  (2007), no.~1, 133--173.

\bibitem{NS-deg-stab-rat}
Johannes Nicaise and Evgeny Shinder,
\newblock The motivic nearby fiber and degeneration of stable rationality,
\newblock {\em Inventiones Mathematicae} \textbf{217} (2019), no.~2, 377--413.


\bibitem{thuong-int-conj}
L{\^e}~Quy Thuong,
\newblock Proofs of the integral identity conjecture over algebraically closed
  fields,
\newblock {\em Duke Mathematical Journal} \textbf{164} (2015), no.~1, 157--194.

\bibitem{Yin:QE:ACVF:min}
Yimu Yin, {Quantifier elimination and minimality conditions in
  algebraically closed valued fields}, arXiv:1006.1393v1, 2009.

% BEGIN REFEREE INSERTION R161
\RefereeRemark{161}{Typo}{Incorrect arXiv year}{%
\textit{Original source line 3130.}\quad The cited version arXiv:1006.1393v1 was submitted on 7 June \emph{2010}, not in 2009. Use 2010 for this arXiv reference, or distinguish an earlier manuscript date explicitly.
}
% END REFEREE INSERTION R161

\bibitem{Yin:special:trans}
\bysame, {Special transformations in algebraically closed valued fields},
  {\em Annals of Pure and Applied Logic} \textbf{161} (2010), no.~12, 1541--1564.

\bibitem{Yin:int:acvf}
\bysame, {Integration in algebraically closed valued fields}, {\em Annals of
  Pure and Applied Logic} \textbf{162} (2011), no.~5, 384--408.

\bibitem{Yin:int:expan:acvf}
\bysame, {Integration in algebraically closed valued fields with
  sections}, {\em Annals of Pure and Applied Logic} \textbf{164} (2013), no.~1,
  1--29.

\bibitem{yin:hk:part:3}
\bysame, {{F}ourier transform of the additive group in algebraically
  closed valued fields}, {\em Selecta Mathematica (New Series)} \textbf{20} (2014), no.~4,
  1111--1157.

\bibitem{Yin:tcon:I}
\bysame, {Generalized {E}uler characteristic in power-bounded {$T$}-convex
  valued fields}, {\em Compositio Mathematica} \textbf{153} (2017), no.~12,
  2591--2642.

\bibitem{Yin:tcon:1.5}
\bysame, \emph{Integration in power-bounded {$T$}-convex valued fields},
  (2018), arXiv:1808.07129 [math.LO].

\end{thebibliography}



\end{document} 
